% ############################################### PREAMBLE AND CONTROL PANEL #

% ########################################################### DOCUMENT CLASS #
\documentclass[  %   draft
             % , demo  % black boxes instead of figures
              , 8pt, twoside, openright]{extbook}

% ########################################################## AND CONTROL PANEL #

% true - hmarginration 1:1, urlcolor - darkgreen, version - for searching
\newif\ifscreen 	   \screentrue
% true - no apendix, no author of grammer, no backmatter ISBN, 
\newif\ifPDF               \PDFtrue
% true - rotate pages of morph verb tables
\newif\ifPDFrotate         \PDFrotatetrue
% true - adds white searchable value !
\newif\ifPDFsearchvalue    \PDFsearchvaluetrue
\newif\ifmakecovers        \makecoverstrue
\newif\ifmakefrontmatter   \makefrontmattertrue
\newif\ifinputintroduction \inputintroductiontrue
\newif\ifinputletters      \inputletterstrue
\newif\ifinputphon         \inputphontrue
\newif\ifinputmorpho       \inputmorphotrue
\newif\ifmakebackmatter    \makebackmattertrue
\newif\ifshowcrop           \showcroptrue
% corrections
\newif\ifC                  \Ctrue

\newif\ifPRINTHOUSE 	   \PRINTHOUSEfalse
\newif\ifPDFSEARCH  	   \PDFSEARCHfalse
\newif\ifPDFPRINT 	   \PDFPRINTfalse
\newif\ifPDFtest	   \PDFtesttrue

%versions
%PRINTHOUSE
\ifPRINTHOUSE
\makecoversfalse
\screenfalse
\PDFfalse
\PDFrotatefalse
\PDFsearchvaluefalse
\showcroptrue
\Ctrue
\fi

%PDFSEARCH
\ifPDFSEARCH
\screentrue
\PDFtrue
\PDFrotatetrue
\PDFsearchvaluetrue
\showcropfalse
\Ctrue
\fi

%PDFPRINT
\ifPDFPRINT
\screenfalse
\PDFtrue
\PDFrotatefalse
\PDFsearchvaluefalse
\showcropfalse
\Ctrue
\fi

%PDFtest
\ifPDFtest
\makecoversfalse
\makefrontmatterfalse
\inputintroductionfalse
\inputletterstrue
\inputphonfalse
\inputmorphofalse
\makebackmatterfalse
\screentrue
\PDFtrue
\PDFrotatefalse
\PDFsearchvaluefalse
\showcropfalse
\Cfalse
\fi

% 1. PDF (for printing)  
% a) searchable value, b) no apendix, c) version= PDF d) no ISBN and no ISBN bar
% e) automatic rotation of landscape pages, f) no author of Icelandic Grammer
% 2. PDF (for viewing offline and searching) + screen (horizontal ratio 1:1) 
% 3. PRINTHOUSE
% a) no searchable value, b) apendix, c) version ISBN d) ISBN bar
% e) no rotation of landscape pages, f)urllinkcolor black
% f) A4 paper dimentions, C5 layout, shows cropping marks

% ########################################################## PAGE GEOMETRIES #

\usepackage%[showframe]
           {geometry}
           

\ifshowcrop
% Basic geometry of document
\geometry
  { headsep    =   \baselineskip
  , textwidth  = 42\baselineskip
  , textheight = 60\baselineskip
  , hmarginratio = \ifscreen 1:1\else 2:3\fi
  , vmarginratio = 2:3
  , bindingoffset = 0cm
  , onecolumn
  , a4paper 
  , layout=c5paper,
  , layouthoffset=\dimexpr(\paperwidth-\csname Gm@layoutwidth\endcsname)/2\relax
   , layoutvoffset=\dimexpr(\paperheight-\csname Gm@layoutheight\endcsname)/2\relax
  , showcrop
  }
\else
\geometry
  { headsep    =   \baselineskip
  , textwidth  = 42\baselineskip
  , textheight = 60\baselineskip
  , hmarginratio = \ifscreen 1:1\else 2:3\fi
  , vmarginratio = 2:3
  , bindingoffset = 0cm
  , onecolumn,
  , c5paper
  }
\fi

\ifshowcrop
% Geometry of the dictionary section
\newcommand\dictionarygeometry{\newgeometry
  { textwidth  = 42\baselineskip
  , textheight = 60\baselineskip
  , hmarginratio = \ifscreen 1:1\else 2:3\fi
  , vmarginratio = 3:2
  , bindingoffset = 0cm
  , twocolumn
  , a4paper 
  , layout=c5paper
  , layouthoffset=\dimexpr(\paperwidth-\csname Gm@layoutwidth\endcsname)/2\relax
   , layoutvoffset=\dimexpr(\paperheight-\csname Gm@layoutheight\endcsname)/2\relax
  , showcrop
  }}
   % Geometry of the cover pages
\newcommand\covergeometry{\newgeometry
  { textwidth  = 42\baselineskip
  , textheight = 60\baselineskip
  , hmarginratio = 1:1
  , vmarginratio = 2:3
  , bindingoffset = 0cm
  , onecolumn
  , layout=c5paper
 , layouthoffset=\dimexpr(\paperwidth-\csname Gm@layoutwidth\endcsname)/2\relax
   , layoutvoffset=\dimexpr(\paperheight-\csname Gm@layoutheight\endcsname)/2\relax
  , showcrop
  }} 
\else
% Geometry of the dictionary section
\newcommand\dictionarygeometry{\newgeometry
  { textwidth  = 42\baselineskip
  , textheight = 60\baselineskip
  , hmarginratio = \ifscreen 1:1\else 2:3\fi
  , vmarginratio = 3:2
  , bindingoffset = 0cm
  , twocolumn,
  , c5paper
  }}
   % Geometry of the cover pages
\newcommand\covergeometry{\newgeometry
  { textwidth  = 42\baselineskip
  , textheight = 60\baselineskip
  , hmarginratio = 1:1
  , vmarginratio = 2:3
  , bindingoffset = 0cm
  , onecolumn
  , c5paper
  }} 
\fi

% ################################################################# LANGUAGES #

\usepackage[icelandic, latin, czech]{babel}
\usepackage{csquotes}


% ################################################################## ENCODING #

\usepackage[utf8]{inputenc}

% Smashing uppercase letters with diacritics
% ------------------ THIS IS EVIL AND SHOULD NEVER BE DONE -------------------
% Don't smash anything if we are on the first run to produce the index
\IfFileExists{figures.idx}{
  % Czech uppercase letters with diacritics:
  \DeclareUnicodeCharacter{00C1}{\strut\smash{\' A}} % Á
  \DeclareUnicodeCharacter{010C}{\strut\smash{\v C}} % Č
  \DeclareUnicodeCharacter{010E}{\strut\smash{\v D}} % Ď
  \DeclareUnicodeCharacter{00C9}{\strut\smash{\' E}} % É
  \DeclareUnicodeCharacter{011A}{\strut\smash{\v E}} % Ě
  \DeclareUnicodeCharacter{00CD}{\strut\smash{\' I}} % Í
  \DeclareUnicodeCharacter{0147}{\strut\smash{\v N}} % Ň
  \DeclareUnicodeCharacter{00D3}{\strut\smash{\' O}} % Ó
  \DeclareUnicodeCharacter{0158}{\strut\smash{\v R}} % Ř
  \DeclareUnicodeCharacter{0160}{\strut\smash{\v S}} % Š
  \DeclareUnicodeCharacter{0164}{\strut\smash{\v T}} % Ť
  \DeclareUnicodeCharacter{00DA}{\strut\smash{\' U}} % Ú
  \DeclareUnicodeCharacter{016E}{\strut\smash{\r U}} % Ů
  \DeclareUnicodeCharacter{00DD}{\strut\smash{\' Y}} % Ý
  \DeclareUnicodeCharacter{017D}{\strut\smash{\v Z}} % Ž
  % Icelandic uppercase letters with diacritics: ÁÉÍÓÚÝ (ÐÞÆ) and
  \DeclareUnicodeCharacter{00D6}{\strut\smash{\" O}} % Ö
}{}
% ------------------ THAT WAS EVIL AND SHOULD NEVER BE DONE ------------------

% Non-breakable space
\DeclareUnicodeCharacter{00A0}{~}

% Non-breakable hyphen
\newcommand*{\nobreakhyphen}{\mbox{-}}
\DeclareUnicodeCharacter{2011}{\nobreakhyphen}


% ############################################################ FONTS, SYMBOLS #

\usepackage[T1]{fontenc}

\usepackage{tgpagella}
\usepackage[scaled=0.88]{helvet}      % relative scale of the two fonts

\def\phvfamily{\fontfamily{phv}\selectfont}

%\usepackage{stmaryrd}
% Let's keep this minimal instead of using the package:
\DeclareSymbolFont{stmry}{U}{stmry}{m}{n}
%\SetSymbolFont{stmry}{bold}{U}{stmry}{b}{n}
\DeclareMathSymbol\shortrightarrow\mathrel{stmry}{"01}
\DeclareMathSymbol\shortuparrow\mathrel{stmry}{"02}

% Methinks these are not necessary
%\usepackage{dingbat}
%\usepackage{manfnt}
%\usepackage{latexsym}

\usepackage{pifont}




% #################################################################### LAYOUT #

% \normalsize should be {8pt}{9.6pt}
\def\HUGE {\fontsize{23.887872pt}{3\baselineskip}\selectfont}
\def\Huge {\fontsize{19.906560pt}{3\baselineskip}\selectfont}
\def\huge {\fontsize{16.588800pt}{3\baselineskip}\selectfont}
\def\LARGE{\fontsize{13.824000pt}{2\baselineskip}\selectfont}
\def\Large{\fontsize{11.520000pt}{2\baselineskip}\selectfont}
\def\large{\fontsize{ 9.600000pt}{2\baselineskip}\selectfont}

% Two columns layout ruler
\setlength\columnsep    {2\baselineskip}
\setlength\columnseprule{0.4pt}

% Necessary for baseline alignment
\topskip=\baselineskip
\raggedbottom
\setlength\parskip{0pt} % it's better to avoid glue

% Temporarily suppress warnings
\hbadness=5000
\vbadness=5000

\setlength\emergencystretch{17pt}

% Allow smaller emergencystretch in several cases
\newenvironment{xtolerant}[2]{%
  \par
  \ifx\empty#1\empty\else\tolerance=#1\relax\fi
  \ifx\empty#2\empty\else\emergencystretch=#2\relax\fi
}{%
  \par
}

% ######################################################### ASSORTED PACKAGES #

\usepackage{tikz}
\usetikzlibrary{calc}

\usepackage{tipa}                      % IPA phonetics

\usepackage{fourier-orns}              % ornaments
\usepackage{amsmath}                   % non-breakable dash
\usepackage{hyphenat}                  % no hyphen in abbreviations

\usepackage{hanging}
\usepackage{fix2col}
\usepackage{fixltx2e}                  % subscript
\usepackage{pagecolor}
\usepackage{afterpage}                 % Used in coloring trick for covers
\usepackage{pdfpages}

\usepackage{multirow}
\usepackage{multicol}
\usepackage{rotating}

\usepackage{etoolbox} % http://ctan.org/pkg/etoolbox
%\usepackage{newunicodechar} 		% new definition of middle dot
\usepackage[ISBN=978-80-260-2385-2,SC0]{ean13isbn}

\usepackage{pgffor}

\usepackage{newunicodechar}

\newunicodechar{·}{\chejnikcdot}
\makeatletter
\newcommand{\chejnikcdot}{%
  \ifnum\lastpenalty=10042
    \kern-0.06667em\relax % dots should always be kerned
  \fi
  \textperiodcentered
  \@ifnextchar|{\kern\chejnikkern\relax}{\penalty10042 }%
}

\newcommand{\closeupdotbar}{%
  \renewcommand\chejnikkern{-0.06667em}%
}
\newcommand\chejnikkern{0pt} % default

% Middle dot + | little bit closer, nicer
%\newunicodechar{·}{\advancedcdot}
%\makeatletter
%\newcommand*\advancedcdot{\textperiodcentered
%  \@ifnextchar|{\hspace{-0.03667em}}{\checkifcdotnext}}
%\newcommand*\checkifcdotnext[2]
%  {\def\tmpa{·}\def\tmpb{#1#2}\ifx\tmpa\tmpb\hspace{-0.06667em}\fi#1#2}
%\makeatother

% ############################################################ BASELINE GRID #

\usepackage{eso-pic}
\usepackage{tikzpagenodes}

% Command to draw a baseline grid
\newcommand\drawbaselinegrid{%
  \begin{tikzpicture}[overlay,remember picture]
    \draw [red!30!white, ultra thin, dashed]
      (0,0) grid [ ystep  = \baselineskip, xstep = \textwidth
                 , shift  = (current page text area.north west)
                 , yshift = -\dp\strutbox
                 ] ++(\textwidth,-\textheight);
    \draw [red!30!white, thin]
      (0,0) grid [ step = \baselineskip, xstep = \textwidth
                 , shift=(current page text area.north)
                 ] ++(0.5\textwidth,-\textheight)
      (0,0) grid [ step = \baselineskip, xstep = \textwidth
                 , shift=(current page text area.north)
                 ] ++(-0.5\textwidth,-\textheight);
    \draw[black!10!white]
      (current page text area.north west)
        rectangle (current page text area.south east);
  \end{tikzpicture}}

% Draw a baseline grid on every page
% \AddToShipoutPicture{\drawbaselinegrid}


% ################################################################## METADATA #

\title{%
  \textbf{Islandsko-český studijní slovník}%
  \thanks{Tato kniha byla vytvořena v \LaTeX{}u pod Ubuntu 14.04.%
    Poděkování patří všem autorům, kteří publikují pod svobodnými licencemi.}}
\author{Aleš Chejn, Ján Zaťko, Jón Gíslason}
\date{říjen 2011}


% ############################################################## BIBLIOGRAPHY #

\usepackage[style=authoryear]{biblatex}
\addbibresource{slovnik.bib}


% ##################################################################### INDEX #

\usepackage{imakeidx}

% Index of authors of photographs
\makeindex
  [ name = figures
  , title = {Seznam autorů fotografií a ilustrací}
  , intoc
  , columnseprule
  , columnsep = \columnsep
  , noautomatic                       % We handle this in make.sh
  ]

% icelandic rules start
\usepackage{filecontents}
\def\mygroup#1{{\bfseries#1}}
\begin{filecontents*}{style.xdy}
;; style.xdy
(markup-letter-group :open-head "~n\mygroup{" :close-head "}")
\end{filecontents*}
% icelandic rules end


% ############################################################### PAGE STYLES #

\usepackage{fancyhdr}

% -------------------------------------- BASIC PAGE STYLE
\fancypagestyle{plain}{
  \fancyhf\relax                       % Clear header/footer
  \renewcommand \headrulewidth {0.0pt} % No header rule
  \renewcommand \footrulewidth {0.0pt} % No footer rule
  \fancyfoot[C]{\thepage}              % Page number in footer, centred
}

% -------------------------------------- INDEX PAGE STYLE  
%\fancypagestyle{indexstyle}{
%  \fancyhf\relax                       % Clear header/footer
%  \renewcommand \headrulewidth {0.4pt} % Thin header rule
%  \renewcommand \footrulewidth {0.0pt} % No footer rule
%  \fancyhead[C]{\thepage}              % Page number in footer, centred
%}

% -------------------------------------- DICTIONARY PAGE STYLE
\fancypagestyle{myheadings}{
  \fancyhf\relax                       % Clear header/footer
  \renewcommand \headrulewidth {0.4pt} % Thin header rule
  \fancyhead[CO,CE]{\thepage}          % Page number in header, centred
  % NOTE: the page numbers will be printed when the dictionary is ready
  \fancyhead[LE,LO]{\phvfamily\bfseries\rightmark}
  \fancyhead[RE,RO]{\phvfamily\bfseries\leftmark}
}

\pagestyle{plain}  

% ##################################################################### LISTS #

\usepackage{expdlist}
\usepackage{enumitem}

\setlist{nosep, topsep=\baselineskip}
\setlist[itemize]{leftmargin=*,label=\scalebox{.6}{\textbullet}}


% #################################################################### FLOATS #

\usepackage{caption}
\captionsetup{labelformat=empty}

% Alter some LaTeX defaults for better treatment of figures:
% See p.105 of "TeX Unbound" for suggested values. 
% See pp. 199-200 of Lamport's "LaTeX" book for details.
% ------------------------------------- Parameters for ALL pages:
\renewcommand \topfraction    {0.9}   %   max fraction of floats at top
\renewcommand \bottomfraction {0.8}   %   max fraction of floats at bottom
% ------------------------------------- Parameters for TEXT (non float) pages:
\setcounter{topnumber}    {1}
\setcounter{bottomnumber} {1}
\setcounter{totalnumber}  {1}         %   2 may work better
\setcounter{dbltopnumber} {2}         %   for 2-column pages
\renewcommand \dbltopfraction   {0.9} %   fit big float above 2-col. text
\renewcommand \textfraction     {0.07}%   allow minimal text w. figs
% ------------------------------------- Parameters for FLOAT (non text) pages:
\renewcommand \floatpagefraction{0.7} %   require fuller float pages
% N.B.: floatpagefraction MUST be less than topfraction !!
\renewcommand{\dblfloatpagefraction}{0.7} % require fuller float pages
% remember to use [htp] or [htpb] for placement

\makeatletter
\setlength{\@fptop}{0pt}
\setlength{\@fpbot}{0pt plus 1fil}
\makeatother

% This kills spacing around floats
\setlength \intextsep        {0pt}
\setlength \floatsep         {0pt}
\setlength \textfloatsep     {0pt}
\setlength \dbltextfloatsep  {0pt}
\setlength \dblfloatsep      {0pt}

% This kills spacing around captions
\setlength \abovecaptionskip {0pt}
\setlength \belowcaptionskip {0pt}


% #################################################################### TABLES #

\usepackage{booktabs}
\usepackage{tabularx}
\usepackage{longtable}
\usepackage{ltxtable}
\usepackage{lscape}

\ifPDFrotate
\else
\makeatletter			       
\renewcommand\PLS@Rotate[1]{}          % this command prevents pdflscape 
\makeatother                           % to rotate pages
\fi

%\LetLtxMacro\oldtabularx\tabularx
%\LetLtxMacro\oldendtabularx\endtabularx
%\renewenvironment{tabularx}[1]
%  {\oldtabularx{\columnwidth}{#2}}
%  {\oldendtabularx}

\newcommand\LTXfw[1]{\LTXtable{\columnwidth}{#1}}

%%% \LetLtxMacro\oldtabularx\tabularx
%%% \LetLtxMacro\oldendtabularx\endtabularx
%%% \newenvironment{longtabularx}[2]
%%%   {\oldtabularx{#1}{#2}}
%%%   {\oldendtabularx}

% ######################################################### czech babel vs cmidrule
% ###################################### http://tex.stackexchange.com/a/112001/8380
\begingroup
  \makeatletter
  \catcode`\-=\active
  \AtBeginDocument{
  \def\@@@cmidrule[#1-#2]#3#4{\global\@cmidla#1\relax
    \global\advance\@cmidla\m@ne
    \ifnum\@cmidla>0\global\let\@gtempa\@cmidrulea\else
    \global\let\@gtempa\@cmidruleb\fi
    \global\@cmidlb#2\relax
    \global\advance\@cmidlb-\@cmidla
    \global\@thisrulewidth=#3
    \@setrulekerning{#4}
    \ifnum\@lastruleclass=\z@\vskip \aboverulesep\fi
    \ifnum0=`{\fi}\@gtempa
    \noalign{\ifnum0=`}\fi\futurenonspacelet\@tempa\@xcmidrule}
  }
\endgroup

% macro that supress cmidrule at page brake
% http://tex.stackexchange.com/questions/110548/using-longtable-with-cmidrule-page-breaking-issues
\makeatletter
\def\LT@output{%
  \ifnum\outputpenalty <-\@Mi
    \ifnum\outputpenalty > -\LT@end@pen
      \LT@err{floats and marginpars not allowed in a longtable}\@ehc
    \else
      \setbox\z@\vbox{\unvbox\@cclv}%
      \ifdim \ht\LT@lastfoot>\ht\LT@foot
        \dimen@\pagegoal
        \advance\dimen@-\ht\LT@lastfoot
        \ifdim\dimen@<\ht\z@
          \setbox\@cclv\vbox{\unvbox\z@
\copy\LT@foot\vss}%
          \@makecol
          \@outputpage
          \setbox\z@\vbox{\box\LT@head}%
        \fi
      \fi
      \global\@colroom\@colht
      \global\vsize\@colht
      \vbox
        {\unvbox\z@
\box\ifvoid\LT@lastfoot\LT@foot\else\LT@lastfoot\fi}%
    \fi
  \else
    \setbox\@cclv\vbox{\unvbox\@cclv

\setbox0\lastbox
\ifdim\ht0=0.29999pt \unskip % hope that was a cline we threw away
\else
\nointerlineskip
\box0     % put it back, whatever it was
\fi

\copy\LT@foot\vss}%
    \@makecol
    \@outputpage
      \global\vsize\@colroom
    \copy\LT@head\nobreak
  \fi}

\makeatother

% These values ensure that for top, mid and bottom rules
% space above + rule width + space below = baselineskip
\setlength \abovetopsep    {0.43\baselineskip}
\setlength \heavyrulewidth {0.10\baselineskip}
\setlength \belowbottomsep {0.43\baselineskip}
\setlength \aboverulesep   {0.47\baselineskip}
\setlength \lightrulewidth {0.06\baselineskip}
\setlength \belowrulesep   {0.47\baselineskip}

%\usepackage{letltxmacro}

%%% \LetLtxMacro\oldtabular\tabular
%%% \LetLtxMacro\oldendtabular\endtabular
%%% \renewenvironment{tabular}[2][t]
%%%   {\oldtabular[#1]{#2}}
%%%   {\oldendtabular}

%%% \LetLtxMacro\oldtable\table
%%% \LetLtxMacro\oldendtable\endtable
%%% \renewenvironment{table}[1][h]
%%%   {\leavevmode\oldtable[#1]\vspace{\dp\strutbox}}
%%%   {\vspace{-\dp\strutbox}\oldendtable}
%%%   % This aligns a [t] aligned tab-envo to baseline grid


%%% % declention and conjugation tables
%%% \usepackage{floatrow}
%%% \DeclareFloatFont{footnotesize}{\footnotesize}
%%% % "scriptsize" is defined by floatrow, "tiny" not
%%% \floatsetup[table]{font=footnotesize}

% ################################################################### FIGURES #

\usepackage{adjustbox} % loads also graphicx

\graphicspath{%
  {/home/chejnik/Dokumenty/hvalur.org/images/biolib/full/}%
  {/home/chejnik/Dokumenty/hvalur.org/images/uploaded_files/}}

\newlength\dicfigheight
\newlength\dicfigskip

\newcommand\dicFigure[3]{ % {filename}{caption}{index entry}
  \setlength\fboxsep{0pt}\setlength\fboxrule{0.5pt}
  \def\dicfigbox{\fbox{%
    \adjincludegraphics [ max height = 0.8\columnwidth
                        , max width  = 0.8\columnwidth ] {#1}}}
  \setlength\abovecaptionskip {\dp\strutbox}
  \setlength\belowcaptionskip {\ht\strutbox} 
  \setlength\dicfigheight {\heightof\dicfigbox+\fboxrule}
  \setlength\dicfigskip
    {\baselineskip*\numexpr1+\dicfigheight/\baselineskip\relax-\dicfigheight}
  \begin{figure}[ht]
    \centering\vspace*{\dicfigskip}\dicfigbox\caption{#2}\index[figures]{#3}
  \end{figure}
}


% #################################################################### COLORS #

\usepackage{color}

% http://www.colorschemer.com/online.html
% main color - grammar color
% #66023C
\definecolor {darkgreen}          {rgb} {0.40, 0.01, 0.24}

\definecolor {royalazure}         {rgb} {0.00, 0.22, 0.66}
\definecolor {brown}              {rgb} {0.40, 0.01, 0.24}

% COLORS FOR THUMB INDEXES
\definecolor {babyblueeyes}       {rgb} {0.63, 0.79, 0.95}
\definecolor {unitednationsblue}  {rgb} {0.36, 0.57, 0.90}
\definecolor {blue(ryb)}          {rgb} {0.01, 0.28, 1.00}
\definecolor {darkblue}           {rgb} {0.00, 0.00, 0.55}
\definecolor {screamingreen}      {rgb} {0.46, 1.00, 0.44}
\definecolor {limegreen}          {rgb} {0.20, 0.80, 0.20}
\definecolor {islamicgreen}       {rgb} {0.00, 0.56, 0.00}
\definecolor {upforestgreen}      {rgb} {0.00, 0.27, 0.13}
\definecolor {icterine}           {rgb} {0.99, 0.97, 0.37}
\definecolor {orange(colorwheel)} {rgb} {1.00, 0.50, 0.00}
\definecolor {orange-red}         {rgb} {1.00, 0.27, 0.00}
\definecolor {oucrimsonred}       {rgb} {0.60, 0.00, 0.00}
\definecolor {cottoncandy}        {rgb} {1.00, 0.74, 0.85}
\definecolor {orchid}             {rgb} {0.85, 0.44, 0.84}
\definecolor {vividcerise}        {rgb} {0.85, 0.11, 0.51}
\definecolor {patriarch}          {rgb} {0.50, 0.00, 0.50}

\definecolor {title}              {RGB} {  16,   13,   32}
\definecolor {golden}             {RGB} { 241,  184,   45}
\definecolor {freqstar}           {RGB} {  91,    3,   99}
\definecolor {newdev}             {RGB} {   3,   11,   99}

\definecolor {color1}             {RGB} { 182,   86,    0}
\definecolor {color2}             {RGB} { 143,    9,    6}
\definecolor {color3}             {RGB} {   3,   23,  118}
\definecolor {color4}             {RGB} {   0,   82,  168}
\definecolor {color5}             {RGB} {   0,   85,  142}
\definecolor {color6}             {RGB} {   0,  115,  162}
\definecolor {color7}             {RGB} {  34,  146,  186}
\definecolor {color8}             {RGB} {  40,  159,  153}
\definecolor {color9}             {RGB} {   0,  125,  111}
\definecolor {color10}            {RGB} {   4,  107,   60}
\definecolor {color11}            {RGB} {  71,  134,   81}
\definecolor {color12}            {RGB} { 109,  134,   42}
\definecolor {color13}            {RGB} { 205,  194,   18}
\definecolor {color14}            {RGB} { 204,  162,   24}

%\definecolor {color1}             {RGB} { 239,  222,  205}
%\definecolor {color2}             {RGB} { 251,  206,  177}
%\definecolor {color3}             {RGB} { 251,  185,  142}
%\definecolor {color4}             {RGB} { 248,  166,  112}
%\definecolor {color5}             {RGB} { 216,  231,  246}
%\definecolor {color6}             {RGB} { 186,  215,  244}
%\definecolor {color7}             {RGB} { 152,  197,  243}
%\definecolor {color8}             {RGB} { 118,  181,  245}
%\definecolor {color9}             {RGB} { 181,  249,  185}
%\definecolor {color10}            {RGB} { 160,  248,  167}
%\definecolor {color11}            {RGB} { 128,  247,  136}
%\definecolor {color12}            {RGB} {  87,  248,   98}
%\definecolor {color13}            {RGB} { 250,  241,  138}
%\definecolor {color14}            {RGB} { 250,  240,  111}
%\definecolor {color15}            {RGB} { 244,  239,   85}
%\definecolor {color16}            {RGB} { 245,  239,   57}

\definecolor {lightgrey}          {RGB} { 105,  105,  105}
\definecolor {darkgreen_real}     {rgb} {0.00, 0.50, 0.00}


% ############################################################# THUMB INDEXES #

% Thumb indexes' colors
%\newcommand\BoxColor{%
%  \ifcase\theletternum darkgreen!30\or babyblueeyes\or unitednationsblue\or blue(ryb)\or screamingreen\or limegreen\or islamicgreen\or upforestgreen\or icterine\or orange(colorwheel)\or orange-red%
%  \or oucrimsonred\or cottoncandy\or orchid\or vividcerise\or patriarch\or babyblueeyes\or unitednationsblue\or blue(ryb)\or screamingreen\or limegreen\or islamicgreen\or upforestgreen\or icterine\or orange(colorwheel)\or orange-red%
%  \or oucrimsonred\or cottoncandy\or orchid\or vividcerise\or patriarch\else darkgreen!30\fi}
  
% Thumb indexes' colors
%\newcommand\BoxColor{%
%  \ifcase\theletternum darkgreen!30\or color1\or color2\or color3\or color4\or color5\or color6\or color7\or color8\or color9\or color10%
%  \or color11\or color12\or color13\or color14\or color1\or color2\or color3\or color4\or color5\or color6\or color7\or color8\or color9%
%  \or color10\or color11\or color12\or color13\or color14\or color1\or color2\or color3\or color4\or color5\or color6\else color7\or color8\or color9\or color10\or color11\or color12\fi}

%\def\BoxColor#1{darkgreen!\the\numexpr103-#1-#1-#1\relax!blue}

\newcommand\BoxColor[1]{%
\ifcase#1 darkgreen!30\or color13\or color14\or color1\or color2\or color3\or color4\or color5\or color6\or color7\or color8\or color9\or color10%
\or color11\or color12\or color13\or color14\or color1\or color2\or color3\or color4\or color5\or color6\or color7\or color8\or color9%
\or color10\or color11\or color12\or color13\or color14\or color1\or color2\or color3\or color4\or color5\or color6\else color7\or color8\or color9\or color10\or color11\or color12\fi}

\ifshowcrop
%%%%%%%%%%%%%%%%%%%%%%%%%%%%%%%%%%%%%%%%%%%%%%%%%%%%%%%%%%%%%%%%%%%%%%%%%%%%%%%%

\newcounter{letternum}
\newcounter{lettersum}                 \setcounter{lettersum}{34}
% these margins separate top thumb from top of page and last thumb from
% bottom of page. BEWARE, it's from top of the (cut out) page, not the folio!
\newlength{\thumbtopmargin}            \setlength{\thumbtopmargin}{1cm}
\newlength{\thumbbottommargin}         \setlength{\thumbbottommargin}{1cm}
\newlength{\thumbwidth}                \setlength{\thumbwidth}{2.20cm}
\newlength{\thumbheight}
\pgfmathsetlength{\thumbheight}%
  {(\csname Gm@layoutheight\endcsname-\thumbtopmargin-\thumbbottommargin)/\value{lettersum}}
% inner xsep of letter from box margin
\newlength{\thumbsep}                  \setlength{\thumbsep}{.30cm}
% overlap over cutting line
\newlength{\thumboverlap}              \setlength{\thumboverlap}{3cm}

\tikzset{
  thumb/.style={
    text=white,
    outer sep=0pt,
    inner xsep=\thumbsep,
    minimum height=\thumbheight,
    text width=\thumbwidth-2\thumbsep,
    font=\sffamily\bfseries,},
  thumb east/.style={thumb,
    xshift=\dimexpr(+\thumboverlap-\paperwidth+\csname Gm@layoutwidth\endcsname)/2\relax,
    align=left,anchor=north east,},
  thumb west/.style={thumb,
    xshift=\dimexpr(-\thumboverlap+\paperwidth-\csname Gm@layoutwidth\endcsname)/2\relax,
    align=right,anchor=north west,},
}

\def\thumbnew{}
\def\thumbold{}
\usepackage{everypage}
\AddEverypageHook{\if\relax\thumbnew\relax\xdef\thumbnew{\thumbold}\fi}

\def\ethumbs#1,#2\relax%
  {\if\relax#1\relax\else\eventhumb{#1}\fi%
   \if\relax#2\relax\else\ethumbs#2\relax\fi%
   \gdef\thumbnew{}%
   \gdef\thumbold{#1,}}

\def\othumbs#1,#2\relax%
  {\if\relax#1\relax\else\oddthumb{#1}\fi%
   \if\relax#2\relax\else\othumbs#2\relax\fi%
   \gdef\thumbold{#1,}%
   \gdef\thumbnew{}}

\newcommand{\drawthumb}[2]{%
  \begin{tikzpicture}[remember picture, overlay]
    \node [thumb #2, fill=\BoxColor{#1}]
       at ($(current page.north #2)-%
            (0,\csname Gm@layoutvoffset\endcsname+\thumbtopmargin-\thumbheight+%
               #1*\thumbheight)$) {\csname Let#1\endcsname};
   \end{tikzpicture}}



%%%%%%%%%%%%%%%%%%%%%%%%%%%%%%%%%%%%%%%%%%%%%%%%%%%%%%%%%%%%%%%%%%%%%%%%%%%%%%%%
\else
% THUMB INDEXES
% new counter to hold the current number of the letter to determine the vertical position
\newcounter{letternum}
% newcounter for the sum of all letters to get the right height of a box
\newcounter{lettersum}
\setcounter{lettersum}{34}
% some margin settings
\newlength{\thumbtopmargin}
\setlength{\thumbtopmargin}{\ifshowcrop 3cm\else 1.5cm\fi} %1cm
\newlength{\thumbbottommargin}
\setlength{\thumbbottommargin}{\ifshowcrop 6cm\else 2cm\fi} %2.5cm
% calculate the box height by dividing the page height
\newlength{\thumbheight}
\pgfmathsetlength{\thumbheight}{%
(\paperheight-\thumbtopmargin-\thumbbottommargin)%
/%
\value{lettersum}
}
% box width
\newlength{\thumbwidth}
\setlength{\thumbwidth}{\ifshowcrop 2cm\else 0.5cm\fi} %0.5cm
% style the boxes
\tikzset{
thumb/.style={
   text=white,
   minimum height=\thumbheight,
   text width=\thumbwidth,
   outer sep=0pt,
   font=\sffamily\bfseries,
 }
 }

\def\thumbnew{}
\def\thumbold{}
\usepackage{everypage}
\AddEverypageHook{\if\relax\thumbnew\relax\xdef\thumbnew{\thumbold}\fi}

\def\ethumbs#1,#2\relax{\if\relax#1\relax\else\eventhumb{#1}\fi%
                        \if\relax#2\relax\else\ethumbs#2\relax\fi%
                        \gdef\thumbnew{}%
                        \gdef\thumbold{#1,}%
}

\def\othumbs#1,#2\relax{\if\relax#1\relax\else\oddthumb{#1}\fi%
                        \if\relax#2\relax\else\othumbs#2\relax\fi%
                        \gdef\thumbold{#1,}%
                        \gdef\thumbnew{}%
}

\newcommand{\drawthumb}[2]{%
  % see pgfmanual.pdf for more information about this part
  \begin{tikzpicture}[remember picture, overlay]
    \node [thumb, fill=\BoxColor{#1}, text centered, anchor=north #2]
       at ($(current page.north #2)-%
%            (0,\thumbtopmargin+\value{letternum}*\thumbheight)%
            (0,\thumbtopmargin+#1*\thumbheight)$) {\csname Let#1\endcsname};
   \end{tikzpicture}}
\fi

\newcommand{\oddthumb} [1]{\drawthumb{#1}{west}}
\newcommand{\eventhumb}[1]{\drawthumb{#1}{east}}

\newcommand{\lettergroup}[1]%
  {\refstepcounter{letternum}%
   \expandafter\gdef\csname Let\theletternum\endcsname{#1}%
   \xdef\thumbnew{\theletternum,\thumbnew}%
   \fancyhead[LO]{\phvfamily\bfseries\rightmark%
     \expandafter\ethumbs\thumbnew\relax\relax\relax}%
   \fancyhead[RE]{\phvfamily\bfseries\leftmark%
     \expandafter\othumbs\thumbnew\relax\relax\relax}}
     
% ############################################################## LOCALIZATION #

% renames the index name
%\addto\captionsczech{%
%  \renewcommand\indexname{{Seznam autorů fotografií a ilustrací}}}

% renames the contents name
\addto\captionsczech{%
  \renewcommand\contentsname{Obsah}}

\renewcommand \contentsname {Obsah}
\renewcommand \chaptername  {Kapitola}


% ################################################################ SECTIONING #
\usepackage[explicit]{titlesec}

\setlength{\beforetitleunit}{\baselineskip}
\setlength{\aftertitleunit} {\baselineskip}

\titlespacing*{\chapter}   {0em}{*0}{*3} % Remembere there's a topskip too
\titlespacing*{\section}   {0em}{*1}{*2}
\titlespacing*{\subsection}{0em}{*1}{*2}

\newcommand\ghostdrop[2]{%
  \raisebox{-#1\baselineskip}[0pt][0pt]{#2}}

\titleformat {\chapter} {} {} {0pt} % {#1}
  {\ghostdrop{1.0}{\parbox[t]{\columnwidth}
     {\LARGE\bfseries\centering\MakeUppercase{#1}}}}

\def\longchapterskip{\vspace{2\baselineskip}}
\def\longsectionskip{\vspace{2\baselineskip}}

\titleformat {\section} {\bfseries} {} {0pt}
  {\ghostdrop{1.0}{\parbox[t]{\columnwidth}
     {\Large\bfseries\centering#1}}}

\titleformat {\subsection} {\bfseries} {} {0pt}
  {\ghostdrop{1.0}{\hspace{1em}\parbox[t]{\columnwidth-1em}
     {\large\bfseries#1}}}


\setlength\multicolsep{0pt}
%\BeforeBeginEnvironment{multicols}{\vspace{-\dp\strutbox}}


% ######################################################## TABLE OF CONTENTS #
\usepackage{titletoc}

\titlecontents {chapter} [0pt] {\addvspace\baselineskip\bfseries}
  {} {} {\titlerule*[2\baselineskip]{}\contentspage}

\titlecontents {section} [\baselineskip] {} {} {}
  {\titlerule*[1\baselineskip]{.}\contentspage}


% #################################################### ASSORTED CUSTOM MACROS #

\def\textIS#1{\foreignlanguage{icelandic}{#1}}
\def\textCS#1{\foreignlanguage{czech}{#1}}
\def\textLA#1{\foreignlanguage{latin}{#1}}

\usepackage{accsupp}

\makeatletter                                                                                                          
  \def\ACCSUPP@bdc{\special{pdf:code \ACCSUPP@span BDC}}                                                               
  \def\ACCSUPP@emc{\special{pdf:code EMC}}                                                                             
\makeatother

\usepackage{pgffor}

\newcommand\blspace[1][1]{\vspace{#1\baselineskip}}

\usepackage{xparse}


\newcommand \dicsymFrequent  {\raisebox{-.2ex}{\color{darkgreen}\ding{167}}}
\newcommand \dicsymSee       {$\shortrightarrow$}
\newcommand \dicsymCompare   {$\shortuparrow$}
\newcommand \dicsymIdiom     {$\diamondsuit$}
\newcommand \dicsymExampleIS {$\triangleright$}
\newcommand \dicsymExampleCS {\guilsinglright}
\newcommand \dicsymProverb   {{\color{newdev}\ding{96}}}
\newcommand \dicsymPhoto
  {\includegraphics[keepaspectratio,height=0.65em]{photo_icon}}

%\NewDocumentCommand\dicTerm{mo}
 % {{\phvfamily\bfseries\textIS{#1}}\IfValueT{#2}{, \phvfamily\bfseries\textIS{#2}}}

\NewDocumentCommand\dicEntry{o}
  {\par\hangpara{0.25\baselineskip}{1}%
   \IfValueT{#1}{\markboth{#1}{#1}}%
    \ifPDFsearchvalue
    \settowidth{\eyja}{#1}%
  \makebox[\eyja]{\color{white}#1}\hspace{-\eyja}%
  \fi
   \ignorespaces}
   
   \NewDocumentCommand\dicEntryGuide{o}
  {\par\hangpara{0.29\baselineskip}{1}%
   \IfValueT{#1}{\markboth{#1}{#1}}%
    \ifPDFsearchvalue
    \settowidth{\eyja}{#1}%
  \makebox[\eyja]{\color{white}#1}\hspace{-\eyja}%
  \fi
   \ignorespaces}
   
\newcommand\dicSubEntry
  {\\\hspace*{-\hangindent}}

\NewDocumentCommand\textAlt{mm}
  {\BeginAccSupp{method=pdfstringdef,ActualText={#2}}#1\EndAccSupp{}}

\newcounter{terms}

\NewDocumentCommand\dicStyleIntTerm{m}{\textIS{{\phvfamily\bfseries#1}}}
\NewDocumentCommand\dicStyleExtTerm{m}{\textIS{\textit{#1}}}

   \NewDocumentCommand\dicTerm{mo}
  {{\closeupdotbar \phvfamily\bfseries\textIS{#1}}\IfValueT{#2}{, \phvfamily\bfseries\textIS{#2}}}

  \newcommand\dicTermList[1]
  {\foreach \t [count=\n] in {#1}{\ifnum\n=1\else, \fi\dicTerm{\t}}}
  
\def\dicIPA#1{\textipa{[#1]}}

\NewDocumentCommand\dicPos{smo}
  {{\color{darkgreen}\textbf{\small#2}%
   \IfValueT{#3}{\nobreak\textsubscript{#3}}}}

%\NewDocumentCommand\dicFlx{smo}
 % {{\color{darkgreen}\footnotesize#1\IfValueT{#2}{\textsubscript{#2}}}}

  \NewDocumentCommand\dicFlx{smo}
  {{\color{darkgreen}\footnotesize#2%
   \IfValueT{#3}{\nobreak\textsubscript{#3}}}}
  
\newcommand\dicExampleCS[1]{\dicsymExampleCS~\textit{#1}}
\newcommand\dicExampleIS[1]{\dicsymExampleIS~\textIS{\textit{#1}}}

\NewDocumentCommand\dicLink{sm}
  {\IfBooleanF#1{\dicsymSee\nobreak\hspace{.08667em plus .03333em}}\nobreak\hspace{.03667em}\dicTerm{#2}}
\NewDocumentCommand\dicSynonym{sm}
  {(\IfBooleanF#1{\dicsymSee\nobreak\hspace{.08667em plus .03333em}}\nobreak\hspace{.03667em}\textIS{\textit{#2}})}
\NewDocumentCommand\dicAntonym{sm}
  {(\IfBooleanF#1{\dicsymCompare\nobreak\hspace{.08667em plus .03333em}}\nobreak\hspace{.03667em}\textIS{\textit{#2}})}
\NewDocumentCommand\dicCompare{sm}
  {(\IfBooleanF#1{\dicsymCompare\nobreak\hspace{.08667em plus .03333em}}\nobreak\hspace{.03667em}\textIS{\textit{#2}})}


\newcommand\dicPhraseIS[1]{\textIS{\textbf{#1}}}

\NewDocumentCommand\dicIdiom{mo}
  {\dicSubEntry{\color{newdev}\dicStyleIntTerm{#1} %
     \IfValueTF{#2}{+ \dicStyleIntTerm{#2} \markboth{#1 + #2}{#1 + #2}}{\markboth{#1}{#1}}}\dicsymIdiom}
     
\newcommand\dicProverb
  {\dicSubEntry\dicsymProverb}



\newcommand\dicLangCat[1]{{\footnotesize#1}}
\newcommand\dicFieldCat[1]{{\footnotesize#1}}


\newcommand\dicDirectTranslationCS[1]{#1}
\newcommand\dicIndirectTranslationCS[1]{{\footnotesize#1}}

\usepackage{placeins}

\newcommand*{\dicLetter}[2]{%
  \FloatBarrier\newpage\phantomsection\pdfbookmark[1]{#1}{#2}%
  \noindent\parbox[b][9\baselineskip][c]{\columnwidth}
    {\centering\HUGE\strut\smash{\MakeUppercase{#1}~\MakeLowercase{#1}}}%
      %\par\xdef\tpd{\the\prevdepth}
      %\par\prevdepth\tpd%
  \expandafter\lettergroup{\MakeLowercase{#1}}}

  
%tabulky
\newcommand{\specialcell}[2][l]{%
\begin{tabular}[#1]{@{}l@{}}#2\end{tabular}}


%\newcommand{\devisionguide}[1]{\hspace*{0.3cm}{{{{\foreignlanguage{icelandic}{\color{newdev}{\phvfamily{\textbf{#1}}}}}}}}}
%\newcommand{\devisionguideindent}[1]{\hspace*{-0.1cm}{{{{\foreignlanguage{icelandic}{\color{newdev}{\phvfamily{\textbf{#1}}}}}}}}}

%\newcommand\n{$n$\nobreakdash-\hspace{0pt}} %nonbreakable dash in grammar endings
%\def \nobreakseq {\nobreak \hskip 0pt \hbox} %nolinebreak in case like (-s, -)

% rule line
\newcommand{\HRule}{\rule{\linewidth}{0.1mm}}

% width of pdf found box
\newlength\eyja

%\newcommand\strutsmash[1]{\strut\smash{#1}}


% Dictionary symbols


% Command synonyms
%\let\freqstar\dicsymFrequent
%\let\photoicon\dicsymPhoto

\newcommand{\vstretch}[1]{\vspace*{\stretch{#1}}}

\newcommand*{\addthin}{\hskip0.16667em\relax}
\newcommand*{\addthinS}{\hskip0.06667em\relax}
\newcommand*{\addthinM}{\hskip0.10007em\relax}
\newcommand*{\addthinSS}{\hskip0.00007em\relax}
% #################################################################### COVERS #

\def\dictnameCZ{islandsko-český studijní slovník}
\def\dictnameIS{íslensk-tékknesk stúdentaorðabók}

\newcommand*{\frontcoverimages}{%
  ds_image_blaberjalyng_0_2.jpg,
  ds_image_lundi_0_1.jpg,
  ds_image_hestur_0_1.jpg,
  ds_image_hrafnafifa_0_2.jpg,
  ds_image_baldursbra_0_1.jpg,
  20948.jpg,
  ds_image_hvalur_0_1.jpg,
  ds_image_blodberg_0_1.jpg}

\newcommand*{\backcoverimages}{%
  ds_image_spoi_0_1.jpg,
  ds_image_mosalyng_0_1.jpg,
  ds_image_blaklukkulyng_0_1.jpg,
  20787.jpg,
  ds_image_ufsi_0_2.jpg,
  %ds_image_gullbra_0_1.jpg,
  ds_image_vatnsberi_0_1.jpg,
  ds_image_islandsfifill_0_2.jpg,
  21137.jpg}

\newcommand\makecoverwith[1]{
 \pagecolor{title}\afterpage\nopagecolor
 \begin{center}
 \vspace*{0.3cm}%{\fill}
 {\Huge\scshape\color{white}%
   \dictnameCZ\\\dictnameIS}
 \vspace*{1em}                        % I think this is necessary
 \begin{figure}[ht]\centering%
   \setlength\fboxsep{0pt}\setlength\fboxrule{0.5pt}%
   \foreach \file in #1{
     \fbox{\includegraphics[width=5.5cm]{\file}}}
 \end{figure}
 \vspace*{\fill}
 \end{center}
}


% ###################################################### HYPERREF & PDF INFO #

% unicode neccessary so that national characters in hypersetup appear ok
\usepackage[pdftex, unicode, hyperfootnotes=false]{hyperref}

% hyperlinks in black
\makeatletter
\let\Hy@linktoc\Hy@linktoc@none
\makeatother

\hypersetup
  { pdftitle={Islandsko-český studijní slovník}
  , pdfauthor={Aleš Chejn, Jón Gíslason, Ján Zaťko}
  , pdfsubject={Islandsko-český studijní slovník}
  , pdfkeywords={Islandsko-český studijní slovník%
               , islandština, čeština, slovník}
  , bookmarks=true
  , colorlinks=true
  , citecolor=black
  , urlcolor=\ifscreen darkgreen\else black\fi
  , linkcolor=black}


% ################################################################# MICROTYPE #

\usepackage[final]{microtype}

% ############################################################### HYPHENATION #

% this allows to include brackets inside hyphenation rules f.e (pro)-dis-ku-to-vat
\lccode`\(`\(
\lccode`\)`\)
\lccode`\[`\[
\lccode`\]`\]

\begin{hyphenrules}{icelandic}
\hyphenation{
  frum-semja
  burm-neskur
  ó-breytan-le-gur
  pla-tó-nskt
  skatta-fram-tal
  skrúf-gang-ur
  strand-rauð-viður
  prent-sver-ta
  grun-nur
  uppi-skrop-pa
  ger-vi-mennska
  sö-lu-stað-ur
  sý-na-gó-ga
  mark-mið
  trau-st-ur
  ís-kyggi-leg-ur
  brjó-sta-höld
  kné-krjú-pa
  nel-li-ka
  kaffi-kvörn
  margs-konar
  skrúf-gang-ur
  uppi-skrop-pa
  veð-rað-ur
  vest-ur-átt
  mis-mun-andi
  Bar-ba-do-si
  tung-u-mál
  au-ga-brún
  há-sker-pu-sjón-varp
  rúnn-stykki
  ger-vi-tungl
  yfir-borðs-menn-ska
  gæ-fu-mun-ur
  hlí-fi-skil-di
  á-vís-un
  ó-um-flýj-an-leg-ur
  san-mar-ín-sk-ur
  gæ-ða-stað-all
}
\end{hyphenrules}
\begin{hyphenrules}{czech}
\hyphenation{
  teo-rie
  mi-liardy
  (z)-od-po-věd-nost
  tan-tié-ma
  alian-ce
  pe-dia-tr(ička)
  pa-cienty
  zá-chra-nář(ka)
  dia-man-to-vá
  poe-zie
  ma-te-riál-ní
  egoi-s-ta
  egoist-ka
  nai-vi-ta
  ne-upřím-nost
  pě-ti-úhel-ník
  še-s-ti-úhel-ník
  za-šmo-dr-chat
  ar-cheo-lo-gic-ký
  ne-omlu-vi-tel-ný
  (vy)-pro-du-ko-vat
  dia-gno-s-ti-ko-vat
  teo-lož-ka
  ateist-ka
  oceá-nu
  im-pe-ria-li-s-ta
  (vy)-zdvih-nout
  ne-úměr-ný
  kon-ti-nuál-ní
  neu-t-rum
  ne-opráv-ně-ně
  vig-vam
  pro-zaic-ký
  geo-lož-ka
  seiz-mic-ká
  teo-rém
  olean-dr
  poe-ta
  (vy)-fo-to-gra-fo-vat
  (vy)-ven-ti-lo-vat
  pneu-mo-nie
  Lu-cem-bur-čan-(ka)
  ma-te-riály
  am-bi-cióz-ní
  Myan-ma(r)
  Myan-ma-řan-(ka)
  mol-dav-šti-na
  si-tua-cí
  ne-ukáz-ně-ný
  he-te-ro-trof-ní
  ne-omlu-vi-tel-ná
  ne-upřím-ný
  dvou-bliz-ný
  ne-úmy-sl-ný
  ne-ustá-lý-mi
  ne-ohle-du-pl-ný
  ne-osob-ní
  ne-oso-le-ný
  ne-uvě-ři-tel-ný
  ne-ohro-že-ný
  ne-oby-čej-ný
  ne-útul-ný
  ne-únav-ný
  ne-ože-nit
  po-ukáz-ka
  za-úto-čit
  aero-dy-na-mic-ký
  hru-bián-ství
  (pro)-di-s-ku-to-vat
  (vy)-ča-lou-nit
  (na)-oč-ko-vá-ní
  Bhú-tán-ka
  (za)-šně-ro-vat
  soud-ní
  re-vo-lu-cio-nář(ka)
  á-zer-báj-d-žán
  čtyř-hran-ný
  ze-m-dle-lý
  klou-zač-ka
  Vest-mann-ské
  tří-stran-ný
  (za)-bu-rá-ce-ní
  os-tud-ný
  pro-du-cent-(ka)
  chvá-s-tal-(ka)
  apar-theid
  asis-ten-ce
  asis-tent-(ka)
  (u)-či-nit
  (u)-za-mknout
  past-vi-na
  (za)-zá-lo-ho-vat
  Af-ri-čan-(ka)
  (za)-an-ga-žo-vá-ní
  se-rióz-ní
  opo-nent-(ka)
  pro-tiv-ník
  zvlášt-ní
  (po)-div-ný
  ho-mo-se-xuál-ní
  ko-alič-ní
  si-tua-ce
  si-tua-cí
  egoi-s-mus
  spe-ciál-ní
  spe-cia-li-zo-va-ný
  spe-cia-li-zo-va-né
  sobě-stač-ný
  sa-mo-zvo-le-ný
  sa-mo-čin-ně
 chvá-s-ta-vost
 klient-(ka)
 ku-rióz-ní
 (na)-šrou-bo-vat
 (za)-šrou-bo-vat
 krea-ti-vi-ta
 pneu-ma-ti-ka
 toa-le-ta
 so-cia-list-ka
 (vy)-dez-in-fi-ko-va-ný
 vše-uměl
 vše-uměl-(kyně)
 (pře)-škrt-nout
 (vy)-svlék-nout
 vy-učo-va-cí
 vy-učo-va-cích
 (vy)-chla-s-tat
 (vy)-te-to-vat
 pa-cient-(ka)
 ne-opod-s-tat-ně-ný
 de-se-ti-úhel-ník
 pe-rio-di-kum
 (za)-mum-lá-ní
 ne-smě-lá
 ne-dů-vě-ři-vá
 tee-na-ger-(ka)
 (z)-de-ge-ne-ro-vá-ní
 mo-ci-chti-vý
 bez-úhon-nost
 zne-uctít
 vo-dě-odol-ný
 we-bo-vý
 web-de-sign
 in-že-nýr-(ka)
 zá-pa-do-is-land-ský
 klient-(ka)
 za-uja-tý
 vy-uží-va-né
 vy-uží-vat
 (za)-re-gi-s-t-ro-vat
 (z)-kom-pri-mo-va-ný
 na-cio-na-li-s-mus
 ma-k-ro-eko-no-mic-ká
 bez-bo-lest-ný
 viol-ka
 troj-úhel-ník
 (za)-bu-rá-ce-ní
 ne-ústup-nost
 (vy)-cvi-če-ný
 zprá-va
 ge-ria-tr-(ička)
 osud-ný
 ná-pra-vy
 vstup-ní
 obec-ně
 re-pub-li-ka
 (po)-div-ný
 pa-slun-ce
 vliv-ný
 ba-ham-ský
 pe-dia-tr-(ička)
 ha-sič-(ka)
 vnuč-ka
 bro-s-kvoň
 lo-di-vod-(ka)
 mi-s-tr-(yně)
 in-že-nýr-(ka)
 dr-ži-tel-(ka)
 to-pič-(ka)
 pro-mi-nent-ní
 takt-ní
 frag-men-tár-ní
 po-zo-ru-hod-ný
 úspě-š-ný
 trans-fu-ze
 bur-žoa-zie
 pře-stro-je-ní
 oh-ni-vzdor-ný
 dis-tan-ční
 de-frau-da-ce
 fi-s-kál-ní
 zběh-lý
 kvap-ný
 střed-ní
 sa-mo-hlá-s-ka
 pře-vo-di-tel-ný
 zban-k-ro-to-va-ný
 zá-clo-na
 zná-mý
 zá-klad-ní
 na-uka
 ro-š-ťa-č-ka
 lí-bán-ky
 sou-du
 zkou-š-ka
 mlý-nek
 prá-šiv-ka
 kul-ma
 la-hod-ný
 zá-meč-ni-ce
 ne-vý-znam-ný
 ovzdu-ší
 mno-ho-úhel-ní-ko-vý
 žra-vý
 zcest-ný
 esen-ce
 (při)-nu-ce-ný
 skr-b-lic-ký
 skle-s-lý
 bez-bož-ný
 ne-opráv-ně-ný
 ošu-mě-lý
 pru-der-ní
 chvá-s-ta-vý
 sou-drž-ný
 van-dal-(ka)
 sa-bo-tér-(ka)
 (po)-div-ně
 mo-k-řad
 plyt-ký
 na-spo-du
 izo-la-ce
 le-toun
 za-han-be-ný
 ří-ze-ní
 kom-plex-ní
 pr-kén-ko
 ba-vič-(ka)
 plyt-ký
 ne-mo-tor-ně
 šéf-(ka)
 přes-čas
 zji-s-ti-tel-ný
 ta-ko-vý
 kňou-ral-(ka)
 se-me-s-tr
 po-pra-s-ka-ný
 han-ba
 vzpří-me-ný
 chys-ta-ný
 lik-vi-da-ce
 or-ga-ni-za-ce
 ri-s-kant-ní
 (za)-pís-kat
 (za)-hvíz-dat
 vrou-cí
  }
\end{hyphenrules}


% ############################################################# CZECH TYPOGRAPHY #
% https://math.feld.cvut.cz/tkadlec/ftp/text/textypo.pdf



% ################################################################# DOCUMENT #

\begin{document}
\pagenumbering{gobble}

% -------------------------------------------------------------- Front cover -

\ifmakecovers
\covergeometry
  \makecoverwith\frontcoverimages
\restoregeometry  
\fi


% ============================================================= FRONT MATTER =

\ifmakefrontmatter

% ------------------------------------------------------------------ Patitul -

\cleardoublepage

\begin{center}
  %\vspace*{2cm}{\fontsize{20}{20}\selectfont\scshape\dictnameCZ\\\dictnameIS}
  \vspace*{\stretch 1}
  {\Huge\scshape\dictnameCZ\\\dictnameIS}
  \vspace{\stretch 4}
\end{center}

% -------------------------------------------------------------------- Titul -

\cleardoublepage

\begin{center}
  %\vspace*{2cm}{\fontsize{24}{24}\selectfont\scshape\dictnameCZ\\\dictnameIS}
  \vspace*{\stretch 1}
  {\HUGE\scshape\dictnameCZ\\\dictnameIS}
  \vspace{\stretch 2}
\end{center}

% ----------------------------------------------------------------- Impresum -

\clearpage

\noindent\begin{minipage}[t][\textheight][t]{\textwidth}
Autoři \textit{Islandsko-českého studijního slovníku}:\\
  \textbf{Aleš Chejn, Ján Zaťko, Jón Gíslason}\\[\baselineskip]
\ifPDF
\else
Autor \textit{Islandské gramatiky}:\\
  \textbf{Vojtěch Kupča}\\[\baselineskip]
\fi
Poděkování:\\
  \textbf{Dorotě Nierychlewské-Chejn, Jiřině Hanzlové, Jiřině Chejnové,
          Simoně Petrásové,\\ Petru Mikešovi, Vojtěchu Kupčovi, 
          Renatě Peškové Emilsson, Amiru Mulahumici, \\ Kristýně Antonové, 
          Martině Kašparové, Lucii Brabcové, Samoriele, Adamu Kožouškovi}\\[\baselineskip]
Finanční příspěvek na výtisk exemplářů určených knihovnám\\ v České republice,
Slovenské republice a na Islandu poskytli:\\
  \textbf{Petr Mikeš, Dorota Nierychlewska-Chejn, Jiřina Chejnová, Otta Chejn}
  \ifPDF
  \\[39\baselineskip]
  \else
  \\[36\baselineskip]
  \fi
\textbf{Hvalur.org - 2016}\\
Licence: \textbf{GNU Free Documentation License}\\
Sazba: \textbf{Aleš Chejn} v prostředí  \textbf\LaTeX\\
  Zvláštní poděkování patří \textbf{Paolu al-Imk\=an\=\i\ Brasolin} a uživatelům  \textbf{\TeX\ - \LaTeX\ Stack Exchange}\\[\baselineskip]
\ifPDF
Verze: \textbf{PDF verze} \textit{Islandsko-českého studijního slovníku} vhodná pro \textbf{\ifscreen offline prohlížení a vyhledávání\else vlastní tisk\fi}\\ volně dostupná na \url{www.hvalur.org}\\[\baselineskip]
\else
\textbf\ISBN\\[\baselineskip]
\fi
\ifPDF
\else
Tisk: \textbf{Tiskárna Melmen, Pardubice}\\
\fi
Heslová část vysázena patkovým fontem \TeX\ Gyre Pagella,
hesla bezpatkovým fontem {\phvfamily Helvetica}.
\end{minipage}

% -------------------------------------------------------- Table of contents -

\cleardoublepage 

\pdfbookmark\contentsname{toc}    
\tableofcontents
\thispagestyle{empty}
\cleardoublepage 

\fi

\pagenumbering{arabic}
\setcounter{page}{3}
% ============================================================= INTRODUCTION =

\ifinputintroduction

% ------------------------------------- The study Icelandic-Czech dictionary -

\chapter{O Islandsko-českém studijním slovníku}
Cílem vytvoření\textit{ Islandsko-českého studijního slovníku} je zpřístupnění islandského jazyka všem zájemcům o~tento jazyk. 

\textit{Islandsko-český studijní slovník} je určen především českým a~slovenským studentům islandštiny, začátečníkům i~pokročilým, kteří zde naleznou ucelené fonetické, morfologické a~syntaktické informace, jakož
i~slovní zásobu potřebnou k~aktivnímu i~pasivnímu užití s~rozsahem vhodným ke studiu islandštiny na vysoké škole.
Slovník je rovněž určen všem zájemcům o~islandštinu a~Island, jimž slovník nabízí encyklopedické informace, úvod do fonetiky a~gramatiky a~velké množství fotografií a~ilustrací. 
Třebaže se jedná o~slovník primárně zaměřený na potřeby českých a~slovenských uživatelů, obsahuje také řadu informací, které usnadní práci islandskému uživateli (např. použití více jak 10 000 synonym a~antonym, překlady příkladů ap.).

Slovní zásoba původně vychází z~\textit{Concise Icelandic-English Dictionary} (\cite {ic_en}), značně však byla rozšířena díky přístupu k~webovému projektu \textit{ISLEX} (\cite {int1}). Velké množství hesel bylo
přidáno na podnět uživatelů online verze \textit{Islandsko-českého studijního slovníku}.

Slovník obsahuje 30~575 hesel s~46~725 významy nebo slovními spojeními. Z~celkového počtu 30~575 je 28~097 svébytných hesel obsahujících alespoň jeden význam nebo slovní spojení. 
Zbylá hesla (v~počtu 2~478) jsou varianty, dublety nebo nepravidelné tvary.

Slovník je primárně online elektronický slovník, volně přístupný na internetových stránkách \url{www.hvalur.org}  (\cite {int14}). 
Tištěná verze slovníku je publikována na vlastní náklady v~počtu 30 kusů, přičemž většina exemplářů je určena knihovnám a~institucím v~České republice, Slovenské republice a~na Islandu
a~několik exemplářů je určeno autorům nebo zájemcům z~řad veřejnosti. Dvě PDF verze (vyhledávání\,/\addthin prohlížení offline a~vlastní tisk) jsou dostupné zdarma na internetových stránkách slovníku (viz výše).

\textit{Islandsko-český studijní slovník} je otevřený svobodný projekt, jehož se v~průběhu tvorby zúčastnilo velké množství spoluautorů (viz Historie tvorby IČSS), kterým tímto jménem autorů velmi děkuji.
Děkuji také všem uživatelům online verze slovníku za trpělivost, dobrou vůli a~veškeré připomínky a~návrhy, jmenovitě především Petru Mikešovi, Adamu Kožouškovi a~Samoriele za neutuchající víru v~náš projekt.
Poděkování patří také našim rodinám a~blízkým, díky kterým jsme se mohli soustavně věnovat tvorbě slovníku.
Srdečné poděkování patří mé ženě Dorotce za podporu a~nadšení potřebné k~tvorbě slovníku.

\blspace[5]

{\centering Aleš Chejn, Brwinów, 24. prosince 2015\par}


% ------------ The history of the formation Icelandic-Czech dictionary study -

\chapter{Historie tvorby Islandsko-českého studijního slovníku}
\longchapterskip
Projekt \textit{Islandsko-českého studijního slovníku} (IČSS) vznikl v~roce 2001 na nepřímý podnět doc. PhDr. Heleny Kadečkové, CSc. z~Univerzity Karlovy,
kdy byl vypracován seznam slov. 
Tento seznam obsahoval islandská slova a~jejich překlady, později byly přidány základní morfologické informace o~slovním druhu a~deklinační\,/\addthin konjugační koncovky podle vzoru z~\textit{Concise Icelandic-English Dictionary} (\cite {ic_en}). 
Seznam v~této době obsahoval okolo 1 200 slov. Byl zveřejněn pomocí statických webových stránek a~později byl umístěn na webu Univerzity Karlovy. 

V~roce 2006 byl slovník rozšířen o~nových 1 500 heslových slov a~stávající hesla byla rozšířena o~nové významy, zvláště významy frázových sloves. 
Ke konci roku 2006  byl slovník převeden z programu Excel do databáze MySQL a~byly zahájeny práce na vytvoření aplikace \textit{Dictionary System} (DS) (\cite {int19}), která by podporovala tvorbu IČSS. 
Obsah slovníku byl zveřejněn pod svobodnou licencí na webových stránkách \url{www.hvalur.org} (\cite {int14}), kde se nachází dodnes. 

V~roce 2007 byla vytvořena první verze aplikace DS. V~této době se k~projektu přidala Renata Pešková Emilsson, která se nyní věnuje podpoře výuky češtiny na Islandu. 
Ta slovník obohatila rozsáhlým seznamem slov z~praktického života na Islandu. Dále se slovník rozšiřoval o~heslová slova podle \textit{Concise Icelandic-English Dictionary} (\cite {ic_en}). 
V průběhu následujících třech let bylo doplněno přibližně 15~000 slov. V této době byla hesla rozšířena o~synonyma, příklady, frekvenci, oborovou a~stylovou charakteristiku. K projektu se na Islandu připojila na krátkou dobu Kristýna Antonová, studentka islandštiny na Islandské univerzitě, která 
nás upozornila na nesrovnalosti v~uživatelském rozhraní aplikace DS.

Během léta 2009 byl vytvořen skript na generování skloňování a~časování heslových slov. Výsledky generování byly porovnány s~\textit{Beygingarlýsing íslensks nútímamáls} (\cite {DalvikVM}) a~během následujících let konzultovány
s~Jónem Gíslasonem. Kristín Bjarnadóttir, vedoucí projektu \textit{Beygingarlýsing íslensks nútímamáls}, zodpověděla řadu našich dotazů týkajících se deklinací a~konjugací.

V~roce 2009 začala spolupráce s~\textit{Biolib.cz} (\cite {int5}), jmenovitě s~Ondřejem Zichou. Obdrželi jsme  databázi s~náhledy obrázků, autory a~licencemi a~latinskými názvy botanických a~zoologických druhů.

V~zimním semestru 2009 jsme požádali s~Amirem Mulahumicem, bosenským studentem islandštiny na Islandské univerzitě, Studentský fond Islandské univerzity o~grant na tvorbu zvukové databáze islandských slov. 
Studentský fond návrh přijal. Peníze z~grantu jsme určili pro islandského rodilého mluvčího, Jóna Gíslasona, lektora islandského jazyka na Islandské univerzitě v~Reykjavíku, který nahrál více jak 22~000 islandských slov do elektronické formy. 
S~Dorotou Nierychlewskou-Chejn jsme připravili dva seznamy heslových slov (jeden pro nahrávajícího a~druhý pro zpracování nahrávek). Na podporu nahrávek jsme vytvořili webové stránky nazvané \textit{Databáze s~výslovností islandských slov} (\cite {int13}) a~tyto stránky jsme přeložili do 6~jazyků (Jón Gíslason do islandštiny, Ján Zaťko do slovenštiny, francouzštiny a~angličtiny, Amir Mulamuhic do srbochorvatštiny a~Dorota Nierychlewska-Chejn do polštiny).

V~roce 2010 jsme přidali do databáze pravidla islandské výslovnosti, která jsou uvedena v~knize \textit{Handbók um íslenskan framburð} (\cite {is2}). Z~těchto pravidel jsme automaticky vygenerovali fonetický zápis výslovnosti všech hesel ve slovníku. 
Jón Gíslason ověřil správnost fonetického zápisu u~všech heslových slov.

Na konci roku 2011 jsme spojili webové stránky slovníku se stránkami aplikace a~sjednotili jsme a~vylepšili vzhled webových stránek. Cennými připomínkami se na vytvoření aplikace podílela
RNDr. Ingrid Nagyová, PhD., vysokoškolská pedagožka Ostravské Univerzity. Ján Zaťko přeložil nové webové stránky, Průvodce po slovníku do slovenštiny, francouzštiny a~angličtiny a~celý manuál aplikace do angličtiny. 
Dorota Nierychlewska-Chejn přeložila stránky do polštiny, Jón Gíslason do islandštiny a~Marcos Helena do portugalštiny.

Od prosince 2011 do března 2014 jsme realizovali první čtení, ve kterém jsme doplnili chybějící informace, přidali nová hesla a~opravili pravopisné chyby. 
Obsah slovníku jsme porovnávali především s~\textit{Íslensk orðabók} (\mbox{\cite {is}}) a~\textit{Islandsko-českým slovníkem} (\cite {is7}). S~Vojtěchem Kupčou jsme konzultovali mnohé nesnáze při překladech.

V~letech 2013 až 2015 byl vytvořen morfologický klíč, který rozšiřuje klasický systém zobrazení morfologické informace v~islandských slovnících o~přesné zařazení daného hesla do určité deklinační\,/\addthin konjugační skupiny, jejíž
skloňování\,/\addthin časování je ukázáno na příkladu. V~roce 2014 jsme obdrželi roční grant od Ostravské univerzity v~Ostravě na tvorbu morfologického klíče. S~realizací grantu nám velmi pomohla doc. PhDr. Diana Svobodová, Ph.D. a~Ing. Zuzana Havrlantová.

Od března 2014 do srpna 2015 jsme realizovali druhé finální čtení, ve kterém jsme porovnávali informace s~webovým slovníkem \textit{ISLEX} (\cite {int1}). 
Na požádání jsme obdrželi od Þórdís Úlfarsdóttir (vedoucí projektu \textit{ISLEX}) seznam víceslovných příslovečných frází. Takto poskytnutý seznam nám usnadnil opravy slovníku. 
S~Þórdís Úlfarsdóttir jsme konzultovali různé gramatické jevy a~záležitosti týkající se deklinací\,/\addthin konjugací.
Martina Kašparová, studentka islandštiny na Islandské univerzitě, prošla celý slovník, přičemž se soustředila mimo jiné na stylistické a~pravopisné korektury překladů. 
Lucie Brabcová udělala jazykové a~typografické korektury celého slovníku.

Vojtěch Kupča vydal v~roce 2014 \textit{Islandskou gramatiku} (\cite {is77}) a~souhlasil s~připojením této publikace ke~knižní verzi IČSS.
V~roce 2015 jsme s~přispěním Paola al-Imk\=an\=\i\ Brasolin vysázeli slovník v~typografickém prostředí {\LaTeX}. V~únoru 2016 byl slovník vytištěn v~Tiskárně MELMEN, s.r.o v~Pardubicích a svázán 
v~Knihařství Vojtíšková Hrubá v~Pardubicích. 

\blspace[5]

{\centering Všem spoluautorům děkujeme za spolupráci.\par}


% -------------------------------------------------  Guide to the dictionary -

\chapter{Průvodce po slovníku}
\begin{multicols}{2}
\def\tableA{\noindent\begin{tabularx}{\columnwidth}{>{\bfseries}lll}
  \toprule
  \textit{přípona} & \textit{přípona s~koncovkou} & \textit{příklad} \\
  \midrule
  -ing-   & -ingur   & rugl··ingur \\
          & -ing     & hreyf··ing \\
          & -ingi    & erf··ingi \\
          & -ingar   & lækn··ingar \\
  -ling-  & -lingur  & vett··lingur \\
  -ning-  & -ningur  & flut··ningur \\
          & -ning    & set··ning \\
          & -ningar  & kos··ningar \\
  -að-    & -aður    & mark··aður \\
  -nað-   & -naður   & skil··naður \\
  -ar-    & -ari     & dóm··ari \\
  -skap-  & -skapur  & kveð··skapur \\
  -dóm-   & -dómur   & lær··dómur \\
  -leik-  & -leikur  & sann··leikur \\
          & -leiki   & mögu··leiki \\
  -lát-   & -látur   & laus··látur \\
  -átt-   & -átta    & víð··átta \\
  -sam-   & -samur   & starf··samur \\
  -sem-   & -semi    & hamingju··semi \\
  -nesk-  & -neskur  & tékk··neskur \\
          & -neska   & tékk··neska \\
  -leg-   & -legur   & blíð··legur \\
          & -lega    & blíð··lega \\
  -ung-   & -ungur   & gler··ungur \\
          & -unga    & snopp··unga \\
          & -ungi    & sveit··ungi \\
          & -ung     & nýj··ung \\
  -ótt-   & -óttur   & blett··óttur \\
  -ræn-   & -rænn    & suð··rænn \\
          & -ræna    & nor··ræna \\
  -ern-   & -erni    & fað··erni \\
  -ind-   & -indi    & vís··indi \\
  -and-   & -andi    & gef··andi \\
  -ug-    & -ugur    & göf··ugur \\
          & -ugt     & stöð··ugt \\
  -ust-   & -usta    & þjón··usta \\
  -ald-   & -aldur   & far··aldur \\
  -uð-    & -uður    & söfn··uður \\
  -úð-    & -úð      & dul··úð \\
  -versk- & -verskur & ung··verskur \\
          & -verska  & ung··verska \\
  -ist-   & -isti    & hass··isti \\
  -ism-   & -ismi    & húman··ismi \\
  -ang-   & -angur   & far··angur \\
  \bottomrule
  \end{tabularx}
  \captionof{table}{Seznam přípon označených ve slovníku}
  \label{table1}}

\def\tableB{\noindent\begin{tabularx}{\columnwidth}{>{\bfseries}lX}
  \toprule
  \textbf{symbol} & \textbf{vysvětlivka} \\
  \midrule
  m                   & označuje slovní druh \\
  \textsubscript{1}   & označuje třídu \\
  (-a, -ar)           & označuje deklinační / konjugační koncovky hesla \\
  \textsubscript{8}   & označuje podtřídu\\
  \bottomrule
  \end{tabularx}
  \captionof{table}{Morfologická informace}
  \label{table2}}

\def\tableC{\noindent\begin{tabularx}{\columnwidth}
    {>{\bfseries}l>{\footnotesize\itshape}lllll}
  \toprule
  \multirow{3}{*}{{{\textbf{f{\textsubscript{1}}} \Large{\textbf{7}}}}}
    & nom & kisa & kisan    & kisur & kisurnar \\
    & acc & kisu & kisuna   & kisur & kisurnar \\
    & dat & kisu & kisunni  & kisum & kisunum \\
    & gen & kisu & kisunnar & kisa  & kisanna \\
  \bottomrule
  \end{tabularx}
  \captionof{table}{Skloňování hesla \textit{kisa}}
  \label{table3}}

\def\tableD{\noindent\begin{tabularx}{\columnwidth}{>{\bfseries}lllll}
  \toprule
  \textit{zkratka} & \textit{slovo} & \textit{pád} & \textit{číslo}
    & \specialcell{\textit{předmět}\\\textit{podmět}} \\
  \midrule
  e-að  & eitthvað   & 1. & sg & neživotný \\
  e-ð   & eitthvað   & 4. & sg & neživotný \\
  e-ir  & einhverjir & 1. & pl & životný \\
  e-ja  & einhverja  & 4. & pl & životný \\
  e-jum & einhverjum & 3. & pl & životný \\
  e-m   & einhverjum & 3. & sg & životný \\
  e-n   & einhvern   & 4. & sg & životný \\
  e-r   & einhver    & 1. & sg & životný \\
  e-rra & einhverra  & 2. & pl & životný \\
  e-rs  & einhvers   & 2. & sg & životný \\
  e-s   & einhvers  & 2. & sg & neživotný \\
  e-u   & einhverju  & 3. & sg & neživotný \\
  \bottomrule
  \end{tabularx}
  \captionof{table}{Seznam použitých zkratek neurčitých zájmen}
  \label{table4}}

%#############################################################################

\section{Heslo}

\subsection*{Řazení slov}

Heslová slova jsou zobrazena tučně a~jsou seřazena podle islandské abecedy:

\blspace
{\bfseries\centering
  a~á b c d ð e é f g h i~í j k~l m n\\
  o~ó p q r s~t u~ú v~w x y ý z~þ æ ö\par}
\blspace

K~snadnější orientaci v~heslové části je použit u~každého písmene tzv. palcový index, který se nachází na hřbetě slovníku. Každé písmeno je v~palcovém indexu barevně odlišeno.

\subsection*{Homonyma}

Homonyma (heslová slova stejného reprezentativního tvaru) jsou označena povýšenými arabskými číslicemi uváděnými za heslovým slovem.

\blspace
  \dicEntry \dicTerm{vor\textsuperscript{1}} \dicPos{n} \dicFlx{(-s, -)}
  \dicEntry \dicTerm{vor\textsuperscript{2}} \dicPos{pron} \dicFlx{pers gen pl}
\blspace

Homonyma jsou seřazena podle českého ustáleného řazení slovních druhů. Poté následují předpony a~zkratky.
V~případě dvou podstatných jmen je pořadí mužský rod, ženský rod, střední rod. V~případě dvou podstatných jmen stejného rodu má přednost slovo se slabým skloňováním před slovem se silným skloňováním. V~případě dvou sloves stejného reprezentativního tvaru má přednost sloveso se slabým časováním před slovesem se silným časováním.

\subsection*{Varianty a~dublety}

U~některých heslových slov jsou v~záhlaví hesla uvedeny a) varianty (obměny)

\blspace
  \dicEntry \dicTerm{ærsla··fenginn}, \dicTerm{ærslafullur} \dicPos{adj}
\blspace

nebo b) dublety (dvojtvary).

\blspace
  \dicEntry \dicTerm{{ævin··týri}, \dicTerm{ævintýr}} \dicPos{n}
\blspace

Varianty a~dublety jsou méně frekventované než samotné heslo. Lze je vyhledat ve slovníku jako hesla obsahující fonetické a~morfologické informace s~odkazem na původní heslové slovo.

\blspace
  \dicEntry \dicTerm{ærsla··|fullur} \dicIPA{{a}{i}{\textsubring{r}}{s}{\textsubring{d}}{l}{a}{f}{\textscy}{\textsubring{d}}{l}{\textscy}{\textsubring{r}}} \dicPos{adj} \dicFlx{(comp -fyllri, sup -fyllstur)} \dicLink{ærslafenginn}
  \dicEntry \dicTerm{ævin··týr} \dicIPA{{a}{i}{\textlengthmark}{v}{\textsci}{n}{t\textsuperscript{h}}{i}{\textsubring{r}}} \dicPos{n} \dicFlx{(-s, -)} \dicLink{ævintýri}
\blspace

\subsection*{Vybrané tvary jako hesla}

Ve slovníku jsou uvedeny vybrané tvary heslových slov, u~kterých by mohl mít uživatel potíže s~identifikací. Jedná se o~a) genitiv jednotného čísla ({\textbf{sg}}) a~nominativ množného čísla ({\textbf{pl}}) podstatných jmen, b) ženský a~střední rod přídavných jmen a~2. a~3. stupeň rodu mužského přídavných jmen, c) oznamovací způsob ({\textbf{ind}}) přítomného času ({\textbf{praes}}) a~minulého času ({\textbf{pf}}) 1. osoby ({\textbf{pers}}) sg a~pl a~spojovací způsob ({\textbf{con}}) minulého času 1. osoby sg a~jiné nepravidelné tvary.

\blspace
  \dicEntry \dicTerm{öttum} \dicPos{v} \dicFlx{ind pf pl 1 pers} \dicLink{etja}
\blspace

\subsection*{Dělení heslových slov}

K~dělení heslových slov je použita střední tečka \textbf{'·'}, která odděluje části složeného slova, přípony a~předpony. Platí následující pravidla:

\renewcommand{\labelenumi}{\alph{enumi})}
\begin{enumerate}
\item jediné dělení heslového slova je označeno dvěma středními tečkami \textbf{'··'}. Může se jednat o~dělení částí složeného slova
\blspace
\dicEntry \dicTerm{ráð··hús}
\blspace

nebo o~dělení mezi kořenem a~předponou / příponou.

\blspace
\dicEntry \dicTerm{þjón··usta}
\blspace
\item V~případě, že se v~hesle vyskytuje více dělení, jsou použity dvě střední tečky \textbf{'··'} pro významově poslední dělení a~jedna střední tečka \textbf{'·'} pro zbylé / zbylá dělení.
\blspace
\dicEntry \dicTerm{við·bótar··líf·eyris·sparn·aður}
\blspace
\item V~hesle s~předponou je označena jedna nebo i~více předpon (každá zvlášť).
\blspace
\dicEntry \dicTerm{ó··að·finnan·legur}
\blspace
\item V~hesle s~příponou je označena jedna nebo i~více přípon (každá zvlášť).
\blspace
\dicEntry \dicTerm{létt·úð··ugur}
\blspace
\item V~hesle s~příponou je označena pouze přípona / přípony poslední části složeného slova. Důvodem pro neoznačení přípon jiných částí složených slov je lepší čitelnost hesla. Příklad označení pouze přípony poslední části složeniny:
\blspace
\dicEntry \dicTerm{flutninga··miðl·ari}
\end{enumerate}

Následuje seznam přípon, které jsou ve slovníku označeny:

\tableA

\section{Frekvence slov}

Symbolem \dicsymFrequent\ je označeno 2~825 hesel patřících mezi velmi frekventovaná islandská slova. Tento seznam vychází z~korpusu \textit{Íslenskur Orðasjóður} (\cite{int7}) a~knihy \textit{Íslensk orðtíðnibók} (\cite{is2}).

\blspace
  \dicEntry \dicTerm{koma}\textsuperscript{2} \dicsymFrequent
\blspace

\section{Fonetický zápis}

V~hranatých závorkách v~záhlaví hesla je uveden fonetický zápis výslovnosti hesla v~mezinárodní fonetické abecedě (IPA, angl. International Phonetic Alphabet).

\blspace
  \dicEntry \dicTerm{öldu··rót} \dicIPA{{\oe}{l}{\textsubring{d}}{\textscy}{r}{ou}{\textsubring{d}}}
\blspace

Znak '{\textlengthmark}' označuje dlouhou samohlásku nebo souhlásku.

\blspace
  \dicEntry \dicTerm{ör}\textsuperscript{2} \dicIPA{{\oe}{\textlengthmark}{\textsubring{r}}}
\blspace

\section{Slovní druhy, předpony a~zkratky}

V~záhlaví hesla je uvedena morfologická informace popisující a) slovní druh hesla, b) předponu nebo c) zkratku. Pokud se heslo řadí do více slovních druhů, jsou jednotlivé zkratky pro slovní druhy odděleny lomítkem.

\blspace
  \dicEntry \dicTerm{af} \dicPos{prep / adv}
\blspace

Informace o~slovním druhu, předponě či zkratce je vysázena barevně a~tučně.

\subsection*{Podstatná jména}

Podstatná jména jsou uvedena v~1. pádě (\textbf{nom}) jednotného čísla (\textbf{sg}) nebo v~1. pádě množného čísla (\textbf{pl}), pokud se heslo v~jednotném čísle nevyskytuje. Podstatná jména jsou označena následovně: \textbf{m} (l. masculinum) - podstatné jméno rodu mužského,
\textbf{f} (l. femininum) - podstatné jméno rodu ženského, \textbf{n} (l. neutrum) - podstatné jméno rodu středního, \textbf{subs} (l. substantivum) - podstatné jméno bez jednoznačného mluvnického rodu.

\blspace
  \dicEntry \dicTerm{hestur} \dicPos{m}
  \dicEntry \dicTerm{Hanoí} \dicPos{subs}
\blspace

\subsection*{Přídavná jména}

Přídavná jména jsou uvedena v~1. pádě (\textbf{nom}) rodu mužského (\textbf{m}) jednotného čísla (\textbf{sg}) 1. stupně (\textbf{pos}) (l. positivus) a~jsou označena zkratkou \textbf{adj} (l. adjectivum).

\blspace
  \dicEntry \dicTerm{fallegur} \dicPos{adj}
\blspace

V~případě, že u~přídavného jména dochází v~rodě ženském (\textbf{f}) nebo středním (\textbf{n}) k~přehlásce anebo je tvar druhého stupně (\textbf{comp} (l. comparativus)) či třetího stupně (\textbf{sup} (l. superlativus)) nepravidelný, pak jsou tyto tvary uvedeny v~nominativu jako heslová slova s~odkazem na hlavní heslový tvar

\blspace
  \dicEntry \dicTerm{yngri} \dicPos{adj} \dicFlx{comp m} \dicLink{ungur}
\blspace

nebo jako svébytná hesla s~významy.

\blspace
  \dicEntry \dicTerm{fremri} \dicPos{adj} \dicFlx{comp (sup fremstur)}
\blspace

\subsection*{Zájmena}

Zájmena jsou označena zkratkou \textbf{pron} (l. pronomen) a~jsou tříděna na zájmena ukazovací \textbf{dem} (angl. demonstrative), přivlastňovací \textbf{poss} (angl. possessive), osobní \textbf{pers} (angl. personal), neurčitá \textbf{indef} (angl. indefinite), zvratná \textbf{refl} (angl. reflexive) a~vztažná \textbf{rel} (angl. relative).

\blspace
  \dicEntry \dicTerm{þessi} \dicPos{pron} \dicFlx{dem}
\blspace

U~zájmen přivlastňovacích je v~3. osobě označen rod.

\blspace
  \dicEntry \dicTerm{hans} \dicPos{pron} \dicFlx{poss m}
\blspace

U~zájmen osobních je označen pád a~číslo.

\blspace
  \dicEntry \dicTerm{mig} \dicPos{pron} \dicFlx{pers acc sg}
\blspace

\subsection*{Číslovky}

Číslovky jsou označeny zkratkou \textbf{num} (l. numerale). V~případě vybraných číslovek (1 až 4) jsou ve slovníku uvedeny jako heslová slova také tvary ženského i~středního rodu. V~takovém případě jsou označeny zkratkou \textbf{m} pro mužský rod, zkratkou \textbf{f} pro ženský rod a~zkratkou \textbf{n} pro střední rod. Zkratkou \textbf{ord} (angl. ordinal number) jsou označeny číslovky řadové.

\blspace
  \dicEntry \dicTerm{tut··tugu} \dicPos{num}
  \dicEntry \dicTerm{tveir} \dicPos{num} \dicFlx{pl m}
  \dicEntry \dicTerm{fyrsti} \dicPos{num} \dicFlx{ord}
\blspace

\subsection*{Slovesa}

Slovesa jsou uvedena v~infinitivu a~jsou označena zkratkou \textbf{v} (l. verbum).

\blspace
  \dicEntry \dicTerm{fara} \dicPos{v}
\blspace

Slovesa ve středním rodě jsou označena jako \textbf{refl} (angl. reflexive).

\blspace
  \dicEntry \dicTerm{nálgast} \dicPos{v} \dicFlx{refl}
\blspace

Slovesa, která jsou vždy neosobní, jsou označena \textbf{impers} (angl. impersonal).

\blspace
  \dicEntry \dicTerm{svima} \dicPos{v} \dicFlx{impers}
\blspace

U~sloves je dále uvedeno, s~jakým pádem se pojí. Zkratka \textbf{acc} označuje 4. pád, \textbf{dat} 3. pád, \textbf{gen} 2. pád a~\textbf{nom} 1. pád. V~případě, že se sloveso pojí s~různými pády, je použito lomítko.

\blspace
  \dicEntry \dicTerm{klóra} \dicPos{v} \dicFlx{acc / dat}
\blspace

V~případě, že se sloveso pojí s~více pády, je použito znaménko plus.

\blspace
  \dicEntry \dicTerm{gefa} \dicPos{v} \dicFlx{dat + acc}
\blspace

\subsection*{Příslovce}

Příslovce jsou označena zkratkou \textbf{adv} (l. adverbium). Pokud chybí informace o~stupni, jedná se o~1. stupeň stupňování (\textbf{pos}). 2. stupeň je vždy označen zkratkou \textbf{comp} a~3. stupeň zkratkou \textbf{sup}.

\blspace
  \dicEntry \dicTerm{þar} \dicPos{adv}
  \dicEntry \dicTerm{ofar} \dicPos{adv} \dicFlx{comp (pos uppi, sup efst)}
  \dicEntry \dicTerm{síðast} \dicPos{adv} \dicFlx{sup (comp síðar)}
\blspace

\subsection*{Předložky}

Předložky jsou označeny zkratkou \textbf{prep} (l. praepositio). Pokud se předložka pojí výhradně s~jedním pádem, jsou použity zkratky pro 4. pád (\textbf{acc}), 3. pád (\textbf{dat}), 2. pád (\textbf{gen}).

\blspace
  \dicEntry \dicTerm{af} \dicPos{prep / adv} \dicFlx{dat}
\blspace

V~případě, že se předložka pojí s~různými pády, je použito lomítko.

\blspace
  \dicEntry \dicTerm{fyrir} \dicPos{prep / adv} \dicFlx{acc / dat}
\blspace

\subsection*{Spojky}

Spojky jsou označeny zkratkou \textbf{conj} (l. conjunctio).

\subsection*{Částice}

Částice jsou označeny zkratkou \textbf{part} (l. particula).

\subsection*{Citoslovce}

Citoslovce jsou označeny zkratkou \textbf{inter} (l. interjectio).

\subsection*{Předpony}

Předpony jsou označeny zkratkou \textbf{predp}.

\subsection*{Zkratky}

Zkratky jsou označeny zkratkou \textbf{zkr}.

\section{Skloňování, časování a~stupňování}

Každé heslo obsahuje zápis morfologické informace, která přesně informuje, jak se heslo skloňuje / časuje nebo stupňuje.

V~slovníku jsou použity dva odlišné systémy zápisu morfologické informace, které se vzájemně doplňují:

\begin{enumerate}
\item tradiční systém - do závorky za slovním druhem se zapisují deklinační / konjugační koncovky. Počet koncovek závisí na slovním druhu (viz dále).
\item numerický systém (vychází ze systému Helgi Haraldssona (\cite {is6})) - číslice odkazující na morfologický klíč, kde jsou podrobně vypsány deklinační / konjugační třídy.
Další informace o~numerickém systému lze nalézt na straně \pageref{sec:dec}.
\end{enumerate}

Deklinační / konjugační tvary nebo koncovky podstatných jmen, přídavných jmen, zájmen, číslovek, sloves a~příslovcí jsou uvedeny v~závorkách za slovním druhem.

Hnízdovací šev \textbf{'|'} označuje místo ve slově, kde se deklinační koncovka připojuje ke slovu.

\blspace
  \dicEntry \dicTerm{heils|a{\textsuperscript{1}}} \dicPos{f} \dicFlx{(-u)}
\blspace

Pokud není hnízdovací šev \textbf{'|'} uveden, tak se koncovka připojuje k~samotnému heslu.

\blspace
  \dicEntry \dicTerm{mynd} \dicPos{f} \dicFlx{(-ar, -ir)}
\blspace

Pokud dochází v~ohýbání k~přehlásce nebo je ohýbání nepravidelné, pak jsou uvedeny celé tvary hesla.

\blspace
  \dicEntry \dicTerm{maður} \dicPos{m} \dicFlx{(manns, menn)}
\blspace

V~případě, že heslové slovo má více tvarů, jsou varianty koncovky odděleny lomítkem.

\blspace
  \dicEntry \dicTerm{beð|ur} \dicPos{m} \dicFlx{(-s / -jar, -ir)}
\blspace

\subsection*{Podstatná jména}

U~podstatných jmen je uveden tvar hesla pro 2.~pád jednotného čísla a~1.~pád množného čísla.

\blspace
  \dicEntry \dicTerm{hval|ur} \dicPos{m} \dicFlx{(-s, -ir)}
\blspace

kde '\textit{hvals}' je 2.~pád jednotného čísla a~'\textit{hvalir}' 1.~pád množného čísla.
Pokud je uveden pouze jeden tvar, jedná se o~2.~pád jednotného čísla a~znamená to rovněž, že heslové slovo se neohýbá v~množném čísle.

\blspace
  \dicEntry \dicTerm{heils|a{\textsuperscript{1}}} \dicPos{f} \dicFlx{(-u)}
\blspace

Pokud je použita zkratka pro množné číslo (\textbf{pl}), znamená to, že se heslo ohýbá pouze v~množném čísle.

\blspace
  \dicEntry \dicTerm{afar··kostir} \dicPos{m} \dicFlx{pl}
\blspace

Pokud je použita zkratka \textbf{indecl}, znamená to, že podstatné jméno je nesklonné.

\blspace
  \dicEntry \dicTerm{október} \dicPos{m} \dicFlx{indecl}
\blspace

\subsection*{Přídavná jména}

U~přídavných jmen jsou uvedeny v~závorkách tvary, které poukazují na nepravidelnosti ve skloňování.

Pokud přídavné jméno v~ženském rodě neztrácí koncové '\textit{-r}', je uveden tvar hesla v~ženském rodě.

\blspace
  \dicEntry \dicTerm{hýr} \dicPos{adj} \dicFlx{(f -)}
\blspace

Dále jsou uvedeny tvary ženského rodu, ve~kterých dochází k~přehlásce:

\blspace
  \dicEntry \dicTerm{reglu··|samur} \dicPos{adj} \dicFlx{(f -söm)}
\blspace

Pokud je přídavné jméno nesklonné, je uvedena zkratka \textbf{indecl}.

\blspace
  \dicEntry \dicTerm{tví··mála} \dicPos{adj} \dicFlx{indecl}
\blspace

Pokud má přídavné jméno nepravidelné tvary 2.~a~3.~stupně stupňování, jsou tyto tvary uvedeny v~závorce spolu se zkratkou pro 2.~stupeň (\textbf{comp}) a~3.~stupeň (\textbf{sup}).

\blspace
  \dicEntry \dicTerm{þunnur} \dicPos{adj} \dicFlx{(comp þynnri, sup þynnstur)}
\blspace

V~případě, že je uveden pouze jeden ze dvou zbývajících stupňů, znamená to, že heslo se stupňuje pouze v~uvedených stupních.

\blspace
  \dicEntry \dicTerm{syðri} \dicPos{adj} \dicFlx{comp (sup syðstur)}
\blspace

\subsection*{Slovesa}

Islandská slovesa se dělí na slabá, silná a~nepravidelná. Vzhledem k~tomuto dělení se liší počet slovních tvarů uvedených v~závorce za zkratkou \textbf{v}. Princip uvádění slovních tvarů spočívá v~tom, že uživatel
je schopen (více méně) z~těchto tvarů odvodit časování slovesa. V~tomto slovníku jsme
záměrně umístili některá slabá slovesa mezi silná (místo dvou slovesných tvarů je uvedeno pět tvarů) kvůli tomu, aby byl uživatel přesněji informován o~časování slovesa.

Slabá slovesa, která tvoří 1.~osobu jednotného čísla minulého času oznamovacího způsobu koncovkou '\textit{-aði}', mají v~závorce uvedenu jen jednu koncovku, jmenovitě '\textit{(-aði)}'.

\blspace
  \dicEntry \dicTerm{ætl|a} \dicPos{v} \dicFlx{(-aði)}
\blspace

U~zbývajících skupin slabých sloves se v~závorce nachází koncovka 1.~osoby jednotného čísla minulého času oznamovacího způsobu a~supinum.

\blspace
  \dicEntry \dicTerm{kenn|a} \dicPos{v} \dicFlx{(-di, -t)}
\blspace

Silná nebo nepravidelná slovesa mají v~závorce vždy pět tvarů, jmenovitě 1.~osobu jednotného čísla přítomného času oznamovacího způsobu, 1.~osobu jednotného čísla minulého času oznamovacího způsobu, 1.~osobu množného čísla minulého času oznamovacího způsobu, 1.~osobu jednotného čísla minulého času spojovacího způsobu a~supinum.

\blspace
  \dicEntry \dicTerm{grípa} \dicPos{v} \dicFlx{(gríp, greip, gripum, gripi, gripið)}
\blspace

Silná nebo nepravidelná slovesa vyskytující se pouze v~neosobním užití mají v~závorce 4 tvary, jmenovitě 3.~osobu jednotného čísla přítomného času oznamovacího způsobu, 3.~osobu jednotného čísla minulého času oznamovacího způsobu, 3.~osobu jednotného čísla minulého času spojovacího způsobu a~supinum.

\blspace
  \dicEntry \dicTerm{kross··|bregða} \dicPos{v} \dicFlx{(-bregður, -brá, -brygði, -brugðið)}
\blspace

U~některých významů v~heslovém slově se vyskytují gramatické informace, které popisují chování daného významu.

Zkratka \textbf{impers} označuje, že se daný význam slovesa vyskytuje v~neosobním užití. Pomocí slovního spojení je vyjádřeno užití slovesa.

\blspace
  \dicEntry \dicTerm{batn|a} \dicPos{v} \dicFlx{(-aði)} \dicTerm{e-m batnar} \dicFlx{impers} (kdo) se uzdravuje, (kdo) se zotavuje, (komu) je lépe
\blspace

Zkratka \textbf{refl} označuje mediopasivní tvary slovesa.

\blspace
  \dicEntry \dicTerm{að··lag|a} \dicPos{v} \dicFlx{(-aði)} \dicIdiom{aðlagast} {\dicPhraseIS{aðlagast e-u}} \dicFlx{refl} {přizpůsobit se (čemu)}
\blspace

\subsection*{Příslovce}

U~příslovcí se v~záhlaví nachází tvary stupňování. Tyto tvary jsou uvedeny v~závorce spolu se zkratkou pro 2.~stupeň (\textbf{comp}) a~3.~stupeň (\textbf{sup}).

\blspace
  \dicEntry \dicTerm{vel\textsuperscript{2}} \dicPos{adv} \dicFlx{(comp betur, sup best)}
\blspace

Nepravidelné tvary 2.~a~3.~stupně jsou uvedeny buď jako hesla odkazující na základní tvar v~1.~stupni (\textbf{pos}), anebo jako svébytná hesla. V~obou případech jsou uvedeny v~závorce zbývající stupně.

\blspace
  \dicEntry \dicTerm{síst} \dicPos{adv} \dicFlx{sup (pos varla, comp síður)}
\blspace

Pokud některý ze stupňů uveden není, znamená to, že pro dané slovo neexistuje.

\blspace
  \dicEntry \dicTerm{síðar} \dicPos{adv} \dicFlx{comp (sup síðast)}
\blspace

\section{Morfologická informace} \label{sec:dec}

V~záhlaví hesla se kromě deklinačních / konjugačních koncovek nachází také morfologické numerické informace, které odkazují na morfologický klíč (viz strana \pageref{sec:morpho}), kde jsou podrobně vypsány kompletní deklinační / konjugační tabulky.

\blspace
  \dicEntry \dicTerm{ap|i} \dicPos{m}[1] \dicFlx{(-a, -ar)}[8]
\blspace

Morfologická informace je zapsána podle následujícího schématu: \textbf{m{\textsubscript{1}} (-a, -ar){\textsubscript{8}}}

\blspace
\tableB
\blspace

Deklinační / konjugační třídy a~jejich počet nejsou v~islandštině ustálené a~slouží pouze pro rozlišení odlišných skloňování / časování.

Názvy podtříd slouží pro rozlišení odlišných skloňování / časování.

\subsection*{Práce s~morfologickým klíčem}

Pro vyhledání kompletních deklinačních / konjugačních tabulek slouží následující postup, který demonstrujeme na ukázce:

\blspace
  \dicEntry \dicTerm{kis|a} \dicPos{f}[1] \dicFlx{(-u, -ur)}[7]
\blspace

Nejdříve určíme slovní druh (v~případě podstatných jmen také rod), v~tomto případě podstatné jméno rodu ženského, a~navštívíme Morfologický klíč (str. \pageref{sec:morpho}) pro podstatná jména rodu ženského (str. \pageref{sec:morpho_f}).
Následně dohledáme třídu (v~tomto případě číslo {\textsubscript{\textbf{1}}}) a~následně dohledáme dodatečné numerické informace za závorkou (v~tomto případě číslo {\textsubscript{\textbf{7}}}).

\tableC

\section{Definice}

\subsection*{Základní tvar}

Základním tvarem definice je, že jednomu islandskému slovu odpovídá jeden nebo více českých významů.

\blspace
  \dicEntry \dicTerm{landa··fræð|i} \dicPos{f} \dicFlx{(-i)} \textCS{zeměpis, geografie}
\blspace

V~některých případech je v~české části použita závorka. V~prvním případě je závorka užívaná ke~zkrácení zápisu dvou slov ('\textit{pomoc}', '\textit{výpomoc}').

\blspace
  \dicEntry \dicTerm{að··stoð} \dicPos{f} \dicFlx{(-ar)} \textCS{(vý)pomoc, asistence}
\blspace

V~druhém případě se v~závorce nachází upřesňující informace, která zabraňuje dvojznačnosti.

\blspace
  \dicEntry \dicTerm{fíkj|a} \dicPos{f} \dicFlx{(-u, -ur)} \textCS{fík} {\footnotesize {\textCS{(plod)}}}
\blspace

V~mnoha případech je závorka v~české části použita pro příklad užití českého slova a~tím vymezení jeho významu.

\blspace
  \dicEntry \dicTerm{að·gengi··legur} \dicPos{adj} {přístupný, dostupný (informace ap.)}
\blspace

Lomítko \textbf{'/'} slouží v~české části významu k~oddělení více významů pojících se s~jedním nebo více slovy a~je použito v~české části následujícím způsobem:

\blspace
  \dicEntry \dicTerm{aðal··|gata} \dicPos{f} \dicFlx{(-götu, -götur)} \textCS{hlavní třída / ulice}
\blspace

Význam přečteme jako '\textit{hlavní třída, hlavní ulice}'.

\subsection*{Opis heslového slova}

V~případech, kdy islandské slovo nemá ekvivalent v~českém jazyce, jsou významy islandských slov opsány delším opisem a~zobrazeny menším písmem. Příkladem jsou slova z~lidových pověstí, kulinářské speciality nebo botanické pojmy.

\blspace
  \dicEntry \dicTerm{til··ber|i} \dicPos{m} \dicFlx{(-a, -ar)} {\footnotesize{pov.}} {\footnotesize{stvoření, které dojí krávy a~ovce hospodářů}}
\blspace

\subsection*{Slovní spojení}

V~definici hesla jsou kromě samotných významů uváděna také slovní spojení, která jsou sázena tučně. Po slovním spojení následuje překlad slovního spojení.

\blspace
  \dicEntry \dicTerm{lær|a} \dicPos{v} \dicFlx{(-ði, -t)} \textbf{læra e-ð utanbókar} {učit se (co) zpaměti, memorovat (co)}
\blspace

Závorky jsou použity na tu část slovního spojení, kterou je možné vypustit.

\blspace
  \dicEntry \dicTerm{ár\textsuperscript{1}} \dicPos{f} \dicFlx{(-ar, -ar)} \textbf{taka (of) djúpt í árinni} {\footnotesize{přen.}} {přehnat (co v~tvrzení)}
\blspace

Příklad můžeme přečíst jako '\textit{taka djúpt í árinni}' nebo '\textit{taka of djúpt í árinni}'.

Mezi slovní spojení v~tomto slovníku řadíme a) tvary heslového slova, b) příkladová spojení, c) ustálená slovní spojení a~d) rčení nebo přísloví.

\begin{enumerate}
\item Tvary heslového slova\\
Do této skupiny řadíme tvar množného čísla hesla (zkratka \textbf{pl}).
\blspace
\dicEntry \dicTerm{bót} \dicPos{f} \dicFlx{(-ar, bætur)} \textbf{bætur} \dicFlx{pl} {dávky, přídavky, příspěvky}\\

Dále do této skupiny řadíme tvar slovesa v~středním slovesném rodě (zkratka \textbf{refl}).\\

\dicEntry \dicTerm{finna} \dicPos{v} \dicFlx{(finn, fann, fundum, fyndi, fundið)} \dicIdiom{finnast} {\dicPhraseIS{finnast}} \dicFlx{refl} {být nalezen, najít se, nalézt se, nacházet se}
\blspace
\item Příkladová spojení\\
Příkladová spojení ilustrují komunikační chování heslových slov v~příslušných významech a~v~různých syntagmatech (např. sloveso s~určitou předložkou ap.). Jsou označena tučně modrou barvou a~odsazená pro snadnou orientaci jako svébytné heslo.
\blspace
\dicEntry \dicTerm{koma\textsuperscript{2}} \dicPos{v} \dicFlx{(kem, kom, komum, kæmi, komið)} \dicIdiom{koma}[að] {\dicPhraseIS{koma að e-m}} {přiblížit se ke (komu), přijít ke (komu)}
\blspace
\item Ustálená slovní spojení\\
Tento druh slovních spojení je charakterizován většinou přeneseným významem. Tato spojení jsou uváděna za posledním významem nebo příkladovým spojením a~jsou označená tučně modrou barvou a~odsazená pro snadnou orientaci jako svébytné heslo.
\blspace
\dicEntry \dicTerm{kast} \dicPos{n} \dicFlx{(-s, köst)} \dicIdiom{kast} {\dicPhraseIS{þegar til kastanna kemur}} {\footnotesize{přen.}} {když jde do tuhého}
\blspace
\item Rčení a~přísloví\\
Rčení a~přísloví se nacházejí vždy na konci definice a~jsou odlišená symbolem \dicsymProverb\ a~odsazená pro snadnou orientaci.
\blspace
\dicEntry \dicTerm{æfing} \dicPos{f} \dicFlx{(-ar, -ar)} \dicProverb{} \dicPhraseIS{Æfing skapar meistara.} {\footnotesize{přís.}} \textCS{Cvičení dělá mistra.}
\end{enumerate}

\section{Vícevýznamovost}

\subsection*{Vícevýznamová hesla}

Značné množství hesel má více než jeden význam. Významy jsou odlišeny arabskými číslicemi a~v~některých případech i~islandskými synonymy.

\blspace
  \dicEntry \dicTerm{land} \dicPos{n} \dicFlx{(-s, lönd)} {\textbf{1.}} \dicSynonym{þurrlendi} {souš, pevnina, země} {\textbf{2.}} \dicSynonym{árbakki} {břeh} {\textbf{3.}} \dicSynonym{ríki} {země, stát, vlast} {\textbf{4.}} \dicSynonym{landareign} {země, půda, pozemek}
\blspace

\subsection*{Vícevýznamová slovní spojení}

V~případě, že slovní spojení uvnitř definice má více významů, jsou jednotlivé významy odlišeny písmeny {\textbf{a., b., c.}} atd.

\blspace
  \dicEntry \dicTerm{drag} {\dicPos{n} \dicFlx{(-s, drög)} {\textbf{1.}} {mokřina, podmoklý terén;} {\textbf{2.}} \textbf{drög} {\footnotesize{pl}} {\textbf{a.}} \dicSynonym{uppsprettur} {prameny (řeky ap.)} {\textbf{b.}} \dicSynonym{undirbúningur} {náčrt, návrh.}}
\blspace

\section{Řazení významů a~slovních spojení v~definici}
\longsectionskip

V~definici jsou nejdříve seřazeny významy slova a~odlišeny arabskými číslicemi.

\subsection*{V rámci významu}

Slovní spojení jsou uváděna v~rámci jednotlivého významu, pokud se s~ním významově vážou.

\blspace
  \dicEntry \dicTerm{tím|i} \dicPos{m} \dicFlx{(-a, -ar)} {\textbf{1.}} \dicSynonym{tíð} {čas, doba}; \textbf{í þann tíma} \dicFlx{adv} \dicSynonym{þá\textsuperscript{2}} {pak} {\textbf{2.}} \dicSynonym{rétt stund} {(správný) čas / moment}
\blspace

\subsection*{Za posledním očíslovaným významem}

Pokud se slovní spojení neváže k~žádnému významu, je řazeno za posledním očíslovaným významem. Ustálená slovní spojení jsou uvedena symbolem \dicsymIdiom.

\blspace
  \dicEntry \dicTerm{finna} \dicPos{v} \dicFlx{(finn, fann, fundum, fyndi, fundið)} {\textbf{3.}} \dicSynonym{skynja} {cítit, vnímat} \dicIdiom{finna}[að] {\dicPhraseIS{finna að e-u}} {kritizovat (co), nacházet chyby na (čem)}
\blspace

\subsection*{Způsob řazení slovních spojení}

Slovní spojení, která se nachází za posledním očíslovaným významem, jsou řazena podle předložky / příslovce / zájmena, se kterou se pojí. Pro snadnou orientaci v~rozsáhlých definicích je uvedeno heslové slovo nebo tvar heslového slova plus předložka / příslovce / zájmeno a~barevně odlišeno modrou barvou.

\blspace
  \dicEntry \dicTerm{koma\textsuperscript{2}} \dicPos{v} \dicFlx{(kem, kom, komum, kæmi, komið)} \dicIdiom{koma}[út] {\dicPhraseIS{koma út}} {vyjít (kniha ap.), objevit se (na trhu ap.)}
\blspace

Poté následují slovní spojení v~slovesném středním rodě řazená opět podle předložky / příslovce / zájmena, se kterou se pojí:

\blspace
  \dicEntry \dicTerm{koma\textsuperscript{2}} \dicPos{v} \dicFlx{(kem, kom, komum, kæmi, komið)} \dicIdiom{komast} {\dicPhraseIS{komast}} \dicFlx{refl} {dostat se, dostávat se}
\blspace

Poté následují ustálená slovní spojení, která nebylo možné zařadit do předchozích položek:

\blspace
  \dicEntry \dicTerm{koma\textsuperscript{2}} \dicPos{v} \dicFlx{(kem, kom, komum, kæmi, komið)} \dicIdiom{koma} {\dicPhraseIS{(komdu) sæll (og blessaður)}} {ahoj, nazdar, dobrý den}
\blspace

Nakonec jsou uvedena přísloví a~rčení označená symbolem \dicsymProverb\:

\blspace
  \dicEntry \dicTerm{eik} \dicPos{n} \dicFlx{(-ar / -ur, -ur)} \dicProverb{} \dicPhraseIS{Ekki / eigi fellur eik við fyrsta högg.} {\footnotesize{přís.}} {Jednou ranou dub nespadne.}
\blspace

\section{Synonyma a~antonyma / hesla ke~srovnání}
\longsectionskip

\subsection*{Synonyma}

Synonyma nebo synonymní slovní spojení jsou uvedeny před opisem významu hesla. Jsou umístěna v~závorkách. Islandská synonyma slouží a) islandskému uživateli k~rozlišení českých významů b) českému uživateli jako dodatečná informace o~významu slova.

\blspace
  \dicEntry \dicTerm{staða} \dicPos{f} \dicFlx{(stöðu, stöður)} {\textbf{1.}} \dicSynonym{ástand} {situace, stav, poměry} {\textbf{2.}} \dicSynonym{embætti} {post, (pracovní) místo / pozice (ve firmě ap.)}
\blspace

V~případě, že synonymum lze nalézt jako heslo ve slovníku, je použita šipka '\dicsymSee\ '.

\blspace
  \dicEntry \dicTerm{bygging} {\dicPos{f} \dicFlx{(-ar, -ar)} {\textbf{1.}} \dicSynonym{hús} {budova, stavba, stavení} {\textbf{2.}} \textit{(það að byggja)} {(vý)stavba} }
\blspace

\subsection*{Antonyma / Hesla ke~srovnání}

Antonyma či hesla ke~srovnání jsou uvedena za opisem významu heslového slova a~jsou označena šipkou '$\shortuparrow$'.

\blspace
  \dicEntry \dicTerm{kaldur} \dicPos{adj} \dicFlx{(f köld)} {\textbf{1.}} {studený, chladný} \textit{($\shortuparrow$ heitur)}
\blspace

\section{Oborové a~stylové charakteristiky}

K~upřesnění významu heslového slova slouží oborové a~stylové charakteristiky.

\subsection*{Oborové charakteristiky}

Oborové charakteristiky zařazují islandské slovo do určitého oboru a~naznačují použití daného slova v~jazyce

\blspace
  \dicEntry \dicTerm{berg} \dicPos{n} \dicFlx{(-s, -)} {\textbf{1.}} {\footnotesize{geol.}} {hornina, kámen} {\textbf{2.}} \dicSynonym{klöpp\textsuperscript{1}} {skalní stěna, skalisko}
\blspace

a~vymezují význam českého slova.

\blspace
  \dicEntry \dicTerm{mús} \dicPos{f} \dicFlx{(-ar, mýs)} {\textbf{1.}} {\footnotesize{zool.}} {myš} (\textit{l. Mus}) {\textbf{2.}} {\footnotesize{poč.}} {myš}
\blspace

\subsection*{Stylové charakteristiky}

Stylové charakteristiky jsou uváděny v~případě, kdy význam heslového slova je možné použít pouze při určité příležitosti nebo je význam jazykově zabarven.

\blspace
  \dicEntry \dicTerm{gaur} \dicPos{m} \dicFlx{(-s, -ar)} {\textbf{1.}} {\footnotesize{hovor.}} {chlápek, maník, člověk}
\blspace

\section{Syntax}

Syntaktické informace popisují:

\begin{enumerate}
\item zdali je sloveso tranzitivní nebo intranzitivní a~s~jakým(i) pádem (pády) se pojí, pokud je tranzitivní,
\item zdali je podmět / předmět životný či neživotný,
\item zdali je sloveso osobní či neosobní a~v~jakém pádu je podmět, pokud je sloveso neosobní,
\item zdali je podmět gramatický,
\item zdali je sloveso použito ve středním rodě,
\item nejčastější příklady, které se pojí se slovesem,
\item se kterými příslovci nebo předložkami se sloveso pojí.
\end{enumerate}

Syntaktické informace se ve slovníku vyskytují ve třech formách: a) informace uvedené zkratkou, b) teoretické příklady a~c) praktické příklady.

\subsection*{Syntaktické informace uvedené zkratkou}

Zde máme na mysli informace o~pádě (\textbf{nom, acc, dat, gen}), slovesu v~středním slovesném rodě (\textbf{refl}), neosobním slovesu (\textbf{impers}), čísle (\textbf{sg, pl}), rozkazovacím způsobu (\textbf{imper}) a~spojovacím způsobu (\textbf{con}).

\subsection*{Teoretické příklady}

Teoretickým příkladem nazýváme příklad, kdy podmět nebo předmět je zastoupen zájmenem neurčitým v~pádě, se~kterým se slovo pojí a~sloveso je zpravidla v~infinitivu. V~případě, že podmět není v~1.~pádě (např. u~sloves neosobních), je sloveso uvedeno v~3.~osobě jednotného čísla přítomného času.
Ve slovních spojeních používáme zkratky pro tvary neurčitého zájmena \textbf{einhver} (někdo) a~\textbf{eitthvað} (něco). Tyto zkratky označují a) pád, b) číslo a~c) životnost nebo neživotnost podmětu / předmětu.\\
Následuje seznam zkratek neurčitých zájmen použitých ve slovníku: \\

\tableD
\blspace
Příklad použití zkratek s~neurčitými zájmeny:

\blspace
  \dicEntry \dicTerm{leit|a} \dicPos{v} \dicFlx{(-aði)} \textbf{leita að e-u / e-m} {hledat (co / koho)} \dicExampleIS{leita að lyklunum} \dicExampleCS{hledat klíče}
\blspace

Příklad tedy přečteme jako '\textit{leita að einhverju}' (\textit{hledat co}) a~'\textit{leita að einhverjum}' (\textit{hledat koho}). V~tomto případě se sloveso pojí s~3.~pádem a~předmět může být životný i~neživotný.

\subsection*{Praktické příklady}

Ve slovníku se vyskytuje velké množství praktických příkladů a~jejich překladů. Praktický příklad se liší od příkladu teoretického tím, že místo zájmen neurčitých jsou uvedena frekventovaná slova z~běžného projevu. Praktické příklady slouží k~ilustraci použití
heslového slova.\\
Praktickým příkladem může být a) (složené) slovo, b) slovní spojení nebo c) věta.

\begin{enumerate}
\item Slovo jako praktický příklad:
\blspace
\dicEntry \dicTerm{hett|a} \dicPos{f} \dicFlx{(-u, -ur)} \textbf{1.} {kapuce, kápě} \dicExampleIS{hettuúlpa} \dicExampleCS{bunda s~kapucí}
\blspace
\item Slovní spojení jako praktický příklad:
\blspace
\dicEntry \dicTerm{tím|i} \dicPos{m} \dicFlx{(-a, -ar)} \textbf{3.} \dicSynonym{klukkustund} {hodina} {\footnotesize{(šedesát minut)}} \dicExampleIS{tveggja tíma gangur} \dicExampleCS{dvouhodinový pochod}
\blspace
\item Věta jako praktický příklad:
\blspace
\dicEntry \dicTerm{mæt|a} \dicPos{v} \dicFlx{(-ti, -t)} \dicIdiom{mætast} \dicPhraseIS{mætast} \dicSynonym*{ná saman} \dicFlx{refl} {setkat se, potkat se} \dicExampleIS{Þau mættust á gangstéttinni.} \dicExampleCS{Potkali se na chodníku.}
\end{enumerate}

\section{Fotografie a~ilustrace}

Ve slovníku je publikováno 630 fotografií rostlin, živočichů a~fotografií, která dokumentují místopisně heslová slova (např. kulinářské speciality).
Fotografie jsou publikovány pod veřejnými licencemi. Jména autorů a~názvy licencí lze nalézt v~Seznamu autorů fotografií a~ilustrací na straně \pageref{sec:photo}. U~hesla a~konkrétně u~daného významu, ke kterému se fotografie vztahuje, je použit symbol fotoaparátu \dicsymPhoto. Fotografie je umístěna buď pod heslem, anebo v~horní části stránky, kde se heslo nachází, anebo na stránce (stránkách) následující(ch).

%\par\begin{center}\setlength\fboxsep{0pt}\setlength\fboxrule{0.5pt}
%\fbox{\includegraphics[width=6cm]{ds_image_posthus_0_2.jpg}}\end{center}
%\par\begin{center}\footnotesize {Pósthús}\end{center}

\end{multicols}

% ---------------------------------------- List of abbreviations and symbols -

\chapter{Seznam použitých zkratek a symbolů}
\begin{multicols}{2}
\begin{description}[leftmargin=!,labelwidth=1.5cm]
%[\breaklabel\setleftmargin{40pt}\setlabelstyle{\bfseries}
%\addtolength{\itemsep}{-0.5\baselineskip}]
\item[{abb}] {zkrácený}
\item[{acc}] {4. pád, akuzativ}
\item[{adj}] {přídavné jméno}
\item[{adv}] {příslovce}
\item[{akt}] {rod činný, aktivum}
\item[{anat.}] {anatomie}
\item[{angl.}] {anglicky}
\item[{astro.}] {astronomie, výzkum vesmíru}
\item[{ap.}] {a podobně}
\item[{básn.}] {básnický, poetický výraz}
\item[{biol.}] {biologie}
\item[{bot.}] {botanika}

\item[{comp}] {2. stupeň stupňování, komparativ}
\item[{con}] {spojovací způsob, konjunktiv}
\item[{conj}] {spojka}
\item[{dat}] {3. pád, dativ}
\item[{def}] {určitý tvar}
\item[{dem}] {zájmeno ukazovací}
\item[{dět.}] {dětsky}

\item[{e-að}] {co (sg nom neživotný) (\textit{eitthvað})}
\item[{e-ð}] {co (sg acc neživotný) (\textit{eitthvað})}
\item[{e-ir}] {kdo (pl nom životný) (\textit{einhverjir})}
\item[{e-ja}] {koho (pl acc životný) (\textit{einhverja})}
\item[{e-jum}] {komu (pl dat životný) (\textit{einhverjum})}
\item[{e-m}] {komu (sg dat životný) (\textit{einhverjum})}
\item[{e-n}] {koho (sg acc životný) (\textit{einhvern})}
\item[{e-r}] {kdo (sg nom životný) (\textit{einhver})}
\item[{e-rra}] {koho (pl gen životný) (\textit{einhverra})}
\item[{e-rs}] {koho (sg gen životný) (\textit{einhvers})}
\item[{e-s}] {čeho (sg gen neživotný) (\textit{einhvers})}
\item[{e-u}] {čemu (sg dat neživotný) (\textit{einhverju})}

\item[{ekon.}] {ekonomika, obchod}
\item[{elek.}] {elektřina}
\item[{f}] {rod ženský}
\item[{filos.}] {filosofie, logika}
\item[{form.}] {formálně}
\item[{fyz.}] {fyzika}
\item[{gen}] {2. pád, genitiv}
\item[{geog.}] {zeměpis, geografie}
\item[{geol.}] {geologie}
\item[{han.}] {hanlivě, pejorativně}
\item[{hist.}] {historie}
\item[{hovor.}] {hovorově}
\item[{hrub.}] {hrubě}
\item[{hud.}] {hudba}
\item[{imper}] {rozkazovací způsob, imperativ}
\item[{impers}] {sloveso neosobní}
\item[{ind}] {oznamovací způsob, indikativ}
\item[{indecl}] {nesklonný}
\item[{indef}] {neurčitý tvar, (pron +) zájmeno neurčité}
\item[{inf}] {infinitiv}
\item[{int}] {tázací zájmeno}
\item[{inter}] {citoslovce}
\item[{chem.}] {chemie}
\item[{jaz.}] {jazykověda}
\item[{kulin.}] {kulinářství, vaření}
\item[{l.}] {latinsky}
\item[{let.}] {letectví}
\item[{lit.}] {literatura, vydavatelství}
\item[{m}] {rod mužský}
\item[{mat.}] {matematika}
\item[{med. }] {lékařství}
\item[{meteo.}] {meteorologie}
\item[{myt.}] {mytologie}
\item[{n}] {rod střední}
\item[{náb.}] {náboženství}
\item[{nám.}] {námořnictví, rybolov}
\item[{nom}] {1. pád, nominativ}
\item[{num}] {číslovka}
\item[{ord}] {řadová číslovka}
\item[{p}] {osoba}
\item[{part}] {částice}
\item[{pers}] {(v +) osoba, (pron +) zájmeno osobní}
\item[{pf}] {minulý čas}
\item[{pl}] {množné číslo, plurál}
\item[{poč.}] {informatika}
\item[{pol.}] {politika, politologie}
\item[{poss}] {zájmeno přivlastňovací}
\item[{pos}] {1. stupeň stupňování, pozitiv}
\item[{pov.}] {lidové pověsti, folkloristika}
\item[{pp}] {příčestí minulé}
\item[{praes}] {přítomný čas}
\item[{práv.}] {právnictví, soudnictví}
\item[{predp}] {předpona}
\item[{prep}] {předložka}
\item[{presp}] {příčestí přítomné}
\item[{pron}] {zájmeno}
\item[{prop}] {vlastní jméno, proprium}
\item[{přen.}] {přeneseně}
\item[{přís.}] {přísloví}
\item[{psych.}] {psychologie}
\item[{refl}] {(v +) médium, (pron +) zvratné zájmeno}
\item[{rel}] {vztažné zájmeno}
\item[{sg}] {jednotné číslo, singulár}
\item[{slang.}] {slang}
\item[{sport.}] {sport}
\item[{stav.}] {stavebnictví, architektura}
\item[{subs}] {podstatné jméno (bez rodu)}
\item[{sup}] {3. stupeň stupňování, superlativ}
\item[{supin}] {supinum}
\item[{škol.}] {školství}
\item[{techn.}] {technika, mechanika}
\item[{v}] {sloveso}
\item[{voj.}] {vojenství}
\item[{zast.}] {zastarale}
\item[{zkr}] {zkratka}
\item[{zool.}] {zoologie}

\item[{;}] {středník uprostřed definice odděluje věty; středník na konci definice odděluje významy}

\item[{(...)}] {v kulatých závorkách jsou uvedeny doplňkové informace}
\item[{[...]}] {v hranatých závorkách je uveden fonetický zápis výslovnosti}\footnote{Vysvětlení fonetických symbolů je uvedeno v kapitole Fonetika \ref{sec:phon_phonems} na straně \pageref{sec:phon_phonems}}
\item[{|}] {naznačuje místo, kde se deklinační koncovka připojuje ke slovnímu základu}
\item[{·}] {vedlejší dělení složeného slova}
\item[{··}] {hlavní dělení složeného slova}
\item[{+}] {plus}
\item[{/}] {lomítko označuje více možností}
\item[{\dicsymSee}] {šipka odkazuje na heslo ve slovníku}
\item[{\dicsymCompare}] {šipka odkazuje na heslo ve slovníku ke srovnání}
\item[{\dicsymExampleIS}] {uvozuje islandský příklad}
\item[{\dicsymExampleCS}] {uvozuje český překlad příkladu}
\item[{\dicsymIdiom}] {uvozuje slovní spojení}
\item[{\dicsymFrequent}] {heslo patří mezi frekventované výrazy}

\item[{\dicsymPhoto}] {fotografie, ilustrace}
\item[{\dicsymProverb}] {přísloví, rčení}

\end{description}
\end{multicols}

\chapter{Skrá yfir skammstafanir og tákn}
\begin{multicols}{2}
\begin{description}[leftmargin=!,labelwidth=1.5cm]
%[\breaklabel\setleftmargin{40pt}\setlabelstyle{\bfseries}
%\addtolength{\itemsep}{-0.5\baselineskip}]
\item[{abb}] {stýfður}
\item[{acc}] {þolfall}
\item[{adj}] {lýsingarorð}
\item[{adv}] {atviksorð}
\item[{akt}] {germynd}
\item[{anat.}] {líffærafræði}
\item[{angl.}] {enska}
\item[{astro.}] {geimvísindi, stjörnufræði}
\item[{ap.}] {e.þ.h.}
\item[{básn.}] {ljóðrænn}
\item[{biol.}] {líffræði}
\item[{bot.}] {grasafræði}

\item[{comp}] {miðstig}
\item[{con}] {viðtengingarháttur}
\item[{conj}] {samtenging}
\item[{dat}] {þágufall}
\item[{def}] {með greini}
\item[{dem}] {ábendingarfornafn}
\item[{dět.}] {barnamál}

\item[{e-að}] {(\textit{eitthvað})}
\item[{e-ð}] {(\textit{eitthvað})}
\item[{e-ir}] {(\textit{einhverjir})}
\item[{e-ja}] {(\textit{einhverja})}
\item[{e-jum}] {(\textit{einhverjum})}
\item[{e-m}] {(\textit{einhverjum})}
\item[{e-n}] {(\textit{einhvern})}
\item[{e-r}] { (\textit{einhver})}
\item[{e-rra}] {(\textit{einhverra})}
\item[{e-rs}] {(\textit{einhvers})}
\item[{e-s}] {(\textit{einhvers})}
\item[{e-u}] { (\textit{einhverju})}

\item[{ekon.}] {hagfræði, viðskipti}
\item[{elek.}] {rafmagn}
\item[{f}] {kvenkyn}
\item[{filos.}] {heimspeki, rökfræði}
\item[{form.}] {formlegt mál}
\item[{fyz.}] {eðlisfræði}
\item[{gen}] {eignarfall}
\item[{geog.}] {landafræði}
\item[{geol.}] {jarðfræði}
\item[{han.}] {niðrandi}
\item[{hist.}] {sagnfræði}
\item[{hovor.}] {talmál}
\item[{hrub.}] {gróft}
\item[{hud.}] {tónlist}
\item[{imper}] {boðháttur}
\item[{impers}] {ópersónuleg sögn}
\item[{ind}] {framsöguháttur}
\item[{indecl}] {óbeygjanlegur}
\item[{indef}] {óákveðinn, (pron +) óákveðið fornafn}
\item[{inf}] {nafnháttur}
\item[{int}] {spurnigarfornafn}
\item[{inter}] {upphrópun}
\item[{chem.}] {efnafræði}
\item[{jaz.}] {málfræði, málvísindi}
\item[{kulin.}] {eldhús}
\item[{l.}] {latína}
\item[{let.}] {flug, flugmál}
\item[{lit.}] {skáldskapur}
\item[{m}] {nafnorð, karlkyn}
\item[{mat.}] {stærðfræði}
\item[{med. }] {læknisfræði}
\item[{meteo.}] {veðurfræði}
\item[{myt.}] {goðafræði}
\item[{n}] {nafnorð, hvorugkyn}
\item[{náb.}] {trúfræði}
\item[{nám.}] {sjómennska}
\item[{nom}] {nefnifall}
\item[{num}] {töluorð}
\item[{ord}] {raðtala}
\item[{p}] {persóna}
\item[{part}] {ögn}
\item[{pers}] {(v +) persóna, (pron +) persónufornafn}
\item[{pf}] {þátíð}
\item[{pl}] {fleirtala}
\item[{poč.}] {tölvufræði, tölvur}
\item[{pol.}] {stjórnmál, stjórnmálafræði}
\item[{poss}] {eignarfornafn}
\item[{pos}] {frumstig}
\item[{pov.}] {þjóðtrú}
\item[{pp}] {lýsingarháttur þátíðar}
\item[{praes}] {nútíð}
\item[{práv.}] {lögfræði, dómsmál}
\item[{predp}] {forskeyti}
\item[{prep}] {forsetning}
\item[{presp}] {lýsingarháttur nútíðar}
\item[{pron}] {fornafn}
\item[{prop}] {sérnafn}
\item[{přen.}] {afleidd merking}
\item[{přís.}] {málsháttur}
\item[{psych.}] {sálfræði}
\item[{refl}] {(v +) miðmynd, (pron +) afturbeygt fornafn}
\item[{rel}] {tilvísunarfornafn}
\item[{sg}] {eintala}
\item[{slang.}] {slangur}
\item[{sport.}] {íþróttir}
\item[{stav.}] {húsasmíði, byggingarlist}
\item[{subs}] {nafnorð (án kyns)}
\item[{sup}] {efstastig}
\item[{supin}] {sagnbót}
\item[{škol.}] {skólamál}
\item[{techn.}] {tækni, verkfræði}
\item[{v}] {sagnorð}
\item[{voj.}] {her}
\item[{zast.}] {úrelt mál}
\item[{zkr}] {skammstöfun}
\item[{zool.}] {dýrafræði}

\item[{;}] {semíkomma í miðri skýringu aðgreinir setningu; semíkomma í enda skýringar aðgreinir mismunandi merkingarsvið}

\item[{(...)}] {innan sviga eru viðbótarupplýsingar}
\item[{[...]}] {innan hornklofa er hljóðritaður framburður}
\item[{|}] {merkir stað, þar sem beygingarending bætist við stofn}
\item[{·}] {aukasamskeyti samsetts orðs}
\item[{··}] {aðalsamskeyti samsetts orðs}
\item[{+}] {plús}
\item[{/}] {skástrik táknar fleiri möguleika}
\item[{\dicsymSee}] {ör vísar til uppflettiorðs í orðabókinni}
\item[{\dicsymCompare}] {ör vísar til uppflettiorðs í bókinni til samanburðar}
\item[{\dicsymExampleIS}] {kemur á undan dæmi á íslensku}
\item[{\dicsymExampleCS}] {kemur á undan tékkneskri þýðingu á dæminu}
\item[{\dicsymIdiom}] {orð með háa tíðni}

\item[{\dicsymPhoto}] {ljósmynd}
\item[{\dicsymProverb}] {málsháttur}

\end{description}
\end{multicols}

% ------------------------------------------------- Headword part and alphabet

\chapter{Heslová část}
{\renewcommand{\arraystretch}{2}%
%\topskip0pt
\vstretch{1}
{\setlength{\tabcolsep}{1.5em}
\begin{table}[htpb]
\centering

\begin{tabular}{ccccccc} \hline
A a &  & G g &  & O o &  & V v \\
Á á  &  & H h &  & Ó ó &  & W w \\
B b &  & I i &  & P p &  & X x \\
C c &  & Í í &  & Q q &  & Y y \\
D d &  & J j &  & R r &  & Ý ý \\
Ð ð &  & K k &  & S s &  & Z z \\
E e &  & L l &  & T t &  & Þ þ \\
É é  &  & M m &  & U u &  & Æ æ \\
F f &  & N n &  & Ú ú &  & Ö ö \\
\hline
\end{tabular}

%\caption{\textbf{Islandská abeceda}}
\label{my-islandska-abeceda}
\end{table}
}
\vstretch{1}
{\renewcommand{\arraystretch}{1}%

\clearpage

\fi

% =========================================================== THE DICTIONARY =

\ifinputletters

\cleardoublepage

\dictionarygeometry
\pagestyle{myheadings}

\newcommand*{\alphabetF}{%
  a,aa,b,c,d,e,ee,f,g,h,i,ii,j,k,l,m,n,o,oo,p,r,s1,s2,t,u,uu,v,w,x,y,yy,th,ae,oe}
\ifPDFtest
\dicLetter{r}{letter21}
\begin{xtolerant}{}{1pt}
\dicEntry[rabarbari] \dicTerm{rabar··bar|i} \dicIPA{{r}{a}{\textlengthmark}{\textsubring{b}}{a}{r}{\textsubring{b}}{a}{r}{\textsci}} \dicPos{m}[1] \dicFlx{(‑a, ‑ar)}[8] \dicFieldCat{bot.} \dicDirectTranslationCS{rebarbora} \textit{(l.~{\textLA{Rhe\-um x cultorum}})}
\end{xtolerant}
\dicEntry[rabb] \dicTerm{rabb} \dicIPA{{r}{a}{\textsubring{b}}{\textlengthmark}} \dicPos{n}[2] \dicFlx{(‑s)}[2] \dicSynonym{samtal} \dicDirectTranslationCS{povídání, klábosení, pokec} \dicExampleIS{rabb um alla heima og geima} \dicExampleCS{povídání o~všem možném}
\dicEntry[rabba] \dicTerm{rabb|a} \dicIPA{{r}{a}{\textsubring{b}}{\textlengthmark}{a}} \dicPos{v}[1] \dicFlx{(‑aði)}[50] \dicSynonym{spjalla} \dicDirectTranslationCS{povídat, klábosit} \dicExampleIS{rabba við e‑n um e‑ð} \dicExampleCS{klábosit s~(kým) o~(čem)}
\dicEntry[radar] \dicTerm{radar} \dicIPA{{r}{a}{\textlengthmark}{\textsubring{d}}{a}{\textsubring{r}}} \dicPos{m}[4] \dicFlx{(‑s, ‑ar)}[16] \dicFieldCat{techn.} \dicSynonym{ratsjá} \dicDirectTranslationCS{radar}
\dicEntry[radda] \dicTerm{radd|a} \dicIPA{{r}{a}{\textsubring{d}}{\textlengthmark}{a}} \dicPos{v}[1] \dicFlx{(‑aði)}[13] \dicFieldCat{jaz.} \dicDirectTranslationCS{vyslovit zněle}
\dicEntry[raddaður] \dicTerm{radd··aður} \dicIPA{{r}{a}{\textsubring{d}}{\textlengthmark}{a}{ð}{\textscy}{\textsubring{r}}} \dicPos{adj}[3] \dicFlx{(f rödduð)}[1] \dicFieldCat{jaz.} \dicDirectTranslationCS{znělý} \dicAntonym{óraddaður}
\dicEntry[raddar] \dicTerm{raddar} \dicIPA{{r}{a}{\textsubring{d}}{\textlengthmark}{a}{\textsubring{r}}} \dicPos{f} \dicFlx{sg gen} \dicLink{rödd}
\dicEntry[raddband] \dicTerm{radd··|band} \dicIPA{{r}{a}{\textsubring{d}}{\textsubring{b}}{a}{n}{\textsubring{d}}} \dicPos{n}[2] \dicFlx{(‑bands, ‑bönd)}[8] \dicPhraseIS{raddbönd} \dicFlx{pl} \dicFieldCat{anat.} \dicDirectTranslationCS{hlasivky}
\dicEntry[raddbandahljóð] \dicTerm{radd·banda··hljóð} \dicIPA{{r}{a}{\textsubring{d}}{\textsubring{b}}{a}{n}{\textsubring{d}}{a}{\textsubring{l}}{j}{ou}{\texttheta}} \dicPos{n}[2] \dicFlx{(‑s, ‑)}[5] \dicFieldCat{jaz.} \dicDirectTranslationCS{hlasivková souhláska, glotála}
\dicEntry[raddbandalokhljóð] \dicTerm{radd·banda··lok·hljóð} \dicIPA{{r}\-{a}\-{\textsubring{d}}\-{\textsubring{b}}\-{a}\-{n}\-{\textsubring{d}}\-{a}\-{l}\-{\textopeno}\-{\r{g}}\-{\textsubring{l}}\-{j}\-{ou}\-{\texttheta}\-} \dicPos{n}[2] \dicFlx{(‑s, ‑)}[5] \dicFieldCat{jaz.} \dicDirectTranslationCS{hlasivková ražená souhláska, glotální ploziva, ráz}
\dicEntry[raddbandamæltur] \dicTerm{radd·banda··mæltur} \dicIPA{{r}{a}{\textsubring{d}}{\textsubring{b}}{a}{n}{\textsubring{d}}{a}{m}{a}{i}{\textsubring{l}}{\textsubring{d}}{\textscy}{\textsubring{r}}} \dicPos{adj}[1]\dicFlx{}[-13] \dicFieldCat{jaz.} \dicDirectTranslationCS{hlasivkový, glotální}
\dicEntry[raddbandaönghljóð] \dicTerm{radd·banda··öng·hljóð} \dicIPA{{r}\-{a}\-{\textsubring{d}}\-{\textsubring{b}}\-{a}\-{n}\-{\textsubring{d}}\-{a}\-{\oe i}\-{\ng}\-{\r{g}}\-{\textsubring{l}}\-{j}\-{ou}\-{\texttheta}\-} \dicPos{n}[2] \dicFlx{(‑s, ‑)}[5] \dicFieldCat{jaz.} \dicDirectTranslationCS{hlasivková třená souhláska, glotální frikativa}
\dicEntry[raddir] \dicTerm{raddir} \dicIPA{{r}{a}{\textsubring{d}}{\textlengthmark}{\textsci}{\textsubring{r}}} \dicPos{f} \dicFlx{pl nom} \dicLink{rödd}
\dicEntry[raddlaus] \dicTerm{radd··laus} \dicIPA{{r}{a}{\textsubring{d}}{l}{\oe i}{s}} \dicPos{adj}[5]\dicFlx{}[-1] \textbf{1.} \dicDirectTranslationCS{(jsoucí) bez hlasu}  \textbf{2.} \dicFieldCat{jaz.} \dicSynonym{óraddaður} \dicDirectTranslationCS{neznělý}
\dicEntry[raddleysi] \dicTerm{radd··leysi} \dicIPA{{r}{a}{\textsubring{d}}{l}{ei}{s}{\textsci}} \dicPos{n}[2] \dicFlx{(‑s)}[20] \dicFieldCat{jaz.} \dicDirectTranslationCS{neznělost}
\dicEntry[radín] \dicTerm{radín} \dicIPA{{r}{a}{\textlengthmark}{\textsubring{d}}{i}{\textsubring{n}}} \dicPos{n}[2] \dicFlx{(‑s)}[2] \dicFieldCat{chem.} \dicDirectTranslationCS{radium} \textit{(l.~{\textLA{Ra, Radium}})}
\dicEntry[radísa] \dicTerm{radís|a} \dicIPA{{r}{a}{\textlengthmark}{\textsubring{d}}{i}{s}{a}} \dicPos{f}[1] \dicFlx{(‑u, ‑ur)}[7] \dicFieldCat{bot.} \dicSynonym{hreðka} \dicDirectTranslationCS{ředkvička, ředkev setá} \textit{(l.~{\textLA{Raphanus sativus}})}
\dicEntry[radíus] \dicTerm{radíus} \dicIPA{{r}{a}{\textlengthmark}{\textsubring{d}}{i}{j}{\textscy}{s}} \dicPos{m}[4] \dicFlx{(‑s, ‑ar)}[25] \dicFieldCat{mat.} \dicDirectTranslationCS{poloměr, rádius}
\dicEntry[radon] \dicTerm{radon} \dicIPA{{r}{a}{\textlengthmark}{\textsubring{d}}{\textopeno}{\textsubring{n}}} \dicPos{n}[2] \dicFlx{(‑s)}[2] \dicFieldCat{chem.} \dicDirectTranslationCS{radon} \textit{(l.~{\textLA{Rn, Radon}})}
\dicEntry[raða] \dicTerm{rað|a} \dicsymFrequent\  \dicIPA{{r}{a}{\textlengthmark}{ð}{a}} \dicPos{v}[1] \dicFlx{(‑aði)}[13] \dicFlx{dat} \textbf{1.} \dicSynonym*{skipa niður} \dicDirectTranslationCS{uspořádat, seřadit, utřídit} \dicExampleIS{raða e‑u} \dicExampleCS{uspořádat (co)}  \textbf{2.} \dicDirectTranslationCS{(se)řadit} \dicExampleIS{raða nöfnunum í stafrófsröð} \dicExampleCS{seřadit jména podle abecedy};  \dicIdiom{raða}[upp]{ \dicPhraseIS{raða sér upp}} \dicDirectTranslationCS{seřadit se (do řady)} \dicExampleIS{Nemendurnir röðuðu sér upp til myndatöku.} \dicExampleCS{Žáci se seřadili kvůli focení.};  \dicIdiom{raðast}{ \dicPhraseIS{raðast}} \dicFlx{refl} \dicDirectTranslationCS{seřadit se, uspořádat se}
\dicEntry[raðar] \dicTerm{raðar} \dicIPA{{r}{a}{\textlengthmark}{ð}{a}{\textsubring{r}}} \dicPos{f} \dicFlx{sg gen} \dicLink{röð}
\dicEntry[raðhús] \dicTerm{rað··hús} \dicIPA{{r}{a}{\textlengthmark}{\texttheta}{h}{u}{s}} \dicPos{n}[2] \dicFlx{(‑s, ‑)}[5] \dicDirectTranslationCS{řadový dům}
\dicEntry[raðir] \dicTerm{raðir} \dicIPA{{r}{a}{\textlengthmark}{ð}{\textsci}{\textsubring{r}}} \dicPos{f} \dicFlx{pl nom} \dicLink{röð}
\dicEntry[raðsmíða] \dicTerm{rað··smíð|a} \dicIPA{{r}{a}{ð}{s}{m}{i}{ð}{a}} \dicPos{v}[1] \dicFlx{(‑aði)}[1] \dicFlx{acc} \dicFieldCat{techn.} \dicDirectTranslationCS{(sériově) vyrábět\,/\addthin konstruovat} \dicExampleIS{raðsmíða skip} \dicExampleCS{sériově vyrábět lodě}
\dicEntry[raðsmíði] \dicTerm{rað··smíð|i} \dicIPA{{r}{a}{ð}{s}{m}{i}{ð}{\textsci}} \dicPos{f}[2] \dicFlx{(‑i, ‑ar)}[1] \dicDirectTranslationCS{sériová výroba} \dicExampleIS{raðsmíði skipa} \dicExampleCS{sériová výroba lodí}
\dicEntry[raðtala] \dicTerm{rað··|tala} \dicIPA{{r}{a}{ð}{t\smash{\textsuperscript{h}}}{a}{l}{a}} \dicPos{f}[1] \dicFlx{(‑tölu, ‑tölur)}[14] \dicFieldCat{jaz.} \dicDirectTranslationCS{řadová číslovka}
\dicEntry[raf] \dicTerm{raf\smash{\textsuperscript{1}}} \dicIPA{{r}{a}{\textlengthmark}{f}} \dicPos{n}[2] \dicFlx{(‑s, röf)}[8] \dicDirectTranslationCS{jantar}
\dicEntry[raf] \dicTerm{raf-\smash{\textsuperscript{2}}} \dicIPA{{r}{a}{\textlengthmark}{f}} \dicPos{predp} \dicDirectTranslationCS{elektrický, elektro-}
\dicEntry[rafali] \dicTerm{rafal|i}\dicTerm{, rafall} \dicIPA{{r}\-{a}\-{\textlengthmark}\-{v}\-{a}\-{l}\-{\textsci}\-} \dicPos{m}[1] \dicFlx{(‑a, ‑ar)}[11] \dicFieldCat{elek.} \dicDirectTranslationCS{generátor, dynamo}
\dicEntry[rafall] \dicTerm{rafal|l} \dicIPA{{r}{a}{\textlengthmark}{v}{a}{\textsubring{d}}{\textsubring{l}}} \dicPos{m}[6] \dicFlx{(‑s, ‑ar)}[48] \dicLink{rafali}
\dicEntry[rafbók] \dicTerm{raf··|bók} \dicIPA{{r}{a}{v}{\textsubring{b}}{ou}{\r{g}}} \dicPos{f}[8] \dicFlx{(‑bókar, ‑bækur)}[5] \dicDirectTranslationCS{elektronická kniha, e-kniha, e-book}
\dicEntry[rafefnafræði] \dicTerm{raf··efna·fræð|i} \dicIPA{{r}{a}{\textlengthmark}{v}{\textepsilon}{\textsubring{b}}{n}{a}{f}{r}{a}{i}{ð}{\textsci}} \dicPos{f}[3] \dicFlx{(‑i)}[3] \dicDirectTranslationCS{elektrochemie}
\dicEntry[rafeind] \dicTerm{raf··eind} \dicIPA{{r}{a}{\textlengthmark}{v}{ei}{n}{\textsubring{d}}} \dicPos{f}[7] \dicFlx{(‑ar, ‑ir)}[1] \dicFieldCat{fyz.} \dicDirectTranslationCS{elektron}
\dicEntry[rafeindafræði] \dicTerm{raf·einda··fræð|i} \dicIPA{{r}{a}{\textlengthmark}{v}{ei}{n}{\textsubring{d}}{a}{f}{r}{a}{i}{ð}{\textsci}} \dicPos{f}[3] \dicFlx{(‑i)}[3] \dicDirectTranslationCS{elektronika} \dicIndirectTranslationCS{(obor elektrotechniky)}
\dicEntry[rafeindasmásjá] \dicTerm{raf·einda··smá·sjá} \dicIPA{{r}{a}{\textlengthmark}{v}{ei}{n}{\textsubring{d}}{a}{s}{m}{au}{s}{j}{au}} \dicPos{f}[4] \dicFlx{(‑r, ‑r)}[18] \dicFieldCat{elek.} \dicDirectTranslationCS{elektronový mikroskop}
\dicEntry[raffræði] \dicTerm{raf··fræð|i} \dicIPA{{r}{a}{f}{r}{a}{i}{ð}{\textsci}} \dicPos{f}[3] \dicFlx{(‑i)}[3] \dicDirectTranslationCS{elektrotechnika}
\dicEntry[raffræðingur] \dicTerm{raf·fræð··ing|ur} \dicIPA{{r}{a}{f}{r}{a}{i}{ð}{i}{\ng}{\r{g}}{\textscy}{\textsubring{r}}} \dicPos{m}[6] \dicFlx{(‑s, ‑ar)}[8] \dicSynonym{rafmagnsverkfræðingur} \dicDirectTranslationCS{elektroinženýr(ka)}
\dicEntry[rafgeymir] \dicTerm{raf··geym|ir} \dicIPA{{r}{a}{v}{\r{\textObardotlessj}}{ei}{m}{\textsci}{\textsubring{r}}} \dicPos{m}[7] \dicFlx{(‑is, ‑ar)}[1] \dicFieldCat{elek.} \dicDirectTranslationCS{akumulátor, akumulační baterie}
\dicEntry[rafgreining] \dicTerm{raf··grein·ing} \dicIPA{{r}{a}{v}{\r{g}}{r}{ei}{n}{i}{\ng}{\r{g}}} \dicPos{f}[4] \dicFlx{(‑ar)}[7] \dicFieldCat{elek.} \dicDirectTranslationCS{elektrolýza}
\dicEntry[rafhlaða] \dicTerm{raf··|hlaða} \dicIPA{{r}{a}{f}{\textsubring{l}}{a}{ð}{a}} \dicPos{f}[1] \dicFlx{(‑hlöðu, ‑hlöður)}[20] \dicDirectTranslationCS{baterie, baterka} \dicIndirectTranslationCS{(suchý galvanický článek)}
\dicEntry[rafhleðsla] \dicTerm{raf··hleðsl|a} \dicIPA{{r}{a}{f}{\textsubring{l}}{\textepsilon}{ð}{s}{\textsubring{d}}{l}{a}} \dicPos{f}[1] \dicFlx{(‑u)}[5] \dicFieldCat{elek.} \dicDirectTranslationCS{elektrický náboj}
\dicEntry[rafiðnfræðingur] \dicTerm{raf·iðn·fræð··ing|ur} \dicIPA{{r}{a}{\textlengthmark}{v}{\textsci}{ð}{n}{f}{r}{a}{i}{ð}{i}{\ng}{\r{g}}{\textscy}{\textsubring{r}}} \dicPos{m}[6] \dicFlx{(‑s, ‑ar)}[8] \dicDirectTranslationCS{elektrotechnik, elektrotechnička}
\dicEntry[rafknúinn] \dicTerm{raf··knúinn} \dicIPA{{r}{a}{f}{k\smash{\textsuperscript{h}}}{n}{u}{\textsci}{\textsubring{n}}} \dicPos{adj}[6]\dicFlx{}[-6] \dicDirectTranslationCS{elektrický} \dicIndirectTranslationCS{(napájený, poháněný elektřinou)}
\dicEntry[raflagnar] \dicTerm{raf··lagnar} \dicIPA{{r}{a}{v}{l}{a}{\r{g}}{n}{a}{\textsubring{r}}} \dicPos{f} \dicFlx{sg gen} \dicLink{raflögn}
\dicEntry[rafleiðsla] \dicTerm{raf··leiðsl|a} \dicIPA{{r}{a}{v}{l}{ei}{ð}{s}{\textsubring{d}}{l}{a}} \dicPos{f}[1] \dicFlx{(‑u, ‑ur)}[19] \dicFieldCat{elek.} \dicDirectTranslationCS{elektrické vedení}
\dicEntry[raflesari] \dicTerm{raf··les·ar|i} \dicIPA{{r}{a}{v}{l}{\textepsilon}{s}{a}{r}{\textsci}} \dicPos{m}[1] \dicFlx{(‑a, ‑ar)}[13] \dicFieldCat{poč.} \dicDirectTranslationCS{čtečka (elektronických knih ap.), e‑reader}
\dicEntry[raflýsing] \dicTerm{raf··lýs·ing} \dicIPA{{r}{a}{v}{l}{i}{s}{i}{\ng}{\r{g}}} \dicPos{f}[4] \dicFlx{(‑ar)}[7] \dicDirectTranslationCS{elektrické osvětlení}
\dicEntry[raflögn] \dicTerm{raf··|lögn} \dicIPA{{r}{a}{v}{l}{\oe}{\r{g}}{\textsubring{n}}} \dicPos{f}[7] \dicFlx{(‑lagnar, ‑lagnir)}[16] \dicFieldCat{elek.} \dicDirectTranslationCS{elektroinstalace}
\dicEntry[rafmagn] \dicTerm{raf··magn\smash{\textsuperscript{1}}} \dicsymFrequent\  \dicIPA{{r}{a}{v}{m}{a}{\r{g}}{\textsubring{n}}{}{/}{r}{a}{\textlengthmark}{m}{a}{\r{g}}{\textsubring{n}}} \dicPos{n}[2] \dicFlx{(‑s)}[2] \dicFieldCat{fyz.} \dicDirectTranslationCS{elektřina} \dicExampleIS{leggja rafmagn í húsið} \dicExampleCS{zavést do domu elektřinu}
\dicEntry[rafmagn] \dicTerm{raf··magn-\smash{\textsuperscript{2}}} \dicIPA{{r}{a}{v}{m}{a}{\r{g}}{\textsubring{n}}{}{/}{r}{a}{\textlengthmark}{m}{a}{\r{g}}{\textsubring{n}}} \dicPos{predp} \dicDirectTranslationCS{elektrický, elek\-t\-ro-}
\dicEntry[rafmagnsleiðsla] \dicTerm{raf·magns··leiðsl|a} \dicIPA{{r}{a}{\textlengthmark}{m}{a}{x}{s}{l}{ei}{ð}{s}{\textsubring{d}}{l}{a}} \dicPos{f}[1] \dicFlx{(‑u, ‑ur)}[19] \dicFieldCat{elek.} \dicDirectTranslationCS{elektrické vedení}
\dicEntry[rafmagnsofn] \dicTerm{raf·magns··ofn} \dicIPA{{r}{a}{\textlengthmark}{m}{a}{x}{s}{\textopeno}{\textsubring{b}}{\textsubring{n}}} \dicPos{m}[4] \dicFlx{(‑s, ‑ar)}[3] \dicDirectTranslationCS{elektrické topení, elektrický ohřívač}
\dicEntry[rafmagnstæki] \dicTerm{raf·magns··tæki} \dicIPA{{r}{a}{\textlengthmark}{m}{a}{x}{s}{t\smash{\textsuperscript{h}}}{a}{i}{\r{\textObardotlessj}}{\textsci}} \dicPos{n}[2] \dicFlx{(‑s, ‑)}[16] \dicDirectTranslationCS{elektrický přístroj, elektrické zařízení}
\begin{xtolerant}{}{1pt}
\dicEntry[rafmagnsverkfræðingur] \textls[-8]{\dicTerm{raf·magns··verk·fræð·ing|ur}\addthinS\dicIPA{{r}\-{a}\-{\textlengthmark}\-{m}\-{a}\-{x}\-{s}\-{v}\-{\textepsilon}\-{\textsubring{r}}\-{\r{g}}\-{f}\-{r}\-{a}\-{i}\-{ð}\-{i}\-{\ng}\-{\r{g}}\-{\textscy}\-{\textsubring{r}}}} \dicPos{m}[6] \dicFlx{(‑s, ‑ar)}[8] \dicSynonym{raffræðingur} \dicDirectTranslationCS{elektroinženýr(ka)}
\end{xtolerant}
\dicEntry[raforka] \dicTerm{raf··ork|a} \dicIPA{{r}{a}{\textlengthmark}{v}{\textopeno}{\textsubring{r}}{\r{g}}{a}} \dicPos{f}[1] \dicFlx{(‑u)}[5] \dicFieldCat{fyz.} \dicSynonym{rafmagn\smash{\textsuperscript{1}}} \dicDirectTranslationCS{elektrická energie}
\dicEntry[raforkuver] \dicTerm{raf·orku··ver} \dicIPA{{r}{a}{\textlengthmark}{v}{\textopeno}{\textsubring{r}}{\r{g}}{\textscy}{v}{\textepsilon}{\textsubring{r}}} \dicPos{n}[2] \dicFlx{(‑s, ‑)}[5] \dicDirectTranslationCS{elektrárna}
\dicEntry[rafrænn] \dicTerm{raf··rænn} \dicIPA{{r}{a}{v}{r}{a}{i}{\textsubring{d}}{\textsubring{n}}} \dicPos{adj}[7]\dicFlx{}[-2] \dicDirectTranslationCS{elektronický} \dicExampleIS{rafræn útgáfa} \dicExampleCS{elektronické vydání}
\dicEntry[rafsegulfræði] \dicTerm{raf·segul··fræð|i} \dicIPA{{r}{a}{f}{s}{\textepsilon}{\textbabygamma}{\textscy}{l}{f}{r}{a}{i}{ð}{\textsci}} \dicPos{f}[3] \dicFlx{(‑i)}[3] \dicDirectTranslationCS{elektrodynamika}
\dicEntry[rafsegull] \dicTerm{raf··seg|ull} \dicIPA{{r}{a}{f}{s}{\textepsilon}{\textbabygamma}{\textscy}{\textsubring{d}}{\textsubring{l}}} \dicPos{m}[6] \dicFlx{(‑uls, ‑lar)}[31] \dicFieldCat{elek.} \dicDirectTranslationCS{elektromagnet}
\dicEntry[rafsegulsvið] \dicTerm{raf·segul··svið} \dicIPA{{r}{a}{f}{s}{\textepsilon}{\textbabygamma}{\textscy}{\textsubring{l}}{s}{v}{\textsci}{\texttheta}} \dicPos{n}[2] \dicFlx{(‑s)}[2] \dicFieldCat{fyz.} \dicDirectTranslationCS{elektromagnetické pole}
\dicEntry[rafskaut] \dicTerm{raf··skaut} \dicIPA{{r}{a}{f}{s}{\r{g}}{\oe i}{\textsubring{d}}} \dicPos{n}[2] \dicFlx{(‑s, ‑)}[5] \dicFieldCat{elek.} \dicDirectTranslationCS{elektroda}
\dicEntry[rafstraumur] \dicTerm{raf··straum|ur} \dicIPA{{r}{a}{f}{s}{\textsubring{d}}{r}{\oe i}{m}{\textscy}{\textsubring{r}}} \dicPos{m}[6] \dicFlx{(‑s, ‑ar)}[22] \dicFieldCat{elek.} \dicDirectTranslationCS{elektrický proud}
\dicEntry[rafstöð] \dicTerm{raf··stöð} \dicIPA{{r}{a}{f}{s}{\textsubring{d}}{\oe}{\texttheta}} \dicPos{f}[6] \dicFlx{(‑var, ‑var)}[1] \dicSynonym{orkuver} \dicDirectTranslationCS{elektrárna}
\dicEntry[rafsuða] \dicTerm{raf··suð|a} \dicIPA{{r}{a}{f}{s}{\textscy}{ð}{a}} \dicPos{f}[1] \dicFlx{(‑u)}[5] \dicFieldCat{techn.} \dicDirectTranslationCS{obloukové svařování, svařování elektrickým obloukem}
\dicEntry[raftæki] \dicTerm{raf··tæki} \dicIPA{{r}{a}{f}{t\smash{\textsuperscript{h}}}{a}{i}{\r{\textObardotlessj}}{\textsci}} \dicPos{n}[2] \dicFlx{(‑s, ‑)}[16] \dicDirectTranslationCS{elektrický přístroj\,/\addthin spotřebič}
\dicEntry[rafveita] \dicTerm{raf··veit|a} \dicIPA{{r}{a}{v}{\textlengthmark}{ei}{\textsubring{d}}{a}} \dicPos{f}[1] \dicFlx{(‑u, ‑ur)}[13] \dicDirectTranslationCS{elektrárenská společnost}
\dicEntry[rafvirki] \dicTerm{raf··virk|i} \dicIPA{{r}{a}{v}{\textlengthmark}{\textsci}{\textsubring{r}}{\r{\textObardotlessj}}{\textsci}} \dicPos{m}[1] \dicFlx{(‑ja, ‑jar)}[14] \dicDirectTranslationCS{elektrikář(ka)}
\dicEntry[rafvæðing] \dicTerm{raf··væð·ing} \dicIPA{{r}{a}{v}{\textlengthmark}{a}{i}{ð}{i}{\ng}{\r{g}}} \dicPos{f}[4] \dicFlx{(‑ar)}[7] \dicDirectTranslationCS{elektrifikace, zavádění elektřiny} \dicExampleIS{rafvæðing sveitanna} \dicExampleCS{elektrifikace venkova}
\dicEntry[raggeit] \dicTerm{rag··geit} \dicIPA{{r}{a}{\r{\textObardotlessj}}{\textlengthmark}{ei}{\textsubring{d}}} \dicPos{f}[8] \dicFlx{(‑ar, ‑ur)}[1] \dicSynonym{bleyða} \dicDirectTranslationCS{zbabělec, zbabělkyně, strašpytel}
\dicEntry[ragmennska] \dicTerm{rag··mennsk|a} \dicIPA{{r}{a}{\textbabygamma}{m}{\textepsilon}{n}{s}{\r{g}}{a}} \dicPos{f}[1] \dicFlx{(‑u)}[5] \dicSynonym{hugleysi} \dicDirectTranslationCS{zbabělost}
\dicEntry[ragna] \dicTerm{ragn|a} \dicIPA{{r}{a}{\r{g}}{n}{a}} \dicPos{v}[1] \dicFlx{(‑aði)}[13] \dicSynonym{bölva} \dicDirectTranslationCS{nadávat, klít}
\dicEntry[ragnarök] \dicTerm{ragna··rök} \dicIPA{{r}{a}{\r{g}}{n}{a}{r}{\oe}{\r{g}}} \dicPos{n}[2] \dicFlx{pl}[9] \dicFieldCat{myt.} \dicDirectTranslationCS{Ragnarök} \dicIndirectTranslationCS{(v~severské mytologii označení pro konec světa a~zánik bohů)}
\dicEntry[ragur] \dicTerm{ragur} \dicIPA{{r}{a}{\textlengthmark}{\textbabygamma}{\textscy}{\textsubring{r}}} \dicPos{adj}[1] \dicFlx{(f rög)}[2] \dicSynonym{huglaus} \dicDirectTranslationCS{zbabělý, bázlivý} \dicAntonym{óragur}
\dicEntry[ragú] \dicTerm{ragú} \dicIPA{{r}{a}{\textlengthmark}{\textbabygamma}{u}} \dicPos{n}[2] \dicFlx{(‑s)}[2] \dicFieldCat{kulin.} \dicDirectTranslationCS{ragú}
\dicEntry[rak] \dicTerm{rak\smash{\textsuperscript{1}}} \dicIPA{{r}{a}{\textlengthmark}{\r{g}}} \dicPos{n}[2] \dicFlx{(‑s, rök)}[8] \dicSynonym{kveikur} \dicDirectTranslationCS{knot}
\dicEntry[rak] \dicTerm{rak\smash{\textsuperscript{2}}} \dicIPA{{r}{a}{\textlengthmark}{\r{g}}} \dicPos{v} \dicFlx{ind pf sg 1 pers} \dicLink{reka\smash{\textsuperscript{2}}}
\dicEntry[raka] \dicTerm{rak|a} \dicIPA{{r}{a}{\textlengthmark}{\r{g}}{a}} \dicPos{v}[1] \dicFlx{(‑aði)}[13] \dicFlx{dat\,/\addthin acc} \textbf{1.} \dicFlx{dat} \dicSynonym{skara} \dicDirectTranslationCS{hrabat, shrabovat} \dicExampleIS{raka hey} \dicExampleCS{shrabovat seno}  \textbf{2.} \dicFlx{acc} \dicDirectTranslationCS{(o)holit} \dicExampleIS{raka skeggið af sér} \dicExampleCS{oholit si vousy};  \dicPhraseIS{raka sig} \dicDirectTranslationCS{(o)holit se};  \dicIdiom{raka}[saman]{ \dicPhraseIS{raka e‑u saman}} \dicSynonym{safna} \dicDirectTranslationCS{sbírat (co), shromáždit (co)} \dicExampleIS{raka saman fé} \dicExampleCS{shromažďovat ovce}
\dicEntry[rakalaus] \dicTerm{raka··laus} \dicIPA{{r}{a}{\textlengthmark}{\r{g}}{a}{l}{\oe i}{s}} \dicPos{adj}[5]\dicFlx{}[-1] \textbf{1.} \dicDirectTranslationCS{suchý, (jsoucí) bez vlhkosti} \dicExampleIS{rakalaus íbúð} \dicExampleCS{suchý byt}  \textbf{2.} \dicSynonym{órökstuddur} \dicDirectTranslationCS{nepodložený, nedoložený}
\dicEntry[rakarastofa] \dicTerm{rakara··stof|a} \dicIPA{{r}{a}{\textlengthmark}{\r{g}}{a}{r}{a}{s}{\textsubring{d}}{\textopeno}{v}{a}} \dicPos{f}[1] \dicFlx{(‑u, ‑ur)}[7] \dicDirectTranslationCS{holičství, pánské kadeřnictví}
\dicEntry[rakari] \dicTerm{rak··ar|i} \dicIPA{{r}{a}{\textlengthmark}{\r{g}}{a}{r}{\textsci}} \dicPos{m}[1] \dicFlx{(‑a, ‑ar)}[10] \dicSynonym{hárskeri} \dicDirectTranslationCS{holič(ka), kadeřník, kadeřnice}
\dicEntry[rakastig] \dicTerm{raka··stig} \dicIPA{{r}{a}{\textlengthmark}{\r{g}}{a}{s}{\textsubring{d}}{\textsci}{x}} \dicPos{n}[2] \dicFlx{(‑s, ‑)}[5] \dicFieldCat{meteo.} \dicDirectTranslationCS{vlhkost, humidita}
\dicEntry[rakblað] \textls[15]{\dicTerm{rak··|blað} \dicIPA{{r}{a}{\textlengthmark}{\r{g}}{\textsubring{b}}{l}{a}{\texttheta}} \dicPos{n}[2] \dicFlx{(‑blaðs, ‑blöð)}[8] \dicDirectTranslationCS{žiletka, holicí čepelka}}
\dicEntry[rakbursti] \dicTerm{rak··burst|i} \dicIPA{{r}{a}{\textlengthmark}{\r{g}}{\textsubring{b}}{\textscy}{\textsubring{r}}{s}{\textsubring{d}}{\textsci}} \dicPos{m}[1] \dicFlx{(‑a, ‑ar)}[1] \dicDirectTranslationCS{štětka na holení}
\dicEntry[raketta] \dicTerm{rakett|a} \dicIPA{{r}{a}{\textlengthmark}{\r{\textObardotlessj}}{\textepsilon}{h}{\textsubring{d}}{a}} \dicPos{f}[1] \dicFlx{(‑u, ‑ur)}[7] \textbf{1.} \dicSynonym{flugeldur} \dicDirectTranslationCS{raketa} \dicIndirectTranslationCS{(pyrotechnický náboj)}  \textbf{2.} \dicLangCat{zast.} \dicSynonym{eldflaug} \dicDirectTranslationCS{raketa (kosmická ap.)}
\dicEntry[rakhnífur] \dicTerm{rak··hníf|ur} \dicIPA{{r}{a}{\textlengthmark}{\r{g}}{\textsubring{n}}{i}{v}{\textscy}{\textsubring{r}}} \dicPos{m}[6] \dicFlx{(‑s, ‑ar)}[22] \dicDirectTranslationCS{břitva}
\dicEntry[raki] \dicTerm{rak|i} \dicIPA{{r}{a}{\textlengthmark}{\r{\textObardotlessj}}{\textsci}} \dicPos{m}[1] \dicFlx{(‑a)}[3] \dicDirectTranslationCS{vlhkost, vlhko} \dicExampleIS{raki í lofti} \dicExampleCS{vlhkost ve vzduchu}
\dicEntry[rakið] \dicTerm{rakið} \dicIPA{{r}{a}{\textlengthmark}{\r{\textObardotlessj}}{\textsci}{\texttheta}} \dicPos{v} \dicFlx{supin} \dicLink{rekja}
\dicEntry[rakka] \dicTerm{rakk|a} \dicIPA{{r}{a}{h}{\r{g}}{a}} \dicPos{v}[1] \dicFlx{(‑aði)}[13] \dicPhraseIS{rakka e‑n niður} \dicLangCat{hovor.} \dicSynonym{niða} \dicDirectTranslationCS{zkritizovat (koho), osočit (koho)}
\dicEntry[rakki] \dicTerm{rakk|i} \dicIPA{{r}{a}{h}{\r{\textObardotlessj}}{\textsci}} \dicPos{m}[1] \dicFlx{(‑a, ‑ar)}[8] \dicLangCat{han.} \dicSynonym{hundur} \dicDirectTranslationCS{psisko, čokl}
\dicEntry[rakna] \dicTerm{rakn|a} \dicIPA{{r}{a}{h}{\r{g}}{n}{a}} \dicPos{v}[1] \dicFlx{(‑aði)}[43] \dicSynonym{losna} \dicDirectTranslationCS{rozmotat se, rozplést se} \dicExampleIS{Hnúturinn raknaði.} \dicExampleCS{Uzel se rozpletl.};  \dicIdiom{rakna}[af]{ \dicPhraseIS{láta e‑ð af hendi rakna}} \dicDirectTranslationCS{dát (co), přispět (čím) (na charitu ap.)};  \dicIdiom{rakna}[upp]{ \dicPhraseIS{rakna upp}} \dicDirectTranslationCS{(roz)párat se} \dicExampleIS{Saumurinn er að rakna upp.} \dicExampleCS{Šev se páře.};  \dicIdiom{rakna}[við\,/\addthin úr]{ \dicPhraseIS{rakna við, rakna úr roti}} \dicDirectTranslationCS{přijít k~sobě, probrat se (z~bezvědomí)}
\dicEntry[raksápa] \dicTerm{rak··sáp|a} \dicIPA{{r}{a}{\textlengthmark}{\r{g}}{s}{au}{\textsubring{b}}{a}} \dicPos{f}[1] \dicFlx{(‑u, ‑ur)}[19] \dicDirectTranslationCS{krém na holení}
\dicEntry[rakspíri] \dicTerm{rak··spír|i} \dicIPA{{r}{a}{\textlengthmark}{\r{g}}{s}{\textsubring{b}}{i}{r}{\textsci}} \dicPos{m}[1] \dicFlx{(‑a)}[3] \dicDirectTranslationCS{voda po holení}
\dicEntry[rakstrarvél] \dicTerm{rakstrar··vél} \dicIPA{{r}{a}{x}{s}{\textsubring{d}}{r}{a}{r}{v}{j}{\textepsilon}{\textsubring{l}}} \dicPos{f}[4] \dicFlx{(‑ar, ‑ar)}[1] \dicDirectTranslationCS{hrabačka, stroj na hrabání}
\dicEntry[rakstraumur] \dicTerm{rak··straum|ur} \dicIPA{{r}{a}{\textlengthmark}{\r{g}}{s}{\textsubring{d}}{r}{\oe i}{m}{\textscy}{\textsubring{r}}} \dicPos{m}[6] \dicFlx{(‑s, ‑ar)}[22] \dicFieldCat{elek.} \dicDirectTranslationCS{stejnosměrný proud}
\dicEntry[rakstur] \dicTerm{rakst|ur} \dicIPA{{r}{a}{x}{s}{\textsubring{d}}{\textscy}{\textsubring{r}}} \dicPos{m}[5] \dicFlx{(‑urs\,/\addthin ‑rar)}[10] \textbf{1.} \dicSynonym*{það að raka skegg} \dicDirectTranslationCS{(o)holení}  \textbf{2.} \dicSynonym*{það að raka hey} \dicDirectTranslationCS{(s)hrabání, shrabování (sena ap.)}
\dicEntry[rakti] \dicTerm{rakti} \dicIPA{{r}{a}{x}{\textsubring{d}}{\textsci}} \dicPos{v} \dicFlx{ind pf sg 1 pers} \dicLink{rekja}
\dicEntry[rakur] \dicTerm{rakur} \dicIPA{{r}{a}{\textlengthmark}{\r{g}}{\textscy}{\textsubring{r}}} \dicPos{adj}[1] \dicFlx{(f rök)}[2] \textbf{1.} \dicSynonym{deigur} \dicDirectTranslationCS{vlhký, provlhlý, humidní} \dicExampleIS{rakt loft} \dicExampleCS{vlhký vzduch}  \textbf{2.} \dicSynonym{hreifur} \dicDirectTranslationCS{mírně podnapilý}
\dicEntry[rakvél] \dicTerm{rak··vél} \dicIPA{{r}{a}{\textlengthmark}{\r{g}}{v}{j}{\textepsilon}{\textsubring{l}}} \dicPos{f}[4] \dicFlx{(‑ar, ‑ar)}[1] \dicDirectTranslationCS{holicí strojek}
\dicEntry[rall] \dicTerm{rall} \dicIPA{{r}{a}{l}{\textlengthmark}} \dicPos{n}[2] \dicFlx{(‑s, röll)}[8] \textbf{1.} \dicSynonym{svall\smash{\textsuperscript{1}}} \dicDirectTranslationCS{pitka, flám} \dicExampleIS{vera á ralli} \dicExampleCS{flámovat}  \textbf{2.} \dicSynonym*{bílakeppni} \dicDirectTranslationCS{rallye} \dicIndirectTranslationCS{(automobilový závod)}
\dicEntry[ramba] \dicTerm{ramb|a} \dicIPA{{r}{a}{m}{\textsubring{b}}{a}} \dicPos{v}[1] \dicFlx{(‑aði)}[50] \dicSynonym{reika} \dicDirectTranslationCS{bloumat, chodit bezcílně};  \dicIdiom{ramba}[á]{ \dicPhraseIS{ramba á e‑ð}} \dicDirectTranslationCS{narazit na (co), najít (co) náhodou}
\dicEntry[ramma] \dicTerm{ramm|a} \dicIPA{{r}{a}{m}{\textlengthmark}{a}} \dicPos{v}[1] \dicFlx{(‑aði)}[13] \dicFlx{acc} \dicDirectTranslationCS{(o)rámovat};  \dicPhraseIS{ramma e‑ð inn} \dicDirectTranslationCS{zarámovat (co)}
\dicEntry[rammasamningur] \dicTerm{ramma··samn·ing|ur} \dicIPA{{r}{a}{m}{\textlengthmark}{a}{s}{a}{m}{n}{i}{\ng}{\r{g}}{\textscy}{\textsubring{r}}} \dicPos{m}[6] \dicFlx{(‑s, ‑ar)}[8] \dicDirectTranslationCS{rámcová dohoda}
\dicEntry[rammbyggður] \dicTerm{ramm··byggður} \dicIPA{{r}{a}{m}{\textsubring{b}}{\textsci}{\textbabygamma}{ð}{\textscy}{\textsubring{r}}} \dicPos{adj}[2]\dicFlx{}[-1] \dicSynonym{traustur} \dicDirectTranslationCS{pevný, solidní, bytelný}
\dicEntry[rammi] \dicTerm{ramm|i} \dicIPA{{r}{a}{m}{\textlengthmark}{\textsci}} \dicPos{m}[1] \dicFlx{(‑a, ‑ar)}[8] \textbf{1.} \dicSynonym{umgerð} \dicDirectTranslationCS{rám} \dicExampleIS{setja mynd í ramma} \dicExampleCS{zarámovat obraz}  \textbf{2.} \dicDirectTranslationCS{rámec, sféra} \dicExampleIS{rammi skáldsögunnar} \dicExampleCS{rámec příběhu}
\dicEntry[rammlegur] \dicTerm{ramm··legur} \dicIPA{{r}{a}{m}{l}{\textepsilon}{\textbabygamma}{\textscy}{\textsubring{r}}} \dicPos{adj}[1]\dicFlx{}[-8] \dicSynonym{sterklegur} \dicDirectTranslationCS{masivní, robustní, solidní}
\dicEntry[rammleikur] \dicTerm{ramm··leik|ur} \dicIPA{{r}{a}{m}{l}{ei}{\r{g}}{\textscy}{\textsubring{r}}} \dicPos{m}[6] \dicFlx{(‑s)}[17] \dicPhraseIS{(gera e‑ð) af eigin rammleik} \dicDirectTranslationCS{(udělat (co)) bez cizí pomoci}
\dicEntry[rammur] \dicTerm{rammur} \dicIPA{{r}{a}{m}{\textlengthmark}{\textscy}{\textsubring{r}}} \dicPos{adj}[1] \dicFlx{(f römm)}[2] \textbf{1.} \dicSynonym{sterkur} \dicDirectTranslationCS{silný, mocný};  \dicPhraseIS{eiga við ramman reip að draga} \dicLangCat{přen.} \dicDirectTranslationCS{ocitnout se v~těžké situaci};  \dicPhraseIS{vera rammur að afli} \dicDirectTranslationCS{mít velkou sílu (člověk ap.)}  \textbf{2.} \dicSynonym{beiskur} \dicDirectTranslationCS{hořký} \dicIndirectTranslationCS{(chuťově)} \dicExampleIS{rammur bjór} \dicExampleCS{hořké pivo}
\dicEntry[ranabjalla] \dicTerm{rana··|bjalla} \dicIPA{{r}{a}{\textlengthmark}{n}{a}{\textsubring{b}}{j}{a}{\textsubring{d}}{l}{a}} \dicPos{f}[1] \dicFlx{(‑bjöllu, ‑bjöllur)}[20] \dicPhraseIS{ranabjöllur} \dicFlx{pl} \dicFieldCat{zool.} \dicDirectTranslationCS{nosatcovití} \textit{(l.~{\textLA{Curculionidae}})}
\dicEntry[randar] \dicTerm{randar} \dicIPA{{r}{a}{n}{\textsubring{d}}{a}{\textsubring{r}}} \dicPos{f} \dicFlx{sg gen} \dicLink{rönd}
\dicEntry[ranga] \dicTerm{ranga} \dicIPA{{r}{au}{\ng}{\r{g}}{a}} \dicPos{f}[1] \dicFlx{(röngu)}[2] \dicDirectTranslationCS{opačná strana (oděvu ap.), rub} \dicAntonym{rétta\smash{\textsuperscript{3}}}
\dicEntry[rangali] \dicTerm{rang··al|i} \dicIPA{{r}{au}{\ng}{\r{g}}{a}{l}{\textsci}} \dicPos{m}[1] \dicFlx{(‑a, ‑ar)}[8] \dicDirectTranslationCS{dlouhý úzký (točitý) průchod\,/\addthin průjezd}
\dicEntry[rangeygður] \dicTerm{rang··eygður} \dicIPA{{r}{au}{\ng}{\r{g}}{ei}{\textbabygamma}{ð}{\textscy}{\textsubring{r}}} \dicPos{adj}[2]\dicFlx{}[-1] \dicLink{rangeygur}
\dicEntry[rangeygur] \dicTerm{rang··eygur}\dicTerm{, rangeygður} \dicIPA{{r}\-{au}\-{\ng}\-{\r{g}}\-{ei}\-{\textbabygamma}\-{\textscy}\-{\textsubring{r}}\-} \dicPos{adj}[1]\dicFlx{}[-1] \dicDirectTranslationCS{šilhavý, šilhající}
\dicEntry[rangfæra] \dicTerm{rang··fær|a} \dicIPA{{r}{au}{\ng}{\r{g}}{f}{a}{i}{r}{a}} \dicPos{v}[2] \dicFlx{(‑ði, ‑t)}[99] \dicFlx{acc} \dicDirectTranslationCS{překroutit, překrucovat, (z)komolit, zkreslit, zkreslovat} \dicExampleIS{rangfæra ummæli e‑rs} \dicExampleCS{překroutit (čí) námitku}
\dicEntry[rangfærsla] \dicTerm{rang··færsl|a} \dicIPA{{r}{au}{\ng}{\r{g}}{f}{a}{i}{\textsubring{r}}{s}{\textsubring{d}}{l}{a}} \dicPos{f}[1] \dicFlx{(‑u, ‑ur)}[13] \dicDirectTranslationCS{překroucení, zkomolení, zkreslení}
\dicEntry[ranghvolfa] \dicTerm{rang··hvolf|a} \dicIPA{{r}{au}{\ng}{\r{g}}{k\smash{\textsuperscript{h}}}{v}{\textopeno}{l}{v}{a}} \dicPos{v}[2] \dicFlx{(‑di, ‑t)}[154] \dicFlx{dat} \dicPhraseIS{ranghvolfa augunum} \dicDirectTranslationCS{obrátit oči v~sloup}
\dicEntry[rangindi] \dicTerm{rang··indi} \dicIPA{{r}{au}{\textltailn}{\r{\textObardotlessj}}{\textsci}{n}{\textsubring{d}}{\textsci}} \dicPos{n}[2] \dicFlx{pl}[19] \dicSynonym{ranglæti} \dicDirectTranslationCS{příkoří, nespravedlnost, křivda};  \dicPhraseIS{beita e‑n rangindum} \dicDirectTranslationCS{ukřivdit (komu)}
\dicEntry[ranglátur] \dicTerm{rang··látur} \dicIPA{{r}{au}{\ng}{\r{g}}{l}{au}{\textsubring{d}}{\textscy}{\textsubring{r}}} \dicPos{adj}[1]\dicFlx{}[-1] \dicSynonym{óréttlátur} \dicDirectTranslationCS{nespravedlivý, křivdící} \dicExampleIS{ranglátur dómur} \dicExampleCS{nespravedlivý rozsudek}
\dicEntry[ranglæti] \dicTerm{rang··læti} \dicIPA{{r}{au}{\ng}{\r{g}}{l}{a}{i}{\textsubring{d}}{\textsci}} \dicPos{n}[2] \dicFlx{(‑s)}[20] \dicSynonym{rangindi} \dicDirectTranslationCS{příkoří, bezpráví, nespravedlnost, křivda} \dicExampleIS{þola ranglæti af e‑m} \dicExampleCS{snášet od (koho) příkoří}
\dicEntry[rangsleitinn] \dicTerm{rangs··leitinn} \dicIPA{{r}{au}{\ng}{\r{g}}{s}{l}{ei}{\textsubring{d}}{\textsci}{\textsubring{n}}} \dicPos{adj}[6]\dicFlx{}[-2] \dicSynonym{ranglátur} \dicDirectTranslationCS{nespravedlivý, předpojatý, (jsoucí) nefér}
\dicEntry[rangsleitni] \dicTerm{rangs··leitn|i} \dicIPA{{r}{au}{\ng}{\r{g}}{s}{l}{ei}{h}{\textsubring{d}}{n}{\textsci}} \dicPos{f}[3] \dicFlx{(‑i)}[3] \dicSynonym{rangindi} \dicDirectTranslationCS{nespravedlivost, neférovost}
\dicEntry[rangstaða] \dicTerm{rang··|staða} \dicIPA{{r}{au}{\ng}{\r{g}}{s}{\textsubring{d}}{a}{ð}{a}} \dicPos{f}[1] \dicFlx{(‑stöðu)}[2] \dicFieldCat{sport.} \dicDirectTranslationCS{postavení mimo hru, ofsajd}
\dicEntry[rangstæður] \dicTerm{rang··stæður} \dicIPA{{r}{au}{\ng}{\r{g}}{s}{\textsubring{d}}{a}{i}{ð}{\textscy}{\textsubring{r}}} \dicPos{adj}[2]\dicFlx{}[-6] \dicFieldCat{sport.} \dicDirectTranslationCS{(o~hráči) (jsoucí) v~postavení mimo hru, (jsoucí) v~ofsajdu}
\dicEntry[rangtúlka] \dicTerm{rang··túlk|a} \dicIPA{{r}{au}{\ng}{\r{g}}{t\smash{\textsuperscript{h}}}{u}{\textsubring{l}}{\r{g}}{a}} \dicPos{v}[1] \dicFlx{(‑aði)}[1] \dicFlx{acc} \dicDirectTranslationCS{nesprávně\,/\addthin myl\-ně si vyložit\,/\addthin vykládat} \dicExampleIS{rangtúlka orð hans} \dicExampleCS{nesprávně si vykládat jeho slova}
\dicEntry[rangtúlkun] \dicTerm{rang··túlk|un} \dicIPA{{r}{au}{\ng}{\r{g}}{t\smash{\textsuperscript{h}}}{u}{\textsubring{l}}{\r{g}}{\textscy}{\textsubring{n}}} \dicPos{f}[7] \dicFlx{(‑unar, ‑anir)}[8] \dicDirectTranslationCS{nesprávná\,/\addthin mylná interpretace, překroucení (významu ap.)}
\dicEntry[rangur] \dicTerm{rangur} \dicsymFrequent\  \dicIPA{{r}{au}{\ng}{\r{g}}{\textscy}{\textsubring{r}}} \dicPos{adj}[1] \dicFlx{(f röng)}[2] \textbf{1.} \dicSynonym{ósannur} \dicDirectTranslationCS{špatný, chybný, mylný} \dicExampleIS{rangt svar} \dicExampleCS{špatná odpověď};  \dicPhraseIS{hafa á röngu að standa, hafa rangt fyrir sér} \dicSynonym*{fara ekki með rétt mál} \dicDirectTranslationCS{(z)mýlit se, (s)plést se}  \textbf{2.} \dicSynonym{óréttur} \dicDirectTranslationCS{špatný, nesprávný} \dicExampleIS{röng hegðun} \dicExampleCS{špatné chování};  \dicPhraseIS{gera e‑m rangt til} \dicDirectTranslationCS{(u)křivdit (komu)}
\dicEntry[rani] \dicTerm{ran|i} \dicIPA{{r}{a}{\textlengthmark}{n}{\textsci}} \dicPos{m}[1] \dicFlx{(‑a, ‑ar)}[8] \textbf{1.} \dicSynonym*{fílsrani} \dicDirectTranslationCS{chobot}  \textbf{2.} \dicSynonym*{fjallsnöf} \dicDirectTranslationCS{výběžek (hory ap.)}  \textbf{3.} \dicDirectTranslationCS{rampa} \dicIndirectTranslationCS{(šikmá komunikace pro pěší nebo vozidla spojující různé výškové úrovně)}  \textbf{4.} \dicLangCat{slang.} \dicDirectTranslationCS{čumák, nos};  \dicPhraseIS{fá sér í ranann} \dicLangCat{slang.} \dicDirectTranslationCS{šňupnout si (kokain ap.)}
\dicEntry[ranka] \dicTerm{rank|a} \dicIPA{{r}{au}{\r{\ng}}{\r{g}}{a}} \dicPos{v}[1] \dicFlx{(‑aði)}[50] \dicPhraseIS{ranka við sér} \dicSynonym*{átta sig} \dicDirectTranslationCS{probrat se, vzpamatovat se}
\dicEntry[rann] \dicTerm{rann} \dicIPA{{r}{a}{n}{\textlengthmark}} \dicPos{v} \dicFlx{ind pf sg 1 pers} \dicLink{renna\smash{\textsuperscript{3}}}
\dicEntry[rannsaka] \dicTerm{rann··sak|a} \dicsymFrequent\  \dicIPA{{r}{a}{n}{s}{a}{\r{g}}{a}} \dicPos{v}[1] \dicFlx{(‑aði)}[13] \dicFlx{acc} \textbf{1.} \dicSynonym{kanna\smash{\textsuperscript{2}}} \dicDirectTranslationCS{vyšetřit, vyšetřovat, (pro)zkoumat, (pro)bádat (záležitost ap.)} \dicExampleIS{rannsaka málið ýtarlega} \dicExampleCS{vyšetřovat podrobně záležitost}  \textbf{2.} \dicDirectTranslationCS{vyšetřit, vyšetřovat (pacienta ap.)} \dicExampleIS{rannsaka sjúklinginn} \dicExampleCS{vyšetřit pacienta}
\dicEntry[rannsakandi] \dicTerm{rann··sak·andi} \dicsymFrequent\  \dicIPA{{r}{a}{n}{s}{a}{\r{g}}{a}{n}{\textsubring{d}}{\textsci}} \dicPos{adj}[13] \dicFlx{indecl}[1] \dicDirectTranslationCS{zkoumavý, zkoumající, vyšetřující} \dicExampleIS{horfa rannsakandi á e‑n} \dicExampleCS{dívat se na (koho) zkoumavě}
\dicEntry[rannsókn] \dicTerm{rann··sókn} \dicsymFrequent\  \dicIPA{{r}{a}{n}{s}{ou}{h}{\r{g}}{\textsubring{n}}} \dicPos{f}[7] \dicFlx{(‑ar, ‑ir)}[1] \textbf{1.} \dicSynonym*{það að rannsaka} \dicDirectTranslationCS{vyšetřování, šetření}  \textbf{2.} \dicSynonym*{fræðileg athugun} \dicDirectTranslationCS{výzkum, zkoumání} \dicExampleIS{rannsóknir í læknisfræði} \dicExampleCS{výzkumy v~medicíně}
\dicEntry[rannsóknaraðferð] \dicTerm{rann·sóknar··að·ferð} \dicIPA{{r}{a}{n}{s}{ou}{h}{\r{g}}{n}{a}{r}{a}{ð}{f}{\textepsilon}{r}{\texttheta}} \dicPos{f}[7] \dicFlx{(‑ar, ‑ir)}[1] \textbf{1.} \dicDirectTranslationCS{výzkumná metoda, metoda zkoumání}  \textbf{2.} \dicDirectTranslationCS{vyšetřovací postup\,/\addthin metoda}
\dicEntry[rannsóknardeild] \dicTerm{rann·sóknar··deild} \dicIPA{{r}{a}{n}{s}{ou}{h}{\r{g}}{n}{a}{r}{\textsubring{d}}{ei}{l}{\textsubring{d}}} \dicPos{f}[7] \dicFlx{(‑ar, ‑ir)}[1] \textbf{1.} \dicDirectTranslationCS{výzkumné oddělení}  \textbf{2.} \dicDirectTranslationCS{vyšetřovací oddělení}
\dicEntry[rannsóknardómari] \textls[15]{\dicTerm{rann·sóknar··dóm·ar|i} \dicIPA{{r}\-{a}\-{n}\-{s}\-{ou}\-{h}\-{\r{g}}\-{n}\-{a}\-{r}\-{\textsubring{d}}\-{ou}\-{m}\-{a}\-{r}\-{\textsci}\-} \dicPos{m}[1] \dicFlx{(‑a, ‑ar)}[13]} \dicDirectTranslationCS{vyšetřující soudce\,/\addthin soudkyně}
\dicEntry[rannsóknarleiðangur] \dicTerm{rann·sóknar··leið·ang|ur} \dicIPA{{r}\-{a}\-{n}\-{s}\-{ou}\-{h}\-{\r{g}}\-{n}\-{a}\-{r}\-{l}\-{ei}\-{ð}\-{au}\-{\ng}\-{\r{g}}\-{\textscy}\-{\textsubring{r}}\-} \dicPos{m}[5] \dicFlx{(‑urs, ‑rar)}[3] \dicDirectTranslationCS{výzkumná expedice\,/\addthin výprava}
\dicEntry[rannsóknarlögregla] \dicTerm{rann·sóknar··lög·regl|a} \dicIPA{{r}\-{a}\-{n}\-{s}\-{ou}\-{h}\-{\r{g}}\-{n}\-{a}\-{r}\-{l}\-{\oe}\-{\textbabygamma}\-{r}\-{\textepsilon}\-{\r{g}}\-{l}\-{a}\-} \dicPos{f}[1] \dicFlx{(‑u, ‑ur)}[19] \dicDirectTranslationCS{kriminální policie}
\begin{xtolerant}{}{1pt}
\dicEntry[rannsóknarlögreglumaður] \dicTerm{rann·sóknar··lög·reglu·|maður} \dicIPA{{r}\-{a}\-{n}\-{s}\-{ou}\-{h}\-{\r{g}}\-{n}\-{a}\-{r}\-{l}\-{\oe}\-{\textbabygamma}\-{r}\-{\textepsilon}\-{\r{g}}\-{l}\-{\textscy}\-{m}\-{a}\-{ð}\-{\textscy}\-{\textsubring{r}}\-} \dicPos{m}[13] \dicFlx{(‑manns, ‑menn)}[2] \dicDirectTranslationCS{detektiv, komisař(ka) kriminální policie}
\end{xtolerant}
\dicEntry[rannsóknarmaður] \dicTerm{rann·sóknar··|maður} \dicIPA{{r}{a}{n}{s}{ou}{h}{\r{g}}{n}{a}{r}{m}{a}{ð}{\textscy}{\textsubring{r}}} \dicPos{m}[13] \dicFlx{(‑manns, ‑menn)}[2] \dicDirectTranslationCS{výzkumník, výzkumnice, výzkumný pracovník, výzkumná pracovnice}
\dicEntry[rannsóknarráð] \dicTerm{rann·sóknar··ráð} \dicIPA{{r}{a}{n}{s}{ou}{h}{\r{g}}{n}{a}{r}{au}{\texttheta}} \dicPos{n}[2] \dicFlx{(‑s, ‑)}[5] \dicDirectTranslationCS{výzkumná rada\,/\addthin komise}
\dicEntry[rannsóknarstofa] \dicTerm{rann·sóknar··stof|a}\textls[-20]{\dicTerm{, rannsóknastofa} \dicIPA{{r}\-{a}\-{n}\-{s}\-{ou}\-{h}\-{\r{g}}\-{n}\-{a}\-{\textsubring{r}}\-{s}\-{\textsubring{d}}\-{\textopeno}\-{v}\-{a}}} \dicPos{f}[1] \dicFlx{(‑u, ‑ur)}[7] \dicDirectTranslationCS{laboratoř} \dicExampleIS{Hún er ungur efnafræðingur og vinnur á rannsóknarstofu.} \dicExampleCS{Ona je mladá chemička a~pracuje v~laboratoři.}
\dicEntry[rannsóknastofa] \dicTerm{rann·sókna··stof|a} \dicIPA{{r}{a}{n}{s}{ou}{h}{\r{g}}{n}{a}{s}{\textsubring{d}}{\textopeno}{v}{a}} \dicPos{f}[1] \dicFlx{(‑u, ‑ur)}[7] \dicLink{rannsóknarstofa}
\dicEntry[rannsóknastofnun] \dicTerm{rann·sókna··stofn|un} \dicIPA{{r}{a}{n}{s}{ou}{h}{\r{g}}{n}{a}{s}{\textsubring{d}}{\textopeno}{\textsubring{b}}{n}{\textscy}{\textsubring{n}}} \dicPos{f}[7] \dicFlx{(‑unar, ‑an\-ir)}[8] \dicDirectTranslationCS{vědecký institut, výzkumný ústav}
\dicEntry[rappa] \dicTerm{rapp|a} \dicIPA{{r}{a}{h}{\textsubring{b}}{a}} \dicPos{v}[1] \dicFlx{(‑aði)}[13] \dicFieldCat{hud.} \dicDirectTranslationCS{rapovat}
\dicEntry[rasa] \dicTerm{ras|a} \dicIPA{{r}{a}{\textlengthmark}{s}{a}} \dicPos{v}[1] \dicFlx{(‑aði)}[13] \textbf{1.} \dicSynonym{hrasa} \dicDirectTranslationCS{klopýtnout, klopýtat}  \textbf{2.} \dicSynonym*{fara hratt} \dicDirectTranslationCS{pospíchat, chvátat};  \dicIdiom{rasa}[um]{ \dicPhraseIS{rasa um ráð fram}} \dicSynonym{flana} \dicDirectTranslationCS{unáhlit se, uspíšit rozhodnutí};  \dicIdiom{rasa}[út]{ \dicPhraseIS{rasa út}} \dicLangCat{hovor.} \dicDirectTranslationCS{vyřádit se, vydovádět se, vyblbnout se}
\dicEntry[rasismi] \dicTerm{ras··ism|i} \dicIPA{{r}{a}{\textlengthmark}{s}{\textsci}{s}{m}{\textsci}} \dicPos{m}[1] \dicFlx{(‑a)}[3] \dicSynonym{kynþáttahatur} \dicDirectTranslationCS{rasismus}
\dicEntry[rasisti] \dicTerm{ras··ist|i} \dicIPA{{r}{a}{\textlengthmark}{s}{\textsci}{s}{\textsubring{d}}{\textsci}} \dicPos{m}[1] \dicFlx{(‑a, ‑ar)}[1] \dicDirectTranslationCS{rasista, rasistka}
\dicEntry[rask] \dicTerm{rask} \dicIPA{{r}{a}{s}{\r{g}}} \dicPos{n}[2] \dicFlx{(‑s)}[2] \textbf{1.} \dicSynonym{umrót} \dicDirectTranslationCS{narušení, porušení (půdy ap.)}  \textbf{2.} \dicSynonym{ónæði} \dicDirectTranslationCS{rozrušení, vyrušení}
\dicEntry[raska] \dicTerm{rask|a} \dicIPA{{r}{a}{s}{\r{g}}{a}} \dicPos{v}[1] \dicFlx{(‑aði)}[13] \dicFlx{dat} \textbf{1.} \dicSynonym{trufla} \dicDirectTranslationCS{(vy)rušit, vytrhnout (z~práce ap.)} \dicExampleIS{raska ró e‑rs} \dicExampleCS{rušit (čí) klid}  \textbf{2.} \dicSynonym{róta} \dicDirectTranslationCS{narušit, porušit (půdu ap.)} \dicExampleIS{raska jarðvegi} \dicExampleCS{narušit půdu};  \dicIdiom{raskast}{ \dicPhraseIS{raskast}} \dicFlx{refl} \dicSynonym*{haggast} \dicDirectTranslationCS{být narušen\,/\addthin porušen} \dicExampleIS{Friður raskaðist.} \dicExampleCS{Mír byl narušen.}
\dicEntry[rasp] \dicTerm{rasp} \dicIPA{{r}{a}{s}{\textsubring{b}}} \dicPos{n}[2] \dicFlx{(‑s)}[2] \dicDirectTranslationCS{strouhanka}
\dicEntry[raspa] \dicTerm{rasp|a} \dicIPA{{r}{a}{s}{\textsubring{b}}{a}} \dicPos{v}[1] \dicFlx{(‑aði)}[13] \dicFlx{acc} \dicDirectTranslationCS{(na)strouhat} \dicExampleIS{raspa súkkulaði yfir kökuna} \dicExampleCS{nastrouhat čokoládu na dort}
\dicEntry[raspur] \dicTerm{rasp|ur} \dicIPA{{r}{a}{s}{\textsubring{b}}{\textscy}{\textsubring{r}}} \dicPos{m}[6] \dicFlx{(‑s, ‑ar)}[5] \textbf{1.} \dicSynonym*{gróf þjöl} \dicDirectTranslationCS{rašple, hrubý pilník}  \textbf{2.} \dicSynonym{rasp} \dicDirectTranslationCS{strouhanka}
\dicEntry[rass] \dicTerm{rass} \dicsymFrequent\  \dicIPA{{r}{a}{s}{\textlengthmark}} \dicPos{m}[4] \dicFlx{(‑, ‑ar)}[29] \dicSynonym{sitjandi} \dicDirectTranslationCS{pozadí, zadnice, zadek};  \dicPhraseIS{falla eins og flís við rass} \dicLangCat{přen.} \dicDirectTranslationCS{padnout jako ulitý};  \dicPhraseIS{farðu í rass og rófu} \dicLangCat{hovor.} \dicDirectTranslationCS{běž někam};  \dicPhraseIS{lenda\,/\addthin vera í rassi með e‑ð} \dicLangCat{hovor.} \dicDirectTranslationCS{být s~(čím) pozadu};  \dicPhraseIS{það er sami rassinn undir þeim öllum} \dicLangCat{přen.} \dicDirectTranslationCS{všichni jsou stejní (lháři ap.)};  \dicPhraseIS{það er seint í rassinn gripið} \dicLangCat{hovor.} \dicDirectTranslationCS{pozdě bycha honit}
\dicEntry[rassgat] \dicTerm{rass··|gat} \dicIPA{{r}{a}{s}{\r{g}}{a}{\textsubring{d}}} \dicPos{n}[2] \dicFlx{(‑gats, ‑göt)}[8] \textbf{1.} \dicLangCat{hovor.} \dicSynonym*{endaþarmsop} \dicDirectTranslationCS{řiť, řitní otvor};  \dicPhraseIS{rífa e‑n upp á rassgatinu} \dicLangCat{hrub.} \dicDirectTranslationCS{vyhnat (koho) z~postele};  \dicPhraseIS{farðu í rassgat!} \dicFlx{imper} \dicLangCat{hrub.} \dicSynonym*{hypja sig} \dicDirectTranslationCS{jdi do prdele!};  \dicPhraseIS{ekki rassgat (í bala)} \dicLangCat{hovor.} \dicSynonym*{ekki neitt} \dicDirectTranslationCS{ani prd}  \textbf{2.} \dicLangCat{slang.} \dicDirectTranslationCS{prdelka} \dicIndirectTranslationCS{(pěkné, zdravé dítě)} \dicExampleIS{Hún er óttalegt rassgat.} \dicExampleCS{Je to roztomilá prdelka.}  \textbf{3.} \dicLangCat{han.} \dicDirectTranslationCS{prdel, řiť} \dicIndirectTranslationCS{(bezvýznamné místo\,/\addthin město)} \dicExampleIS{í rassgati} \dicExampleCS{v~prdeli}
\dicEntry[rasskinn] \dicTerm{rass··kinn} \dicIPA{{r}{a}{s}{c\smash{\textsuperscript{h}}}{\textsci}{\textsubring{n}}} \dicPos{f}[4] \dicFlx{(‑ar, ‑ar)}[1] \dicDirectTranslationCS{hýždě}
\dicEntry[rassskella] \dicTerm{rass··skell|a} \dicIPA{{r}{a}{s}{\textlengthmark}{\r{\textObardotlessj}}{\textepsilon}{\textsubring{d}}{l}{a}} \dicPos{v}[2] \dicFlx{(‑ti, ‑t)}[77] \dicFlx{acc} \dicDirectTranslationCS{naplácat\,/\addthin dát na zadek} \dicExampleIS{rassskella e‑n duglega} \dicExampleCS{dát (komu) pořádně na zadek}
\dicEntry[rassvasi] \dicTerm{rass··vas|i} \dicIPA{{r}{a}{s}{v}{a}{s}{\textsci}} \dicPos{m}[1] \dicFlx{(‑a, ‑ar)}[8] \dicDirectTranslationCS{kapsa na zadku}
\dicEntry[rata] \dicTerm{rat|a} \dicIPA{{r}{a}{\textlengthmark}{\textsubring{d}}{a}} \dicPos{v}[1] \dicFlx{(‑aði)}[29] \textbf{1.} \dicSynonym*{þekkja leið} \dicDirectTranslationCS{orientovat se, znát cestu} \dicExampleIS{rata í myrkrinu} \dicExampleCS{orientovat se ve tmě}  \textbf{2.} \dicSynonym{finna} \dicDirectTranslationCS{trefit, nalézt se};  \dicIdiom{rata}[í]{ \dicPhraseIS{rata í raunir}} \dicDirectTranslationCS{dostat se do potíží};  \dicIdiom{ratast}{ \dicPhraseIS{e‑m ratast satt á munn}} \dicFlx{refl impers} \dicDirectTranslationCS{(kdo) má pravdu}
\dicEntry[rati] \dicTerm{rat|i} \dicIPA{{r}{a}{\textlengthmark}{\textsubring{d}}{\textsci}} \dicPos{m}[1] \dicFlx{(‑a, ‑ar)}[8] \dicSynonym{kjáni} \dicDirectTranslationCS{magor, pako}
\dicEntry[ratleikur] \dicTerm{rat··leik|ur} \dicIPA{{r}{a}{\textlengthmark}{\textsubring{d}}{l}{ei}{\r{g}}{\textscy}{\textsubring{r}}} \dicPos{m}[9] \dicFlx{(‑s, ‑ir)}[15] \dicFieldCat{sport.} \dicDirectTranslationCS{orientační běh}
\dicEntry[ratsjá] \dicTerm{rat··sjá} \dicIPA{{r}{a}{\textlengthmark}{\textsubring{d}}{s}{j}{au}} \dicPos{f}[4] \dicFlx{(‑r, ‑r)}[18] \dicFieldCat{techn.} \dicDirectTranslationCS{radar}
\dicEntry[ratvísi] \dicTerm{rat··vís|i} \dicIPA{{r}{a}{\textlengthmark}{\textsubring{d}}{v}{i}{s}{\textsci}} \dicPos{f}[3] \dicFlx{(‑i)}[3] \dicDirectTranslationCS{orientační smysl}
\dicEntry[rauð] \dicTerm{rauð} \dicIPA{{r}{\oe i}{\textlengthmark}{\texttheta}} \dicPos{v} \dicFlx{ind pf sg 1 pers} \dicLink{rjóða}
\dicEntry[rauða] \dicTerm{rauð|a} \dicIPA{{r}{\oe i}{\textlengthmark}{ð}{a}} \dicPos{f}[1] \dicFlx{(‑u, ‑ur)}[19] \dicDirectTranslationCS{žloutek}
\dicEntry[rauðabíti] \dicTerm{rauða··bíti} \dicIPA{{r}{\oe i}{\textlengthmark}{ð}{a}{\textsubring{b}}{i}{\textsubring{d}}{\textsci}} \dicPos{n}[2] \dicFlx{(‑s)}[20] \dicSynonym{dögun} \dicDirectTranslationCS{rozbřesk dne, svítání} \dicExampleIS{vakna í rauðabítið} \dicExampleCS{probudit se za svítání}
\dicEntry[Rauðahaf] \dicTerm{Rauða··haf} \dicIPA{{r}{\oe i}{\textlengthmark}{ð}{a}{h}{a}{f}} \dicPos{n}[2] \dicFlx{(‑s)}[2] \dicFieldCat{geog.} \dicDirectTranslationCS{Rudé moře}
\dicEntry[rauðber] \dicTerm{rauð··ber} \dicIPA{{r}{\oe i}{ð}{\textsubring{b}}{\textepsilon}{\textsubring{r}}} \dicPos{n}[2] \dicFlx{(‑s, ‑)}[13] \dicDirectTranslationCS{brusinka} \dicIndirectTranslationCS{(plod)}
\dicEntry[rauðberjalyng] \dicTerm{rauð·berja··lyng} \dicIPA{{r}{\oe i}{ð}{\textsubring{b}}{\textepsilon}{r}{j}{a}{l}{i}{\ng}{\r{g}}} \dicPos{n}[2] \dicFlx{(‑s)}[2] \dicFieldCat{bot.} \dicDirectTranslationCS{brusnice brusinka} \textit{(l.~{\textLA{Vaccinium vitis-idaea}})}  \dicsymPhoto\ 
\dicFigure{69465.jpg}{Rauðberjalyng}{Rauðberjalyng - Krejčík Stanislav, Biolib, Copyright/CC-BY}
\dicEntry[rauðbrúnn] \dicTerm{rauð··brúnn} \dicIPA{{r}{\oe i}{ð}{\textsubring{b}}{r}{u}{\textsubring{d}}{\textsubring{n}}} \dicPos{adj}[7]\dicFlx{}[-1] \dicDirectTranslationCS{červenohnědý, rezavý}
\begin{xtolerant}{}{1pt}
\dicEntry[rauðbrystingur] \dicTerm{rauð··bryst·ing|ur} \dicIPA{{r}{\oe i}{ð}{\textsubring{b}}{r}{\textsci}{s}{\textsubring{d}}{i}{\ng}{\r{g}}{\textscy}{\textsubring{r}}} \dicPos{m}[6] \dicFlx{(‑s, ‑ar)}[8] \dicFieldCat{zool.} \dicDirectTranslationCS{jespák rezavý} \textit{(l.~{\textLA{Calidris canutus}})}  \dicsymPhoto\
\end{xtolerant}
\dicFigure{ds_image_raudbrystingur_0_2.jpg}{Rauðbrystingur}{Rauðbrystingur - Gregory Breese/U. S. Fish and Wildlife Servic, PD}
\dicEntry[rauðeygður] \dicTerm{rauð··eygður}\dicTerm{, rauðeygur} \dicIPA{{r}\-{\oe i}\-{\textlengthmark}\-{ð}\-{ei}\-{\textbabygamma}\-{ð}\-{\textscy}\-{\textsubring{r}}\-} \dicPos{adj}[2]\dicFlx{}[-1] \dicDirectTranslationCS{(jsoucí) s~červenýma a~opuchlýma očima (od pláče ap.)} \dicExampleIS{Hún var rauðeygð af gráti.} \dicExampleCS{Oči měla rudé od pláče.}
\dicEntry[rauðeygur] \dicTerm{rauð··eygur} \dicIPA{{r}{\oe i}{\textlengthmark}{ð}{ei}{\textbabygamma}{\textscy}{\textsubring{r}}} \dicPos{adj}[1]\dicFlx{}[-6] \dicLink{rauðeygður}
\dicEntry[rauðgreni] \dicTerm{rauð··greni} \dicIPA{{r}{\oe i}{ð}{\r{g}}{r}{\textepsilon}{n}{\textsci}} \dicPos{n}[2] \dicFlx{(‑s)}[20] \dicFieldCat{bot.} \dicDirectTranslationCS{smrk ztepilý} \textit{(l.~{\textLA{Picea abies}})}  \dicsymPhoto\ 
\dicFigure{67575.jpg}{Rauðgreni}{Rauðgreni - Marian Huc, Biolib, CC-BY-SA/CC-BY}
\dicEntry[Rauðhetta] \dicTerm{Rauð··hett|a} \dicIPA{{r}{\oe i}{\textlengthmark}{\texttheta}{h}{\textepsilon}{h}{\textsubring{d}}{a}} \dicPos{f}[1] \dicFlx{(‑u)}[6] \dicFlx{prop} \dicDirectTranslationCS{Červená karkulka}
\dicEntry[rauðhærður] \dicTerm{rauð··hærður} \dicIPA{{r}{\oe i}{\textlengthmark}{\texttheta}{h}{a}{i}{r}{ð}{\textscy}{\textsubring{r}}} \dicPos{adj}[2]\dicFlx{}[-1] \dicDirectTranslationCS{rudovlasý, zrzavý}
\dicEntry[rauðhöfðaönd] \dicTerm{rauð·höfða··|önd} \dicIPA{{r}{\oe i}{\textlengthmark}{\texttheta}{h}{\oe}{v}{ð}{a}{\oe}{n}{\textsubring{d}}} \dicPos{f}[8] \dicFlx{(‑andar, ‑endur)}[2] \dicFieldCat{zool.} \dicDirectTranslationCS{hvízdák euroasijský} \textit{(l.~{\textLA{Anas penelope}})}
\dicEntry[rauðkál] \dicTerm{rauð··kál} \dicIPA{{r}{\oe i}{ð}{k\smash{\textsuperscript{h}}}{au}{\textsubring{l}}} \dicPos{n}[2] \dicFlx{(‑s)}[2] \dicFieldCat{bot.} \dicDirectTranslationCS{zelí hlávkové} \textit{(l.~{\textLA{Brassica oleracea var. capitata}})}
\dicEntry[rauðkorn] \dicTerm{rauð··korn} \dicIPA{{r}{\oe i}{ð}{k\smash{\textsuperscript{h}}}{\textopeno}{r}{\textsubring{d}}{\textsubring{n}}} \dicPos{n}[2] \dicFlx{(‑s, ‑)}[5] \dicFieldCat{biol.} \dicDirectTranslationCS{červená krvinka, erytrocyt}
\dicEntry[rauðlaukur] \dicTerm{rauð··lauk|ur} \dicIPA{{r}{\oe i}{ð}{l}{\oe i}{\r{g}}{\textscy}{\textsubring{r}}} \dicPos{m}[6] \dicFlx{(‑s, ‑ar)}[22] \dicFieldCat{bot.} \dicSynonym{laukur} \dicDirectTranslationCS{(červená) cibule, cibule kuchyňská červená} \textit{(l.~{\textLA{Allium cepa}})}
\dicEntry[rauðleitur] \dicTerm{rauð··leitur} \dicIPA{{r}{\oe i}{ð}{l}{ei}{\textsubring{d}}{\textscy}{\textsubring{r}}} \dicPos{adj}[1]\dicFlx{}[-1] \dicDirectTranslationCS{načervenalý, narudlý}
\dicEntry[rauðmagi] \dicTerm{rauð··mag|i} \dicIPA{{r}{\oe i}{ð}{m}{a}{i}{j}{\textsci}} \dicPos{m}[1] \dicFlx{(‑a, ‑ar)}[8] \dicFieldCat{zool.} \dicDirectTranslationCS{hranáč šedý} \textit{(l.~{\textLA{Cyclopterus lumpus}})}  \dicsymPhoto\ 
\dicFigure{26717.jpg}{Rauðmagi}{Rauðmagi - Dr. Aleksey Zuyev, Biolib, PD}
\dicEntry[rauðrófa] \dicTerm{rauð··róf|a} \dicIPA{{r}{\oe i}{ð}{r}{ou}{a}} \dicPos{f}[1] \dicFlx{(‑u, ‑ur)}[19] \dicFieldCat{bot.} \dicDirectTranslationCS{řepa červená} \textit{(l.~{\textLA{Beta vulgaris var. vulgaris}})}
\dicEntry[Rauðskinna] \dicTerm{Rauð··skinn|a} \dicIPA{{r}{\oe i}{ð}{s}{\r{\textObardotlessj}}{\textsci}{n}{a}} \dicPos{f}[1] \dicFlx{(‑u)}[6] \dicFieldCat{myt.} \dicDirectTranslationCS{Rauðskinna} \dicIndirectTranslationCS{(kniha černé magie)}
\dicEntry[rauðskinni] \dicTerm{rauð··skinn|i} \dicIPA{{r}{\oe i}{ð}{s}{\r{\textObardotlessj}}{\textsci}{n}{\textsci}} \dicPos{m}[1] \dicFlx{(‑a, ‑ar)}[1] \dicLangCat{han.} \dicDirectTranslationCS{rudoch, rudoška}
\dicEntry[rauðsmári] \dicTerm{rauð··smár|i} \dicIPA{{r}{\oe i}{ð}{s}{m}{au}{r}{\textsci}} \dicPos{m}[1] \dicFlx{(‑a, ‑ar)}[1] \dicFieldCat{bot.} \dicDirectTranslationCS{jetel luční} \textit{(l.~{\textLA{Trifolium pratense}})}  \dicsymPhoto\ 
\dicFigure{ds_image_raudsmari_0_1.jpg}{Rauðsmári}{Rauðsmári - Masaki Ikeda, CC BY-SA 3.0}
\dicEntry[rauðsokka] \dicTerm{rauð··sokk|a} \dicIPA{{r}{\oe i}{ð}{s}{\textopeno}{h}{\r{g}}{a}} \dicPos{f}[1] \dicFlx{(‑u, ‑ur)}[7] \dicFieldCat{hist.} \dicDirectTranslationCS{feministka} \dicIndirectTranslationCS{(v~60. a~70. letech 20. století)}
\dicEntry[rauðspretta] \dicTerm{rauð··sprett|a} \dicIPA{{r}{\oe i}{ð}{s}{\textsubring{b}}{r}{\textepsilon}{h}{\textsubring{d}}{a}} \dicPos{f}[1] \dicFlx{(‑u, ‑ur)}[7] \dicFieldCat{zool.} \dicSynonym{skarkoli} \dicDirectTranslationCS{platýs velký} \textit{(l.~{\textLA{Pleuronectes platessa}})}
\dicEntry[rauður] \dicTerm{rauður} \dicsymFrequent\  \dicIPA{{r}{\oe i}{\textlengthmark}{ð}{\textscy}{\textsubring{r}}} \dicPos{adj}[2]\dicFlx{}[-6] \textbf{1.} \dicDirectTranslationCS{červený, rudý};  \dicPhraseIS{Rauði krossinn} \dicFlx{prop} \dicDirectTranslationCS{Červený kříž}  \textbf{2.} \dicSynonym{kommúnískur} \dicDirectTranslationCS{rudý, komunistický};  \dicPhraseIS{Rauði herinn} \dicFieldCat{hist.} \dicDirectTranslationCS{Rudá armáda}  \textbf{3.} \dicSynonym{snjólaus} \dicDirectTranslationCS{(jsoucí) bez sněhu, nezasněžený};  \dicPhraseIS{rauð jól} \dicDirectTranslationCS{Vánoce bez sněhu, Vánoce na blátě};  \dicIdiom{rauður}{ \dicPhraseIS{gjalda e‑m rauðan belg fyrir gráan}} \dicLangCat{přen.} \dicDirectTranslationCS{vrátit to (komu) i~s~úroky}; { \dicPhraseIS{mála bæinn rauðan}} \dicLangCat{přen.} \dicDirectTranslationCS{jít se bavit (do města), obrátit město vzhůru nohama}
\dicEntry[rauðvín] \dicTerm{rauð··vín} \dicIPA{{r}{\oe i}{ð}{v}{i}{\textsubring{n}}} \dicPos{n}[2] \dicFlx{(‑s, ‑)}[5] \dicDirectTranslationCS{červené víno}
\dicEntry[rauf] \dicTerm{rauf\smash{\textsuperscript{1}}} \dicIPA{{r}{\oe i}{\textlengthmark}{f}} \dicPos{f}[4] \dicFlx{(‑ar, ‑ar)}[1] \textbf{1.} \dicSynonym{op} \dicDirectTranslationCS{díra, otvor}  \textbf{2.} \dicSynonym{rifa} \dicDirectTranslationCS{štěrbina, prasklina}
\dicEntry[rauf] \dicTerm{rauf\smash{\textsuperscript{2}}} \dicIPA{{r}{\oe i}{\textlengthmark}{f}} \dicPos{v} \dicFlx{ind pf sg 1 pers} \dicLink{rjúfa}
\dicEntry[rauk] \dicTerm{rauk} \dicIPA{{r}{\oe i}{\textlengthmark}{\r{g}}} \dicPos{v} \dicFlx{ind pf sg 1 pers} \dicLink{rjúka}
\dicEntry[raula] \dicTerm{raul|a} \dicIPA{{r}{\oe i}{\textlengthmark}{l}{a}} \dicPos{v}[1] \dicFlx{(‑aði)}[1] \dicFlx{acc} \dicSynonym{söngla} \dicDirectTranslationCS{broukat si, potichu si prozpěvovat} \dicExampleIS{raula lag fyrir munni sér} \dicExampleCS{broukat si pro sebe písničku}
\dicEntry[raun] \dicTerm{raun} \dicsymFrequent\  \dicIPA{{r}{\oe i}{\textlengthmark}{\textsubring{n}}} \dicPos{f}[7] \dicFlx{(‑ar, ‑ir)}[1] \textbf{1.} \dicSynonym{reynsla} \dicDirectTranslationCS{zkušenost, zážitek};  \dicPhraseIS{af eigin raun} \dicFlx{adv} \dicDirectTranslationCS{z~vlastní zkušenosti};  \dicPhraseIS{í rauninni, í raun og veru} \dicFlx{adv} \dicDirectTranslationCS{ve skutečnosti, skutečně};  \dicPhraseIS{komast að raun um e‑ð} \dicFlx{refl} \dicDirectTranslationCS{přijít na (co), dozvědět se (co)}  \textbf{2.} \dicSynonym{þraut} \dicDirectTranslationCS{(těžká) zkouška, utrpení};  \dicPhraseIS{standast raun} \dicFlx{refl} \dicDirectTranslationCS{přestát těžkou zkoušku};  \dicPhraseIS{það er raun að} \dicDirectTranslationCS{je utrpení (být svědkem něčeho ap.)};  \dicProverb\  \dicPhraseIS{Sá er vinur sem í raun reynist.} \dicLangCat{přís.} \dicDirectTranslationCS{V~nouzi poznáš přítele.}
\dicEntry[raunalegur] \dicTerm{rauna··legur} \dicIPA{{r}{\oe i}{\textlengthmark}{n}{a}{l}{\textepsilon}{\textbabygamma}{\textscy}{\textsubring{r}}} \dicPos{adj}[1]\dicFlx{}[-8] \textbf{1.} \dicSynonym{sorglegur} \dicDirectTranslationCS{žalostný, tragický} \dicExampleIS{raunaleg saga} \dicExampleCS{tragický příběh}  \textbf{2.} \dicSynonym{dapur} \dicDirectTranslationCS{smutný, zarmoucený} \dicExampleIS{vera raunalegur á svipinn} \dicExampleCS{mít smutný výraz}
\dicEntry[raunar] \dicTerm{raunar} \dicsymFrequent\  \dicIPA{{r}{\oe i}{\textlengthmark}{n}{a}{\textsubring{r}}} \dicPos{adv} \dicSynonym*{að sönnu} \dicDirectTranslationCS{vlastně, ve skutečnosti} \dicExampleIS{Ég hef raunar aldrei efast um það.} \dicExampleCS{Ve skutečnosti jsem o~tom nikdy nepochyboval.}
\dicEntry[raungildi] \dicTerm{raun··gildi} \dicIPA{{r}{\oe i}{\textltailn}{\r{\textObardotlessj}}{\textsci}{l}{\textsubring{d}}{\textsci}} \dicPos{n}[2] \dicFlx{(‑s)}[20] \dicDirectTranslationCS{reálná\,/\addthin skutečná hodnota}
\dicEntry[raunhyggja] \dicTerm{raun··hyggj|a} \dicIPA{{r}{\oe i}{\textlengthmark}{\textsubring{n}}{h}{\textsci}{\r{\textObardotlessj}}{a}} \dicPos{f}[1] \dicFlx{(‑u)}[5] \dicFieldCat{filos.} \dicDirectTranslationCS{empirismus}
\dicEntry[raunhæfur] \dicTerm{raun··hæfur} \dicIPA{{r}{\oe i}{\textlengthmark}{\textsubring{n}}{h}{a}{i}{v}{\textscy}{\textsubring{r}}} \dicPos{adj}[1]\dicFlx{}[-1] \dicSynonym{hagnýtur} \dicDirectTranslationCS{praktický} \dicExampleIS{raunhæft starf} \dicExampleCS{praktická práce} \dicAntonym{fræðilegur}
\dicEntry[raunspeki] \dicTerm{raun··spek|i} \dicIPA{{r}{\oe i}{n}{s}{\textsubring{b}}{\textepsilon}{\r{\textObardotlessj}}{\textsci}} \dicPos{f}[3] \dicFlx{(‑i)}[3] \dicFieldCat{filos.} \dicDirectTranslationCS{pozitivismus}
\dicEntry[raunsæisstefna] \dicTerm{raun·sæis··stefn|a} \dicIPA{{r}{\oe i}{n}{s}{a}{i}{j}{\textsci}{s}{\textsubring{d}}{\textepsilon}{\textsubring{b}}{n}{a}} \dicPos{f}[1] \dicFlx{(‑u)}[5] \dicFieldCat{lit.} \dicDirectTranslationCS{realismus}
\dicEntry[raunsær] \dicTerm{raun··sær} \dicIPA{{r}{\oe i}{n}{s}{a}{i}{\textsubring{r}}} \dicPos{adj}[4]\dicFlx{}[-7] \dicDirectTranslationCS{realistický, věcný} \dicExampleIS{raunsæjar bókmenntir} \dicExampleCS{literatura realismu}
\dicEntry[raunverulega] \dicTerm{raun·veru··lega} \dicsymFrequent\  \dicIPA{{r}{\oe i}{n}{v}{\textepsilon}{r}{\textscy}{l}{\textepsilon}{\textbabygamma}{a}} \dicPos{adv} \dicSynonym*{í rauninni} \dicDirectTranslationCS{skutečně, opravdu} \dicExampleIS{rifja upp hvað raunverulega gerðist} \dicExampleCS{připomenout, co se skutečně stalo}
\dicEntry[raunverulegur] \dicTerm{raun·veru··legur} \dicsymFrequent\  \dicIPA{{r}{\oe i}{n}{v}{\textepsilon}{r}{\textscy}{l}{\textepsilon}{\textbabygamma}{\textscy}{\textsubring{r}}} \dicPos{adj}[1]\dicFlx{}[-8] \dicSynonym{áþreifanlegur} \dicDirectTranslationCS{skutečný, reálný, opravdový} \dicExampleIS{raunveruleg ástæða} \dicExampleCS{skutečný důvod} \dicAntonym{óraunverulegur}
\dicEntry[raunveruleiki] \dicTerm{raun·veru··leik|i} \dicIPA{{r}{\oe i}{n}{v}{\textepsilon}{r}{\textscy}{l}{ei}{\r{\textObardotlessj}}{\textsci}} \dicPos{m}[1] \dicFlx{(‑a)}[3] \dicSynonym{veruleiki} \dicDirectTranslationCS{skutečnost, realita}
\dicEntry[raunvextir] \dicTerm{raun··vextir} \dicIPA{{r}{\oe i}{n}{v}{\textepsilon}{x}{s}{\textsubring{d}}{\textsci}{\textsubring{r}}} \dicPos{m}[11] \dicFlx{pl}[1] \dicFieldCat{ekon.} \dicDirectTranslationCS{reálná úroková sazba}
\dicEntry[raunvísindi] \dicTerm{raun··vís·indi} \dicIPA{{r}{\oe i}{n}{v}{i}{s}{\textsci}{n}{\textsubring{d}}{\textsci}} \dicPos{n}[2] \dicFlx{pl}[19] \dicDirectTranslationCS{přírodní vědy}
\dicEntry[raup] \dicTerm{raup} \dicIPA{{r}{\oe i}{\textlengthmark}{\textsubring{b}}} \dicPos{n}[2] \dicFlx{(‑s)}[2] \dicSynonym{gort} \dicDirectTranslationCS{chvástání, chlubení se, vytahování se}
\dicEntry[raupa] \dicTerm{raup|a} \dicIPA{{r}{\oe i}{\textlengthmark}{\textsubring{b}}{a}} \dicPos{v}[1] \dicFlx{(‑aði)}[1] \dicSynonym{gorta} \dicDirectTranslationCS{chvástat se, vychloubat se, vytahovat se} \dicExampleIS{raupa af e‑u} \dicExampleCS{vytahovat se (čím)}
\dicEntry[raupsamur] \dicTerm{raup··|samur} \dicIPA{{r}{\oe i}{\textlengthmark}{\textsubring{b}}{s}{a}{m}{\textscy}{\textsubring{r}}} \dicPos{adj}[1] \dicFlx{(f ‑söm)}[2] \dicSynonym*{grobbinn} \dicDirectTranslationCS{chvástavý, chlubivý}
\dicEntry[raus] \dicTerm{raus} \dicIPA{{r}{\oe i}{\textlengthmark}{s}} \dicPos{n}[2] \dicFlx{(‑s, ‑)}[5] \dicDirectTranslationCS{žvást, blábol}
\dicEntry[rausa] \dicTerm{raus|a} \dicIPA{{r}{\oe i}{\textlengthmark}{s}{a}} \dicPos{v}[1] \dicFlx{(‑aði)}[44] \dicSynonym{þusa} \dicDirectTranslationCS{blábolit, mlít nesmysly}
\dicEntry[rausn] \dicTerm{rausn} \dicIPA{{r}{\oe i}{s}{\textsubring{d}}{\textsubring{n}}} \dicPos{f}[4] \dicFlx{(‑ar)}[3] \dicSynonym{örlæti} \dicDirectTranslationCS{štědrost, velkorysost}
\dicEntry[rausnarlegur] \dicTerm{rausnar··legur} \dicIPA{{r}{\oe i}{s}{\textsubring{d}}{n}{a}{r}{l}{\textepsilon}{\textbabygamma}{\textscy}{\textsubring{r}}} \dicPos{adj}[1]\dicFlx{}[-8] \dicSynonym{örlátur} \dicDirectTranslationCS{štědrý, velkorysý} \dicExampleIS{rausnarlegt boð} \dicExampleCS{velkorysá nabídka}
\dicEntry[raust] \dicTerm{raust} \dicIPA{{r}{\oe i}{s}{\textsubring{d}}} \dicPos{f}[7] \dicFlx{(‑ar, ‑ir)}[6] \dicSynonym{rödd} \dicDirectTranslationCS{hlas};  \dicPhraseIS{brýna raustina} \dicSynonym*{hækka róminn} \dicDirectTranslationCS{zvýšit hlas};  \dicPhraseIS{kveða við raust} \dicDirectTranslationCS{hlasitě si notovat}
\dicEntry[rá] \dicTerm{rá} \dicIPA{{r}{au}{\textlengthmark}} \dicPos{f}[4] \dicFlx{(‑r\,/\addthin ‑ar, ‑r\,/\addthin rær)}[28] \dicFieldCat{nám.} \dicDirectTranslationCS{ráhno}
\dicEntry[rádýr] \dicTerm{rá··dýr} \dicIPA{{r}{au}{\textlengthmark}{\textsubring{d}}{i}{\textsubring{r}}} \dicPos{n}[2] \dicFlx{(‑s, ‑)}[5] \dicFieldCat{zool.} \dicDirectTranslationCS{srnec, srnec obecný} \textit{(l.~{\textLA{Capreolus capreolus}})}  \dicsymPhoto\ 
\dicFigure{ds_image_radyr_0_1.jpg}{Rádýr}{Rádýr - Marek Szczepanek, GFDL}
\dicEntry[rádýrstarfur] \dicTerm{rá·dýrs··tarf|ur} \dicIPA{{r}{au}{\textlengthmark}{\textsubring{d}}{i}{\textsubring{r}}{s}{t\smash{\textsuperscript{h}}}{a}{r}{v}{\textscy}{\textsubring{r}}} \dicPos{m}[6] \dicFlx{(‑s, ‑ar)}[5] \dicDirectTranslationCS{srnec}
\dicEntry[ráð] \dicTerm{ráð} \dicsymFrequent\  \dicIPA{{r}{au}{\textlengthmark}{\texttheta}} \dicPos{n}[2] \dicFlx{(‑s, ‑)}[5] \textbf{1.} \dicSynonym{ráðlegging} \dicDirectTranslationCS{rada, doporučení};  \dicPhraseIS{fara að ráði\,/\addthin ráðum e‑rs} \dicDirectTranslationCS{postupovat podle (čí) rady};  \dicPhraseIS{með ráðum og dáð} \dicFlx{adv} \dicLangCat{přen.} \dicDirectTranslationCS{slovem i~skutkem};  \dicPhraseIS{spyrja e‑n ráða} \dicDirectTranslationCS{požádat (koho) o~radu}  \textbf{2.} \dicSynonym{úrræði} \dicDirectTranslationCS{řešení, východisko} \dicExampleIS{Mér kemur ráð í hug.} \dicExampleCS{Napadlo mě řešení.};  \dicPhraseIS{bregða á það ráð að (gera e‑ð)} \dicDirectTranslationCS{uchýlit se k~(udělání (čeho))};  \dicPhraseIS{hafa ráð undir rifi hverju} \dicLangCat{přen.} \dicDirectTranslationCS{mít na všechno odpověď};  \dicPhraseIS{hvað er til ráða?} \dicDirectTranslationCS{a~co teď?, co dělat?};  \dicPhraseIS{kunna ráð við e‑u} \dicDirectTranslationCS{najít řešení (čeho)}  \textbf{3.} \dicSynonym{samþykki} \dicDirectTranslationCS{souhlas} \dicExampleIS{gert með ráði e‑rs} \dicExampleCS{uděláno s~(čím) souhlasem}  \textbf{4.} \dicSynonym{skynsemi} \dicDirectTranslationCS{úsudek, rozum};  \dicPhraseIS{vera ekki með réttu ráði} \dicLangCat{přen.} \dicDirectTranslationCS{nebýt při smyslech};  \dicPhraseIS{vita ekki sitt rjúkandi ráð} \dicLangCat{přen.} \dicDirectTranslationCS{nevědět si rady, být (naprosto) bezradný}  \textbf{5.} \dicSynonym{áætlun} \dicDirectTranslationCS{plán, záměr, úmysl};  \dicPhraseIS{af ásettu ráði} \dicFlx{adv} \dicDirectTranslationCS{úmyslně, záměrně};  \dicPhraseIS{leggja á ráðin} \dicDirectTranslationCS{hodlat, plánovat};  \dicPhraseIS{það varð að ráði að (gera e‑ð)} \dicDirectTranslationCS{bylo rozhodnuto (udělat (co))}  \textbf{6.} \dicSynonym{nefnd} \dicDirectTranslationCS{rada, komise (bankovní ap.)} \dicExampleIS{bankaráð} \dicExampleCS{bankovní rada};  \dicIdiom{ráð}{ \dicPhraseIS{bæta ráð sitt}} \dicDirectTranslationCS{polepšit se, napravit se};  \dicPhraseIS{e‑að að yfirlögðu ráði} \dicDirectTranslationCS{(co) předem naplánované (vražda ap.)};  \dicPhraseIS{eins og lög gera ráð fyrir} \dicFlx{adv} \dicDirectTranslationCS{v~souladu se zákonem};  \dicPhraseIS{gera ráð fyrir e‑u} \dicDirectTranslationCS{domnívat se (co), mít za (co)};  \dicPhraseIS{festa\,/\addthin staðfesta ráð sitt} \dicDirectTranslationCS{uzavřít sňatek};  \dicPhraseIS{nú eru góð ráð dýr} \dicLangCat{přen.} \dicDirectTranslationCS{teď je každá rada drahá};  \dicPhraseIS{rasa (ekki) um ráð fram} \dicDirectTranslationCS{(ne)pospíchat}
\dicEntry[ráða] \dicTerm{ráða} \dicsymFrequent\  \dicIPA{{r}{au}{\textlengthmark}{ð}{a}} \dicPos{v}[6] \dicFlx{(ræð, réð, réðum, réði, ráðið)}[85] \dicFlx{acc\,/\addthin dat} \textbf{1.} \dicFlx{acc} \dicDirectTranslationCS{zaměstnat, zaměstnávat, najmout si (pracovníka ap.)} \dicExampleIS{ráða e‑n í vinnu} \dicExampleCS{zaměstnat (koho)}  \textbf{2.} \dicFlx{acc} \dicSynonym{leysa} \dicDirectTranslationCS{vyřešit, rozluštit} \dicExampleIS{ráða gátu} \dicExampleCS{vyřešit hádanku}  \textbf{3.} \dicFlx{dat} \dicSynonym{stjórna} \dicDirectTranslationCS{řídit, vést, spravovat (firmu ap.)};  \dicPhraseIS{ráða ríkjum, ráða lögum og lofum} \dicDirectTranslationCS{dominovat, vládnout (všemu)}  \textbf{4.} \dicFlx{dat} \dicSynonym{ákveða} \dicDirectTranslationCS{rozhodnout, rozhodovat} \dicExampleIS{ráða öllu} \dicExampleCS{rozhodnout vše}  \textbf{5.} \dicFlx{dat} \dicSynonym{ráðleggja} \dicDirectTranslationCS{(po)radit, doporučit} \dicExampleIS{Ég ræð henni eindregið að leita læknis.} \dicExampleCS{Rozhodně jí doporučuji vyhledat lékaře.};  \dicPhraseIS{ráða e‑m heilt} \dicDirectTranslationCS{dát (komu) dobrou radu};  \dicIdiom{ráða}[af]{ \dicPhraseIS{ráða e‑ð af e‑u}} \dicDirectTranslationCS{vyvozovat (co) z~(čeho)}; { \dicPhraseIS{ráða e‑n af lífi\,/\addthin dögum}} \dicDirectTranslationCS{zavraždit (koho), sprovodit (koho) ze světa};  \dicIdiom{ráða}[fram úr]{ \dicPhraseIS{ráða fram úr e‑u}} \dicDirectTranslationCS{nalézt východisko z~(čeho)} \dicExampleIS{ráða fram úr erfiðleikunum} \dicExampleCS{nalézt východisko z~potíží};  \dicIdiom{ráða}[frá]{ \dicPhraseIS{ráða e‑m frá e‑u}} \dicDirectTranslationCS{odradit\,/\addthin odrazovat (koho) od (čeho)};  \dicIdiom{ráða}[fyrir]{ \dicPhraseIS{ráða fyrir e‑u}} \dicDirectTranslationCS{vládnout (čemu), řídit (co)};  \dicIdiom{ráða}[við]{ \dicPhraseIS{ráða við e‑ð}} \dicDirectTranslationCS{poradit si s~(čím), zvládnout (co)}; { \dicPhraseIS{ráða við e‑n}} \dicDirectTranslationCS{poradit si s~(kým)};  \dicIdiom{ráða}[yfir]{ \dicPhraseIS{ráða yfir e‑u\,/\addthin e‑m}} \dicDirectTranslationCS{rozhodovat o~(čem\,/\addthin kom), ovládat (co\,/\addthin koho)};  \dicIdiom{ráðast}{ \dicPhraseIS{ráðast}} \dicFlx{refl} {\textbf{a.}} \dicDirectTranslationCS{rozhodnout se, rozřešit se (osud ap.)};  {\textbf{b.}} \dicDirectTranslationCS{zaútočit, atakovat, udeřit}; { \dicPhraseIS{það verður að ráðast}} \dicFlx{refl} \dicDirectTranslationCS{to se uvidí, to se ukáže};  \dicIdiom{ráðast}[að]{ \dicPhraseIS{ráðast að e‑m}} \dicFlx{refl} \dicDirectTranslationCS{napadnout (koho), zaútočit na (koho)};  \dicIdiom{ráðast}[á]{ \dicPhraseIS{ráðast á e‑n}} \dicFlx{refl} \dicDirectTranslationCS{napadnout (koho), zaútočit na (koho)} \dicExampleIS{ráðast á manninn úr launsátri} \dicExampleCS{napadnout muže ze zálohy};  \dicIdiom{ráðast}[í]{ \dicPhraseIS{ráðast í e‑ð}} \dicFlx{refl} \dicDirectTranslationCS{pustit se do (čeho), vrhnout se do (čeho)};  \dicIdiom{ráða}{ \dicPhraseIS{e‑að ræður úrslitum}} \dicDirectTranslationCS{(co) rozhoduje, (co) je rozhodující};  \dicPhraseIS{ráða ferðinni} \dicDirectTranslationCS{vybírat cestu, volit směr}; { \dicPhraseIS{ráða sér ekki (fyrir gleði)}} \dicDirectTranslationCS{být bez sebe (štěstím)}
\dicEntry[ráðabrugg] \dicTerm{ráða··brugg} \dicIPA{{r}{au}{\textlengthmark}{ð}{a}{\textsubring{b}}{r}{\textscy}{\r{g}}} \dicPos{n}[2] \dicFlx{(‑s)}[2] \dicSynonym{brugg} \dicDirectTranslationCS{komplot, intrika, pleticha}
\dicEntry[ráðagerð] \dicTerm{ráða··gerð} \dicIPA{{r}{au}{\textlengthmark}{ð}{a}{\r{\textObardotlessj}}{\textepsilon}{r}{\texttheta}} \dicPos{f}[7] \dicFlx{(‑ar, ‑ir)}[1] \dicSynonym{ætlun} \dicDirectTranslationCS{plán, návrh, záměr} \dicExampleIS{ráðagerðir um nýja virkjun} \dicExampleCS{plán nové elektrárny}
\dicEntry[ráðagóður] \dicTerm{ráða··|góður} \dicIPA{{r}{au}{\textlengthmark}{ð}{a}{\r{g}}{ou}{ð}{\textscy}{\textsubring{r}}} \dicPos{adj}[11] \dicFlx{(comp ‑betri, sup ‑bestur)}[1] \dicSynonym{hygginn} \dicDirectTranslationCS{důmyslný, vynalézavý}
\dicEntry[ráðalaus] \dicTerm{ráða··laus} \dicIPA{{r}{au}{\textlengthmark}{ð}{a}{l}{\oe i}{s}} \dicPos{adj}[5]\dicFlx{}[-1] \dicSynonym{úrræðalaus} \dicDirectTranslationCS{bezradný, rozpačitý};  \dicPhraseIS{deyja ekki ráðalaus} \dicLangCat{přen.} \dicDirectTranslationCS{dokázat\,/\addthin umět si poradit}
\dicEntry[ráðaleysi] \dicTerm{ráða··leysi}\dicTerm{, ráðleysi} \dicIPA{{r}\-{au}\-{\textlengthmark}\-{ð}\-{a}\-{l}\-{ei}\-{s}\-{\textsci}\-} \dicPos{n}[2] \dicFlx{(‑s)}[20] \dicSynonym{úrræðaleysi} \dicDirectTranslationCS{bezradnost, rozpačitost}
\dicEntry[ráðamaður] \dicTerm{ráða··|maður} \dicIPA{{r}{au}{\textlengthmark}{ð}{a}{m}{a}{ð}{\textscy}{\textsubring{r}}} \dicPos{m}[13] \dicFlx{(‑manns, ‑menn)}[2] \dicDirectTranslationCS{vlivný činitel, vlivná činitelka, autorita}
\dicEntry[ráðandi] \dicTerm{ráð··|andi\smash{\textsuperscript{1}}} \dicIPA{{r}{au}{\textlengthmark}{ð}{a}{n}{\textsubring{d}}{\textsci}} \dicPos{m}[2] \dicFlx{(‑anda, ‑endur)}[1] \textbf{1.} \dicSynonym{stjórnandi} \dicDirectTranslationCS{správce, správkyně, správcová, administrátor(ka)}  \textbf{2.} \dicSynonym*{sá sem ræður hugþraut} \dicDirectTranslationCS{řešitel(ka)}
\dicEntry[ráðandi] \dicTerm{ráð··andi\smash{\textsuperscript{2}}} \dicIPA{{r}{au}{\textlengthmark}{ð}{a}{n}{\textsubring{d}}{\textsci}} \dicPos{adj}[13] \dicFlx{indecl}[1] \dicDirectTranslationCS{dominantní, převládající, převažující} \dicExampleIS{ráðandi öfl} \dicExampleCS{dominující síly};  \dicPhraseIS{vera fjár síns ráðandi} \dicDirectTranslationCS{disponovat svým majetkem}
\dicEntry[ráðdeild] \dicTerm{ráð··deild} \dicIPA{{r}{au}{ð}{\textsubring{d}}{ei}{l}{\textsubring{d}}} \dicPos{f}[7] \dicFlx{(‑ar)}[3] \dicSynonym{hagsýni} \dicDirectTranslationCS{hospodárnost, šetrnost, ekonomičnost} \dicExampleIS{sýna ráðdeild} \dicExampleCS{vykazovat hospodárnost}
\dicEntry[ráðfæra] \dicTerm{ráð··fær|a} \dicIPA{{r}{au}{ð}{f}{a}{i}{r}{a}} \dicPos{v}[2] \dicFlx{(‑ði, ‑t)}[99] \dicFlx{acc} \dicPhraseIS{ráðfæra sig við e‑n} \dicSynonym*{spyrja ráða} \dicDirectTranslationCS{poradit se s~(kým)} \dicExampleIS{ráðfæra sig við lögfræðing} \dicExampleCS{poradit se s~právníkem}
\dicEntry[ráðgáta] \dicTerm{ráð··gát|a} \dicIPA{{r}{au}{ð}{\r{g}}{au}{\textsubring{d}}{a}} \dicPos{f}[1] \dicFlx{(‑u, ‑ur)}[19] \dicSynonym{gáta\smash{\textsuperscript{1}}} \dicDirectTranslationCS{záhada, hádanka} \dicExampleIS{Þetta er mér ráðgáta.} \dicExampleCS{Tak to je mi hádankou.}
\dicEntry[ráðgefandi] \dicTerm{ráð··gef·andi} \dicIPA{{r}{au}{ð}{\r{\textObardotlessj}}{\textepsilon}{v}{a}{n}{\textsubring{d}}{\textsci}} \dicPos{adj}[13] \dicFlx{indecl}[1] \dicDirectTranslationCS{poradní, poradenský, konzultační} \dicExampleIS{ráðgefandi aðili} \dicExampleCS{poradní strana}
\dicEntry[ráðgera] \dicTerm{ráð··ger|a} \dicIPA{{r}{au}{ð}{\r{\textObardotlessj}}{\textepsilon}{r}{a}} \dicPos{v}[2] \dicFlx{(‑ði, ‑t)}[105] \dicFlx{acc} \dicSynonym*{hafa í hyggju} \dicDirectTranslationCS{plánovat, chystat, zamýšlet} \dicExampleIS{ráðgera ferð} \dicExampleCS{plánovat cestu}
\dicEntry[ráðgjafi] \dicTerm{ráð··gjaf|i} \dicIPA{{r}{au}{ð}{\r{\textObardotlessj}}{a}{v}{\textsci}} \dicPos{m}[1] \dicFlx{(‑a, ‑ar)}[8] \dicDirectTranslationCS{poradce, poradkyně, konzultant(ka)} \dicExampleIS{ráðgjafi e‑rs í fjármálum} \dicExampleCS{(čí) konzultant ve finančních záležitostech}
\dicEntry[ráðgjöf] \dicTerm{ráð··|gjöf} \dicIPA{{r}{au}{ð}{\r{\textObardotlessj}}{\oe}{f}} \dicPos{f}[7] \dicFlx{(‑gjafar)}[19] \dicDirectTranslationCS{poradenství, konzultace} \dicExampleIS{fá ráðgjöf} \dicExampleCS{konzultovat}
\dicEntry[ráðherra] \dicTerm{ráð··herr|a} \dicsymFrequent\  \dicIPA{{r}{au}{\textlengthmark}{\texttheta}{h}{\textepsilon}{r}{a}} \dicPos{m}[1] \dicFlx{(‑a, ‑ar)}[17] \dicDirectTranslationCS{ministr(yně)} \dicExampleIS{hæstvirtur ráðherra} \dicExampleCS{vážený pane ministře}
\dicEntry[ráðhús] \dicTerm{ráð··hús} \dicIPA{{r}{au}{\textlengthmark}{\texttheta}{h}{u}{s}} \dicPos{n}[2] \dicFlx{(‑s, ‑)}[5] \dicDirectTranslationCS{radnice} \dicExampleIS{Á myndinni er ráðhús í Reykjavík.} \dicExampleCS{Na obrázku je radnice v~Reykjavíku.}
\dicEntry[ráðið] \dicTerm{ráðið} \dicIPA{{r}{au}{\textlengthmark}{ð}{\textsci}{\texttheta}} \dicPos{v} \dicFlx{supin} \dicLink{ráða}
\dicEntry[ráðlegging] \dicTerm{ráð··legg·ing} \dicIPA{{r}{au}{ð}{l}{\textepsilon}{\r{\textObardotlessj}}{i}{\ng}{\r{g}}} \dicPos{f}[4] \dicFlx{(‑ar, ‑ar)}[5] \textbf{1.} \dicSynonym{ráð} \dicDirectTranslationCS{rada, instrukce, ponaučení} \dicExampleIS{ráðlegging um notkun lyfsins} \dicExampleCS{návod k~použití léku}  \textbf{2.} \dicSynonym*{það að ráðleggja} \dicDirectTranslationCS{poradění, doporučení}
\dicEntry[ráðleggja] \dicTerm{ráð··|leggja} \dicIPA{{r}{au}{ð}{l}{\textepsilon}{\r{\textObardotlessj}}{a}} \dicPos{v}[4] \dicFlx{(‑legg, ‑lagði, ‑lögðum, ‑legði, ‑lagt)}[21] \dicFlx{dat + acc} \dicSynonym*{gefa ráð} \dicDirectTranslationCS{(po)radit, doporučit} \dicExampleIS{ráðleggja e‑m e‑ð} \dicExampleCS{poradit (co komu)}
\dicEntry[ráðlegur] \dicTerm{ráð··legur} \dicIPA{{r}{au}{ð}{l}{\textepsilon}{\textbabygamma}{\textscy}{\textsubring{r}}} \dicPos{adj}[1]\dicFlx{}[-8] \dicSynonym{heppilegur} \dicDirectTranslationCS{doporučeníhodný, vhodný, záhodný}
\dicEntry[ráðleysi] \dicTerm{ráð··leysi} \dicIPA{{r}{au}{ð}{l}{ei}{s}{\textsci}} \dicPos{n}[2] \dicFlx{(‑s)}[20] \dicLink{ráðaleysi}
\dicEntry[ráðning] \dicTerm{ráð··ning} \dicIPA{{r}{au}{ð}{n}{i}{\ng}{\r{g}}} \dicPos{f}[4] \dicFlx{(‑ar, ‑ar)}[5] \textbf{1.} \dicSynonym{lausn} \dicDirectTranslationCS{(vy)řešení} \dicExampleIS{ráðning gátu} \dicExampleCS{řešení hádanky}  \textbf{2.} \dicDirectTranslationCS{nábor, přijímání, najímání} \dicExampleIS{ráðning í vist} \dicExampleCS{nábor na brigádu}  \textbf{3.} \dicSynonym{refsing} \dicDirectTranslationCS{potrestání, lekce};  \dicPhraseIS{veita e‑m ráðningu} \dicLangCat{přen.} \dicDirectTranslationCS{uštědřit\,/\addthin udělit (komu) lekci}
\dicEntry[ráðningarsamningur] \dicTerm{ráðningar··samn·ing|ur} \dicIPA{{r}\-{au}\-{ð}\-{n}\-{i}\-{\ng}\-{\r{g}}\-{a}\-{\textsubring{r}}\-{s}\-{a}\-{m}\-{n}\-{i}\-{\ng}\-{\r{g}}\-{\textscy}\-{\textsubring{r}}\-} \dicPos{m}[6] \dicFlx{(‑s, ‑ar)}[8] \dicFieldCat{práv.} \dicDirectTranslationCS{dohoda o~práci}
\dicEntry[ráðningarskilyrði] \dicTerm{ráðningar··skil·yrði} \dicIPA{{r}{au}{ð}{n}{i}{\ng}{\r{g}}{a}{\textsubring{r}}{s}{\r{\textObardotlessj}}{\textsci}{l}{\textsci}{r}{ð}{\textsci}} \dicPos{n}[2] \dicFlx{(‑s, ‑)}[14] \dicDirectTranslationCS{podmínky\,/\addthin požadavky pro přijetí do zaměstnání}
\dicEntry[ráðningarskrifstofa] \dicTerm{ráðningar··skrif·stof|a} \dicIPA{{r}\-{au}\-{ð}\-{n}\-{i}\-{\ng}\-{\r{g}}\-{a}\-{\textsubring{r}}\-{s}\-{\r{g}}\-{r}\-{\textsci}\-{f}\-{s}\-{\textsubring{d}}\-{\textopeno}\-{v}\-{a}\-} \dicPos{f}[1] \dicFlx{(‑u, ‑ur)}[7] \dicDirectTranslationCS{personální agentura}
\dicEntry[ráðríkur] \dicTerm{ráð··ríkur} \dicIPA{{r}{au}{ð}{r}{i}{\r{g}}{\textscy}{\textsubring{r}}} \dicPos{adj}[1]\dicFlx{}[-1] \dicSynonym{drottnunargjarn} \dicDirectTranslationCS{panovačný, pánovitý, despotický}
\dicEntry[ráðrúm] \dicTerm{ráð··rúm} \dicIPA{{r}{au}{ð}{r}{u}{\textsubring{m}}} \dicPos{n}[2] \dicFlx{(‑s)}[2] \dicDirectTranslationCS{(časový) prostor, čas};  \dicPhraseIS{hafa ráðrúm til e‑s} \dicLangCat{přen.} \dicDirectTranslationCS{mít prostor na (co)}
\begin{xtolerant}{}{1pt}
\dicEntry[ráðsettur] \dicTerm{ráð··settur}\addthinS\dicIPA{{r}{au}{ð}{s}{\textepsilon}{h}{\textsubring{d}}{\textscy}{\textsubring{r}}}\addthinS\dicPos{adj}[1]\dicFlx{}[-10]\addthinS\dicSynonym{stilltur}\addthinS\dicDirectTranslationCS{usedlý, klidný, konzervativní}
\end{xtolerant}
\dicEntry[ráðskona] \textls[20]{\dicTerm{ráðs··kon|a} \dicIPA{{r}{au}{ð}{s}{k\smash{\textsuperscript{h}}}{\textopeno}{n}{a}} \dicPos{f}[1] \dicFlx{(‑u, ‑ur)}[27] \dicDirectTranslationCS{hospodářka, do\nobreak movnice, školnice}}
\dicEntry[ráðsmaður] \dicTerm{ráðs··|maður} \dicIPA{{r}{au}{ð}{s}{m}{a}{ð}{\textscy}{\textsubring{r}}} \dicPos{m}[13] \dicFlx{(‑manns, ‑menn)}[2] \dicDirectTranslationCS{hospodář(ka), správce, správkyně (kolejí ap.), domovník, domovnice} \dicExampleIS{ráðsmaður heimavistar} \dicExampleCS{správce internátu}
\dicEntry[ráðstafa] \dicTerm{ráð··staf|a} \dicIPA{{r}{au}{ð}{s}{\textsubring{d}}{a}{v}{a}} \dicPos{v}[1] \dicFlx{(‑aði)}[28] \dicFlx{dat} \dicSynonym{skipa\smash{\textsuperscript{1}}} \dicDirectTranslationCS{disponovat, mít k~dispozici, (volně) nakládat} \dicExampleIS{ráðstafa peningunum} \dicExampleCS{disponovat penězi}
\dicEntry[ráðstefna] \dicTerm{ráð··stefn|a} \dicIPA{{r}{au}{ð}{s}{\textsubring{d}}{\textepsilon}{\textsubring{b}}{n}{a}} \dicPos{f}[1] \dicFlx{(‑u, ‑ur)}[7] \dicDirectTranslationCS{konference, sjezd, kongres} \dicExampleIS{ráðstefna um málfræði} \dicExampleCS{konference o~gramatice}
\dicEntry[ráðstefnutúlkun] \dicTerm{ráð·stefnu··túlk|un} \dicIPA{{r}{au}{ð}{s}{\textsubring{d}}{\textepsilon}{\textsubring{b}}{n}{\textscy}{t\smash{\textsuperscript{h}}}{u}{\textsubring{l}}{\r{g}}{\textscy}{\textsubring{n}}} \dicPos{f}[7] \dicFlx{(‑unar, ‑anir)}[8] \dicDirectTranslationCS{konferenční tlumočení}
\dicEntry[ráðstöfun] \dicTerm{ráð··|stöfun} \dicsymFrequent\  \dicIPA{{r}{au}{ð}{s}{\textsubring{d}}{\oe}{v}{\textscy}{\textsubring{n}}} \dicPos{f}[7] \dicFlx{(‑stöfunar, ‑stafanir)}[11] \textbf{1.} \dicSynonym*{það að ráðstafa} \dicDirectTranslationCS{dispozice, disponování};  \dicPhraseIS{hafa e‑ð til ráðstöfunar} \dicDirectTranslationCS{mít (co) k~dispozici}  \textbf{2.} \dicSynonym{framkvæmd} \dicDirectTranslationCS{opatření, krok} \dicExampleIS{tímabundin ráðstöfun} \dicExampleCS{dočasné opatření};  \dicPhraseIS{gera ráðstafanir} \dicDirectTranslationCS{učinit opatření}
\dicEntry[ráðstöfunarfé] \dicTerm{ráð·stöfunar··|fé} \dicIPA{{r}{au}{ð}{s}{\textsubring{d}}{\oe}{v}{\textscy}{n}{a}{\textsubring{r}}{f}{j}{\textepsilon}} \dicPos{n}[3] \dicFlx{(‑fjár)}[5] \dicDirectTranslationCS{disponibilní prostředky}
\dicEntry[ráðunautur] \dicTerm{ráðu··naut|ur} \dicIPA{{r}{au}{\textlengthmark}{ð}{\textscy}{n}{\oe i}{\textsubring{d}}{\textscy}{\textsubring{r}}} \dicPos{m}[6] \dicFlx{(‑s\,/\addthin ‑ar, ‑ar)}[65] \dicDirectTranslationCS{poradce, poradkyně} \dicExampleIS{ráðunautur forsetans} \dicExampleCS{poradkyně prezidenta}
\dicEntry[ráðuneyti] \dicTerm{ráðu··neyti} \dicsymFrequent\  \dicIPA{{r}{au}{\textlengthmark}{ð}{\textscy}{n}{ei}{\textsubring{d}}{\textsci}} \dicPos{n}[2] \dicFlx{(‑s, ‑)}[14] \dicDirectTranslationCS{ministerstvo} \dicExampleIS{Þau heyra undir fleiri en eitt ráðuneyti.} \dicExampleCS{Spadají pod více než jedno ministerstvo.}
\dicEntry[ráðuneytisstjóri] \dicTerm{ráðu·neytis··stjór|i} \dicIPA{{r}{au}{\textlengthmark}{ð}{\textscy}{n}{ei}{\textsubring{d}}{\textsci}{s}{\textsubring{d}}{j}{ou}{r}{\textsci}} \dicPos{m}[1] \dicFlx{(‑a, ‑ar)}[1] \dicFieldCat{pol.} \dicDirectTranslationCS{náměstek\,/\addthin náměstkyně ministra\,/\addthin ministryně}
\dicEntry[ráðvandur] \dicTerm{ráð··|vandur} \dicIPA{{r}{au}{ð}{v}{a}{n}{\textsubring{d}}{\textscy}{\textsubring{r}}} \dicPos{adj}[2] \dicFlx{(f ‑vönd)}[15] \dicSynonym{heiðarlegur} \dicDirectTranslationCS{čestný, poctivý}
\dicEntry[ráðvilltur] \dicTerm{ráð··villtur} \dicIPA{{r}{au}{ð}{v}{\textsci}{\textsubring{l}}{\textsubring{d}}{\textscy}{\textsubring{r}}} \dicPos{adj}[1]\dicFlx{}[-10] \dicSynonym{ruglaður} \dicDirectTranslationCS{nevědoucí si rady, bezradný}
\dicEntry[ráðþrota] \dicTerm{ráð··þrota} \dicIPA{{r}{au}{ð}{\texttheta}{r}{\textopeno}{\textsubring{d}}{a}} \dicPos{adj}[13] \dicFlx{indecl}[1] \dicSynonym{ráðalaus} \dicDirectTranslationCS{bezradný, bezmocný}
\dicEntry[ráf] \dicTerm{ráf} \dicIPA{{r}{au}{\textlengthmark}{f}} \dicPos{n}[2] \dicFlx{(‑s)}[2] \dicSynonym{reik} \dicDirectTranslationCS{bloumání, potloukání se}
\dicEntry[ráfa] \dicTerm{ráf|a} \dicIPA{{r}{au}{\textlengthmark}{v}{a}} \dicPos{v}[1] \dicFlx{(‑aði)}[44] \dicSynonym{rölta} \dicDirectTranslationCS{bloumat, potulovat se} \dicExampleIS{ráfa fram og til baka um göturnar} \dicExampleCS{potulovat se po ulicích}
\dicEntry[rák] \dicTerm{rák} \dicIPA{{r}{au}{\textlengthmark}{\r{g}}} \dicPos{f}[7] \dicFlx{(‑ar, ‑ir)}[1] \textbf{1.} \dicSynonym{rönd} \dicDirectTranslationCS{proužek, pruh}  \textbf{2.} \dicSynonym{skora\smash{\textsuperscript{1}}} \dicDirectTranslationCS{vrub, rýha}
\dicEntry[rákum] \dicTerm{rákum} \dicIPA{{r}{au}{\textlengthmark}{\r{g}}{\textscy}{\textsubring{m}}} \dicPos{v} \dicFlx{ind pf pl 1 pers} \dicLink{reka\smash{\textsuperscript{2}}}
\dicEntry[rámur] \dicTerm{rámur} \dicIPA{{r}{au}{\textlengthmark}{m}{\textscy}{\textsubring{r}}} \dicPos{adj}[1]\dicFlx{}[-1] \dicSynonym{hás} \dicDirectTranslationCS{zastřený, chraplavý} \dicExampleIS{rám rödd} \dicExampleCS{chraplavý hlas}
\dicEntry[rán] \dicTerm{rán} \dicIPA{{r}{au}{\textlengthmark}{\textsubring{n}}} \dicPos{n}[2] \dicFlx{(‑s, ‑)}[5] \dicDirectTranslationCS{loupež, oloupení};  \dicPhraseIS{fremja rán} \dicDirectTranslationCS{spáchat loupež};  \dicPhraseIS{vopnað rán} \dicDirectTranslationCS{ozbrojená loupež}
\dicEntry[rándýr] \dicTerm{rán··dýr\smash{\textsuperscript{1}}} \dicIPA{{r}{au}{n}{\textsubring{d}}{i}{\textsubring{r}}} \dicPos{n}[2] \dicFlx{(‑s, ‑)}[5] \dicFieldCat{zool.} \dicDirectTranslationCS{šelma, dravec, predátor} \textit{(l.~{\textLA{Carnivora}})}
\dicEntry[rándýr] \dicTerm{rán··dýr\smash{\textsuperscript{2}}} \dicIPA{{r}{au}{n}{\textsubring{d}}{i}{\textsubring{r}}} \dicPos{adj}[5] \dicFlx{(f ‑)}[10] \dicSynonym*{fokdýr} \dicDirectTranslationCS{velmi drahý}
\dicEntry[ránfugl] \dicTerm{rán··fugl} \dicIPA{{r}{au}{n}{f}{\textscy}{\r{g}}{\textsubring{l}}} \dicPos{m}[4] \dicFlx{(‑s, ‑ar)}[3] \dicFieldCat{zool.} \dicDirectTranslationCS{dravec} \dicIndirectTranslationCS{(masožravý pták)} \textit{(l.~{\textLA{Falconiformes}})}
\dicEntry[ránsfengur] \dicTerm{ráns··feng|ur} \dicIPA{{r}{au}{n}{s}{f}{ei}{\ng}{\r{g}}{\textscy}{\textsubring{r}}} \dicPos{m}[9] \dicFlx{(‑s\,/\addthin ‑jar)}[27] \dicDirectTranslationCS{lup, kořist}
\dicEntry[rányrkja] \dicTerm{rán··yrkj|a} \dicIPA{{r}{au}{\textlengthmark}{n}{\textsci}{\textsubring{r}}{\r{\textObardotlessj}}{a}} \dicPos{f}[1] \dicFlx{(‑u)}[5] \dicDirectTranslationCS{bezohledné užívání přírodních zdrojů}
\dicEntry[rápa] \dicTerm{ráp|a} \dicIPA{{r}{au}{\textlengthmark}{\textsubring{b}}{a}} \dicPos{v}[1] \dicFlx{(‑aði)}[44] \dicSynonym*{flækjast} \dicDirectTranslationCS{potloukat se, toulat se}
\dicEntry[rás] \dicTerm{rás} \dicsymFrequent\  \dicIPA{{r}{au}{\textlengthmark}{s}} \dicPos{f}[7] \dicFlx{(‑ar, ‑ir)}[1] \textbf{1.} \dicSynonym{hlaup} \dicDirectTranslationCS{běh, sprint};  \dicPhraseIS{taka á rás} \dicSynonym*{taka til  fótanna} \dicDirectTranslationCS{rozběhnout se}  \textbf{2.} \dicSynonym{rennsli} \dicDirectTranslationCS{běh, tok (řeky, času ap.)};  \dicPhraseIS{í tímans rás} \dicFlx{adv} \dicDirectTranslationCS{v~průběhu času}  \textbf{3.} \dicSynonym{renna\smash{\textsuperscript{1}}} \dicDirectTranslationCS{koryto, kanál}  \textbf{4.} \dicDirectTranslationCS{kanál (rádiový ap.)} \dicExampleIS{útvarpsrás} \dicExampleCS{rozhlasový kanál};  \dicPhraseIS{Rás 1, Rás 2} \dicDirectTranslationCS{Rás 1, Rás 2} \dicIndirectTranslationCS{(islandské rozhlasové stanice)}
\dicEntry[rásegl] \dicTerm{rá··segl} \dicIPA{{r}{au}{\textlengthmark}{s}{\textepsilon}{\r{g}}{\textsubring{l}}} \dicPos{n}[2] \dicFlx{(‑s, ‑)}[5] \dicFieldCat{nám.} \dicDirectTranslationCS{ráhnová plachta}
\dicEntry[redda] \dicTerm{redd|a} \dicIPA{{r}{\textepsilon}{\textsubring{d}}{\textlengthmark}{a}} \dicPos{v}[1] \dicFlx{(‑aði)}[22] \dicFlx{dat} \dicLangCat{hovor.} \dicSynonym{bjarga} \dicDirectTranslationCS{helfnout, vypomoct, zachránit} \dicExampleIS{redda honum um peninga} \dicExampleCS{vypomoct mu finančně};  \dicIdiom{reddast}{ \dicPhraseIS{reddast}} \dicFlx{refl} \dicLangCat{hovor.} \dicDirectTranslationCS{srovnat se, vyřešit se} \dicExampleIS{þetta reddast} \dicExampleCS{to se vyřeší}
\dicEntry[reður] \dicTerm{reð|ur} \dicIPA{{r}{\textepsilon}{\textlengthmark}{ð}{\textscy}{\textsubring{r}}} \dicPos{m}[5] \dicFlx{(‑urs, ‑rar)}[1] \dicSynonym{getnaðarlimur} \dicDirectTranslationCS{penis, pyj}
\dicEntry[reðurhúfa] \dicTerm{reður··húf|a} \dicIPA{{r}{\textepsilon}{\textlengthmark}{ð}{\textscy}{\textsubring{r}}{h}{u}{v}{a}} \dicPos{f}[1] \dicFlx{(‑u, ‑ur)}[19] \dicFieldCat{anat.} \dicDirectTranslationCS{žalud} \dicIndirectTranslationCS{(konečná část penisu nebo poštěváčku)}
\dicEntry[reffilegur] \dicTerm{reffi··legur} \dicIPA{{r}{\textepsilon}{f}{\textlengthmark}{\textsci}{l}{\textepsilon}{\textbabygamma}{\textscy}{\textsubring{r}}} \dicPos{adj}[1]\dicFlx{}[-8] \dicSynonym{myndarlegur} \dicDirectTranslationCS{vznešený, velkolepý}
\dicEntry[refhvörf] \dicTerm{ref··hvörf} \dicIPA{{r}{\textepsilon}{v}{k\smash{\textsuperscript{h}}}{v}{\oe}{r}{f}} \dicPos{n}[2] \dicFlx{pl}[9] \dicFieldCat{lit.} \dicDirectTranslationCS{oxymóron, protimluv}
\dicEntry[refjar] \dicTerm{refjar} \dicIPA{{r}{\textepsilon}{v}{j}{a}{\textsubring{r}}} \dicPos{f}[12] \dicFlx{pl}[1] \dicSynonym{undanbrögð} \dicDirectTranslationCS{vytáčky, vykrucování};  \dicPhraseIS{vera með refjar} \dicDirectTranslationCS{vytáčet se, vykrucovat se}
\dicEntry[refsa] \dicTerm{refs|a} \dicIPA{{r}{\textepsilon}{f}{s}{a}} \dicPos{v}[1] \dicFlx{(‑aði)}[22] \dicFlx{dat} \dicSynonym{hegna} \dicDirectTranslationCS{(po)trestat, ztrestat} \dicExampleIS{refsa innbrotsþjófi} \dicExampleCS{potrestat lupiče};  \dicPhraseIS{refsa e‑m fyrir e‑ð} \dicDirectTranslationCS{potrestat (koho) za (co)}
\begin{xtolerant}{}{1pt}
\dicEntry[refsbróðir] \dicTerm{refs··|bróðir}\addthinS\dicIPA{{r}{\textepsilon}{f}{s}{\textsubring{b}}{r}{ou}{ð}{\textsci}{\textsubring{r}}}\addthinS\dicPos{m}[13]\addthinS\dicFlx{(‑bróður, ‑bræður)}[4] \dicFieldCat{zool.} \dicSynonym{sjak\-ali} \dicDirectTranslationCS{šakal, šakal obecný} \textit{(l.~{\textLA{Canis aureus}})}
\end{xtolerant}
\dicEntry[refsiaðgerð] \dicTerm{refsi··að·gerð} \dicIPA{{r}{\textepsilon}{f}{s}{\textsci}{a}{ð}{\r{\textObardotlessj}}{\textepsilon}{r}{\texttheta}} \dicPos{f}[7] \dicFlx{(‑ar, ‑ir)}[1] \dicFieldCat{práv.} \dicDirectTranslationCS{sankce}
\dicEntry[refsiákvæði] \dicTerm{refsi··á·kvæði} \dicIPA{{r}{\textepsilon}{f}{s}{\textsci}{au}{k\smash{\textsuperscript{h}}}{v}{a}{i}{ð}{\textsci}} \dicPos{n}[2] \dicFlx{(‑s, ‑)}[14] \dicFieldCat{práv.} \dicDirectTranslationCS{trestní ustanovení}
\dicEntry[refsidómur] \dicTerm{refsi··dóm|ur} \dicIPA{{r}{\textepsilon}{f}{s}{\textsci}{\textsubring{d}}{ou}{m}{\textscy}{\textsubring{r}}} \dicPos{m}[6] \dicFlx{(‑s, ‑ar)}[10] \dicFieldCat{práv.} \dicDirectTranslationCS{trestní výrok}
\dicEntry[refsilöggjöf] \dicTerm{refsi··lög·|gjöf} \dicIPA{{r}{\textepsilon}{f}{s}{\textsci}{l}{\oe}{\textbabygamma}{\r{\textObardotlessj}}{\oe}{f}} \dicPos{f}[7] \dicFlx{(‑gjafar)}[19] \dicFieldCat{práv.} \dicDirectTranslationCS{trestní zákon}
\dicEntry[refsing] \dicTerm{refs··ing} \dicIPA{{r}{\textepsilon}{f}{s}{i}{\ng}{\r{g}}} \dicPos{f}[4] \dicFlx{(‑ar, ‑ar)}[5] \dicSynonym{hegning} \dicDirectTranslationCS{trest, potrestání} \dicExampleIS{líkamleg refsing} \dicExampleCS{tělesný trest}
\dicEntry[refsiréttur] \dicTerm{refsi··rétt|ur} \dicIPA{{r}{\textepsilon}{f}{s}{\textsci}{r}{j}{\textepsilon}{h}{\textsubring{d}}{\textscy}{\textsubring{r}}} \dicPos{m}[10] \dicFlx{(‑ar)}[7] \dicFieldCat{práv.} \dicDirectTranslationCS{trestní právo}
\dicEntry[refsiverður] \dicTerm{refsi··verður} \dicIPA{{r}{\textepsilon}{f}{s}{\textsci}{v}{\textepsilon}{r}{ð}{\textscy}{\textsubring{r}}} \dicPos{adj}[2]\dicFlx{}[-1] \dicFieldCat{práv.} \dicDirectTranslationCS{trestný} \dicExampleIS{refsiverður verknaður} \dicExampleCS{trestná činnost}
\dicEntry[refur] \dicTerm{ref|ur} \dicIPA{{r}{\textepsilon}{\textlengthmark}{v}{\textscy}{\textsubring{r}}} \dicPos{m}[6] \dicFlx{(‑s, ‑ir\,/\addthin ‑ar)}[70] \textbf{1.} \dicFieldCat{zool.} \dicDirectTranslationCS{liška, liška obecná} \textit{(l.~{\textLA{Vulpes vulpes}})}  \dicsymPhoto\   \textbf{2.} \dicLangCat{přen.} \dicSynonym*{bragðvís maður} \dicDirectTranslationCS{lišák} \dicIndirectTranslationCS{(vychytralý člověk)};  \dicPhraseIS{til þess eru refirnir\,/\addthin refarnir skornir} \dicDirectTranslationCS{záměrem je, aby};  \dicProverb\  \dicPhraseIS{Illt er að ginna gamlan ref.} \dicLangCat{přís.} \dicDirectTranslationCS{Lišáka rohlíkem neopiješ.}
\dicFigure{21297.jpg}{Refur}{Refur - Haggblom, Lisa, Biolib, PD}
\dicEntry[regla] \dicTerm{regl|a} \dicsymFrequent\  \dicIPA{{r}{\textepsilon}{\r{g}}{l}{a}} \dicPos{f}[1] \dicFlx{(‑u, ‑ur)}[19] \textbf{1.} \dicSynonym{fyrirmæli} \dicDirectTranslationCS{pravidlo, zásada} \dicExampleIS{fara eftir reglum} \dicExampleCS{držet se pravidel}  \textbf{2.} \dicSynonym{lögmál} \dicDirectTranslationCS{zákon, norma} \dicExampleIS{Gangur tunglanna lýtur föstum reglum.} \dicExampleCS{Pohyb měsíců se řídí pevnými zákony.}  \textbf{3.} \dicSynonym{félagsskapur} \dicDirectTranslationCS{řád (náboženský ap.)}
\dicEntry[reglubundinn] \dicTerm{reglu··bundinn} \dicIPA{{r}{\textepsilon}{\r{g}}{l}{\textscy}{\textsubring{b}}{\textscy}{n}{\textsubring{d}}{\textsci}{\textsubring{n}}} \dicPos{adj}[6]\dicFlx{}[-2] \dicDirectTranslationCS{pravidelný, regulérní} \dicExampleIS{reglubundið viðhald} \dicExampleCS{pravidelná údržba}
\dicEntry[reglufastur] \dicTerm{reglu··|fastur} \dicIPA{{r}{\textepsilon}{\r{g}}{l}{\textscy}{f}{a}{s}{\textsubring{d}}{\textscy}{\textsubring{r}}} \dicPos{adj}[1] \dicFlx{(f ‑föst)}[11] \dicDirectTranslationCS{ukázněný, disciplinovaný}
\dicEntry[reglugerð] \dicTerm{reglu··gerð} \dicIPA{{r}{\textepsilon}{\r{g}}{l}{\textscy}{\r{\textObardotlessj}}{\textepsilon}{r}{\texttheta}} \dicPos{f}[7] \dicFlx{(‑ar, ‑ir)}[1] \dicFieldCat{práv.} \dicDirectTranslationCS{ustanovení, nařízení, předpis, směrnice}
\dicEntry[reglugerðasöfn] \dicTerm{reglu·gerða··söfn} \dicIPA{{r}{\textepsilon}{\r{g}}{l}{\textscy}{\r{\textObardotlessj}}{\textepsilon}{r}{ð}{a}{s}{\oe}{\textsubring{b}}{\textsubring{n}}} \dicPos{n}[2] \dicFlx{pl}[9] \dicFieldCat{práv.} \dicDirectTranslationCS{sbírka předpisů\,/\addthin nařízení}
\dicEntry[reglulega] \dicTerm{reglu··lega} \dicsymFrequent\  \dicIPA{{r}{\textepsilon}{\r{g}}{l}{\textscy}{l}{\textepsilon}{\textbabygamma}{a}} \dicPos{adv} \dicDirectTranslationCS{pravidelně} \dicExampleIS{borða reglulega} \dicExampleCS{jíst pravidelně}
\dicEntry[reglulegur] \dicTerm{reglu··legur} \dicsymFrequent\  \dicIPA{{r}{\textepsilon}{\r{g}}{l}{\textscy}{l}{\textepsilon}{\textbabygamma}{\textscy}{\textsubring{r}}} \dicPos{adj}[1]\dicFlx{}[-8] \textbf{1.} \dicSynonym*{kerfislegur} \dicDirectTranslationCS{pravidelný, souměrný} \dicIndirectTranslationCS{(stejnoměrně se opakující)} \dicExampleIS{reglulegur hjartsláttur} \dicExampleCS{pravidelný tep}  \textbf{2.} \dicSynonym{stöðugur} \dicDirectTranslationCS{pravidelný, regulérní, běžný}
\dicEntry[reglumaður] \dicTerm{reglu··|maður} \dicIPA{{r}{\textepsilon}{\r{g}}{l}{\textscy}{m}{a}{ð}{\textscy}{\textsubring{r}}} \dicPos{m}[13] \dicFlx{(‑manns, ‑menn)}[2] \textbf{1.} \dicSynonym*{hófsmaður} \dicDirectTranslationCS{umírněný člověk (v~pití alkoholu ap.)}  \textbf{2.} \dicSynonym*{reglu\-fastur maður} \dicDirectTranslationCS{spořádaný člověk}
\dicEntry[reglusamur] \dicTerm{reglu··|samur} \dicIPA{{r}{\textepsilon}{\r{g}}{l}{\textscy}{s}{a}{m}{\textscy}{\textsubring{r}}} \dicPos{adj}[1] \dicFlx{(f ‑söm)}[2] \textbf{1.} \dicSynonym{reglufastur} \dicDirectTranslationCS{spořádaný, řádný}  \textbf{2.} \dicSynonym{bindindissamur} \dicDirectTranslationCS{střídmý, umírněný}
\dicEntry[reglusemi] \dicTerm{reglu··sem|i} \dicIPA{{r}{\textepsilon}{\r{g}}{l}{\textscy}{s}{\textepsilon}{m}{\textsci}} \dicPos{f}[2] \dicFlx{(‑i)}[2] \textbf{1.} \dicSynonym*{reglufesta} \dicDirectTranslationCS{spořádanost}  \textbf{2.} \dicSynonym*{bindindissemi} \dicDirectTranslationCS{střídmost, zdrženlivost}
\dicEntry[reglustika] \dicTerm{reglu··stik|a}\dicTerm{, reglustrika} \dicIPA{{r}\-{\textepsilon}\-{\r{g}}\-{l}\-{\textscy}\-{s}\-{\textsubring{d}}\-{\textsci}\-{\r{g}}\-{a}\-} \dicPos{f}[1] \dicFlx{(‑u, ‑ur)}[19] \dicDirectTranslationCS{pravítko}
\dicEntry[reglustrika] \dicTerm{reglu··strik|a} \dicIPA{{r}{\textepsilon}{\r{g}}{l}{\textscy}{s}{\textsubring{d}}{r}{\textsci}{\r{g}}{a}} \dicPos{f}[1] \dicFlx{(‑u, ‑ur)}[7] \dicLink{reglustika}
\dicEntry[regn] \dicTerm{regn} \dicsymFrequent\  \dicIPA{{r}{\textepsilon}{\r{g}}{\textsubring{n}}} \dicPos{n}[2] \dicFlx{(‑s)}[2] \dicSynonym{rigning} \dicDirectTranslationCS{déšť} \dicExampleIS{mjúkt regn} \dicExampleCS{slabý déšť}
\dicEntry[regnbogasilungur] \dicTerm{regn·boga··sil·ung|ur} \dicIPA{{r}\-{\textepsilon}\-{\r{g}}\-{\textsubring{n}}\-{\textsubring{b}}\-{\textopeno}\-{\textbabygamma}\-{a}\-{s}\-{\textsci}\-{l}\-{u}\-{\ng}\-{\r{g}}\-{\textscy}\-{\textsubring{r}}\-} \dicPos{m}[6] \dicFlx{(‑s, ‑ar)}[8] \dicFieldCat{zool.} \dicDirectTranslationCS{pstruh duhový} \textit{(l.~{\textLA{Oncorhynchus mykiss}})}  \dicsymPhoto\ 
\dicFigure{15999.jpg}{Regnbogasilungur}{Regnbogasilungur - Zapletal Tomáš, Biolib, PD}
\dicEntry[regnbogi] \dicTerm{regn··bog|i} \dicIPA{{r}{\textepsilon}{\r{g}}{\textsubring{n}}{\textsubring{b}}{\textopeno i}{j}{\textsci}} \dicPos{m}[1] \dicFlx{(‑a, ‑ar)}[1] \dicDirectTranslationCS{duha}
\dicEntry[regngalli] \dicTerm{regn··gall|i} \dicIPA{{r}{\textepsilon}{\r{g}}{\textsubring{n}}{\r{g}}{a}{l}{\textsci}} \dicPos{m}[1] \dicFlx{(‑a, ‑ar)}[8] \dicDirectTranslationCS{nepromokavé oblečení}
\dicEntry[regnhlíf] \dicTerm{regn··hlíf} \dicIPA{{r}{\textepsilon}{\r{g}}{\textsubring{n}}{\textsubring{l}}{i}{f}} \dicPos{f}[4] \dicFlx{(‑ar, ‑ar)}[1] \dicDirectTranslationCS{deštník}
\dicEntry[regnkápa] \dicTerm{regn··káp|a} \dicIPA{{r}{\textepsilon}{\r{g}}{\textsubring{n}}{k\smash{\textsuperscript{h}}}{au}{\textsubring{b}}{a}} \dicPos{f}[1] \dicFlx{(‑u, ‑ur)}[19] \dicDirectTranslationCS{pláštěnka}
\dicEntry[regnmælir] \dicTerm{regn··mæl|ir} \dicIPA{{r}{\textepsilon}{\r{g}}{n}{m}{a}{i}{l}{\textsci}{\textsubring{r}}} \dicPos{m}[7] \dicFlx{(‑is, ‑ar)}[1] \dicFieldCat{meteo.} \dicDirectTranslationCS{srážkoměr}
\dicEntry[regnskógur] \dicTerm{regn··skóg|ur} \dicIPA{{r}{\textepsilon}{\r{g}}{\textsubring{n}}{s}{\r{g}}{ou}{\textscy}{\textsubring{r}}} \dicPos{m}[6] \dicFlx{(‑ar, ‑ar)}[56] \dicDirectTranslationCS{deštný prales}
\dicEntry[regnskúr] \dicTerm{regn··skúr} \dicIPA{{r}{\textepsilon}{\r{g}}{\textsubring{n}}{s}{\r{g}}{u}{\textsubring{r}}} \dicPos{f}[7] \dicFlx{(‑ar, ‑ir)}[1] \dicDirectTranslationCS{přeháňka, přeprška}
\dicEntry[regnvatn] \dicTerm{regn··vatn} \dicIPA{{r}{\textepsilon}{\r{g}}{n}{v}{a}{h}{\textsubring{d}}{\textsubring{n}}} \dicPos{n}[2] \dicFlx{(‑s)}[2] \dicDirectTranslationCS{dešťová voda}
\dicEntry[reið] \dicTerm{reið\smash{\textsuperscript{1}}} \dicIPA{{r}{ei}{\textlengthmark}{\texttheta}} \dicPos{f}[7] \dicFlx{(‑ar, ‑ir)}[1] \textbf{1.} \dicSynonym*{það að ríða} \dicDirectTranslationCS{jízda (na koni ap.)}  \textbf{2.} \dicLangCat{hrub.} \dicSynonym{samræði} \dicDirectTranslationCS{jízda, číslo} \dicIndirectTranslationCS{(soulož)}
\dicEntry[reið] \dicTerm{reið\smash{\textsuperscript{2}}} \dicIPA{{r}{ei}{\textlengthmark}{\texttheta}} \dicPos{v} \dicFlx{ind pf sg 1 pers} \dicLink{ríða}
\dicEntry[reiða] \dicTerm{reið|a\smash{\textsuperscript{1}}} \dicIPA{{r}{ei}{\textlengthmark}{ð}{a}} \dicPos{f}[1] \dicFlx{(‑u, ‑ur)}[19] \dicDirectTranslationCS{pořádek, uspořádání};  \dicPhraseIS{hafa e‑ð til reiðu} \dicDirectTranslationCS{mít (co) k~dispozici};  \dicPhraseIS{henda reiður á e‑u} \dicDirectTranslationCS{získat přehled o~(čem)};  \dicPhraseIS{vera til reiðu} \dicDirectTranslationCS{být k~dispozici}
\dicEntry[reiða] \dicTerm{rei|ða\smash{\textsuperscript{2}}} \dicsymFrequent\  \dicIPA{{r}{ei}{\textlengthmark}{ð}{a}} \dicPos{v}[2] \dicFlx{(‑ddi, ‑tt)}[178] \dicFlx{acc} \textbf{1.} \dicSynonym{flytja} \dicDirectTranslationCS{dopravit\,/\addthin vést na koni\,/\addthin kole} \dicExampleIS{reiða bagga á hjólinu} \dicExampleCS{vést balík na kole};  \dicPhraseIS{reiða e‑n} \dicDirectTranslationCS{svézt (koho) (na koni, kole ap.)}  \textbf{2.} \dicSynonym{greiða\smash{\textsuperscript{2}}} \dicDirectTranslationCS{vydat, vyplatit (peníze ap.)};  \dicPhraseIS{reiða fé af hendi} \dicDirectTranslationCS{zaplatit hotově};  \dicIdiom{reiða}[af]{ \dicPhraseIS{e‑u\,/\addthin e‑m reiðir vel af}} \dicFlx{impers} \dicDirectTranslationCS{(čemu\,/\addthin komu) se dobře daří};  \dicIdiom{reiða}[á]{ \dicPhraseIS{reiða sig á e‑n\,/\addthin e‑ð}} \dicDirectTranslationCS{spolehnout\,/\addthin spoléhat se na (koho\,/\addthin co)};  \dicIdiom{reiða}[fram]{ \dicPhraseIS{reiða e‑ð fram}} {\textbf{a.}} \dicDirectTranslationCS{podat (co), předložit (co) (jídlo na stůl ap.)} \dicExampleIS{reiða kaffi fram} \dicExampleCS{podat kávu};  {\textbf{b.}} \dicSynonym{borga} \dicDirectTranslationCS{vydat (co), zaplatit (co)};  \dicIdiom{reiða}[upp]{ \dicPhraseIS{reiða e‑ð upp}} \dicDirectTranslationCS{napřáhnout (co)} \dicExampleIS{reiða kylfu upp} \dicExampleCS{napřáhnout obušek}
\dicEntry[reiðast] \dicTerm{rei|ðast} \dicIPA{{r}{ei}{\textlengthmark}{ð}{a}{s}{\textsubring{d}}} \dicPos{v}[2] \dicFlx{(‑ddist, ‑ðst)}[210] \dicFlx{refl dat} \dicSynonym*{verða reiður} \dicDirectTranslationCS{rozzlobit se, rozezlít se, rozčílit se} \dicExampleIS{reiðast e‑m} \dicExampleCS{rozzlobit se na (koho)};  \dicPhraseIS{reiðast við e‑n} \dicDirectTranslationCS{rozzlobit se na (koho)};  \dicPhraseIS{reiðast yfir e‑u, reiðast út af e‑u} \dicDirectTranslationCS{rozčílit se kvůli (čemu)}
\dicEntry[reiðgata] \dicTerm{reið··|gata} \dicIPA{{r}{ei}{ð}{\r{g}}{a}{\textsubring{d}}{a}} \dicPos{f}[1] \dicFlx{(‑götu, ‑götur)}[14] \dicDirectTranslationCS{jezdecká stezka}
\dicEntry[reiðhestur] \dicTerm{reið··hest|ur} \dicIPA{{r}{ei}{\textlengthmark}{\texttheta}{h}{\textepsilon}{s}{\textsubring{d}}{\textscy}{\textsubring{r}}} \dicPos{m}[6] \dicFlx{(‑s, ‑ar)}[4] \dicDirectTranslationCS{jezdecký kůň}
\dicEntry[reiðhjól] \dicTerm{reið··hjól} \dicIPA{{r}{ei}{\texttheta}{\c{c}}{ou}{\textsubring{l}}} \dicPos{n}[2] \dicFlx{(‑s, ‑)}[5] \dicDirectTranslationCS{(jízdní) kolo, bicykl}
\dicEntry[reiði] \dicTerm{reið|i\smash{\textsuperscript{1}}} \dicIPA{{r}{ei}{\textlengthmark}{ð}{\textsci}} \dicPos{m}[1] \dicFlx{(‑a, ‑ar)}[1] \textbf{1.} \dicFieldCat{nám.} \dicDirectTranslationCS{takeláž}  \textbf{2.} \dicDirectTranslationCS{podocasník} \dicIndirectTranslationCS{(řemen postroje)}
\dicEntry[reiði] \dicTerm{reið|i\smash{\textsuperscript{2}}} \dicsymFrequent\  \dicIPA{{r}{ei}{\textlengthmark}{ð}{\textsci}} \dicPos{f}[3] \dicFlx{(‑i)}[3] \dicDirectTranslationCS{hněv, zlost, vztek} \dicExampleIS{finna fyrir mikilli reiði} \dicExampleCS{cítit velký vztek};  \dicPhraseIS{reiði út í e‑n} \dicDirectTranslationCS{vztek na (koho)};  \dicPhraseIS{reita e‑n til reiði} \dicDirectTranslationCS{rozhněvat (koho)};  \dicIdiom{reiði}{ \dicPhraseIS{reiðinnar býsn\,/\addthin ósköp af e‑u}} \dicDirectTranslationCS{obrovské množství (čeho)}
\dicEntry[reiðikast] \dicTerm{reiði··|kast} \dicIPA{{r}{ei}{\textlengthmark}{ð}{\textsci}{k\smash{\textsuperscript{h}}}{a}{s}{\textsubring{d}}} \dicPos{n}[2] \dicFlx{(‑kasts, ‑köst)}[8] \dicDirectTranslationCS{záchvat\,/\addthin výbuch vzteku}
\dicEntry[reiðilegur] \dicTerm{reiði··legur} \dicIPA{{r}{ei}{\textlengthmark}{ð}{\textsci}{l}{\textepsilon}{\textbabygamma}{\textscy}{\textsubring{r}}} \dicPos{adj}[1]\dicFlx{}[-8] \dicDirectTranslationCS{rozzlobený, rozhněvaný, rozčilený}
\dicEntry[reiðileysi] \dicTerm{reiði··leysi} \dicIPA{{r}{ei}{\textlengthmark}{ð}{\textsci}{l}{ei}{s}{\textsci}} \dicPos{n}[2] \dicFlx{(‑s)}[20] \dicSynonym*{umhirðuleysi} \dicDirectTranslationCS{opominutí, zanedbání}
\dicEntry[reiðisvipur] \dicTerm{reiði··svip|ur} \dicIPA{{r}{ei}{\textlengthmark}{ð}{\textsci}{s}{v}{\textsci}{\textsubring{b}}{\textscy}{\textsubring{r}}} \dicPos{m}[9] \dicFlx{(‑s, ‑ir)}[8] \dicDirectTranslationCS{rozzlobený výraz}
\dicEntry[reiðmaður] \dicTerm{reið··|maður} \dicIPA{{r}{ei}{ð}{m}{a}{ð}{\textscy}{\textsubring{r}}} \dicPos{m}[13] \dicFlx{(‑manns, ‑menn)}[2] \dicDirectTranslationCS{jezdec, jezdkyně (na koni ap.)} \dicExampleIS{góður reiðmaður} \dicExampleCS{dobrý jezdec}
\dicEntry[reiðskjóti] \dicTerm{reið··skjót|i} \dicIPA{{r}{ei}{ð}{s}{\r{\textObardotlessj}}{ou}{\textsubring{d}}{\textsci}} \dicPos{m}[1] \dicFlx{(‑a, ‑ar)}[1] \dicSynonym{reiðhestur} \dicDirectTranslationCS{jezdecké zvíře}
\dicEntry[reiðubúinn] \dicTerm{reiðu··búinn} \dicsymFrequent\  \dicIPA{{r}{ei}{\textlengthmark}{ð}{\textscy}{\textsubring{b}}{u}{\textsci}{\textsubring{n}}} \dicPos{adj}[6]\dicFlx{}[-2] \dicSynonym{tilbúinn} \dicDirectTranslationCS{připravený, přichystaný} \dicExampleIS{reiðubúinn til aðstoðar} \dicExampleCS{připravený pomoct}
\dicEntry[reiðufé] \dicTerm{reiðu··|fé} \dicIPA{{r}{ei}{\textlengthmark}{ð}{\textscy}{f}{j}{\textepsilon}} \dicPos{n}[3] \dicFlx{(‑fjár)}[5] \dicSynonym{lausafé} \dicDirectTranslationCS{(peněžní) hotovost};  \dicPhraseIS{í reiðufé} \dicFlx{adv} \dicDirectTranslationCS{v~hotovosti, hotově}
\begin{xtolerant}{}{1pt}
\dicEntry[reiður] \dicTerm{reiður} \dicsymFrequent\addthinS\dicIPA{{r}{ei}{\textlengthmark}{ð}{\textscy}{\textsubring{r}}}\addthinS\dicPos{adj}[2]\dicFlx{}[-6]\addthinS\dicDirectTranslationCS{rozzlobený, rozhněvaný, vzteklý} \dicExampleIS{Ég veit að þú ert reið.} \dicExampleCS{Vím, že jsi rozzlobená.}
\end{xtolerant}
\dicEntry[reif] \dicTerm{reif} \dicIPA{{r}{ei}{\textlengthmark}{f}} \dicPos{v} \dicFlx{ind pf sg 1 pers} \dicLink{rífa}
\dicEntry[reifa] \dicTerm{reif|a} \dicIPA{{r}{ei}{\textlengthmark}{v}{a}} \dicPos{v}[1] \dicFlx{(‑aði)}[1] \dicFlx{acc} \textbf{1.} \dicSynonym*{binda um sár} \dicDirectTranslationCS{obvázat, zabandážovat, zafačovat (koleno ap.)}  \textbf{2.} \dicSynonym*{gera grein fyrir} \dicDirectTranslationCS{objasnit, vyjasnit (názory ap.)}
\dicEntry[reifar] \dicTerm{reifar} \dicIPA{{r}{ei}{\textlengthmark}{v}{a}{\textsubring{r}}} \dicPos{f}[12] \dicFlx{pl}[2] \dicLangCat{zast.} \dicSynonym{bleia} \dicDirectTranslationCS{(dětské) pleny}
\dicEntry[reigingur] \dicTerm{reig··ing|ur} \dicIPA{{r}{ei}{j}{\textlengthmark}{\textsci}{\ng}{\r{g}}{\textscy}{\textsubring{r}}} \dicPos{m}[6] \dicFlx{(‑s)}[9] \dicDirectTranslationCS{povýšenost, namyšlenost, nadutost};  \dicPhraseIS{það er reigingur í e‑m} \dicFlx{impers} \dicDirectTranslationCS{(kdo) je namyšlený}
\dicEntry[reik] \dicTerm{reik} \dicIPA{{r}{ei}{\textlengthmark}{\r{g}}} \dicPos{n}[2] \dicFlx{(‑s)}[2] \dicSynonym{flakk} \dicDirectTranslationCS{potulka, toulání};  \dicIdiom{reik}{ \dicPhraseIS{e‑að er á reiki}} \dicDirectTranslationCS{(co) se mění, (co) není jisté};  \dicPhraseIS{vera illa til reika} \dicDirectTranslationCS{špatně vypadat, být špatně oblečen}
\dicEntry[reika] \dicTerm{reik|a} \dicsymFrequent\  \dicIPA{{r}{ei}{\textlengthmark}{\r{g}}{a}} \dicPos{v}[1] \dicFlx{(‑aði)}[44] \dicSynonym{ráfa} \dicDirectTranslationCS{bloumat, chodit bez cíle, (za)toulat se} \dicExampleIS{reika niður að sjó} \dicExampleCS{zatoulat se k~moři};  \dicPhraseIS{láta hugann reika} \dicLangCat{přen.} \dicDirectTranslationCS{nechat myšlenky plynout}
\dicEntry[reikistjarna] \dicTerm{reiki··|stjarna} \dicIPA{{r}{ei}{\textlengthmark}{\r{\textObardotlessj}}{\textsci}{s}{\textsubring{d}}{j}{a}{r}{\textsubring{d}}{n}{a}} \dicPos{f}[1] \dicFlx{(‑stjörnu, ‑stjörnur)}[8] \dicFieldCat{astro.} \dicSynonym{pláneta} \dicDirectTranslationCS{planeta}
\dicEntry[reikna] \dicTerm{reikn|a} \dicsymFrequent\  \dicIPA{{r}{ei}{h}{\r{g}}{n}{a}} \dicPos{v}[1] \dicFlx{(‑aði)}[1] \dicFlx{acc} \dicDirectTranslationCS{(vy)počítat, spočítat} \dicExampleIS{reikna dæmið nákvæmlega} \dicExampleCS{vypočítat příklad přesně};  \dicIdiom{reikna}[með]{ \dicPhraseIS{reikna með e‑u}} \dicDirectTranslationCS{počítat s~(čím), spoléhat na (co)};  \dicIdiom{reikna}[út]{ \dicPhraseIS{reikna e‑ð út}} \dicDirectTranslationCS{vypočítat (co)};  \dicIdiom{reiknast}{ \dicPhraseIS{reiknast}} \dicFlx{refl} \dicDirectTranslationCS{(vy)počítat se}
\dicEntry[reikningshald] \dicTerm{reiknings··hald} \dicIPA{{r}{ei}{h}{\r{g}}{n}{i}{\ng}{s}{h}{a}{l}{\textsubring{d}}} \dicPos{n}[2] \dicFlx{(‑s)}[2] \dicFieldCat{ekon.} \dicSynonym{bókhald} \dicDirectTranslationCS{účetnictví}
\dicEntry[reikningsjöfnuður] \dicTerm{reiknings··|jöfn·uður} \dicIPA{{r}\-{ei}\-{h}\-{\r{g}}\-{n}\-{i}\-{\ng}\-{s}\-{j}\-{\oe}\-{\textsubring{b}}\-{n}\-{\textscy}\-{ð}\-{\textscy}\-{\textsubring{r}}\-} \dicPos{m}[10] \dicFlx{(‑jafnaðar\,/\addthin ‑jöfnuðar)}[43] \dicFieldCat{ekon.} \dicDirectTranslationCS{účetní bilance}
\dicEntry[reikningsnúmer] \dicTerm{reiknings··númer} \dicIPA{{r}{ei}{h}{\r{g}}{n}{i}{\ng}{s}{n}{u}{m}{\textepsilon}{\textsubring{r}}} \dicPos{n}[2] \dicFlx{(‑s, ‑)}[5] \dicDirectTranslationCS{číslo účtu}
\dicEntry[reikningsskil] \dicTerm{reiknings··skil} \dicIPA{{r}{ei}{h}{\r{g}}{n}{i}{\ng}{s}{\r{\textObardotlessj}}{\textsci}{\textsubring{l}}} \dicPos{n}[2] \dicFlx{pl}[1] \textbf{1.} \dicFieldCat{ekon.} \dicSynonym{skilagrein} \dicDirectTranslationCS{vyúčtování, zúčtování}  \textbf{2.} \dicLangCat{přen.} \dicSynonym*{lokkauppgjör} \dicDirectTranslationCS{vyúčtování, zúčtování}
\dicEntry[reikningsuppgjör] \dicTerm{reiknings··upp·gjör} \dicIPA{{r}{ei}{h}{\r{g}}{n}{i}{\ng}{s}{\textscy}{h}{\textsubring{b}}{\r{\textObardotlessj}}{\oe}{\textsubring{r}}} \dicPos{n}[2] \dicFlx{(‑s)}[2] \dicFieldCat{ekon.} \dicDirectTranslationCS{(účetní) výkaz}
\dicEntry[reikningsverð] \dicTerm{reiknings··verð} \dicIPA{{r}{ei}{h}{\r{g}}{n}{i}{\ng}{s}{v}{\textepsilon}{r}{\texttheta}} \dicPos{n}[2] \dicFlx{(‑s)}[2] \dicFieldCat{ekon.} \dicDirectTranslationCS{fakturovaná cena}
\dicEntry[reikningsyfirlit] \dicTerm{reiknings··yfir·lit} \dicIPA{{r}{ei}{h}{\r{g}}{n}{i}{\ng}{s}{\textsci}{v}{\textsci}{r}{l}{\textsci}{\textsubring{d}}} \dicPos{n}[2] \dicFlx{(‑s, ‑)}[5] \dicFieldCat{ekon.} \dicDirectTranslationCS{výpis z~účtu}
\dicEntry[reikningur] \dicTerm{reikn··ing|ur} \dicsymFrequent\  \dicIPA{{r}{ei}{h}{\r{g}}{n}{i}{\ng}{\r{g}}{\textscy}{\textsubring{r}}} \dicPos{m}[6] \dicFlx{(‑s, ‑ar)}[8] \textbf{1.} \dicFieldCat{ekon.} \dicSynonym{kvittun} \dicDirectTranslationCS{účet, účtenka, stvrzenka, faktura} \dicExampleIS{reikningur fyrir vörunni} \dicExampleCS{účtenka ke zboží}  \textbf{2.} \dicFieldCat{ekon.} \dicDirectTranslationCS{účet, konto  (v~bance ap.)} \dicExampleIS{opna reikning í bankanum} \dicExampleCS{otevřít si účet v~bance}  \textbf{3.} \dicSynonym{talnafræði} \dicDirectTranslationCS{aritmetika, počítání} \dicExampleIS{læra reikning í skólanum} \dicExampleCS{naučit se aritmetiku ve škole};  \dicIdiom{reikningur}{ \dicPhraseIS{e‑að gerir strik í reikninginn}} \dicLangCat{přen.} \dicDirectTranslationCS{(co) dělá čáru přes rozpočet}; { \dicPhraseIS{gera upp reikningana við e‑n}} \dicDirectTranslationCS{vyřídit si účty s~(kým), zúčtovat s~(kým)};  \dicPhraseIS{skrifa e‑ð á e‑rs reikning} \dicLangCat{přen.} \dicDirectTranslationCS{připsat (co komu) na účet};  \dicPhraseIS{taka e‑ð með í reikninginn} \dicDirectTranslationCS{vzít (co) v~úvahu, brát ohled na (co)}
\dicEntry[reiknirit] \dicTerm{reikni··rit} \dicIPA{{r}{ei}{h}{\r{g}}{n}{\textsci}{r}{\textsci}{\textsubring{d}}} \dicPos{n}[2] \dicFlx{(‑s, ‑)}[5] \dicFieldCat{poč.} \dicDirectTranslationCS{algoritmus}
\dicEntry[reiknivél] \dicTerm{reikni··vél} \dicIPA{{r}{ei}{h}{\r{g}}{n}{\textsci}{v}{j}{\textepsilon}{\textsubring{l}}} \dicPos{f}[4] \dicFlx{(‑ar, ‑ar)}[1] \dicDirectTranslationCS{kalkulačka}
\dicEntry[reikull] \dicTerm{reikull} \dicIPA{{r}{ei}{\textlengthmark}{\r{g}}{\textscy}{\textsubring{d}}{\textsubring{l}}} \dicPos{adj}[8]\dicFlx{}[-4] \textbf{1.} \dicSynonym*{óstöðuglyndur} \dicDirectTranslationCS{nejistý, nepevný, vratký (chůze ap.)} \dicExampleIS{vera reikull í göngulagi} \dicExampleCS{mít nepevnou chůzi};  \dicPhraseIS{vera reikull í spori} \dicDirectTranslationCS{potácet se, vrávorat}  \textbf{2.} \dicDirectTranslationCS{rozpolcený, nerozhodný}
\dicEntry[reim] \dicTerm{reim} \dicIPA{{r}{ei}{\textlengthmark}{\textsubring{m}}} \dicPos{f}[4] \dicFlx{(‑ar, ‑ar)}[1] \textbf{1.} \dicSynonym{festi} \dicDirectTranslationCS{šňůrka, řemínek, pásek}  \textbf{2.} \dicDirectTranslationCS{tkanička (u~boty ap.)} \dicExampleIS{skóreim} \dicExampleCS{tkanička u~boty}
\dicEntry[reima] \dicTerm{reim|a} \dicIPA{{r}{ei}{\textlengthmark}{m}{a}} \dicPos{v}[1] \dicFlx{(‑aði)}[1] \dicFlx{acc} \dicDirectTranslationCS{zavázat, zavazovat, (za)šněrovat} \dicExampleIS{reima skóna} \dicExampleCS{zašněrovat boty}
\dicEntry[reimleiki] \dicTerm{reim··leik|i} \dicIPA{{r}{ei}{m}{l}{ei}{\r{\textObardotlessj}}{\textsci}} \dicPos{m}[1] \dicFlx{(‑a, ‑ar)}[1] \dicSynonym{draugagangur} \dicDirectTranslationCS{strašení (strašidla ap.)}
\dicEntry[reip] \dicTerm{reip} \dicIPA{{r}{ei}{\textlengthmark}{\textsubring{b}}} \dicPos{n}[2] \dicFlx{(‑s, ‑)}[5] \dicPhraseIS{það er við ramman reip að draga} \dicLangCat{přen.} \dicDirectTranslationCS{je to dlouhá a~obtížná cesta (zlepšit úroveň vzdělání ap.)}
\dicEntry[reipi] \dicTerm{reipi} \dicIPA{{r}{ei}{\textlengthmark}{\textsubring{b}}{\textsci}} \dicPos{n}[2] \dicFlx{(‑s, ‑)}[14] \dicSynonym{kaðall} \dicDirectTranslationCS{provaz, lano, šňůra};  \dicPhraseIS{hæla e‑m á hvert reipi} \dicLangCat{přen.} \dicDirectTranslationCS{vychvalovat (koho) do nebes}
\dicEntry[reiprennandi] \dicTerm{reip··renn·andi} \dicIPA{{r}{ei}{\textlengthmark}{\textsubring{b}}{r}{\textepsilon}{n}{a}{n}{\textsubring{d}}{\textsci}} \dicPos{adv} \dicDirectTranslationCS{plynule, plynně} \dicExampleIS{tala íslensku reiprennandi} \dicExampleCS{mluvit plynně islandsky}
\dicEntry[reis] \dicTerm{reis} \dicIPA{{r}{ei}{\textlengthmark}{s}} \dicPos{v} \dicFlx{ind pf sg 1 pers} \dicLink{rísa}
\dicEntry[reisa] \dicTerm{reis|a\smash{\textsuperscript{1}}} \dicIPA{{r}{ei}{\textlengthmark}{s}{a}} \dicPos{f}[1] \dicFlx{(‑u, ‑ur)}[7] \dicLangCat{hovor.} \dicSynonym{ferð} \dicDirectTranslationCS{cesta, cestování}
\dicEntry[reisa] \dicTerm{reis|a\smash{\textsuperscript{2}}} \dicsymFrequent\  \dicIPA{{r}{ei}{\textlengthmark}{s}{a}} \dicPos{v}[2] \dicFlx{(‑ti, ‑t)}[63] \dicFlx{acc} \textbf{1.} \dicSynonym{byggja} \dicDirectTranslationCS{(vy)budovat, postavit, stavět} \dicExampleIS{reisa tjald} \dicExampleCS{postavit stan};  \dicPhraseIS{reisa e‑ð úr rústum} \dicDirectTranslationCS{přestavět (co)}  \textbf{2.} \dicSynonym*{rétta við} \dicDirectTranslationCS{postavit, narovnat};  \dicIdiom{reisast}{ \dicPhraseIS{reisast}} \dicFlx{refl} \dicSynonym*{rísa upp} \dicDirectTranslationCS{pozvednout se, zvedat se};  \dicIdiom{reisa}{ \dicPhraseIS{e‑að er ekki á rökum reist}} \dicLangCat{přen.} \dicDirectTranslationCS{(co) se nezakládá na skutečnosti}
\dicEntry[reisn] \dicTerm{reisn} \dicIPA{{r}{ei}{s}{\textsubring{d}}{\textsubring{n}}} \dicPos{f}[7] \dicFlx{(‑ar)}[3] \dicSynonym*{mikilleiki} \dicDirectTranslationCS{(morální) velikost, ušlechtilost} \dicExampleIS{Það er reisn yfir honum.} \dicExampleCS{Je to velmi ušlechtilý člověk.}
\dicEntry[reisugildi] \dicTerm{reisu··gildi} \dicIPA{{r}{ei}{\textlengthmark}{s}{\textscy}{\r{\textObardotlessj}}{\textsci}{l}{\textsubring{d}}{\textsci}} \dicPos{n}[2] \dicFlx{(‑s, ‑)}[14] \dicDirectTranslationCS{večírek na oslavu dokončení stavby}
\begin{xtolerant}{}{1pt}
\dicEntry[reisulegur] \dicTerm{reisu··legur}\addthinS\dicIPA{{r}{ei}{\textlengthmark}{s}{\textscy}{l}{\textepsilon}{\textbabygamma}{\textscy}{\textsubring{r}}}\addthinS\dicPos{adj}[1]\dicFlx{}[-8]\addthinS\dicSynonym{gerðarlegur} \dicDirectTranslationCS{majestátní, impozantní, vznešený} \dicExampleIS{reisuleg bygging} \dicExampleCS{majestátní budova}
\end{xtolerant}
\dicEntry[reitur] \dicTerm{reit|ur} \dicsymFrequent\  \dicIPA{{r}{ei}{\textlengthmark}{\textsubring{d}}{\textscy}{\textsubring{r}}} \dicPos{m}[6] \dicFlx{(‑s, ‑ir\,/\addthin ‑ar)}[70] \textbf{1.} \dicSynonym{beð} \dicDirectTranslationCS{pole, záhon}  \textbf{2.} \dicDirectTranslationCS{pole, kolonka} \dicExampleIS{fylla hvern reit með viðeigandi upplýsingum} \dicExampleCS{vyplnit každou kolonku příslušnými informacemi}  \textbf{3.} \dicDirectTranslationCS{pole (na šachovnici ap.)}
\dicEntry[rek] \dicTerm{rek\smash{\textsuperscript{1}}} \dicIPA{{r}{\textepsilon}{\textlengthmark}{\r{g}}} \dicPos{n}[2] \dicFlx{(‑s)}[2] \textbf{1.} \dicSynonym*{það að reka} \dicDirectTranslationCS{unášení (vodou ap.)};  \dicPhraseIS{vera á reki} \dicDirectTranslationCS{být unášený (proudem, větrem ap.)}  \textbf{2.} \dicDirectTranslationCS{unášený předmět (led ap.)}  \textbf{3.} \dicSynonym{aldur} \dicDirectTranslationCS{věk, stáří} \dicExampleIS{Þeir eru á svipuðu reki.} \dicExampleCS{Jsou podobného věku.}
\dicEntry[rek] \dicTerm{rek\smash{\textsuperscript{2}}} \dicIPA{{r}{\textepsilon}{\textlengthmark}{\r{g}}} \dicPos{v} \dicFlx{ind praes sg 1 pers} \textbf{1.} \dicLink{rekja}  \textbf{2.} \dicLink{reka\smash{\textsuperscript{2}}}
\dicEntry[reka] \dicTerm{rek|a\smash{\textsuperscript{1}}} \dicIPA{{r}{\textepsilon}{\textlengthmark}{\r{g}}{a}} \dicPos{f}[1] \dicFlx{(‑u, ‑ur)}[19] \dicLangCat{zast.} \dicSynonym{skófla\smash{\textsuperscript{1}}} \dicDirectTranslationCS{lopata} \dicExampleIS{páll og reka} \dicExampleCS{rýč a~lopata}
\dicEntry[reka] \dicTerm{reka\smash{\textsuperscript{2}}} \dicsymFrequent\  \dicIPA{{r}{\textepsilon}{\textlengthmark}{\r{g}}{a}} \dicPos{v}[6] \dicFlx{(rek, rak, rákum, ræki, rekið)}[6] \dicFlx{acc} \textbf{1.} \dicSynonym{stugga} \dicDirectTranslationCS{hnát, zahnat, odehnat, vyhnat} \dicExampleIS{reka e‑n burt} \dicExampleCS{vyhnat (koho)}  \textbf{2.} \dicSynonym*{segja upp} \dicDirectTranslationCS{vyhodit, vyrazit (ze zaměstnání ap.)} \dicExampleIS{reka e‑n úr starfi} \dicExampleCS{vyhodit (koho) z~práce}  \textbf{3.} \dicSynonym{slá\smash{\textsuperscript{3}}} \dicDirectTranslationCS{vrazit, zabodnout} \dicExampleIS{reka naglana á kaf} \dicExampleCS{zarazit hřebík};  \dicPhraseIS{reka e‑m löðrung} \dicDirectTranslationCS{vrazit (komu) facku}  \textbf{4.} \dicSynonym{starfrækja} \dicDirectTranslationCS{provozovat, vést (firmu ap.)} \dicExampleIS{reka verslunina} \dicExampleCS{provozovat obchod}  \textbf{5.} \dicPhraseIS{e‑ð rekur} \dicFlx{impers} \dicDirectTranslationCS{(co) je unášeno (proudem, větrem ap.)} \dicExampleIS{Bát rekur upp á land.} \dicExampleCS{Člun je unášen k~pevnině.};  \dicIdiom{reka}[á]{ \dicPhraseIS{reka e‑n á dyr}} \dicDirectTranslationCS{vyrazit (koho) ze dveří};  \dicIdiom{reka}[á eftir]{ \dicPhraseIS{reka á eftir e‑m}} \dicDirectTranslationCS{pobízet (koho), popohánět (koho)};  \dicIdiom{reka}[upp]{ \dicPhraseIS{reka upp öskur\,/\addthin óp}} \dicDirectTranslationCS{zařvat, zakřičet}; { \dicPhraseIS{reka upp stór augu}} \dicDirectTranslationCS{vytřeštit zrak, vyvalit oči};  \dicIdiom{reka}[út]{ \dicPhraseIS{reka út úr sér tunguna}} \dicDirectTranslationCS{vypláznout jazyk};  \dicIdiom{reka}[við]{ \dicPhraseIS{reka við}} \dicSynonym*{leysa vind} \dicDirectTranslationCS{prdnout si, prdět, pšouknout si, pšoukat};  \dicIdiom{rekast}{ \dicPhraseIS{rekast}} \dicFlx{refl} \dicDirectTranslationCS{srazit se, narazit (do sebe) (auta ap.)};  \dicIdiom{rekast}[á]{ \dicPhraseIS{rekast á e‑n}} \dicFlx{refl} \dicDirectTranslationCS{narazit na (koho), potkat (koho)} \dicExampleIS{rekast á e‑n af tilviljun} \dicExampleCS{narazit na (koho) náhodou}; { \dicPhraseIS{rekast á e‑ð}} \dicFlx{refl} \dicDirectTranslationCS{narazit na (co), objevit (co)} \dicExampleIS{rekast á villu} \dicExampleCS{narazit na chybu};  \dicIdiom{reka}{ \dicPhraseIS{reka lestina}} \dicDirectTranslationCS{uzavírat pořadí, být poslední (v~závodě ap.)}
\dicEntry[rekald] \dicTerm{rek··|ald} \dicIPA{{r}{\textepsilon}{\textlengthmark}{\r{g}}{a}{l}{\textsubring{d}}} \dicPos{n}[2] \dicFlx{(‑alds, ‑öld)}[8] \textbf{1.} \dicSynonym*{hlutur á reki} \dicDirectTranslationCS{plovoucí\,/\addthin vyplavená věc} \dicIndirectTranslationCS{(předmět vyplavený mořem nebo plovoucí po hladině)}  \textbf{2.} \dicSynonym{ræfill} \dicDirectTranslationCS{(lidská) troska}
\dicEntry[rekaviður] \dicTerm{reka··við|ur} \dicIPA{{r}{\textepsilon}{\textlengthmark}{\r{g}}{a}{v}{\textsci}{ð}{\textscy}{\textsubring{r}}} \dicPos{m}[10] \dicFlx{(‑ar, ‑ir)}[19] \dicDirectTranslationCS{vyplavené\,/\addthin naplavené dříví}
\dicEntry[reki] \dicTerm{rek|i} \dicIPA{{r}{\textepsilon}{\textlengthmark}{\r{\textObardotlessj}}{\textsci}} \dicPos{m}[1] \dicFlx{(‑a, ‑ar)}[1] \dicDirectTranslationCS{plovoucí\,/\addthin vyplavená věc (dříví ap.)} \dicIndirectTranslationCS{(předmět vyplavený mořem nebo plovoucí po hladině)}
\dicEntry[rekið] \dicTerm{rekið} \dicIPA{{r}{\textepsilon}{\textlengthmark}{\r{\textObardotlessj}}{\textsci}{\texttheta}} \dicPos{v} \dicFlx{supin} \dicLink{reka\smash{\textsuperscript{2}}}
\dicEntry[rekja] \dicTerm{rekja} \dicsymFrequent\  \dicIPA{{r}{\textepsilon}{\textlengthmark}{\r{\textObardotlessj}}{a}} \dicPos{v}[4] \dicFlx{(rek, rakti, röktum, rekti, rakið)}[46] \dicFlx{acc} \textbf{1.} \dicSynonym{fylgja\smash{\textsuperscript{2}}} \dicDirectTranslationCS{sledovat, stopovat} \dicExampleIS{rekja sporin eftir e‑n} \dicExampleCS{sledovat (čí) stopu}  \textbf{2.} \dicSynonym*{leysa úr vafningum} \dicDirectTranslationCS{rozmotat, rozplést, rozpárat};  \dicPhraseIS{rekja e‑ð upp} \dicDirectTranslationCS{rozplést (co) (svetr ap.)}  \textbf{3.} \dicSynonym*{gera grein fyrir} \dicDirectTranslationCS{vylíčit, vystihnout} \dicExampleIS{rekja feril hennar ýtarlega} \dicExampleCS{vylíčit podrobně její životní dráhu}  \textbf{4.} \dicPhraseIS{rekja bolta} \dicFieldCat{sport.} \dicDirectTranslationCS{driblovat}
\dicEntry[rekkja] \dicTerm{rekkj|a} \dicIPA{{r}{\textepsilon}{h}{\r{\textObardotlessj}}{a}} \dicPos{f}[1] \dicFlx{(‑u, ‑ur)}[17] \dicSynonym{rúm} \dicDirectTranslationCS{lůžko, lože};  \dicPhraseIS{ganga til rekkju} \dicLangCat{kniž.} \dicSynonym{hátta} \dicDirectTranslationCS{jít spát, jít si lehnout}
\dicEntry[rekkjunautur] \dicTerm{rekkju··naut|ur} \dicIPA{{r}{\textepsilon}{h}{\r{\textObardotlessj}}{\textscy}{n}{\oe i}{\textsubring{d}}{\textscy}{\textsubring{r}}} \dicPos{m}[6] \dicFlx{(‑s\,/\addthin ‑ar, ‑ar)}[65] \textbf{1.} \dicSynonym{elskhugi} \dicDirectTranslationCS{sexuální partner(ka)}  \textbf{2.} \dicSynonym*{rúmfélagi} \dicDirectTranslationCS{společník\,/\addthin společnice sdílící lůžko}
\dicEntry[reknet] \dicTerm{rek··net} \dicIPA{{r}{\textepsilon}{\textlengthmark}{\r{g}}{n}{\textepsilon}{\textsubring{d}}} \dicPos{n}[2] \dicFlx{(‑s, ‑)}[5] \dicFieldCat{nám.} \dicDirectTranslationCS{tenatová síť}
\dicEntry[rekspölur] \dicTerm{rek··|spölur} \dicIPA{{r}{\textepsilon}{\textlengthmark}{\r{g}}{s}{\textsubring{b}}{\oe}{l}{\textscy}{\textsubring{r}}} \dicPos{m}[11] \dicFlx{(‑spalar)}[8] \dicSynonym{skrið} \dicDirectTranslationCS{pohyb, chod};  \dicPhraseIS{koma e‑u á rekspölur} \dicDirectTranslationCS{uvést (co) do chodu}
\dicEntry[rekstrarafkoma] \dicTerm{rekstrar··af·kom|a} \dicIPA{{r}{\textepsilon}{x}{s}{\textsubring{d}}{r}{a}{r}{a}{f}{k\smash{\textsuperscript{h}}}{\textopeno}{m}{a}} \dicPos{f}[1] \dicFlx{(‑u)}[5] \dicFieldCat{ekon.} \dicDirectTranslationCS{provozní výsledek}
\dicEntry[rekstraráætlun] \dicTerm{rekstrar··á·ætl|un} \dicIPA{{r}{\textepsilon}{x}{s}{\textsubring{d}}{r}{a}{r}{au}{a}{i}{h}{\textsubring{d}}{l}{\textscy}{\textsubring{n}}} \dicPos{f}[7] \dicFlx{(‑unar, ‑anir)}[8] \dicFieldCat{ekon.} \dicDirectTranslationCS{provozní harmonogram}
\dicEntry[rekstrarfé] \dicTerm{rekstrar··|fé} \dicIPA{{r}{\textepsilon}{x}{s}{\textsubring{d}}{r}{a}{\textsubring{r}}{f}{j}{\textepsilon}} \dicPos{n}[3] \dicFlx{(‑fjár)}[5] \dicFieldCat{ekon.} \dicDirectTranslationCS{provozní kapitál}
\dicEntry[rekstrarhagfræði] \dicTerm{rekstrar··hag·fræð|i} \dicIPA{{r}{\textepsilon}{x}{s}{\textsubring{d}}{r}{a}{\textsubring{r}}{h}{a}{\textbabygamma}{f}{r}{a}{i}{ð}{\textsci}} \dicPos{f}[3] \dicFlx{(‑i)}[3] \dicFieldCat{ekon.} \dicDirectTranslationCS{mikroekonomie}
\dicEntry[rekstrarhagnaður] \dicTerm{rekstrar··hag·nað|ur} \dicIPA{{r}{\textepsilon}{x}{s}{\textsubring{d}}{r}{a}{\textsubring{r}}{h}{a}{\r{g}}{n}{a}{ð}{\textscy}{\textsubring{r}}} \dicPos{m}[10] \dicFlx{(‑ar)}[9] \dicFieldCat{ekon.} \dicDirectTranslationCS{provozní zisk}
\dicEntry[rekstrarhalli] \dicTerm{rekstrar··hall|i} \dicIPA{{r}{\textepsilon}{x}{s}{\textsubring{d}}{r}{a}{\textsubring{r}}{h}{a}{\textsubring{d}}{l}{\textsci}} \dicPos{m}[1] \dicFlx{(‑a)}[3] \dicFieldCat{ekon.} \dicDirectTranslationCS{provozní schodek\,/\addthin deficit}
\dicEntry[rekstrarkostnaður] \textls[15]{\dicTerm{rekstrar··kost·nað|ur} \dicIPA{{r}\-{\textepsilon}\-{x}\-{s}\-{\textsubring{d}}\-{r}\-{a}\-{\textsubring{r}}\-{k\smash{\textsuperscript{h}}}\-{\textopeno}\-{s}\-{\textsubring{d}}\-{n}\-{a}\-{ð}\-{\textscy}\-{\textsubring{r}}} \dicPos{m}[10] \dicFlx{(‑ar)}[9]} \dicFieldCat{ekon.} \dicDirectTranslationCS{provozní náklady}
\dicEntry[rekstrarlán] \dicTerm{rekstrar··lán} \dicIPA{{r}{\textepsilon}{x}{s}{\textsubring{d}}{r}{a}{r}{l}{au}{\textsubring{n}}} \dicPos{n}[2] \dicFlx{(‑s, ‑)}[5] \dicFieldCat{ekon.} \dicDirectTranslationCS{provozní půjčka}
\dicEntry[rekstrarleyfi] \dicTerm{rekstrar··leyfi} \dicIPA{{r}{\textepsilon}{x}{s}{\textsubring{d}}{r}{a}{r}{l}{ei}{v}{\textsci}} \dicPos{n}[2] \dicFlx{(‑s, ‑)}[14] \dicFieldCat{ekon.} \dicDirectTranslationCS{provozní povolení, koncese}
\dicEntry[rekstrarreikningur] \dicTerm{rekstrar··reikn·ing|ur} \dicIPA{{r}\-{\textepsilon}\-{x}\-{s}\-{\textsubring{d}}\-{r}\-{a}\-{r}\-{ei}\-{h}\-{\r{g}}\-{n}\-{i}\-{\ng}\-{\r{g}}\-{\textscy}\-{\textsubring{r}}\-} \dicPos{m}[6] \dicFlx{(‑s, ‑ar)}[8] \dicFieldCat{ekon.} \dicDirectTranslationCS{výkaz zisků a~ztrát}
\dicEntry[rekstraryfirlit] \dicTerm{rekstrar··yfir·lit} \dicIPA{{r}{\textepsilon}{x}{s}{\textsubring{d}}{r}{a}{r}{\textsci}{v}{\textsci}{r}{l}{\textsci}{\textsubring{d}}} \dicPos{n}[2] \dicFlx{(‑s, ‑)}[5] \dicFieldCat{ekon.} \dicDirectTranslationCS{provozní přehled\,/\addthin souhrn}
\dicEntry[rekstur] \dicTerm{rekst|ur} \dicsymFrequent\  \dicIPA{{r}{\textepsilon}{x}{s}{\textsubring{d}}{\textscy}{\textsubring{r}}} \dicPos{m}[5] \dicFlx{(‑rar\,/\addthin ‑urs, ‑rar)}[9] \textbf{1.} \dicSynonym{starfræksla} \dicDirectTranslationCS{provozování, podnikání} \dicExampleIS{einkarekstur} \dicExampleCS{soukromé podnikání}  \textbf{2.} \dicSynonym*{það að reka} \dicDirectTranslationCS{hnaní, pohánění} \dicExampleIS{rekstur fjár} \dicExampleCS{hnaní ovcí}  \textbf{3.} \dicSynonym*{fjárhópur} \dicDirectTranslationCS{stádo (ovcí ap.)}
\dicEntry[rekti] \dicTerm{rekti} \dicIPA{{r}{\textepsilon}{x}{\textsubring{d}}{\textsci}} \dicPos{v} \dicFlx{con pf sg 1 pers} \dicLink{rekja}
\dicEntry[rektor] \dicTerm{rektor} \dicIPA{{r}{\textepsilon}{x}{\textsubring{d}}{\textopeno}{\textsubring{r}}} \dicPos{m}[4] \dicFlx{(‑s, ‑ar)}[14] \dicFieldCat{škol.} \dicSynonym{skólastjóri} \dicDirectTranslationCS{rektor(ka)}
\dicEntry[rella] \dicTerm{rell|a\smash{\textsuperscript{1}}} \dicIPA{{r}{\textepsilon}{\textsubring{d}}{l}{a}} \dicPos{f}[1] \dicFlx{(‑u, ‑ur)}[19] \textbf{1.} \dicSynonym{kvabb} \dicDirectTranslationCS{povyk, poprask};  \dicPhraseIS{gera sér rellu út af e‑u} \dicDirectTranslationCS{udělat scénu\,/\addthin dusno kvůli (čemu), vyvádět kvůli (čemu)}  \textbf{2.} \dicDirectTranslationCS{větrník} \dicIndirectTranslationCS{(hračka)}  \textbf{3.} \dicSynonym*{lítil flugvél} \dicDirectTranslationCS{malé letadlo, letadélko}
\dicEntry[rella] \dicTerm{rell|a\smash{\textsuperscript{2}}} \dicIPA{{r}{\textepsilon}{\textsubring{d}}{l}{a}} \dicPos{v}[1] \dicFlx{(‑aði)}[1] \dicDirectTranslationCS{obtěžovat, otravovat} \dicExampleIS{rella í e‑m} \dicExampleCS{otravovat (koho)}
\dicEntry[rembast] \dicTerm{remb|ast} \dicIPA{{r}{\textepsilon}{m}{\textsubring{b}}{a}{s}{\textsubring{d}}} \dicPos{v}[2] \dicFlx{(‑dist, ‑st)}[205] \dicFlx{refl} \textbf{1.} \dicSynonym{streitast} \dicDirectTranslationCS{dřít se, pachtit se} \dicExampleIS{rembast við e‑ð} \dicExampleCS{pachtit se s~(čím)}  \textbf{2.} \dicDirectTranslationCS{tlačit (při porodu)}
\dicEntry[rembingur] \dicTerm{remb··ing|ur} \dicIPA{{r}{\textepsilon}{m}{\textsubring{b}}{i}{\ng}{\r{g}}{\textscy}{\textsubring{r}}} \dicPos{m}[6] \dicFlx{(‑s)}[9] \textbf{1.} \dicSynonym{áreynsla} \dicDirectTranslationCS{dřina, lopota}  \textbf{2.} \dicSynonym{hroki} \dicDirectTranslationCS{povýšenost, domýšlivost}
\dicEntry[remúlaði] \dicTerm{remúlaði} \dicIPA{{r}{\textepsilon}{\textlengthmark}{m}{u}{l}{a}{ð}{\textsci}} \dicPos{n}[2] \dicFlx{(‑s)}[20] \dicFieldCat{kulin.} \dicDirectTranslationCS{remuláda}
\dicEntry[rendur] \dicTerm{rendur} \dicIPA{{r}{\textepsilon}{n}{\textsubring{d}}{\textscy}{\textsubring{r}}} \dicPos{f} \dicFlx{pl nom} \dicLink{rönd}
\dicEntry[rengja] \dicTerm{reng|ja} \dicIPA{{r}{ei}{\textltailn}{\r{\textObardotlessj}}{a}} \dicPos{v}[2] \dicFlx{(‑di, ‑t)}[129] \dicFlx{acc} \dicSynonym*{vefengja} \dicDirectTranslationCS{zpochybnit, zpochybňovat (pravdivost ap.)} \dicIndirectTranslationCS{(nejčastěji se záporem)} \dicExampleIS{rengja ekki þessar tölur} \dicExampleCS{nezpochybňovat tato čísla}
\dicEntry[renglulegur] \dicTerm{renglu··legur} \dicIPA{{r}{ei}{\ng}{\r{g}}{l}{\textscy}{l}{\textepsilon}{\textbabygamma}{\textscy}{\textsubring{r}}} \dicPos{adj}[1]\dicFlx{}[-8] \dicSynonym{krangalegur} \dicDirectTranslationCS{vytáhlý, vysoký a~štíhlý}
\dicEntry[renn] \dicTerm{renn} \dicIPA{{r}{\textepsilon}{n}{\textlengthmark}} \dicPos{v} \dicFlx{ind praes sg 1 pers} \dicLink{renna\smash{\textsuperscript{3}}}
\dicEntry[renna] \dicTerm{renn|a\smash{\textsuperscript{1}}} \dicIPA{{r}{\textepsilon}{n}{\textlengthmark}{a}} \dicPos{f}[1] \dicFlx{(‑u, ‑ur)}[7] \textbf{1.} \dicSynonym{stokkur} \dicDirectTranslationCS{koryto, strouha, žlab}  \textbf{2.} \dicSynonym{þakrenna} \dicDirectTranslationCS{okap}
\dicEntry[renna] \dicTerm{renn|a\smash{\textsuperscript{2}}} \dicsymFrequent\  \dicIPA{{r}{\textepsilon}{n}{\textlengthmark}{a}} \dicPos{v}[2] \dicFlx{(‑di, ‑t)}[141] \dicFlx{dat} \textbf{1.} \dicDirectTranslationCS{klouzat (se), sklouznout (se)} \dicExampleIS{renna sér eftir ísnum} \dicExampleCS{klouzat se po ledu}  \textbf{2.} \dicDirectTranslationCS{(pro)jet} \dicExampleIS{renna upp í sveit} \dicExampleCS{vyjet na venkov}  \textbf{3.} \dicDirectTranslationCS{zapnout (zip)} \dicExampleIS{renna upp peysunni} \dicExampleCS{zapnout svetr (na zip)};  \dicIdiom{renna}[niður]{ \dicPhraseIS{renna e‑u niður}} \dicDirectTranslationCS{spláchnout (co), zapít (co)} \dicExampleIS{renna niður matnum með mjólk} \dicExampleCS{spláchnout jídlo mlékem};  \dicIdiom{renna}[yfir]{ \dicPhraseIS{renna yfir e‑ð}} \dicDirectTranslationCS{prolistovat (co)};  \dicIdiom{renna}{ \dicPhraseIS{e‑n rennir e‑ð í grun, e‑r rennir e‑ð í grun}} \dicFlx{impers} \dicDirectTranslationCS{(kdo) má ponětí o~(čem)}
\dicEntry[renna] \dicTerm{renna\smash{\textsuperscript{3}}} \dicsymFrequent\  \dicIPA{{r}{\textepsilon}{n}{\textlengthmark}{a}} \dicPos{v}[6] \dicFlx{(renn, rann, runnum, rynni, runnið)}[23] \textbf{1.} \dicSynonym{streyma} \dicDirectTranslationCS{téct, plynout, proudit} \dicExampleIS{Það rennur úr krananum.} \dicExampleCS{Teče z~kohoutku.}  \textbf{2.} \dicDirectTranslationCS{vyprchat, pominout, opadnout (pocit hněvu ap.)}  \textbf{3.} \dicSynonym{hlaupa} \dicDirectTranslationCS{běžet, utíkat, pádit} \dicExampleIS{renna í kapp við e‑n} \dicExampleCS{běžet s~(kým) o~závod};  \dicIdiom{renna}[af]{ \dicPhraseIS{það rennur af e‑m e‑að}} \dicFlx{impers} \dicSynonym{hverfa\smash{\textsuperscript{2}}} \dicDirectTranslationCS{(co) na (koho) padá (ospalost ap.)}; { \dicPhraseIS{það rennur af e‑m}} \dicFlx{impers} \dicDirectTranslationCS{(kdo) střízliví (z~opojení ap.)};  \dicIdiom{renna}[saman]{ \dicPhraseIS{renna saman}} \dicSynonym*{blandast} \dicDirectTranslationCS{slít se, sjednotit se, splynout};  \dicIdiom{renna}[upp]{ \dicPhraseIS{renna upp}} \dicDirectTranslationCS{vyjít, vycházet (slunce ap.)} \dicExampleIS{Sólin rennur upp.} \dicExampleCS{Slunce vychází.};  \dicIdiom{renna}[út]{ \dicPhraseIS{renna út}} {\textbf{a.}} \dicSynonym*{enda} \dicDirectTranslationCS{vypršet, uplynout (čas ap.)} \dicExampleIS{Fresturinn er runninn út.} \dicExampleCS{Lhůta vypršela.};  {\textbf{b.}} \dicDirectTranslationCS{vypršet, skončit (záruční doba ap.)};  \dicIdiom{renna}{ \dicPhraseIS{láta renna í e‑ð}} \dicDirectTranslationCS{nechat napustit (co)}; { \dicPhraseIS{dagur rennur}} \dicDirectTranslationCS{rozednívá se, vychází slunce}
\dicEntry[rennibekkur] \dicTerm{renni··bekk|ur} \dicIPA{{r}{\textepsilon}{n}{\textlengthmark}{\textsci}{\textsubring{b}}{\textepsilon}{h}{\r{g}}{\textscy}{\textsubring{r}}} \dicPos{m}[9] \dicFlx{(‑jar\,/\addthin ‑s, ‑ir)}[26] \dicDirectTranslationCS{soustruh}
\dicEntry[rennibraut] \dicTerm{renni··braut} \dicIPA{{r}{\textepsilon}{n}{\textlengthmark}{\textsci}{\textsubring{b}}{r}{\oe i}{\textsubring{d}}} \dicPos{f}[7] \dicFlx{(‑ar, ‑ir)}[1] \dicDirectTranslationCS{skluzavka, klouzačka}
\dicEntry[rennihurð] \dicTerm{renni··hurð} \dicIPA{{r}{\textepsilon}{n}{\textlengthmark}{\textsci}{h}{\textscy}{r}{\texttheta}} \dicPos{f}[7] \dicFlx{(‑ar, ‑ir)}[1] \dicDirectTranslationCS{posuvné dveře}
\dicEntry[rennilás] \dicTerm{renni··lás} \dicIPA{{r}{\textepsilon}{n}{\textlengthmark}{\textsci}{l}{au}{s}} \dicPos{m}[4] \dicFlx{(‑s, ‑ar)}[25] \dicDirectTranslationCS{zip, zdrhovadlo} \dicExampleIS{renna upp rennilás} \dicExampleCS{zapnout zip}
\begin{xtolerant}{}{1pt}
\dicEntry[rennilegur] \dicTerm{renni··legur}\addthinS\dicIPA{{r}{\textepsilon}{n}{\textlengthmark}{\textsci}{l}{\textepsilon}{\textbabygamma}{\textscy}{\textsubring{r}}}\addthinS\dicPos{adj}[1]\dicFlx{}[-8]\addthinS\dicSynonym{straumlínulagaður}\addthinS\dicDirectTranslationCS{aerodynamický, zeštíhlený} \dicExampleIS{rennilegur sportbíll} \dicExampleCS{aerodynamické sportovní auto}
\end{xtolerant}
\dicEntry[renningur] \dicTerm{renn··ing|ur} \dicIPA{{r}{\textepsilon}{n}{\textlengthmark}{i}{\ng}{\r{g}}{\textscy}{\textsubring{r}}} \dicPos{m}[6] \dicFlx{(‑s, ‑ar)}[8] \textbf{1.} \dicSynonym{ræma} \dicDirectTranslationCS{úzký pruh, proužek}  \textbf{2.} \dicSynonym{dregill} \dicDirectTranslationCS{běhoun}
\dicEntry[rennisléttur] \textls[15]{\dicTerm{renni··sléttur} \dicIPA{{r}{\textepsilon}{n}{\textlengthmark}{\textsci}{s}{\textsubring{d}}{l}{j}{\textepsilon}{h}{\textsubring{d}}{\textscy}{\textsubring{r}}} \dicPos{adj}[1]\dicFlx{}[-13] \dicDirectTranslationCS{(zcela) rovný\,/\addthin hladký}}
\dicEntry[rennistigi] \dicTerm{renni··stig|i} \dicIPA{{r}{\textepsilon}{n}{\textlengthmark}{\textsci}{s}{\textsubring{d}}{i}{j}{\textsci}} \dicPos{m}[1] \dicFlx{(‑a, ‑ar)}[1] \dicDirectTranslationCS{eskalátor, pohyblivé\,/\addthin jezdicí schody}
\dicEntry[rennsli] \dicTerm{rennsli} \dicIPA{{r}{\textepsilon}{n}{s}{\textsubring{d}}{l}{\textsci}} \dicPos{n}[2] \dicFlx{(‑s, ‑)}[14] \textbf{1.} \dicSynonym{streymi} \dicDirectTranslationCS{(prů)tok}  \textbf{2.} \dicDirectTranslationCS{(prů)běh (představení ap.)}
\dicEntry[renta] \dicTerm{rent|a\smash{\textsuperscript{1}}} \dicIPA{{r}{\textepsilon}{\textsubring{n}}{\textsubring{d}}{a}} \dicPos{f}[1] \dicFlx{(‑u, ‑ur)}[7] \dicSynonym{tekjur} \dicDirectTranslationCS{renta};  \dicPhraseIS{bera nafn með rentu} \dicLangCat{přen.} \dicDirectTranslationCS{dostát svému jménu}
\dicEntry[renta] \dicTerm{rent|a\smash{\textsuperscript{2}}} \dicIPA{{r}{\textepsilon}{\textsubring{n}}{\textsubring{d}}{a}} \dicPos{v}[1] \dicFlx{(‑aði)}[1] \dicFlx{acc} \dicPhraseIS{renta sig} \dicSynonym*{borga sig} \dicDirectTranslationCS{vyplatit se, vyplácet se, přinést\,/\addthin přinášet zisk}
\dicEntry[repja] \dicTerm{repj|a} \dicIPA{{r}{\textepsilon}{\textlengthmark}{\textsubring{b}}{j}{a}} \dicPos{f}[1] \dicFlx{(‑u, ‑ur)}[7] \dicFieldCat{bot.} \dicDirectTranslationCS{řepka, řepka olejka} \textit{(l.~{\textLA{Brassica napus}})}
\dicEntry[reri] \dicTerm{reri} \dicIPA{{r}{\textepsilon}{\textlengthmark}{r}{\textsci}} \dicPos{v} \dicFlx{ind\,/\addthin con pf sg 1 pers} \dicLink{róa\smash{\textsuperscript{2}}}
\dicEntry[rerum] \dicTerm{rerum} \dicIPA{{r}{\textepsilon}{\textlengthmark}{r}{\textscy}{\textsubring{m}}} \dicPos{v} \dicFlx{ind pf pl 1 pers} \dicLink{róa\smash{\textsuperscript{2}}}
\dicEntry[rest] \dicTerm{rest} \dicIPA{{r}{\textepsilon}{s}{\textsubring{d}}} \dicPos{f}[7] \dicFlx{(‑ar, ‑ir)}[1] \dicLangCat{hovor.} \dicSynonym{afgangur} \dicDirectTranslationCS{zbytek (jídla ap.)};  \dicPhraseIS{fyrir rest} \dicFlx{adv} \dicSynonym*{að lokum} \dicDirectTranslationCS{nakonec}
\dicEntry[revía] \dicTerm{reví|a} \dicIPA{{r}{\textepsilon}{\textlengthmark}{v}{i}{j}{a}} \dicPos{f}[1] \dicFlx{(‑u, ‑ur)}[7] \dicDirectTranslationCS{revue, (zábavné) pásmo, kabaret}
\dicEntry[rexa] \dicTerm{rex|a} \dicIPA{{r}{\textepsilon}{x}{s}{a}} \dicPos{v}[1] \dicFlx{(‑aði)}[1] \dicPhraseIS{rexa í e‑m} \dicSynonym*{skipta sér af} \dicDirectTranslationCS{otravovat (koho), rýt do (koho)}
\dicEntry[reyðarhvalur] \dicTerm{reyðar··hval|ur} \dicIPA{{r}{ei}{\textlengthmark}{ð}{a}{r}{k\smash{\textsuperscript{h}}}{v}{a}{l}{\textscy}{\textsubring{r}}} \dicPos{m}[9] \dicFlx{(‑s, ‑ir)}[9] \dicFieldCat{zool.} \dicDirectTranslationCS{kytovec z~čeledi plejtvákovitých} \textit{(l.~{\textLA{Balaenopteridae}})}
\dicEntry[reyfarakaup] \textls[15]{\dicTerm{reyfara··kaup} \dicIPA{{r}{ei}{\textlengthmark}{v}{a}{r}{a}{k\smash{\textsuperscript{h}}}{\oe i}{\textsubring{b}}} \dicPos{n}[2] \dicFlx{pl}[1] \dicLangCat{hovor.} \dicDirectTranslationCS{výhodná kou\nobreak pě} \dicExampleIS{gera reyfarakaup} \dicExampleCS{výhodně nakoupit}}
\dicEntry[reyfari] \dicTerm{reyf··ar|i} \dicIPA{{r}{ei}{\textlengthmark}{v}{a}{r}{\textsci}} \dicPos{m}[1] \dicFlx{(‑a, ‑ar)}[13] \dicFieldCat{lit.} \dicDirectTranslationCS{detektivní\,/\addthin napínavý román, thriller}
\dicEntry[reyfi] \dicTerm{reyfi} \dicIPA{{r}{ei}{\textlengthmark}{v}{\textsci}} \dicPos{n}[2] \dicFlx{(‑s, ‑)}[14] \dicDirectTranslationCS{ovčí rouno}
\dicEntry[reykelsi] \dicTerm{reyk··elsi} \dicIPA{{r}{ei}{\textlengthmark}{\r{\textObardotlessj}}{\textepsilon}{l}{s}{\textsci}} \dicPos{n}[2] \dicFlx{(‑s, ‑)}[14] \dicDirectTranslationCS{kadidlo, vonná tyčinka}
\dicEntry[reykháfur] \dicTerm{reyk··háf|ur} \dicIPA{{r}{ei}{\textlengthmark}{\r{g}}{h}{au}{v}{\textscy}{\textsubring{r}}} \dicPos{m}[6] \dicFlx{(‑s, ‑ar)}[22] \dicSynonym{skorsteinn} \dicDirectTranslationCS{komín}
\dicEntry[reyking] \dicTerm{reyk··ing} \dicIPA{{r}{ei}{\textlengthmark}{\r{\textObardotlessj}}{i}{\ng}{\r{g}}} \dicPos{f}[4] \dicFlx{(‑ar, ‑ar)}[5] \textbf{1.} \dicDirectTranslationCS{uzení}  \textbf{2.} \dicPhraseIS{reykingar} \dicFlx{pl} \dicDirectTranslationCS{kouření} \dicExampleIS{spilla heilsunni með reykingum} \dicExampleCS{ničit si zdraví kouřením};  \dicPhraseIS{reykingar bannaðar!} \dicDirectTranslationCS{zákaz kouření!}
\dicEntry[reykingamaður] \dicTerm{reykinga··|maður} \dicIPA{{r}{ei}{\textlengthmark}{\r{\textObardotlessj}}{i}{\ng}{\r{g}}{a}{m}{a}{ð}{\textscy}{\textsubring{r}}} \dicPos{m}[13] \dicFlx{(‑manns, ‑menn)}[2] \dicDirectTranslationCS{kuřák, kuřačka} \dicExampleIS{Hann er mikill reykingamaður.} \dicExampleCS{Je to velký kuřák.}
\dicEntry[reykja] \dicTerm{reyk|ja} \dicsymFrequent\  \dicIPA{{r}{ei}{\textlengthmark}{\r{\textObardotlessj}}{a}} \dicPos{v}[2] \dicFlx{(‑ti, ‑t)}[25] \dicFlx{acc} \textbf{1.} \dicDirectTranslationCS{(vy)kouřit} \dicExampleIS{reykja sígarettur} \dicExampleCS{kouřit cigarety}  \textbf{2.} \dicDirectTranslationCS{(vy)udit} \dicExampleIS{reykja fisk} \dicExampleCS{udit rybu}  \textbf{3.} \dicSynonym{rjúka} \dicDirectTranslationCS{kouřit, dýmat, čadit} \dicExampleIS{Kertið reykir.} \dicExampleCS{Svíčka kouří.}
\dicEntry[Reykjavík] \dicTerm{Reykja··vík} \dicsymFrequent\  \dicIPA{{r}{ei}{\textlengthmark}{\r{\textObardotlessj}}{a}{v}{i}{\r{g}}} \dicPos{f}[10] \dicFlx{(‑ur)}[3] \dicFieldCat{geog.} \dicDirectTranslationCS{Reykjavík} \dicIndirectTranslationCS{(hlavní město Islandu)} \dicExampleIS{Reykjavík er höfuðborg Íslands.} \dicExampleCS{Reykjavík je hlavní město Islandu.}
\dicEntry[reyklaus] \dicTerm{reyk··laus} \dicIPA{{r}{ei}{\r{g}}{l}{\oe i}{s}} \dicPos{adj}[5]\dicFlx{}[-1] \dicDirectTranslationCS{nekuřácký} \dicExampleIS{reyklaus stofa} \dicExampleCS{nekuřácká místnost}
\dicEntry[reykskynjari] \dicTerm{reyk··skynj·ar|i} \dicIPA{{r}{ei}{\r{g}}{s}{\r{\textObardotlessj}}{\textsci}{n}{j}{a}{r}{\textsci}} \dicPos{m}[1] \dicFlx{(‑a, ‑ar)}[13] \dicDirectTranslationCS{detektor kouře}
\dicEntry[reyktóbak] \dicTerm{reyk··tóbak} \dicIPA{{r}{ei}{\r{g}}{t\smash{\textsuperscript{h}}}{ou}{\textsubring{b}}{a}{\r{g}}} \dicPos{n}[2] \dicFlx{(‑s)}[2] \dicDirectTranslationCS{dýmkový tabák}
\dicEntry[reyktur] \dicTerm{reyktur} \dicIPA{{r}{ei}{x}{\textsubring{d}}{\textscy}{\textsubring{r}}} \dicPos{adj}[1]\dicFlx{}[-10] \dicDirectTranslationCS{uzený} \dicExampleIS{reykt kjöt} \dicExampleCS{uzené maso}
\dicEntry[reykur] \dicTerm{reyk|ur} \dicsymFrequent\  \dicIPA{{r}{ei}{\textlengthmark}{\r{g}}{\textscy}{\textsubring{r}}} \dicPos{m}[9] \dicFlx{(‑jar\,/\addthin ‑s, ‑ir)}[26] \dicDirectTranslationCS{kouř, dým} \dicExampleIS{reykur af eldi} \dicExampleCS{kouř ohně};  \dicPhraseIS{vaða reyk} \dicLangCat{přen.} \dicDirectTranslationCS{být vedle, mýlit se}
\dicEntry[reykvískur] \dicTerm{reyk··vískur} \dicIPA{{r}{ei}{\textlengthmark}{\r{g}}{v}{i}{s}{\r{g}}{\textscy}{\textsubring{r}}} \dicPos{adj}[1]\dicFlx{}[-1] \dicDirectTranslationCS{(jsoucí) z~Reykjavíku, reykjavický}
\dicEntry[reyna] \dicTerm{reyn|a} \dicsymFrequent\  \dicIPA{{r}{ei}{\textlengthmark}{n}{a}} \dicPos{v}[2] \dicFlx{(‑di, ‑t)}[141] \dicFlx{acc} \textbf{1.} \dicSynonym{þola} \dicDirectTranslationCS{zkusit, zakusit, prožít} \dicExampleIS{reyna margt á langri ævi} \dicExampleCS{zkusit mnoho během dlouhého života}  \textbf{2.} \dicSynonym{prófa} \dicDirectTranslationCS{zkusit, (vy)zkoušet, (o)testovat} \dicExampleIS{reyna ýmsar aðferðir} \dicExampleCS{vyzkoušet různé metody}  \textbf{3.} \dicSynonym{freista} \dicDirectTranslationCS{pokusit se, pokoušet se} \dicExampleIS{reyndu að sofa} \dicExampleCS{pokus se usnout};  \dicIdiom{reyna}[á]{ \dicPhraseIS{reyna á e‑ð}} {\textbf{a.}} \dicDirectTranslationCS{zkoušet (co) (výdrž ap.)};  {\textbf{b.}} \dicDirectTranslationCS{namáhat (co) (svaly ap.)}; { \dicPhraseIS{reyna á sig}} \dicDirectTranslationCS{vyvíjet fyzickou aktivitu, namáhat se}; { \dicPhraseIS{láta á e‑ð reyna}} \dicDirectTranslationCS{zkusit (co), risknout (co)};  \dicIdiom{reyna}[fyrir]{ \dicPhraseIS{reyna fyrir sér (sem leikari)}} \dicLangCat{přen.} \dicDirectTranslationCS{zkusit štěstí (jako herec)};  \dicIdiom{reyna}[við]{ \dicPhraseIS{reyna við e‑ð}} \dicDirectTranslationCS{zkusit (co), vyzkoušet (co) (vyluštit křížovku ap.)}; { \dicPhraseIS{reyna við e‑n}} \dicDirectTranslationCS{balit (koho), zkoušet to na (koho)};  \dicIdiom{reynast}{ \dicPhraseIS{reynast}} \dicFlx{refl} \dicDirectTranslationCS{ukázat se, jevit se, zdát se} \dicExampleIS{Þessi aðferð hefur reynst vel.} \dicExampleCS{Tato metoda se ukázala být dobrá.}; { \dicPhraseIS{reynast e‑m vel}} \dicFlx{refl} \dicDirectTranslationCS{osvědčit\,/\addthin osvědčovat se (komu)}
\dicEntry[reynd] \dicTerm{reynd} \dicIPA{{r}{ei}{n}{\textsubring{d}}} \dicPos{f}[7] \dicFlx{(‑ar, ‑ir)}[1] \dicSynonym{veruleiki} \dicDirectTranslationCS{skutečnost, realita, praxe};  \dicPhraseIS{í reynd} \dicFlx{adv} \dicDirectTranslationCS{ve skutečnosti};  \dicPhraseIS{reyndin er sú að} \dicDirectTranslationCS{skutečnost je taková, že}
\dicEntry[reyndar] \dicTerm{reyndar} \dicsymFrequent\  \dicIPA{{r}{ei}{n}{\textsubring{d}}{a}{\textsubring{r}}} \dicPos{adv} \dicSynonym*{að vísu} \dicDirectTranslationCS{vlastně, ve skutečnosti} \dicExampleIS{Ég held að það sé reyndar ekki tímabundið.} \dicExampleCS{Myslím, že ve skutečnosti to není dočasné.}
\dicEntry[reyndur] \dicTerm{reyndur} \dicIPA{{r}{ei}{n}{\textsubring{d}}{\textscy}{\textsubring{r}}} \dicPos{adj}[2]\dicFlx{}[-14] \dicDirectTranslationCS{zkušený, zběhlý, znalý} \dicExampleIS{vera reyndur skipstjóri} \dicExampleCS{být zkušený kapitán}
\dicEntry[reyniber] \dicTerm{reyni··ber} \dicIPA{{r}{ei}{\textlengthmark}{n}{\textsci}{\textsubring{b}}{\textepsilon}{\textsubring{r}}} \dicPos{n}[2] \dicFlx{(‑s, ‑)}[13] \dicDirectTranslationCS{jeřabina}
\dicEntry[reynir] \dicTerm{reyn|ir} \dicIPA{{r}{ei}{\textlengthmark}{n}{\textsci}{\textsubring{r}}} \dicPos{m}[7] \dicFlx{(‑is)}[2] \dicFieldCat{bot.} \dicDirectTranslationCS{jeřáb} \textit{(l.~{\textLA{Sorbus}})}  \dicsymPhoto\ 
\dicFigure{15952.jpg}{Reynir}{Reynir - Hanzlík Václav, Biolib, Copyright/CC-BY-NC}
\dicEntry[reyniviður] \dicTerm{reyni··við|ur} \dicIPA{{r}{ei}{\textlengthmark}{n}{\textsci}{v}{\textsci}{ð}{\textscy}{\textsubring{r}}} \dicPos{m}[10] \dicFlx{(‑ar, ‑ir)}[19] \dicFieldCat{bot.} \dicDirectTranslationCS{jeřáb ptačí} \textit{(l.~{\textLA{Sorbus aucuparia}})}  \dicsymPhoto\ 
\dicFigure{1343.jpg}{Reyniviður}{Reyniviður - Zicha Ondřej, Biolib, Copyright/CC-BY-NC}
\dicEntry[reynsla] \dicTerm{reynsl|a} \dicsymFrequent\  \dicIPA{{r}{ei}{n}{s}{\textsubring{d}}{l}{a}} \dicPos{f}[1] \dicFlx{(‑u)}[5] \textbf{1.} \dicSynonym{upplifun} \dicDirectTranslationCS{zkušenost, zážitek}  \textbf{2.} \dicSynonym{æfing} \dicDirectTranslationCS{cvik, zběhlost} \dicExampleIS{hafa reynslu í e‑u} \dicExampleCS{mít cvik v~(čem)}  \textbf{3.} \dicSynonym{prófun} \dicDirectTranslationCS{zkouška, zkušená} \dicExampleIS{til reynslu} \dicExampleCS{na zkušenou}
\dicEntry[reynslulaus] \dicTerm{reynslu··laus} \dicIPA{{r}{ei}{n}{s}{\textsubring{d}}{l}{\textscy}{l}{\oe i}{s}} \dicPos{adj}[5]\dicFlx{}[-1] \dicSynonym{óreyndur} \dicDirectTranslationCS{nezkušený}
\dicEntry[reynsluleysi] \dicTerm{reynslu··leysi} \dicIPA{{r}{ei}{n}{s}{\textsubring{d}}{l}{\textscy}{l}{ei}{s}{\textsci}} \dicPos{n}[2] \dicFlx{(‑s)}[20] \dicDirectTranslationCS{nezkušenost}
\dicEntry[reynslutími] \dicTerm{reynslu··tím|i} \dicIPA{{r}{ei}{n}{s}{\textsubring{d}}{l}{\textscy}{t\smash{\textsuperscript{h}}}{i}{m}{\textsci}} \dicPos{m}[1] \dicFlx{(‑a)}[3] \dicDirectTranslationCS{zkušební doba}
\dicEntry[reyr] \dicTerm{reyr} \dicIPA{{r}{ei}{\textlengthmark}{\textsubring{r}}} \dicPos{m}[4] \dicFlx{(‑s, ‑ar)}[14] \dicFieldCat{bot.} \dicDirectTranslationCS{tomka} \textit{(l.~{\textLA{Anthoxanthum}})}  \dicsymPhoto\ 
\dicFigure{51143.jpg}{Reyr}{Reyr - Forest \& Kim [[p:2684;Starr]], Biolib, CC-BY}
\dicEntry[reyra] \dicTerm{reyr|a} \dicIPA{{r}{ei}{\textlengthmark}{r}{a}} \dicPos{v}[2] \dicFlx{(‑ði, ‑t)}[99] \dicFlx{acc} \dicSynonym*{binda fast} \dicDirectTranslationCS{(pevně) uvázat\,/\addthin svázat}
\dicEntry[reyta] \dicTerm{reyt|a} \dicIPA{{r}{ei}{\textlengthmark}{\textsubring{d}}{a}} \dicPos{v}[2] \dicFlx{(‑ti, ‑t)}[54] \dicFlx{acc} \textbf{1.} \dicSynonym{plokka} \dicDirectTranslationCS{(o)škubat, vyškubat} \dicExampleIS{reyta fugl} \dicExampleCS{oškubat ptáka}  \textbf{2.} \dicSynonym*{lúa} \dicDirectTranslationCS{(vy)plít} \dicExampleIS{reyta arfa} \dicExampleCS{vyplít plevel};  \dicIdiom{reyta}[af]{ \dicPhraseIS{reyta af sér brandarana}} \dicLangCat{přen.} \dicDirectTranslationCS{sypat ze sebe vtipy}
\dicEntry[reytingur] \dicTerm{reyt··ing|ur} \dicIPA{{r}{ei}{\textlengthmark}{\textsubring{d}}{i}{\ng}{\r{g}}{\textscy}{\textsubring{r}}} \dicPos{m}[6] \dicFlx{(‑s)}[9] \dicDirectTranslationCS{hrstka, drobet} \dicExampleIS{reytingur af fiski} \dicExampleCS{drobet ryby}
\dicEntry[réð] \dicTerm{réð} \dicIPA{{r}{j}{\textepsilon}{\textlengthmark}{\texttheta}} \dicPos{v} \dicFlx{ind pf sg 1 pers} \dicLink{ráða}
\dicEntry[réði] \dicTerm{réði} \dicIPA{{r}{j}{\textepsilon}{\textlengthmark}{ð}{\textsci}} \dicPos{v} \dicFlx{con pf sg 1 pers} \dicLink{ráða}
\dicEntry[réðum] \dicTerm{réðum} \dicIPA{{r}{j}{\textepsilon}{\textlengthmark}{ð}{\textscy}{\textsubring{m}}} \dicPos{v} \dicFlx{ind pf pl 1 pers} \dicLink{ráða}
\dicEntry[réna] \dicTerm{rén|a} \dicIPA{{r}{j}{\textepsilon}{\textlengthmark}{n}{a}} \dicPos{v}[1] \dicFlx{(‑aði)}[45] \textbf{1.} \dicSynonym{minnka} \dicDirectTranslationCS{klesnout, klesat, opadnout, opadat (voda ap.)} \dicExampleIS{Áin rénar.} \dicExampleCS{Řeka klesá.}  \textbf{2.} \dicSynonym{þverra} \dicDirectTranslationCS{ubývat, menšit se} \dicExampleIS{Jökullinn rénar.} \dicExampleCS{Ubývá ledovce.};  \dicPhraseIS{e‑m rénar reiði} \dicFlx{impers} \dicDirectTranslationCS{(koho) přechází vztek}
\dicEntry[rénun] \dicTerm{rén|un} \dicIPA{{r}{j}{\textepsilon}{\textlengthmark}{n}{\textscy}{\textsubring{n}}} \dicPos{f}[7] \dicFlx{(‑unar)}[9] \dicSynonym{minnkun} \dicDirectTranslationCS{opadnutí, pokles, úbytek};  \dicPhraseIS{e‑að er í rénun} \dicDirectTranslationCS{(co) opadá, (co) je na ústupu} \dicExampleIS{Flóðið er í rénun.} \dicExampleCS{Záplavy již opadávají.}
\dicEntry[réri] \dicTerm{réri} \dicIPA{{r}{j}{\textepsilon}{\textlengthmark}{r}{\textsci}} \dicPos{v} \dicFlx{ind\,/\addthin con pf sg 1 pers} \dicLink{róa\smash{\textsuperscript{2}}}
\dicEntry[rérum] \dicTerm{rérum} \dicIPA{{r}{j}{\textepsilon}{\textlengthmark}{r}{\textscy}{\textsubring{m}}} \dicPos{v} \dicFlx{ind pf pl 1 pers} \dicLink{róa\smash{\textsuperscript{2}}}
\dicEntry[rétt] \dicTerm{rétt\smash{\textsuperscript{1}}} \dicIPA{{r}{j}{\textepsilon}{h}{\textsubring{d}}} \dicPos{f}[7] \dicFlx{(‑ar, ‑ir)}[1] \textbf{1.} \dicDirectTranslationCS{ohrada} \dicExampleIS{Það er margt fé inni í réttinni.} \dicExampleCS{V~ohradě je mnoho ovcí.}  \textbf{2.} \dicPhraseIS{réttir} \dicFlx{pl} \dicDirectTranslationCS{shánění ovcí} \dicIndirectTranslationCS{(každý podzim shánějí farmáři ovce z~hor na určené místo a~pak je třídí)} \dicExampleIS{fara í réttirnar} \dicExampleCS{shánět ovce}
\dicEntry[rétt] \dicTerm{rétt\smash{\textsuperscript{2}}} \dicIPA{{r}{j}{\textepsilon}{h}{\textsubring{d}}} \dicPos{adv} \dicFlx{(comp ‑ar, sup ‑ast)} \dicDirectTranslationCS{správně, dobře} \dicExampleIS{ef ég man rétt} \dicExampleCS{pokud si dobře pamatuji};  \dicPhraseIS{rétt áðan} \dicFlx{adv} \dicDirectTranslationCS{před chvílí\,/\addthin okamžikem}
\dicEntry[rétta] \dicTerm{rétt|a\smash{\textsuperscript{1}}} \dicIPA{{r}{j}{\textepsilon}{h}{\textsubring{d}}{a}} \dicPos{f}[1] \dicFlx{(‑u)}[5] \dicDirectTranslationCS{líc} \dicIndirectTranslationCS{(svrchní strana)} \dicExampleIS{réttan á peysunni} \dicExampleCS{líc svetru} \dicAntonym{ranga}
\dicEntry[rétta] \dicTerm{rétt|a\smash{\textsuperscript{2}}} \dicIPA{{r}{j}{\textepsilon}{h}{\textsubring{d}}{a}} \dicPos{v}[1] \dicFlx{(‑aði)}[1] \dicDirectTranslationCS{vést soudní spor}
\dicEntry[rétta] \dicTerm{rétt|a\smash{\textsuperscript{3}}} \dicsymFrequent\  \dicIPA{{r}{j}{\textepsilon}{h}{\textsubring{d}}{a}} \dicPos{v}[2] \dicFlx{(‑i, ‑)}[1] \dicFlx{(dat +) acc} \textbf{1.} \dicFlx{dat + acc} \dicSynonym*{fá e‑m e‑ð} \dicDirectTranslationCS{podat, podávat};  \dicPhraseIS{rétta e‑m e‑ð} \dicDirectTranslationCS{podat\,/\addthin podávat (komu co)} \dicExampleIS{rétta henni bréfið} \dicExampleCS{podat jí dopis}  \textbf{2.} \dicFlx{acc} \dicSynonym*{gera beinan} \dicDirectTranslationCS{narovnat, napřímit} \dicExampleIS{rétta stöngina} \dicExampleCS{napřímit stožár}  \textbf{3.} \dicFlx{acc} \dicSynonym*{teygja fram} \dicDirectTranslationCS{natáhnout, napřímit} \dicExampleIS{rétta upp hönd} \dicExampleCS{zvednout ruku};  \dicIdiom{rétta}[af]{ \dicPhraseIS{rétta e‑ð af}} \dicDirectTranslationCS{narovnat (co), srovnat (co), vyrovnat (co)};  \dicIdiom{rétta}[úr]{ \dicPhraseIS{rétta úr kútnum}} \dicLangCat{přen.} \dicDirectTranslationCS{vzpamatovat se, dát se dohromady}; { \dicPhraseIS{rétta úr sér}} \dicDirectTranslationCS{napřímit se, narovnat se};  \dicIdiom{rétta}[við]{ \dicPhraseIS{rétta e‑ð við}} \dicDirectTranslationCS{srovnat (co), napravit (co)};  \dicPhraseIS{rétta við} \dicLangCat{přen.} \dicDirectTranslationCS{sebrat se, dát se dohromady};  \dicIdiom{réttast}{ \dicPhraseIS{réttast}} \dicFlx{refl} \dicDirectTranslationCS{narovnat se, napřímit se}
\dicEntry[réttarfar] \dicTerm{réttar··far} \dicIPA{{r}{j}{\textepsilon}{h}{\textsubring{d}}{a}{\textsubring{r}}{f}{a}{\textsubring{r}}} \dicPos{n}[2] \dicFlx{(‑s)}[2] \dicFieldCat{práv.} \dicDirectTranslationCS{právní proces}
\dicEntry[réttarhald] \dicTerm{réttar··|hald} \dicIPA{{r}{j}{\textepsilon}{h}{\textsubring{d}}{a}{\textsubring{r}}{h}{a}{l}{\textsubring{d}}} \dicPos{n}[2] \dicFlx{(‑halds, ‑höld)}[8] \dicFieldCat{práv.} \dicDirectTranslationCS{přelíčení}
\dicEntry[réttarhlé] \dicTerm{réttar··hlé} \dicIPA{{r}{j}{\textepsilon}{h}{\textsubring{d}}{a}{\textsubring{r}}{\textsubring{l}}{j}{\textepsilon}} \dicPos{n}[2] \dicFlx{(‑s, ‑)}[5] \dicFieldCat{práv.} \dicDirectTranslationCS{soudní prázdniny}
\dicEntry[réttarkerfi] \dicTerm{réttar··kerfi} \dicIPA{{r}{j}{\textepsilon}{h}{\textsubring{d}}{a}{r}{c\smash{\textsuperscript{h}}}{\textepsilon}{r}{v}{\textsci}} \dicPos{n}[2] \dicFlx{(‑s, ‑)}[14] \dicFieldCat{práv.} \dicDirectTranslationCS{právní systém}
\dicEntry[réttarlæknir] \dicTerm{réttar··lækn|ir} \dicIPA{{r}{j}{\textepsilon}{h}{\textsubring{d}}{a}{r}{l}{a}{i}{h}{\r{g}}{n}{\textsci}{\textsubring{r}}} \dicPos{m}[7] \dicFlx{(‑is, ‑ar)}[1] \dicDirectTranslationCS{soudní lékař(ka)}
\begin{xtolerant}{}{1pt}
\dicEntry[réttarlæknisfræði] \dicTerm{réttar··læknis·fræð|i} \dicIPA{{r}\-{j}\-{\textepsilon}\-{h}\-{\textsubring{d}}\-{a}\-{r}\-{l}\-{a}\-{i}\-{h}\-{\r{g}}\-{n}\-{\textsci}\-{s}\-{f}\-{r}\-{a}\-{i}\-{ð}\-{\textsci}\-} \dicPos{f}[3] \dicFlx{(‑i)}[3] \dicDirectTranslationCS{soud\-ní lékařství}
\end{xtolerant}
\dicEntry[réttarmeinafræði] \dicTerm{réttar··meina·fræð|i} \dicIPA{{r}{j}{\textepsilon}{h}{\textsubring{d}}{a}{r}{m}{ei}{n}{a}{f}{r}{a}{i}{ð}{\textsci}} \dicPos{f}[3] \dicFlx{(‑i)}[3] \dicDirectTranslationCS{soudní patologie}
\dicEntry[réttarmeinafræðingur] \dicTerm{réttar··meina·fræð·ing|ur} \dicIPA{{r}\-{j}\-{\textepsilon}\-{h}\-{\textsubring{d}}\-{a}\-{r}\-{m}\-{ei}\-{n}\-{a}\-{f}\-{r}\-{a}\-{i}\-{ð}\-{i}\-{\ng}\-{\r{g}}\-{\textscy}\-{\textsubring{r}}\-} \dicPos{m}[6] \dicFlx{(‑s, ‑ar)}[8] \dicDirectTranslationCS{soudní patolog\,/\addthin patoložka}
\dicEntry[réttarríki] \dicTerm{réttar··ríki} \dicIPA{{r}{j}{\textepsilon}{h}{\textsubring{d}}{a}{r}{i}{\r{\textObardotlessj}}{\textsci}} \dicPos{n}[2] \dicFlx{(‑s, ‑)}[16] \dicFieldCat{pol.} \dicDirectTranslationCS{právní stát}
\dicEntry[réttarstaða] \dicTerm{réttar··|staða} \dicIPA{{r}{j}{\textepsilon}{h}{\textsubring{d}}{a}{\textsubring{r}}{s}{\textsubring{d}}{a}{ð}{a}} \dicPos{f}[1] \dicFlx{(‑stöðu)}[2] \dicFieldCat{práv.} \dicDirectTranslationCS{právní postavení}
\dicEntry[rétthendur] \dicTerm{rétt··hendur} \dicIPA{{r}{j}{\textepsilon}{h}{\textsubring{d}}{h}{\textepsilon}{n}{\textsubring{d}}{\textscy}{\textsubring{r}}} \dicPos{adj}[2]\dicFlx{}[-17] \dicDirectTranslationCS{pravoruký} \dicAntonym{örvhendur}
\dicEntry[rétthyrndur] \dicTerm{rétt··hyrndur} \dicIPA{{r}{j}{\textepsilon}{h}{\textsubring{d}}{h}{\textsci}{r}{n}{\textsubring{d}}{\textscy}{\textsubring{r}}} \dicPos{adj}[2]\dicFlx{}[-17] \dicFieldCat{mat.} \dicDirectTranslationCS{pravoúhlý}
\dicEntry[réttindalaus] \dicTerm{réttinda··laus} \dicIPA{{r}{j}{\textepsilon}{h}{\textsubring{d}}{\textsci}{n}{\textsubring{d}}{a}{l}{\oe i}{s}} \dicPos{adj}[5]\dicFlx{}[-1] \dicDirectTranslationCS{neoprávněný}
\dicEntry[réttindi] \dicTerm{rétt··indi} \dicIPA{{r}{j}{\textepsilon}{h}{\textsubring{d}}{\textsci}{n}{\textsubring{d}}{\textsci}} \dicPos{n}[2] \dicFlx{pl}[19] \textbf{1.} \dicSynonym{réttur\smash{\textsuperscript{1}}} \dicDirectTranslationCS{právo} \dicExampleIS{náttúrleg réttindi} \dicExampleCS{přirozené právo}  \textbf{2.} \dicSynonym{leyfi} \dicDirectTranslationCS{oprávnění, licence} \dicExampleIS{réttindi til skipstjórnar} \dicExampleCS{kapitánská licence}
\dicEntry[réttkjörinn] \dicTerm{rétt··kjörinn} \dicIPA{{r}{j}{\textepsilon}{h}{\textsubring{d}}{c\smash{\textsuperscript{h}}}{\oe}{r}{\textsci}{\textsubring{n}}} \dicPos{adj}[6]\dicFlx{}[-4] \dicDirectTranslationCS{řádně zvolený}
\dicEntry[réttlaus] \dicTerm{rétt··laus} \dicIPA{{r}{j}{\textepsilon}{h}{\textsubring{d}}{l}{\oe i}{s}} \dicPos{adj}[5]\dicFlx{}[-1] \dicDirectTranslationCS{bezprávný, (jsoucí) bez práv} \dicExampleIS{réttlausir þrælar} \dicExampleCS{bezprávní otroci}
\dicEntry[réttlátur] \dicTerm{rétt··látur} \dicIPA{{r}{j}{\textepsilon}{h}{\textsubring{d}}{l}{au}{\textsubring{d}}{\textscy}{\textsubring{r}}} \dicPos{adj}[1]\dicFlx{}[-1] \dicSynonym{réttsýnn} \dicDirectTranslationCS{spravedlivý} \dicExampleIS{réttlátur dómur} \dicExampleCS{spravedlivý rozsudek}
\dicEntry[réttlæta] \dicTerm{rétt··læt|a} \dicIPA{{r}{j}{\textepsilon}{h}{\textsubring{d}}{l}{a}{i}{\textsubring{d}}{a}} \dicPos{v}[2] \dicFlx{(‑ti, ‑t)}[54] \dicDirectTranslationCS{ospravedlnit, ospravedlňovat, omluvit, omlouvat} \dicExampleIS{réttlæta gerðir sínar} \dicExampleCS{ospravedlnit své činy}
\dicEntry[réttlæti] \dicTerm{rétt··læti} \dicIPA{{r}{j}{\textepsilon}{h}{\textsubring{d}}{l}{a}{i}{\textsubring{d}}{\textsci}} \dicPos{n}[2] \dicFlx{(‑s)}[20] \dicSynonym{réttsýni} \dicDirectTranslationCS{spravedlnost} \dicExampleIS{sýna e‑m réttlæti} \dicExampleCS{být ke (komu) spravedlivý}
\dicEntry[réttmæti] \dicTerm{rétt··mæti} \dicIPA{{r}{j}{\textepsilon}{h}{\textsubring{d}}{m}{a}{i}{\textsubring{d}}{\textsci}} \dicPos{n}[2] \dicFlx{(‑s)}[20] \dicDirectTranslationCS{zákonnost, legitimnost}
\dicEntry[réttmætur] \dicTerm{rétt··mætur} \dicIPA{{r}{j}{\textepsilon}{h}{\textsubring{d}}{m}{a}{i}{\textsubring{d}}{\textscy}{\textsubring{r}}} \dicPos{adj}[1]\dicFlx{}[-1] \textbf{1.} \dicSynonym{verðskuldaður} \dicDirectTranslationCS{oprávněný, odůvodněný} \dicExampleIS{réttmæt gagnrýni} \dicExampleCS{oprávněná kritika}  \textbf{2.} \dicSynonym{lögmætur} \dicDirectTranslationCS{zákonný, legitimní}
\dicEntry[réttritun] \dicTerm{rétt··rit|un} \dicIPA{{r}{j}{\textepsilon}{h}{\textsubring{d}}{r}{\textsci}{\textsubring{d}}{\textscy}{\textsubring{n}}} \dicPos{f}[7] \dicFlx{(‑unar)}[9] \dicDirectTranslationCS{pravopis, ortografie}
\dicEntry[réttritunarorðabók] \textls[30]{\dicTerm{rétt·ritunar··orða·|bók} \dicIPA{{r}\-{j}\-{\textepsilon}\-{h}\-{\textsubring{d}}\-{r}\-{\textsci}\-{\textsubring{d}}\-{\textscy}\-{n}\-{a}\-{r}\-{\textopeno}\-{r}\-{ð}\-{a}\-{\textsubring{b}}\-{ou}\-{\r{g}}\-} \dicPos{f}[8] \mbox{\dicFlx{(‑bókar, ‑bækur)}[5]}} \dicDirectTranslationCS{pravopisný slovník}
\dicEntry[réttsýni] \dicTerm{rétt··sýn|i} \dicIPA{{r}{j}{\textepsilon}{h}{\textsubring{d}}{s}{i}{n}{\textsci}} \dicPos{f}[3] \dicFlx{(‑i)}[3] \dicDirectTranslationCS{spravedlnost, nestrannost}
\dicEntry[réttsýnn] \dicTerm{rétt··sýnn} \dicIPA{{r}{j}{\textepsilon}{h}{\textsubring{d}}{s}{i}{\textsubring{d}}{\textsubring{n}}} \dicPos{adj}[7]\dicFlx{}[-1] \dicSynonym{réttlátur} \dicDirectTranslationCS{spravedlivý, nestranný}
\dicEntry[rétttrúaður] \dicTerm{rétt··trú·|aður} \dicIPA{{r}{j}{\textepsilon}{h}{\textsubring{d}}{t\smash{\textsuperscript{h}}}{r}{u}{a}{ð}{\textscy}{\textsubring{r}}} \dicPos{adj}[3] \dicFlx{(f ‑uð)}[3] \textbf{1.} \dicDirectTranslationCS{pravověrný, ortodoxní}  \textbf{2.} \dicDirectTranslationCS{pravoslavný}
\dicEntry[rétttrúnaður] \dicTerm{rétt··trú·nað|ur} \dicIPA{{r}{j}{\textepsilon}{h}{\textsubring{d}}{t\smash{\textsuperscript{h}}}{r}{u}{n}{a}{ð}{\textscy}{\textsubring{r}}} \dicPos{m}[10] \dicFlx{(‑ar)}[9] \textbf{1.} \dicFieldCat{náb.} \dicDirectTranslationCS{pravověrnost, ortodoxie}  \textbf{2.} \dicFieldCat{náb.} \dicDirectTranslationCS{pravoslaví}
\dicEntry[réttur] \dicTerm{rétt|ur\smash{\textsuperscript{1}}} \dicsymFrequent\  \dicIPA{{r}{j}{\textepsilon}{h}{\textsubring{d}}{\textscy}{\textsubring{r}}} \dicPos{m}[10] \dicFlx{(‑ar, ‑ir)}[4] \textbf{1.} \dicFieldCat{práv.} \dicSynonym{dómstóll} \dicDirectTranslationCS{soud, tribunál};  \dicPhraseIS{draga e‑n fyrir rétt} \dicDirectTranslationCS{postavit (koho) před soud}  \textbf{2.} \dicSynonym{réttindi} \dicDirectTranslationCS{právo, oprávnění, nárok};  \dicPhraseIS{áskilja sér rétt til þess að (gera e‑ð)} \dicDirectTranslationCS{vyhradit si právo ((co) udělat)};  \dicPhraseIS{e‑að á rétt á sér} \dicDirectTranslationCS{(co) je odůvodněné};  \dicPhraseIS{eiga rétt á e‑u} \dicDirectTranslationCS{mít právo na (co)};  \dicPhraseIS{hafa rétt til að (gera e‑ð)} \dicDirectTranslationCS{mít právo na (udělání (čeho))};  \dicPhraseIS{standa á rétti sínum} \dicDirectTranslationCS{trvat na svých právech};  \dicPhraseIS{vera í rétti} \dicDirectTranslationCS{mít přednost} \dicIndirectTranslationCS{(v~dopravním ruchu)}  \textbf{3.} \dicSynonym{matur} \dicDirectTranslationCS{jídlo, pokrm, chod} \dicExampleIS{Hver er réttur dagsins?} \dicExampleCS{Jaké jídlo je dnes v~nabídce?}
\dicEntry[réttur] \dicTerm{réttur\smash{\textsuperscript{2}}} \dicsymFrequent\  \dicIPA{{r}{j}{\textepsilon}{h}{\textsubring{d}}{\textscy}{\textsubring{r}}} \dicPos{adj}[1]\dicFlx{}[-10] \textbf{1.} \dicSynonym{villulaus} \dicDirectTranslationCS{správný, bezchybný} \dicExampleIS{rétt svar} \dicExampleCS{správná odpověď}  \textbf{2.} \dicSynonym{beinn} \dicDirectTranslationCS{přímý, rovný} \dicExampleIS{rétt lína} \dicExampleCS{rovná čára};  \dicIdiom{réttur}{ \dicPhraseIS{hafa á réttu að standa, hafa rétt fyrir sér}} \dicDirectTranslationCS{mít pravdu}
\dicEntry[réttvísi] \dicTerm{rétt··vís|i} \dicIPA{{r}{j}{\textepsilon}{h}{\textsubring{d}}{v}{i}{s}{\textsci}} \dicPos{f}[3] \dicFlx{(‑i)}[3] \textbf{1.} \dicSynonym{réttlæti} \dicDirectTranslationCS{spravedlnost} \dicIndirectTranslationCS{(nestranné hodnocení, uznávání práv druhého)}  \textbf{2.} \dicFieldCat{práv.} \dicSynonym*{ákæruvald} \dicDirectTranslationCS{spravedlnost} \dicIndirectTranslationCS{(rozhodování podle práva v~nějakém sporu)};  \dicPhraseIS{hindra framgang réttvísinnar} \dicDirectTranslationCS{bránit výkonu spravedlnosti}
\dicEntry[Réunion] \dicTerm{Réunion} \dicIPA{{r}{j}{\textepsilon}{\textscy}{n}{\textsci}{\textopeno}{\textsubring{n}}} \dicPos{n}[4] \dicFlx{indecl}[2] \dicFieldCat{geog.} \dicDirectTranslationCS{Réunion}
\dicEntry[riddarakross] \dicTerm{riddara··kross} \dicIPA{{r}{\textsci}{\textsubring{d}}{\textlengthmark}{a}{r}{a}{k\smash{\textsuperscript{h}}}{r}{\textopeno}{s}} \dicPos{m}[4] \dicFlx{(‑, ‑ar)}[28] \dicFieldCat{pol.} \dicDirectTranslationCS{Rytířský kříž} \dicIndirectTranslationCS{(nejvyšší islandské vyznamenání)}
\dicEntry[riddaralegur] \dicTerm{riddara··legur} \dicIPA{{r}{\textsci}{\textsubring{d}}{\textlengthmark}{a}{r}{a}{l}{\textepsilon}{\textbabygamma}{\textscy}{\textsubring{r}}} \dicPos{adj}[1]\dicFlx{}[-8] \dicDirectTranslationCS{rytířský, galantní, dvorný} \dicExampleIS{riddaraleg kurteisi} \dicExampleCS{rytířská zdvořilost}
\dicEntry[riddarasaga] \dicTerm{riddara··|saga} \dicIPA{{r}{\textsci}{\textsubring{d}}{\textlengthmark}{a}{r}{a}{s}{a}{\textbabygamma}{a}} \dicPos{f}[1] \dicFlx{(‑sögu, ‑sögur)}[14] \dicFieldCat{lit.} \dicDirectTranslationCS{rytířský román}
\dicEntry[riddari] \dicTerm{ridd··ar|i} \dicIPA{{r}{\textsci}{\textsubring{d}}{\textlengthmark}{a}{r}{\textsci}} \dicPos{m}[1] \dicFlx{(‑a, ‑ar)}[13] \textbf{1.} \dicSynonym*{ríðandi hermaður} \dicDirectTranslationCS{rytíř} \dicIndirectTranslationCS{(feudální šlechtic)}  \textbf{2.} \dicDirectTranslationCS{rytíř} \dicIndirectTranslationCS{(nositel určitého stupně některého řádu)};  \dicPhraseIS{slá e‑n til riddara} \dicDirectTranslationCS{pasovat (koho) na rytíře}  \textbf{3.} \dicDirectTranslationCS{jezdec, kůň} \dicIndirectTranslationCS{(šachová figura)}
\dicEntry[riða] \dicTerm{rið|a\smash{\textsuperscript{1}}} \dicIPA{{r}{\textsci}{\textlengthmark}{ð}{a}} \dicPos{f}[1] \dicFlx{(‑u)}[5] \textbf{1.} \dicSynonym{skjálfti} \dicDirectTranslationCS{třesení, chvění se, zachvívání se}  \textbf{2.} \dicFieldCat{med.} \dicDirectTranslationCS{scrapie, klusavka} \dicIndirectTranslationCS{(nemoc ovcí a~koz)}
\dicEntry[riða] \dicTerm{rið|a\smash{\textsuperscript{2}}} \dicIPA{{r}{\textsci}{\textlengthmark}{ð}{a}} \dicPos{v}[1] \dicFlx{(‑aði)}[1] \dicSynonym{titra} \dicDirectTranslationCS{třást se, chvět se, zachvívat se} \dicExampleIS{riða í hverju spori} \dicExampleCS{zachvívat se při každém kroku};  \dicPhraseIS{riða til falls} \dicLangCat{přen.} \dicDirectTranslationCS{směřovat k~zániku}
\dicEntry[riði] \dicTerm{riði} \dicIPA{{r}{\textsci}{\textlengthmark}{ð}{\textsci}} \dicPos{v} \dicFlx{con pf sg 1 pers} \dicLink{ríða}
\dicEntry[riðið] \dicTerm{riðið} \dicIPA{{r}{\textsci}{\textlengthmark}{ð}{\textsci}{\texttheta}} \dicPos{v} \dicFlx{supin} \dicLink{ríða}
\dicEntry[riðill] \dicTerm{rið|ill} \dicIPA{{r}{\textsci}{\textlengthmark}{ð}{\textsci}{\textsubring{d}}{\textsubring{l}}} \dicPos{m}[6] \dicFlx{(‑ils, ‑lar)}[35] \dicFieldCat{sport.} \dicSynonym{hópur} \dicDirectTranslationCS{skupina (v~turnaji ap.)} \dicExampleIS{skipta mönnum í riðla} \dicExampleCS{rozdělit muže do skupin}
\dicEntry[riðla] \dicTerm{riðl|a} \dicIPA{{r}{\textsci}{ð}{l}{a}} \dicPos{v}[1] \dicFlx{(‑aði)}[1] \dicFlx{dat} \dicSynonym{aflaga\smash{\textsuperscript{2}}} \dicDirectTranslationCS{zpřeházet, pomíchat, způsobit zmatek};  \dicIdiom{riðlast}{ \dicPhraseIS{riðlast}} \dicFlx{refl} \dicSynonym*{ruglast} \dicDirectTranslationCS{zpřeházet se, pomíchat se} \dicExampleIS{Liðið riðlaðist.} \dicExampleCS{Skupina se pomíchala.}
\dicEntry[riðstraumur] \dicTerm{rið··straum|ur} \dicIPA{{r}{\textsci}{ð}{s}{\textsubring{d}}{r}{\oe i}{m}{\textscy}{\textsubring{r}}} \dicPos{m}[6] \dicFlx{(‑s, ‑ar)}[22] \dicFieldCat{elek.} \dicDirectTranslationCS{střídavý proud} \dicAntonym{jafnstraumur}
\dicEntry[riðum] \dicTerm{riðum} \dicIPA{{r}{\textsci}{\textlengthmark}{ð}{\textscy}{\textsubring{m}}} \dicPos{v} \dicFlx{ind pf pl 1 pers} \dicLink{ríða}
\dicEntry[riðvaxinn] \dicTerm{rið··vaxinn} \dicIPA{{r}{\textsci}{ð}{v}{a}{x}{s}{\textsci}{\textsubring{n}}} \dicPos{adj}[6]\dicFlx{}[-3] \dicSynonym{þrekvaxinn} \dicDirectTranslationCS{podsaditý, zavalitý}
\dicEntry[rif] \dicTerm{rif} \dicIPA{{r}{\textsci}{\textlengthmark}{f}} \dicPos{n}[2] \dicFlx{(‑s, ‑)}[6] \textbf{1.} \dicFieldCat{anat.} \dicSynonym{rifbein} \dicDirectTranslationCS{žebro};  \dicPhraseIS{hafa ráð undir rifi hverju} \dicLangCat{přen.} \dicDirectTranslationCS{mít na všechno odpověď};  \dicPhraseIS{e‑m rennur e‑að til rifja} \dicFlx{impers} \dicDirectTranslationCS{(co koho) znepokojuje}  \textbf{2.} \dicSynonym{sandrif} \dicDirectTranslationCS{rif, (podmořský) útes}
\dicEntry[rifa] \dicTerm{rif|a} \dicIPA{{r}{\textsci}{\textlengthmark}{v}{a}} \dicPos{f}[1] \dicFlx{(‑u, ‑ur)}[19] \dicSynonym{glenna\smash{\textsuperscript{1}}} \dicDirectTranslationCS{prasklina, trhlina} \dicExampleIS{rifa á veggnum} \dicExampleCS{prasklina ve zdi}
\dicEntry[rifbein] \dicTerm{rif··bein} \dicIPA{{r}{\textsci}{v}{\textsubring{b}}{ei}{\textsubring{n}}} \dicPos{n}[2] \dicFlx{(‑s, ‑)}[5] \dicFieldCat{anat.} \dicDirectTranslationCS{žebro}
\dicEntry[riffill] \dicTerm{riff|ill} \dicIPA{{r}{\textsci}{f}{\textlengthmark}{\textsci}{\textsubring{d}}{\textsubring{l}}} \dicPos{m}[6] \dicFlx{(‑ils, ‑lar)}[35] \dicDirectTranslationCS{puška, ručnice}
\dicEntry[rifflaður] \dicTerm{riffl··|aður} \dicIPA{{r}{\textsci}{f}{l}{a}{ð}{\textscy}{\textsubring{r}}} \dicPos{adj}[3] \dicFlx{(f ‑uð)}[3] \dicDirectTranslationCS{žebrovaný, žebrový} \dicExampleIS{rifflað flauel} \dicExampleCS{žebrovaný samet}
\dicEntry[rifi] \dicTerm{rifi} \dicIPA{{r}{\textsci}{\textlengthmark}{v}{\textsci}} \dicPos{v} \dicFlx{con pf sg 1 pers} \dicLink{rífa}
\dicEntry[rifið] \dicTerm{rifið} \dicIPA{{r}{\textsci}{\textlengthmark}{v}{\textsci}{\texttheta}} \dicPos{v} \dicFlx{supin} \dicLink{rífa}
\dicEntry[rifja] \dicTerm{rifj|a} \dicsymFrequent\  \dicIPA{{r}{\textsci}{v}{j}{a}} \dicPos{v}[1] \dicFlx{(‑aði)}[1] \dicFlx{acc} \dicSynonym*{snúa heyi} \dicDirectTranslationCS{uhrabat seno (do řádku)};  \dicIdiom{rifja}[upp]{ \dicPhraseIS{rifja e‑ð upp}} \dicSynonym*{muna} \dicDirectTranslationCS{připomenout si (co), osvěžit si (co), oprášit si (co), zopakovat si (co)} \dicExampleIS{rifja upp atburði síðasta árs} \dicExampleCS{připomenout si události minulého roku};  \dicIdiom{rifjast}[upp]{ \dicPhraseIS{e‑að rifjast upp}} \dicFlx{refl} \dicDirectTranslationCS{(co) se vybavuje\,/\addthin připomíná};  \dicPhraseIS{e‑að rifjast upp fyrir e‑m} \dicFlx{refl} \dicDirectTranslationCS{(co) se (komu) vybavuje\,/\addthin připomíná}
\dicEntry[rifjárn] \dicTerm{rif··járn} \dicIPA{{r}{\textsci}{v}{j}{au}{r}{\textsubring{d}}{\textsubring{n}}} \dicPos{n}[2] \dicFlx{(‑s, ‑)}[5] \dicDirectTranslationCS{struhadlo}
\dicEntry[riflás] \dicTerm{rif··lás} \dicIPA{{r}{\textsci}{v}{l}{au}{s}} \dicPos{m}[4] \dicFlx{(‑s, ‑ar)}[25] \dicDirectTranslationCS{suchý zip}
\dicEntry[rifna] \dicTerm{rifn|a} \dicIPA{{r}{\textsci}{\textsubring{b}}{n}{a}} \dicPos{v}[1] \dicFlx{(‑aði)}[34] \dicDirectTranslationCS{roztrhnout se, rozervat se} \dicExampleIS{Blaðið rifnaði í tvennt.} \dicExampleCS{List se roztrhl na dva kusy.};  \dicIdiom{rifna}[upp]{ \dicPhraseIS{e‑að rifnar upp}} {\textbf{a.}} \dicDirectTranslationCS{(co) se znovu otevírá\,/\addthin trhá (rána ap.)} \dicExampleIS{Sárið rifnaði upp.} \dicExampleCS{Rána se znovu otevřela.};  {\textbf{b.}} \dicDirectTranslationCS{(co) je vyvráceno (strom ap.)} \dicExampleIS{Tré hafa rifnað upp með rótum.} \dicExampleCS{Stromy byly vyvráceny i~s~kořeny.};  \dicIdiom{rifna}[úr]{ \dicPhraseIS{rifna úr e‑u}} \dicDirectTranslationCS{svíjet se (čím), moci se potrhat (čím) (smíchy ap.)} \dicExampleIS{vera að rifna úr hlátri} \dicExampleCS{moci se potrhat smíchy}
\dicEntry[rifrildi] \dicTerm{rif··rildi} \dicIPA{{r}{\textsci}{v}{r}{\textsci}{l}{\textsubring{d}}{\textsci}} \dicPos{n}[2] \dicFlx{(‑s, ‑)}[14] \textbf{1.} \dicSynonym{illdeila} \dicDirectTranslationCS{hádka, spor, pře}  \textbf{2.} \dicSynonym{slitur} \dicDirectTranslationCS{útržek, cár}
\dicEntry[rifs] \dicTerm{rifs} \dicIPA{{r}{\textsci}{f}{s}} \dicPos{n}[2] \dicFlx{(‑)}[24] \dicFieldCat{bot.} \dicDirectTranslationCS{červený rybíz} \dicIndirectTranslationCS{(keř)} \textit{(l.~{\textLA{Ribes rubrum}})}  \dicsymPhoto\ 
\dicFigure{ds_image_rifs_0_1.jpg}{Rifs}{Rifs - Vmenkov, CC BY-SA 3.0}
\dicEntry[rifsber] \dicTerm{rifs··ber} \dicIPA{{r}{\textsci}{f}{s}{\textsubring{b}}{\textepsilon}{\textsubring{r}}} \dicPos{n}[2] \dicFlx{(‑s, ‑)}[13] \dicDirectTranslationCS{červený rybíz} \dicIndirectTranslationCS{(plod)}
\dicEntry[rifta] \dicTerm{rift|a} \dicIPA{{r}{\textsci}{f}{\textsubring{d}}{a}} \dicPos{v}[3] \dicFlx{(‑aði) (‑i, ‑)}[5] \dicFlx{dat} \dicSynonym{afturkalla} \dicDirectTranslationCS{zrušit, zneplatnit, anulovat} \dicExampleIS{rifta samningi} \dicExampleCS{anulovat smlouvu}
\dicEntry[rifting] \dicTerm{rift··ing} \dicIPA{{r}{\textsci}{f}{\textsubring{d}}{i}{\ng}{\r{g}}} \dicPos{f}[4] \dicFlx{(‑ar, ‑ar)}[5] \dicLink{riftun}
\dicEntry[riftun] \dicTerm{rift|un}\dicTerm{, rifting} \dicIPA{{r}\-{\textsci}\-{f}\-{\textsubring{d}}\-{\textscy}\-{\textsubring{n}}\-} \dicPos{f}[7] \dicFlx{(‑unar, ‑anir)}[8] \dicDirectTranslationCS{zrušení, zneplatnění, anulace}
\dicEntry[rifum] \dicTerm{rifum} \dicIPA{{r}{\textsci}{\textlengthmark}{v}{\textscy}{\textsubring{m}}} \dicPos{v} \dicFlx{ind pf pl 1 pers} \dicLink{rífa}
\dicEntry[rigna] \dicTerm{rign|a} \dicsymFrequent\  \dicIPA{{r}{\textsci}{\r{g}}{n}{a}} \dicPos{v}[2] \dicFlx{(‑di, ‑t)}[163] \dicFlx{impers} \textbf{1.} \dicPhraseIS{það\,/\addthin hann rignir} \dicDirectTranslationCS{prší, padá déšť};  \dicPhraseIS{það rignir eldi og brennisteini} \dicLangCat{přen.} \dicDirectTranslationCS{padají trakaře}  \textbf{2.} \dicPhraseIS{e‑u rignir} \dicDirectTranslationCS{(co) prší, (co) padá jako déšť}
\dicEntry[rigning] \dicTerm{rign··ing} \dicsymFrequent\  \dicIPA{{r}{\textsci}{\r{g}}{n}{i}{\ng}{\r{g}}} \dicPos{f}[4] \dicFlx{(‑ar, ‑ar)}[5] \dicSynonym{regn} \dicDirectTranslationCS{déšť, deštivé počasí} \dicExampleIS{Næstu helgi er gert ráð fyrir rigningu fyrir norðan.} \dicExampleCS{Příští víkend se na severu očekává déšť.}
\dicEntry[rigningarlegur] \dicTerm{rigningar··legur} \dicIPA{{r}{\textsci}{\r{g}}{n}{i}{\ng}{\r{g}}{a}{r}{l}{\textepsilon}{\textbabygamma}{\textscy}{\textsubring{r}}} \dicPos{adj}[1]\dicFlx{}[-8] \dicDirectTranslationCS{deštivý, dešťový} \dicExampleIS{rigningarleg ský} \dicExampleCS{dešťové mraky}
\dicEntry[rim] \dicTerm{rim} \dicIPA{{r}{\textsci}{\textlengthmark}{\textsubring{m}}} \dicPos{f}[4] \dicFlx{(‑ar, ‑ar)}[1] \textbf{1.} \dicSynonym{rimill} \dicDirectTranslationCS{příčka, příčel (žebříku ap.)}  \textbf{2.} \dicSynonym{þrep} \dicDirectTranslationCS{schod}
\dicEntry[rimill] \dicTerm{rim|ill} \dicIPA{{r}{\textsci}{\textlengthmark}{m}{\textsci}{\textsubring{d}}{\textsubring{l}}} \dicPos{m}[6] \dicFlx{(‑ils, ‑lar)}[35] \dicDirectTranslationCS{příčka, příčel}
\dicEntry[rimlagluggatjöld] \dicTerm{rimla··glugga·tjöld} \dicIPA{{r}{\textsci}{m}{l}{a}{\r{g}}{l}{\textscy}{\r{g}}{a}{t\smash{\textsuperscript{h}}}{j}{\oe}{l}{\textsubring{d}}} \dicPos{n}[2] \dicFlx{pl}[9] \dicSynonym{rimlatjöld} \dicDirectTranslationCS{žaluzie}
\dicEntry[rimlahlið] \dicTerm{rimla··hlið} \dicIPA{{r}{\textsci}{m}{l}{a}{\textsubring{l}}{\textsci}{\texttheta}} \dicPos{n}[2] \dicFlx{(‑s, ‑)}[5] \dicDirectTranslationCS{mřížový vjezd} \dicIndirectTranslationCS{(silniční vjezd plnící funkci brány mezi oploceními; tvořený z~mříže a~z~toho důvodu neschůdný pro dobytek)}
\dicEntry[rimlatjöld] \dicTerm{rimla··tjöld} \dicIPA{{r}{\textsci}{m}{l}{a}{t\smash{\textsuperscript{h}}}{j}{\oe}{l}{\textsubring{d}}} \dicPos{n}[2] \dicFlx{pl}[9] \dicDirectTranslationCS{žaluzie}
\dicEntry[rimpa] \dicTerm{rimp|a} \dicIPA{{r}{\textsci}{\textsubring{m}}{\textsubring{b}}{a}} \dicPos{v}[1] \dicFlx{(‑aði)}[1] \dicFlx{acc} \dicSynonym{staga} \dicDirectTranslationCS{zašít, zašívat, (za)látat} \dicExampleIS{rimpa í sokkana} \dicExampleCS{zalátat ponožky};  \dicIdiom{rimpa}[saman]{ \dicPhraseIS{rimpa saman e‑ð}} \dicDirectTranslationCS{sešít (co), slátat (co)} \dicExampleIS{rimpa rifuna saman} \dicExampleCS{sešít trhlinu}
\dicEntry[ringlaður] \dicTerm{ringl··|aður} \dicsymFrequent\  \dicIPA{{r}{i}{\ng}{\r{g}}{l}{a}{ð}{\textscy}{\textsubring{r}}} \dicPos{adj}[3] \dicFlx{(f ‑uð)}[3] \dicSynonym*{utan við sig} \dicDirectTranslationCS{zmatený, popletený} \dicExampleIS{Ég var ringluð og skildi ekki hvað var að gerast.} \dicExampleCS{Byla jsem zmatená a~nechápala jsem, co se děje.}
\dicEntry[ringulreið] \dicTerm{ringul··reið} \dicIPA{{r}{i}{\ng}{\r{g}}{\textscy}{l}{r}{ei}{\texttheta}} \dicPos{f}[4] \dicFlx{(‑ar)}[31] \dicSynonym{óreiða} \dicDirectTranslationCS{chaos, zmatek} \dicExampleIS{skapa ringulreið} \dicExampleCS{vytvořit zmatek}
\dicEntry[ris] \dicTerm{ris} \dicIPA{{r}{\textsci}{\textlengthmark}{s}} \dicPos{n}[2] \dicFlx{(‑s, ‑)}[5] \textbf{1.} \dicDirectTranslationCS{podkroví} \dicExampleIS{íbúð í risi} \dicExampleCS{byt v~podkroví}  \textbf{2.} \dicFieldCat{lit.} \dicDirectTranslationCS{vyvrcholení, klimax}  \textbf{3.} \dicFieldCat{geol.} \dicSynonym*{landris} \dicDirectTranslationCS{vzedmutí, výzdvih};  \dicIdiom{ris}{ \dicPhraseIS{það er ekki hátt á e‑m risið}} \dicLangCat{přen.} \dicDirectTranslationCS{(kdo) neoplývá zrovna sebejistotou}
\dicEntry[risaeðla] \dicTerm{risa··eðl|a} \dicIPA{{r}{\textsci}{\textlengthmark}{s}{a}{\textepsilon}{ð}{l}{a}} \dicPos{f}[1] \dicFlx{(‑u, ‑ur)}[7] \dicFieldCat{zool.} \dicDirectTranslationCS{dinosaurus} \textit{(l.~{\textLA{Dinosauria}})}
\dicEntry[risapanda] \dicTerm{risa··|panda} \dicIPA{{r}{\textsci}{\textlengthmark}{s}{a}{p\smash{\textsuperscript{h}}}{a}{n}{\textsubring{d}}{a}} \dicPos{f}[1] \dicFlx{(‑pöndu, ‑pöndur)}[20] \dicLink{pandabjörn}
\dicEntry[risastór] \dicTerm{risa··stór} \dicIPA{{r}{\textsci}{\textlengthmark}{s}{a}{s}{\textsubring{d}}{ou}{\textsubring{r}}} \dicPos{adj}[10] \dicFlx{(f ‑)}[18] \dicDirectTranslationCS{obrovský, obrovitý, gigantický}
\dicEntry[risavaxinn] \dicTerm{risa··vaxinn} \dicIPA{{r}{\textsci}{\textlengthmark}{s}{a}{v}{a}{x}{s}{\textsci}{\textsubring{n}}} \dicPos{adj}[6]\dicFlx{}[-3] \dicDirectTranslationCS{obrovský, ohromný, nesmírný, gigantický} \dicExampleIS{risavaxnar skuldir} \dicExampleCS{obrovské dluhy}
\dicEntry[risi] \dicTerm{ris|i\smash{\textsuperscript{1}}} \dicsymFrequent\  \dicIPA{{r}{\textsci}{\textlengthmark}{s}{\textsci}} \dicPos{m}[1] \dicFlx{(‑a, ‑ar)}[1] \textbf{1.} \dicFieldCat{pov.} \dicSynonym{jötunn} \dicDirectTranslationCS{obr} \dicExampleIS{bardagi milli goða og risa} \dicExampleCS{bitva mezi bohy a~obry}  \textbf{2.} \dicDirectTranslationCS{obr, kolos, gigant} \dicIndirectTranslationCS{(tvor\,/\addthin věc mimořádné velikosti)}
\dicEntry[risi] \dicTerm{risi\smash{\textsuperscript{2}}} \dicIPA{{r}{\textsci}{\textlengthmark}{s}{\textsci}} \dicPos{v} \dicFlx{con pf sg 1 pers} \dicLink{rísa}
\dicEntry[risið] \dicTerm{risið} \dicIPA{{r}{\textsci}{\textlengthmark}{s}{\textsci}{\texttheta}} \dicPos{v} \dicFlx{supin} \dicLink{rísa}
\dicEntry[rismæltur] \dicTerm{ris··mæltur} \dicIPA{{r}{\textsci}{\textlengthmark}{s}{m}{a}{i}{\textsubring{l}}{\textsubring{d}}{\textscy}{\textsubring{r}}} \dicPos{adj}[1]\dicFlx{}[-13] \dicFieldCat{jaz.} \dicDirectTranslationCS{retroflexní, cerebrální}
\dicEntry[rispa] \dicTerm{risp|a\smash{\textsuperscript{1}}} \dicIPA{{r}{\textsci}{s}{\textsubring{b}}{a}} \dicPos{f}[1] \dicFlx{(‑u, ‑ur)}[19] \textbf{1.} \dicSynonym{rák} \dicDirectTranslationCS{škrábanec, škrábnutí} \dicExampleIS{Það kom rispa í glerið.} \dicExampleCS{Na skle se objevil škrábanec.}  \textbf{2.} \dicSynonym{törn} \dicDirectTranslationCS{zabíračka, dřina}
\dicEntry[rispa] \dicTerm{risp|a\smash{\textsuperscript{2}}} \dicIPA{{r}{\textsci}{s}{\textsubring{b}}{a}} \dicPos{v}[1] \dicFlx{(‑aði)}[1] \dicFlx{acc} \dicSynonym{klóra} \dicDirectTranslationCS{poškrábat, podrápat} \dicExampleIS{rispa glerið} \dicExampleCS{poškrábat sklo};  \dicIdiom{rispast}{ \dicPhraseIS{rispast}} \dicFlx{refl} \dicDirectTranslationCS{poškrábat se, podrápat se} \dicIndirectTranslationCS{(poškodit se škrábáním)}
\dicEntry[riss] \dicTerm{riss} \dicIPA{{r}{\textsci}{s}{\textlengthmark}} \dicPos{n}[2] \dicFlx{(‑)}[24] \dicDirectTranslationCS{črta, náčrtek, skica} \dicExampleIS{rissmynd} \dicExampleCS{skica}
\dicEntry[rissa] \dicTerm{riss|a} \dicIPA{{r}{\textsci}{s}{\textlengthmark}{a}} \dicPos{v}[1] \dicFlx{(‑aði)}[1] \dicFlx{acc} \dicDirectTranslationCS{načrtnout, črtat, (na)skicovat} \dicExampleIS{rissa upp e‑ð} \dicExampleCS{načrtnout (co)}
\dicEntry[rist] \dicTerm{rist} \dicIPA{{r}{\textsci}{s}{\textsubring{d}}} \dicPos{f}[7] \dicFlx{(‑ar, ‑ir)}[1] \textbf{1.} \dicFieldCat{anat.} \dicDirectTranslationCS{nárt}  \textbf{2.} \dicSynonym*{járnrimlagrind} \dicDirectTranslationCS{rošt}  \textbf{3.} \dicSynonym{brauðrist} \dicDirectTranslationCS{toustovač, touster, toaster}
\dicEntry[rista] \dicTerm{rist|a\smash{\textsuperscript{1}}} \dicIPA{{r}{\textsci}{s}{\textsubring{d}}{a}} \dicPos{f}[1] \dicFlx{(‑u, ‑ur)}[7] \dicSynonym{skora\smash{\textsuperscript{1}}} \dicDirectTranslationCS{zářez, (za)říznutí}
\dicEntry[rista] \dicTerm{rist|a\smash{\textsuperscript{2}}} \dicIPA{{r}{\textsci}{s}{\textsubring{d}}{a}} \dicPos{v}[1] \dicFlx{(‑aði)}[1] \dicFlx{acc} \dicDirectTranslationCS{opéct, opékat} \dicIndirectTranslationCS{(v~toustovači)} \dicExampleIS{rista brauð} \dicExampleCS{opékat chleba (v~toustovači)}
\dicEntry[rista] \dicTerm{rist|a\smash{\textsuperscript{3}}} \dicIPA{{r}{\textsci}{s}{\textsubring{d}}{a}} \dicPos{v}[2] \dicFlx{(‑i, ‑)}[16] \dicFlx{acc} \textbf{1.} \dicSynonym{skera} \dicDirectTranslationCS{(po)krájet, (po)řezat} \dicExampleIS{rista tréð í sundur} \dicExampleCS{rozřezat dřevo na kousky}  \textbf{2.} \dicDirectTranslationCS{nořit se, mít ponor} \dicExampleIS{Báturinn ristir djúpt.} \dicExampleCS{Člun má hluboký ponor.};  \dicPhraseIS{e‑að ristir ekki djúpt} \dicLangCat{přen.} \dicDirectTranslationCS{(co) není moc zakořeněné (vztek ap.)}
\dicEntry[ristill] \dicTerm{rist|ill} \dicIPA{{r}{\textsci}{s}{\textsubring{d}}{\textsci}{\textsubring{d}}{\textsubring{l}}} \dicPos{m}[6] \dicFlx{(‑ils, ‑lar)}[35] \textbf{1.} \dicFieldCat{anat.} \dicDirectTranslationCS{tračník}  \textbf{2.} \dicFieldCat{med.} \dicDirectTranslationCS{pásový opar} \textit{(l.~{\textLA{Herpes zoster}})}
\dicEntry[risum] \dicTerm{risum} \dicIPA{{r}{\textsci}{\textlengthmark}{s}{\textscy}{\textsubring{m}}} \dicPos{v} \dicFlx{ind pf pl 1 pers} \dicLink{rísa}
\dicEntry[rit] \dicTerm{rit} \dicsymFrequent\  \dicIPA{{r}{\textsci}{\textlengthmark}{\textsubring{d}}} \dicPos{n}[2] \dicFlx{(‑s, ‑)}[5] \dicDirectTranslationCS{(psaná) práce, (literární) dílo, spis} \dicExampleIS{ódauðlegt rit} \dicExampleCS{nesmrtelné dílo}
\dicEntry[rita] \dicTerm{rit|a\smash{\textsuperscript{1}}} \dicIPA{{r}{\textsci}{\textlengthmark}{\textsubring{d}}{a}} \dicPos{f}[1] \dicFlx{(‑u, ‑ur)}[7] \dicFieldCat{zool.} \dicDirectTranslationCS{racek tříprstý} \textit{(l.~{\textLA{Rissa tridactyla}})}  \dicsymPhoto\ 
\dicFigure{21009.jpg}{Rita}{Rita - Matthews, Sue, Biolib, PD}
\dicEntry[rita] \dicTerm{rit|a\smash{\textsuperscript{2}}} \dicsymFrequent\  \dicIPA{{r}{\textsci}{\textlengthmark}{\textsubring{d}}{a}} \dicPos{v}[1] \dicFlx{(‑aði)}[1] \dicFlx{acc} \dicSynonym{skrifa} \dicDirectTranslationCS{napsat, psát} \dicExampleIS{rita grein} \dicExampleCS{psát článek}
\dicEntry[ritari] \dicTerm{rit··ar|i} \dicsymFrequent\  \dicIPA{{r}{\textsci}{\textlengthmark}{\textsubring{d}}{a}{r}{\textsci}} \dicPos{m}[1] \dicFlx{(‑a, ‑ar)}[13] \dicDirectTranslationCS{sekretář(ka), tajemník, tajemnice} \dicExampleIS{Hún var ritari félagsins í sextán ár.} \dicExampleCS{Byla tajemnicí spolku šestnáct let.}
\dicEntry[ritaskrá] \dicTerm{rita··skrá} \dicIPA{{r}{\textsci}{\textlengthmark}{\textsubring{d}}{a}{s}{\r{g}}{r}{au}} \dicPos{f}[4] \dicFlx{(‑r\,/\addthin ‑ar, ‑r)}[21] \dicDirectTranslationCS{bibliografie}
\dicEntry[ritdeila] \dicTerm{rit··deil|a} \dicIPA{{r}{\textsci}{\textlengthmark}{\textsubring{d}}{ei}{l}{a}} \dicPos{f}[1] \dicFlx{(‑u, ‑ur)}[13] \dicDirectTranslationCS{(psaná) polemika\,/\addthin disputace}
\dicEntry[ritdómari] \dicTerm{rit··dóm·ar|i} \dicIPA{{r}{\textsci}{\textlengthmark}{\textsubring{d}}{ou}{m}{a}{r}{\textsci}} \dicPos{m}[1] \dicFlx{(‑a, ‑ar)}[13] \dicDirectTranslationCS{recenzent(ka), (literární) kritik\,/\addthin kritička}
\dicEntry[ritdómur] \dicTerm{rit··dóm|ur} \dicIPA{{r}{\textsci}{\textlengthmark}{\textsubring{d}}{ou}{m}{\textscy}{\textsubring{r}}} \dicPos{m}[6] \dicFlx{(‑s, ‑ar)}[10] \dicDirectTranslationCS{recenze, (literární) kritika}
\dicEntry[ritfangaverslun] \dicTerm{rit·fanga··versl|un} \dicIPA{{r}{\textsci}{\textlengthmark}{\textsubring{d}}{f}{au}{\ng}{\r{g}}{a}{v}{\textepsilon}{\textsubring{r}}{s}{\textsubring{d}}{l}{\textscy}{\textsubring{n}}} \dicPos{f}[7] \dicFlx{(‑unar, ‑anir)}[8] \dicDirectTranslationCS{papírnictví}
\dicEntry[ritfrelsi] \dicTerm{rit··frelsi} \dicIPA{{r}{\textsci}{\textlengthmark}{\textsubring{d}}{f}{r}{\textepsilon}{l}{s}{\textsci}} \dicPos{n}[2] \dicFlx{(‑s)}[20] \dicDirectTranslationCS{svoboda tisku}
\dicEntry[ritfæra] \dicTerm{rit··fær|a} \dicIPA{{r}{\textsci}{\textlengthmark}{\textsubring{d}}{f}{a}{i}{r}{a}} \dicPos{v}[2] \dicFlx{(‑ði, ‑t)}[105] \dicFieldCat{poč.} \dicDirectTranslationCS{editovat, upravit (text ap.)}
\dicEntry[ritfærsla] \dicTerm{rit··færsl|a} \dicIPA{{r}{\textsci}{\textlengthmark}{\textsubring{d}}{f}{a}{i}{\textsubring{r}}{s}{\textsubring{d}}{l}{a}} \dicPos{f}[1] \dicFlx{(‑u)}[5] \dicFieldCat{poč.} \dicDirectTranslationCS{úprava\,/\addthin editace textu}
\dicEntry[ritföng] \dicTerm{rit··föng} \dicIPA{{r}{\textsci}{\textlengthmark}{\textsubring{d}}{f}{\oe i}{\ng}{\r{g}}} \dicPos{n}[2] \dicFlx{pl}[9] \dicDirectTranslationCS{psací potřeby}
\dicEntry[ritgerð] \dicTerm{rit··gerð} \dicIPA{{r}{\textsci}{\textlengthmark}{\textsubring{d}}{\r{\textObardotlessj}}{\textepsilon}{r}{\texttheta}} \dicPos{f}[7] \dicFlx{(‑ar, ‑ir)}[1] \dicDirectTranslationCS{(písemná) práce, esej, kompozice}
\dicEntry[rithöfundur] \dicTerm{rit··höf·und|ur} \dicsymFrequent\  \dicIPA{{r}{\textsci}{\textlengthmark}{\textsubring{d}}{h}{\oe}{v}{\textscy}{n}{\textsubring{d}}{\textscy}{\textsubring{r}}} \dicPos{m}[6] \dicFlx{(‑ar, ‑ar)}[56] \dicDirectTranslationCS{spisovatel(ka), autor(ka) (literárního díla)} \dicExampleIS{Hann var blaðamaður og rithöfundur.} \dicExampleCS{Byl novinář a~spisovatel.}
\dicEntry[rithönd] \dicTerm{rit··|hönd} \dicIPA{{r}{\textsci}{\textlengthmark}{\textsubring{d}}{h}{\oe}{n}{\textsubring{d}}} \dicPos{f}[8] \dicFlx{(‑handar, ‑hendur\,/\addthin ‑handir)}[12] \dicDirectTranslationCS{rukopis, písmo}
\dicEntry[ritill] \dicTerm{rit|ill} \dicIPA{{r}{\textsci}{\textlengthmark}{\textsubring{d}}{\textsci}{\textsubring{d}}{\textsubring{l}}} \dicPos{m}[6] \dicFlx{(‑ils, ‑lar)}[35] \dicFieldCat{poč.} \dicDirectTranslationCS{(textový) editor}
\dicEntry[ritlaun] \dicTerm{rit··laun} \dicIPA{{r}{\textsci}{\textlengthmark}{\textsubring{d}}{l}{\oe i}{\textsubring{n}}} \dicPos{n}[2] \dicFlx{pl}[1] \dicDirectTranslationCS{(autorský) honorář}
\dicEntry[ritmál] \dicTerm{rit··mál} \dicIPA{{r}{\textsci}{\textlengthmark}{\textsubring{d}}{m}{au}{\textsubring{l}}} \dicPos{n}[2] \dicFlx{(‑s, ‑)}[5] \dicFieldCat{jaz.} \dicDirectTranslationCS{psaný jazyk}
\dicEntry[ritmi] \dicTerm{ritm|i}\dicTerm{, rytmi} \dicIPA{{r}\-{\textsci}\-{h}\-{\textsubring{d}}\-{m}\-{\textsci}\-} \dicPos{m}[1] \dicFlx{(‑a, ‑ar)}[1] \dicSynonym*{hljóðfall} \dicDirectTranslationCS{rytmus}
\dicEntry[ritnefnd] \dicTerm{rit··nefnd} \dicIPA{{r}{\textsci}{\textlengthmark}{\textsubring{d}}{n}{\textepsilon}{m}{\textsubring{d}}} \dicPos{f}[7] \dicFlx{(‑ar, ‑ir)}[1] \dicDirectTranslationCS{redakce, redakční rada}
\dicEntry[ritning] \dicTerm{rit··ning} \dicIPA{{r}{\textsci}{h}{\textsubring{d}}{n}{i}{\ng}{\r{g}}} \dicPos{f}[4] \dicFlx{(‑ar, ‑ar)}[5] \textbf{1.} \dicSynonym{rit} \dicDirectTranslationCS{(psané) dílo, spis}  \textbf{2.} \dicSynonym{biblía} \dicDirectTranslationCS{Písmo} \dicIndirectTranslationCS{(bible)};  \dicPhraseIS{heilög ritning} \dicFieldCat{náb.} \dicDirectTranslationCS{Písmo svaté}
\dicEntry[ritröð] \dicTerm{rit··|röð} \dicIPA{{r}{\textsci}{\textlengthmark}{\textsubring{d}}{r}{\oe}{\texttheta}} \dicPos{f}[7] \dicFlx{(‑raðar, ‑raðir)}[16] \dicDirectTranslationCS{série článků\,/\addthin spisů}
\dicEntry[ritsafn] \dicTerm{rit··|safn} \dicIPA{{r}{\textsci}{\textlengthmark}{\textsubring{d}}{s}{a}{\textsubring{b}}{\textsubring{n}}} \dicPos{n}[2] \dicFlx{(‑safns, ‑söfn)}[8] \dicDirectTranslationCS{sebrané dílo, sebraný spis}
\dicEntry[ritsími] \dicTerm{rit··sím|i} \dicIPA{{r}{\textsci}{\textlengthmark}{\textsubring{d}}{s}{i}{m}{\textsci}} \dicPos{m}[1] \dicFlx{(‑a)}[3] \dicDirectTranslationCS{telegraf}
\dicEntry[ritsmíð] \dicTerm{rit··smíð} \dicIPA{{r}{\textsci}{\textlengthmark}{\textsubring{d}}{s}{m}{i}{\texttheta}} \dicPos{f}[4] \dicFlx{(‑ar, ‑ar)}[1] \dicDirectTranslationCS{literární činnost, psaní}
\dicEntry[ritstjóri] \dicTerm{rit··stjór|i} \dicIPA{{r}{\textsci}{\textlengthmark}{\textsubring{d}}{s}{\textsubring{d}}{j}{ou}{r}{\textsci}} \dicPos{m}[1] \dicFlx{(‑a, ‑ar)}[1] \textbf{1.} \dicSynonym*{ritstjóri blaðs} \dicDirectTranslationCS{redaktor(ka)}  \textbf{2.} \dicSynonym*{starfsmaður bókaútgáfu} \dicDirectTranslationCS{editor(ka)}
\dicEntry[ritstjórn] \dicTerm{rit··stjórn} \dicIPA{{r}{\textsci}{\textlengthmark}{\textsubring{d}}{s}{\textsubring{d}}{j}{ou}{r}{\textsubring{d}}{\textsubring{n}}} \dicPos{f}[7] \dicFlx{(‑ar, ‑ir)}[1] \textbf{1.} \dicSynonym*{ritstýring} \dicDirectTranslationCS{redigování (časopisu ap.)}  \textbf{2.} \dicSynonym*{ritstjórar} \dicDirectTranslationCS{redakce, redakční rada}  \textbf{3.} \dicSynonym{skrifstofa} \dicDirectTranslationCS{redakce} \dicIndirectTranslationCS{(místnost, kde pracuje kolektiv redaktorů)}
\begin{xtolerant}{}{1pt}
\dicEntry[ritstjórnargrein] \dicTerm{rit·stjórnar··grein}\addthinS\dicIPA{{r}{\textsci}{\textlengthmark}{\textsubring{d}}{s}{\textsubring{d}}{j}{ou}{r}{\textsubring{d}}{n}{a}{r}{\r{g}}{r}{ei}{\textsubring{n}}}\addthinS\dicPos{f}[4]\addthinS\dicFlx{(‑ar, ‑ar)}[1]\addthinS\dicSynonym{leiðari} \dicDirectTranslationCS{úvodník} \dicIndirectTranslationCS{(úvodní článek novin nebo časopisu)}
\end{xtolerant}
\dicEntry[ritstuldur] \dicTerm{rit··stuld|ur} \dicIPA{{r}{\textsci}{\textlengthmark}{\textsubring{d}}{s}{\textsubring{d}}{\textscy}{l}{\textsubring{d}}{\textscy}{\textsubring{r}}} \dicPos{m}[10] \dicFlx{(‑ar, ‑ir)}[4] \dicDirectTranslationCS{plagiátorství}
\dicEntry[ritstörf] \dicTerm{rit··störf} \dicIPA{{r}{\textsci}{\textlengthmark}{\textsubring{d}}{s}{\textsubring{d}}{\oe}{r}{f}} \dicPos{n}[2] \dicFlx{pl}[9] \dicDirectTranslationCS{psaní} \dicIndirectTranslationCS{(literární tvorba)};  \dicPhraseIS{fást við ritstörf} \dicFlx{refl} \dicDirectTranslationCS{živit se psaním, psát}
\dicEntry[ritun] \dicTerm{rit|un} \dicIPA{{r}{\textsci}{\textlengthmark}{\textsubring{d}}{\textscy}{\textsubring{n}}} \dicPos{f}[7] \dicFlx{(‑unar)}[9] \dicDirectTranslationCS{psaní} \dicIndirectTranslationCS{(písemné zaznamenání něčeho)}
\dicEntry[ritvél] \dicTerm{rit··vél} \dicIPA{{r}{\textsci}{\textlengthmark}{\textsubring{d}}{v}{j}{\textepsilon}{\textsubring{l}}} \dicPos{f}[4] \dicFlx{(‑ar, ‑ar)}[1] \dicDirectTranslationCS{psací stroj}
\dicEntry[ritvélarborði] \dicTerm{rit·vélar··borð|i} \dicIPA{{r}{\textsci}{\textlengthmark}{\textsubring{d}}{v}{j}{\textepsilon}{l}{a}{r}{\textsubring{b}}{\textopeno}{r}{ð}{\textsci}} \dicPos{m}[1] \dicFlx{(‑a, ‑ar)}[1] \dicDirectTranslationCS{průklepová páska}
\dicEntry[ritvinnsla] \dicTerm{rit··vinnsl|a} \dicIPA{{r}{\textsci}{\textlengthmark}{\textsubring{d}}{v}{\textsci}{n}{s}{\textsubring{d}}{l}{a}} \dicPos{f}[1] \dicFlx{(‑u)}[5] \dicFieldCat{poč.} \dicSynonym{ritfærsla} \dicDirectTranslationCS{úprava\,/\addthin editace textu}
\dicEntry[ritvinnsluforrit] \dicTerm{rit·vinnslu··for·rit} \dicIPA{{r}{\textsci}{\textlengthmark}{\textsubring{d}}{v}{\textsci}{n}{s}{\textsubring{d}}{l}{\textscy}{f}{\textopeno}{r}{\textlengthmark}{\textsci}{\textsubring{d}}} \dicPos{n}[2] \dicFlx{(‑s, ‑)}[5] \dicFieldCat{poč.} \dicDirectTranslationCS{textový editor}
\dicEntry[ríð] \dicTerm{ríð} \dicIPA{{r}{i}{\textlengthmark}{\texttheta}} \dicPos{v} \dicFlx{ind praes sg 1 pers} \dicLink{ríða}
\dicEntry[ríða] \dicTerm{ríða} \dicsymFrequent\  \dicIPA{{r}{i}{\textlengthmark}{ð}{a}} \dicPos{v}[6] \dicFlx{(ríð, reið, riðum, riði, riðið)}[75] \dicFlx{dat\,/\addthin acc} \textbf{1.} \dicFlx{dat} \dicDirectTranslationCS{jet (na koni, oslu ap.)} \dicExampleIS{ríða gráum hesti} \dicExampleCS{jet na šedém koni};  \dicPhraseIS{ríða út} \dicDirectTranslationCS{vyjet si (na projížďku)}  \textbf{2.} \dicPhraseIS{e‑að ríður e‑m að fullu} \dicLangCat{přen.} \dicDirectTranslationCS{(co komu) zasadí smrtelnou ránu}  \textbf{3.} \dicFlx{dat} \dicLangCat{hrub.} \dicSynonym*{hafa kynmök við} \dicDirectTranslationCS{šukat, šoustat}  \textbf{4.} \dicFlx{acc} \dicSynonym{flétta\smash{\textsuperscript{2}}} \dicDirectTranslationCS{vázat, plést (síť ap.)};  \dicIdiom{ríða}[á]{ \dicPhraseIS{það ríður á e‑u}} \dicFlx{impers} \dicDirectTranslationCS{(co) je důležité, na (čem) záleží}
\dicEntry[ríf] \dicTerm{ríf} \dicIPA{{r}{i}{\textlengthmark}{f}} \dicPos{v} \dicFlx{ind praes sg 1 pers} \dicLink{rífa}
\dicEntry[rífa] \dicTerm{rífa} \dicsymFrequent\  \dicIPA{{r}{i}{\textlengthmark}{v}{a}} \dicPos{v}[6] \dicFlx{(ríf, reif, rifum, rifi, rifið)}[69] \dicFlx{acc} \textbf{1.} \dicSynonym*{slíta e‑ð sundur með rifum} \dicDirectTranslationCS{roztrhnout, (roz)trhat} \dicExampleIS{rífa bréfið í sundur} \dicExampleCS{roztrhat dopis na kousky}  \textbf{2.} \dicSynonym*{losa sundur smíði} \dicDirectTranslationCS{strhnout, zbořit, zničit} \dicExampleIS{rífa húsið} \dicExampleCS{strhnout dům}  \textbf{3.} \dicSynonym{hrifsa} \dicDirectTranslationCS{vytrhnout, vyrvat} \dicExampleIS{rífa blaðið út úr höndunum á e‑m} \dicExampleCS{vytrhnout (komu) noviny z~rukou}  \textbf{4.} \dicSynonym{klóra} \dicDirectTranslationCS{drápnout, (po)drápat, škrábnout, (po)škrábat} \dicExampleIS{Kötturinn reif mig.} \dicExampleCS{Kočka mě škrábla.};  \dicIdiom{rífa}[í]{ \dicPhraseIS{rífa e‑ð í sig}} \dicSynonym{háma} \dicDirectTranslationCS{cpát se (čím), hltat (co)};  \dicIdiom{rífa}[niður]{ \dicPhraseIS{rífa e‑ð niður}} \dicDirectTranslationCS{(ostře) zkritizovat (co)} \dicExampleIS{rífa sýninguna niður} \dicExampleCS{ostře zkritizovat výstavu};  \dicIdiom{rífa}[upp]{ \dicPhraseIS{rífa upp e‑ð}} \dicSynonym{opna\smash{\textsuperscript{2}}} \dicDirectTranslationCS{roztrhnout (co), otevřít (co) (dopis ap.)} \dicExampleIS{rífa upp bréf} \dicExampleCS{roztrhnout dopis}; { \dicPhraseIS{rífa e‑n upp}} \dicDirectTranslationCS{vytrhnout\,/\addthin vyrvat (koho) ze spánku};  \dicIdiom{rífa}[út]{ \dicPhraseIS{rífa e‑ð út}} \dicDirectTranslationCS{skoupit (co), vykoupit (co)};  \dicIdiom{rífast}{ \dicPhraseIS{rífast}} \dicFlx{refl} \dicDirectTranslationCS{(po)hádat se, přít se} \dicExampleIS{rífast við e‑n} \dicExampleCS{hádat se s~(kým)}
\dicEntry[ríflega] \dicTerm{ríf··lega} \dicIPA{{r}{i}{v}{l}{\textepsilon}{\textbabygamma}{a}} \dicPos{adv} \textbf{1.} \dicSynonym*{rausnarlega} \dicDirectTranslationCS{hojně, velkoryse}  \textbf{2.} \dicSynonym{rúmlega\smash{\textsuperscript{2}}} \dicDirectTranslationCS{přes, více než} \dicExampleIS{ríflega helmingur} \dicExampleCS{více než polovina}
\dicEntry[ríflegur] \dicTerm{ríf··legur} \dicIPA{{r}{i}{v}{l}{\textepsilon}{\textbabygamma}{\textscy}{\textsubring{r}}} \dicPos{adj}[1]\dicFlx{}[-8] \textbf{1.} \dicSynonym{ríkulegur} \dicDirectTranslationCS{hojný, bohatý}  \textbf{2.} \dicSynonym{örlátur} \dicDirectTranslationCS{štědrý, velkorysý}
\dicEntry[Ríga] \dicTerm{Ríga} \dicIPA{{r}{i}{\textlengthmark}{\textbabygamma}{a}} \dicPos{subs} \dicFlx{indecl} \dicFieldCat{geog.} \dicDirectTranslationCS{Riga} \dicIndirectTranslationCS{(hlavní město Lotyšska)}
\dicEntry[rígmontinn] \dicTerm{ríg··montinn} \dicIPA{{r}{i}{\textbabygamma}{m}{\textopeno}{\textsubring{n}}{\textsubring{d}}{\textsci}{\textsubring{n}}} \dicPos{adj}[6]\dicFlx{}[-6] \dicDirectTranslationCS{nafoukaný, pyšný}
\dicEntry[rígur] \dicTerm{ríg|ur} \dicIPA{{r}{i}{\textlengthmark}{\textbabygamma}{\textscy}{\textsubring{r}}} \dicPos{m}[6] \dicFlx{(‑s)}[17] \textbf{1.} \dicSynonym{stirðleiki} \dicDirectTranslationCS{strnulost, ztuhlost}  \textbf{2.} \dicSynonym{spenna\smash{\textsuperscript{1}}} \dicDirectTranslationCS{rivalita, soupeření} \dicExampleIS{Það er rígur á milli þeirra.} \dicExampleCS{Panuje mezi nimi rivalita.}
\dicEntry[ríki] \dicTerm{ríki} \dicsymFrequent\  \dicIPA{{r}{i}{\textlengthmark}{\r{\textObardotlessj}}{\textsci}} \dicPos{n}[2] \dicFlx{(‑s, ‑)}[16] \textbf{1.} \dicSynonym{land} \dicDirectTranslationCS{stát, země} \dicExampleIS{sjálfstætt ríki} \dicExampleCS{svrchovaný stát}  \textbf{2.} \dicSynonym{vald} \dicDirectTranslationCS{moc, vláda} \dicIndirectTranslationCS{(zvláště v~pl)}  \textbf{3.} \dicFieldCat{biol.} \dicDirectTranslationCS{říše}  \textbf{4.} \dicPhraseIS{Ríkið} \dicFlx{def} \dicDirectTranslationCS{obchod s~alkoholem} \dicIndirectTranslationCS{(státní monopol s~prodejem alkoholu)}
\dicEntry[ríkis] \dicTerm{ríkis-} \dicIPA{{r}{i}{\textlengthmark}{\r{\textObardotlessj}}{\textsci}{s}} \dicPos{predp} \dicDirectTranslationCS{státní}
\dicEntry[ríkisábyrgð] \dicTerm{ríkis··á·byrgð} \dicIPA{{r}{i}{\textlengthmark}{\r{\textObardotlessj}}{\textsci}{s}{au}{\textsubring{b}}{\textsci}{r}{\texttheta}} \dicPos{f}[7] \dicFlx{(‑ar, ‑ir)}[1] \dicDirectTranslationCS{státní záruka\,/\addthin garance}
\dicEntry[ríkisborgararéttur] \dicTerm{ríkis·borgara··rétt|ur} \dicIPA{{r}\-{i}\-{\textlengthmark}\-{\r{\textObardotlessj}}\-{\textsci}\-{s}\-{\textsubring{b}}\-{\textopeno}\-{r}\-{\r{g}}\-{a}\-{r}\-{a}\-{r}\-{j}\-{\textepsilon}\-{h}\-{\textsubring{d}}\-{\textscy}\-{\textsubring{r}}\-} \dicPos{m}[10] \dicFlx{(‑ar)}[7] \dicFieldCat{práv.} \dicDirectTranslationCS{(státní) občanství}
\dicEntry[ríkisborgari] \dicTerm{ríkis··borg·ar|i} \dicIPA{{r}{i}{\textlengthmark}{\r{\textObardotlessj}}{\textsci}{s}{\textsubring{b}}{\textopeno}{r}{\r{g}}{a}{r}{\textsci}} \dicPos{m}[1] \dicFlx{(‑a, ‑ar)}[13] \dicDirectTranslationCS{občan(ka)}
\dicEntry[ríkisfang] \dicTerm{ríkis··|fang} \dicIPA{{r}{i}{\textlengthmark}{\r{\textObardotlessj}}{\textsci}{s}{f}{au}{\ng}{\r{g}}} \dicPos{n}[2] \dicFlx{(‑fangs, ‑föng)}[8] \dicFieldCat{práv.} \dicDirectTranslationCS{(státní) občanství}
\dicEntry[ríkisfjármál] \dicTerm{ríkis··fjár·mál} \dicIPA{{r}{i}{\textlengthmark}{\r{\textObardotlessj}}{\textsci}{s}{f}{j}{au}{r}{m}{au}{\textsubring{l}}} \dicPos{n}[2] \dicFlx{pl}[1] \dicFieldCat{ekon.} \dicDirectTranslationCS{státní finance}
\dicEntry[ríkislán] \dicTerm{ríkis··lán} \dicIPA{{r}{i}{\textlengthmark}{\r{\textObardotlessj}}{\textsci}{s}{l}{au}{\textsubring{n}}} \dicPos{n}[2] \dicFlx{(‑s, ‑)}[5] \dicFieldCat{ekon.} \dicDirectTranslationCS{státní půjčka}
\dicEntry[ríkislögreglustjóri] \dicTerm{ríkis··lög·reglu·stjór|i} \dicIPA{{r}\-{i}\-{\textlengthmark}\-{\r{\textObardotlessj}}\-{\textsci}\-{s}\-{l}\-{\oe}\-{\textbabygamma}\-{r}\-{\textepsilon}\-{\r{g}}\-{l}\-{\textscy}\-{s}\-{\textsubring{d}}\-{j}\-{ou}\-{r}\-{\textsci}\-} \dicPos{m}[1] \dicFlx{(‑a, ‑ar)}[1] \dicDirectTranslationCS{policejní prezident(ka)}
\dicEntry[ríkisrekstur] \dicTerm{ríkis··rekst|ur} \dicIPA{{r}{i}{\textlengthmark}{\r{\textObardotlessj}}{\textsci}{s}{r}{\textepsilon}{x}{s}{\textsubring{d}}{\textscy}{\textsubring{r}}} \dicPos{m}[5] \dicFlx{(‑rar\,/\addthin ‑urs)}[10] \dicDirectTranslationCS{státní společnost\,/\addthin podnik}
\dicEntry[ríkissaksóknari] \dicTerm{ríkis··sak·sókn·ar|i} \dicIPA{{r}{i}{\textlengthmark}{\r{\textObardotlessj}}{\textsci}{s}{a}{\r{g}}{s}{ou}{h}{\r{g}}{n}{a}{r}{\textsci}} \dicPos{m}[1] \dicFlx{(‑a, ‑ar)}[13] \dicFieldCat{práv.} \dicDirectTranslationCS{státní zástupce\,/\addthin zástupkyně, prokurátor(ka)}
\dicEntry[ríkissjóður] \dicTerm{ríkis··sjóð|ur} \dicIPA{{r}{i}{\textlengthmark}{\r{\textObardotlessj}}{\textsci}{s}{j}{ou}{ð}{\textscy}{\textsubring{r}}} \dicPos{m}[9] \dicFlx{(‑s, ‑ir)}[6] \dicDirectTranslationCS{státní pokladna}
\dicEntry[ríkisskattstjóri] \dicTerm{ríkis··skatt·stjór|i} \dicIPA{{r}{i}{\textlengthmark}{\r{\textObardotlessj}}{\textsci}{s}{\r{g}}{a}{h}{s}{\textsubring{d}}{j}{ou}{r}{\textsci}} \dicPos{m}[1] \dicFlx{(‑a, ‑ar)}[1] \dicDirectTranslationCS{státní pokladník\,/\addthin pokladnice\,/\addthin pokladní} \dicIndirectTranslationCS{(ředitel integrovaného rozpočtového systému)}
\dicEntry[ríkisskuldabréf] \dicTerm{ríkis··skulda·bréf} \dicIPA{{r}{i}{\textlengthmark}{\r{\textObardotlessj}}{\textsci}{s}{\r{g}}{\textscy}{l}{\textsubring{d}}{a}{\textsubring{b}}{r}{j}{\textepsilon}{f}} \dicPos{n}[2] \dicFlx{(‑s, ‑)}[5] \dicFieldCat{ekon.} \dicDirectTranslationCS{státní dluhopis}

\dicEntry[ríkisstarfsmaður] \textls[19]{\dicTerm{ríkis··starfs·|maður} \dicIPA{{r}{i}{\textlengthmark}{\r{\textObardotlessj}}{\textsci}{s}{\textsubring{d}}{a}{r}{v}{s}{m}{a}{ð}{\textscy}{\textsubring{r}}} \dicPos{m}[13] \dicFlx{(‑manns, ‑menn)}[2]} \dicDirectTranslationCS{státní zaměstnanec\,/\addthin zaměstnankyně}

\dicEntry[ríkisstjóri] \dicTerm{ríkis··stjór|i} \dicIPA{{r}{i}{\textlengthmark}{\r{\textObardotlessj}}{\textsci}{s}{\textsubring{d}}{j}{ou}{r}{\textsci}} \dicPos{m}[1] \dicFlx{(‑a, ‑ar)}[1] \dicDirectTranslationCS{guvernér(ka)}
\dicEntry[ríkisstjórn] \dicTerm{ríkis··stjórn} \dicIPA{{r}{i}{\textlengthmark}{\r{\textObardotlessj}}{\textsci}{s}{\textsubring{d}}{j}{ou}{r}{\textsubring{d}}{\textsubring{n}}} \dicPos{f}[7] \dicFlx{(‑ar, ‑ir)}[1] \dicDirectTranslationCS{(státní\,/\addthin národní) vláda}
\dicEntry[ríkisstyrkur] \dicTerm{ríkis··styrk|ur} \dicIPA{{r}{i}{\textlengthmark}{\r{\textObardotlessj}}{\textsci}{s}{\textsubring{d}}{\textsci}{\textsubring{r}}{\r{g}}{\textscy}{\textsubring{r}}} \dicPos{m}[9] \dicFlx{(‑s, ‑ir)}[15] \dicFieldCat{ekon.} \dicDirectTranslationCS{státní subvence}
\dicEntry[ríkisútvarp] \dicTerm{ríkis··út·|varp} \dicsymFrequent\  \dicIPA{{r}{i}{\textlengthmark}{\r{\textObardotlessj}}{\textsci}{s}{u}{\textsubring{d}}{v}{a}{\textsubring{r}}{\textsubring{b}}} \dicPos{n}[2] \dicFlx{(‑varps, ‑vörp)}[8] \dicDirectTranslationCS{státní rozhlas} \dicExampleIS{afnotagjald ríkisútvarps} \dicExampleCS{poplatek za užívání státního rozhlasu}
\dicEntry[ríkisvald] \dicTerm{ríkis··vald} \dicIPA{{r}{i}{\textlengthmark}{\r{\textObardotlessj}}{\textsci}{s}{v}{a}{l}{\textsubring{d}}} \dicPos{n}[2] \dicFlx{(‑s)}[2] \dicDirectTranslationCS{státní moc}
\dicEntry[ríkja] \dicTerm{rík|ja} \dicsymFrequent\  \dicIPA{{r}{i}{\textlengthmark}{\r{\textObardotlessj}}{a}} \dicPos{v}[2] \dicFlx{(‑ti, ‑t)}[34] \textbf{1.} \dicSynonym{stjórna} \dicDirectTranslationCS{vládnout, panovat} \dicIndirectTranslationCS{(být panovníkem)}  \textbf{2.} \dicDirectTranslationCS{vládnout, panovat} \dicIndirectTranslationCS{(existovat všeobecně nebo ve velké míře)} \dicExampleIS{það ríkti friður e‑s staðar} \dicExampleCS{(kde) vládl mír}
\dicEntry[ríkjandi] \dicTerm{ríkj··andi} \dicIPA{{r}{i}{\textlengthmark}{\r{\textObardotlessj}}{a}{n}{\textsubring{d}}{\textsci}} \dicPos{adj}[13] \dicFlx{indecl}[1] \textbf{1.} \dicSynonym*{sem ríkir} \dicDirectTranslationCS{převládající, dominující}  \textbf{2.} \dicFieldCat{biol.} \dicDirectTranslationCS{dominantní};  \dicPhraseIS{ríkjandi gen} \dicFieldCat{biol.} \dicDirectTranslationCS{dominantní gen}
\dicEntry[ríkulegur] \dicTerm{ríku··legur} \dicIPA{{r}{i}{\textlengthmark}{\r{g}}{\textscy}{l}{\textepsilon}{\textbabygamma}{\textscy}{\textsubring{r}}} \dicPos{adj}[1]\dicFlx{}[-8] \dicSynonym{ríflegur} \dicDirectTranslationCS{hojný, bohatý (pohoštění ap.)}
\dicEntry[ríkur] \dicTerm{ríkur} \dicsymFrequent\  \dicIPA{{r}{i}{\textlengthmark}{\r{g}}{\textscy}{\textsubring{r}}} \dicPos{adj}[1]\dicFlx{}[-1] \textbf{1.} \dicSynonym{auðugur} \dicDirectTranslationCS{bohatý, zámožný} \dicExampleIS{verða ríkur} \dicExampleCS{zbohatnout} \dicAntonym{fátækur}  \textbf{2.} \dicSynonym*{innihaldsríkur} \dicDirectTranslationCS{bohatý, vydatný} \dicExampleIS{ríkur að vítamínum} \dicExampleCS{bohatý na vitamíny}  \textbf{3.} \dicSynonym{voldugur} \dicDirectTranslationCS{mocný, silný} \dicExampleIS{ríkur höfðingi} \dicExampleCS{mocný náčelník}
\dicEntry[rím] \dicTerm{rím} \dicIPA{{r}{i}{\textlengthmark}{\textsubring{m}}} \dicPos{n}[2] \dicFlx{(‑s, ‑)}[5] \dicFieldCat{lit.} \dicDirectTranslationCS{rým};  \dicPhraseIS{langsetis rím} \dicFieldCat{lit.} \dicIndirectTranslationCS{vnitřní rým, ve kterém se rýmují slabiky ve stejném verši};  \dicPhraseIS{veggjað rím} \dicFieldCat{lit.} \dicSynonym{þrírím} \dicDirectTranslationCS{trojslabičný rým};  \dicPhraseIS{þversetis rím} \dicFieldCat{lit.} \dicIndirectTranslationCS{vnitřní rým, ve kterém se rýmují slabiky nacházející se v~různých verších};  \dicIdiom{rím}{ \dicPhraseIS{rugla e‑n í ríminu}} \dicDirectTranslationCS{poplést (koho), splést (koho)}
\dicEntry[ríma] \dicTerm{rím|a\smash{\textsuperscript{1}}} \dicIPA{{r}{i}{\textlengthmark}{m}{a}} \dicPos{f}[1] \dicFlx{(‑u, ‑ur)}[13] \dicFieldCat{lit.} \dicDirectTranslationCS{výpravná báseň}
\dicEntry[ríma] \dicTerm{rím|a\smash{\textsuperscript{2}}} \dicIPA{{r}{i}{\textlengthmark}{m}{a}} \dicPos{v}[1] \dicFlx{(‑aði)}[1] \dicFlx{acc} \textbf{1.} \dicDirectTranslationCS{rýmovat se} \dicExampleIS{Orðin ríma saman.} \dicExampleCS{Slova se rýmují.}  \textbf{2.} \dicDirectTranslationCS{pasovat, hodit se}
\dicEntry[rímorð] \dicTerm{rím··orð} \dicIPA{{r}{i}{\textlengthmark}{m}{\textopeno}{r}{\texttheta}} \dicPos{n}[2] \dicFlx{(‑s, ‑)}[5] \dicFieldCat{lit.} \dicDirectTranslationCS{rým, rýmové slovo}
\dicEntry[Rín] \dicTerm{Rín} \dicIPA{{r}{i}{\textlengthmark}{\textsubring{n}}} \dicPos{f}[4] \dicFlx{(‑ar)}[4] \dicFieldCat{geog.} \dicDirectTranslationCS{Rýn} \dicIndirectTranslationCS{(řeka v~Evropě)}
\dicEntry[rís] \dicTerm{rís} \dicIPA{{r}{i}{\textlengthmark}{s}} \dicPos{v} \dicFlx{ind praes sg 1 pers} \dicLink{rísa}
\dicEntry[rísa] \dicTerm{rísa} \dicsymFrequent\  \dicIPA{{r}{i}{\textlengthmark}{s}{a}} \dicPos{v}[6] \dicFlx{(rís, reis, risum, risi, risið)}[77] \textbf{1.} \dicSynonym*{standa upp} \dicDirectTranslationCS{vstát, povstat, zvednout se} \dicExampleIS{rísa úr sæti} \dicExampleCS{vstát ze židle};  \dicPhraseIS{rísa á fætur} \dicDirectTranslationCS{postavit se, vstát}  \textbf{2.} \dicSynonym*{koma upp} \dicDirectTranslationCS{vycházet, vyjít (slunce ap.)} \dicExampleIS{Sólin rís.} \dicExampleCS{Slunce vychází.}  \textbf{3.} \dicDirectTranslationCS{zvedat se, stoupat, vztyčovat se (vlny ap.)}  \textbf{4.} \dicDirectTranslationCS{povstat, vybudovat se (nový obchod ap.)}  \textbf{5.} \dicDirectTranslationCS{vzniknout, povstat (neshoda ap.)};  \dicIdiom{rísa}[gegn\,/\addthin móti\,/\addthin við]{ \dicPhraseIS{rísa gegn e‑m\,/\addthin e‑u, rísa móti e‑m\,/\addthin e‑u, rísa við e‑m\,/\addthin e‑u}} \dicDirectTranslationCS{povstat proti (komu\,/\addthin čemu), postavit se proti (komu\,/\addthin čemu)};  \dicIdiom{rísa}[undir]{ \dicPhraseIS{rísa (ekki) undir nafni}} \dicLangCat{přen.} \dicDirectTranslationCS{(ne)dostát svému jménu};  \dicIdiom{rísa}[upp]{ \dicPhraseIS{rísa upp}} {\textbf{a.}} \dicSynonym*{standa upp} \dicDirectTranslationCS{postavit se, vstát};  {\textbf{b.}} \dicDirectTranslationCS{povstat, vzbouřit se}
\dicEntry[rjátla] \dicTerm{rjátl|a} \dicIPA{{r}{j}{au}{h}{\textsubring{d}}{l}{a}} \dicPos{v}[1] \dicFlx{(‑aði)}[22] \dicSynonym{rölta} \dicDirectTranslationCS{toulat se, bloumat} \dicExampleIS{rjátla um göturnar} \dicExampleCS{toulat se po ulicích};  \dicIdiom{rjátlast}{ \dicPhraseIS{e‑að rjátlast af e‑m}} \dicFlx{refl} \dicDirectTranslationCS{(co) se vytratí časem u~(koho) (špatná vlastnost ap.)}
\dicEntry[rjóða] \dicTerm{rjóða} \dicIPA{{r}{j}{ou}{\textlengthmark}{ð}{a}} \dicPos{v}[6] \dicFlx{(rýð, rauð, ruðum, ryði, roðið)}[95] \dicFlx{dat} \dicLangCat{zast.} \dicSynonym{smyrja} \dicDirectTranslationCS{potřít, natřít} \dicExampleIS{rjóða smyrsli á andlitið} \dicExampleCS{natřít mast na tvář}
\dicEntry[rjóður] \dicTerm{rjóður\smash{\textsuperscript{1}}} \dicIPA{{r}{j}{ou}{\textlengthmark}{ð}{\textscy}{\textsubring{r}}} \dicPos{n}[2] \dicFlx{(‑s, ‑)}[25] \dicDirectTranslationCS{paseka, mýtina}
\dicEntry[rjóður] \dicTerm{rjóður\smash{\textsuperscript{2}}} \dicIPA{{r}{j}{ou}{\textlengthmark}{ð}{\textscy}{\textsubring{r}}} \dicPos{adj}[2]\dicFlx{}[-6] \dicSynonym{rauður} \dicDirectTranslationCS{červený, načervenalý} \dicExampleIS{rjóður í kinnum} \dicExampleCS{červený v~tváři}
\dicEntry[rjómaís] \dicTerm{rjóma··ís} \dicIPA{{r}{j}{ou}{\textlengthmark}{m}{a}{i}{s}} \dicPos{m}[4] \dicFlx{(‑s, ‑ar)}[25] \dicDirectTranslationCS{(smetanová) zmrzlina}
\dicEntry[rjómakaka] \dicTerm{rjóma··|kaka} \dicIPA{{r}{j}{ou}{\textlengthmark}{m}{a}{k\smash{\textsuperscript{h}}}{a}{\r{g}}{a}} \dicPos{f}[1] \dicFlx{(‑köku, ‑kökur)}[20] \dicFieldCat{kulin.} \dicDirectTranslationCS{šlehačkový\,/\addthin krémový dort}
\dicEntry[rjómi] \dicTerm{rjóm|i} \dicIPA{{r}{j}{ou}{\textlengthmark}{m}{\textsci}} \dicPos{m}[1] \dicFlx{(‑a)}[3] \dicDirectTranslationCS{smetana, šlehačka, krém};  \dicPhraseIS{sýrður rjómi} \dicDirectTranslationCS{zakysaná smetana};  \dicPhraseIS{þeyta rjóma} \dicDirectTranslationCS{(u)šlehat smetanu};  \dicPhraseIS{þeyttur rjómi} \dicDirectTranslationCS{šlehačka}
\dicEntry[rjúfa] \dicTerm{rjúfa} \dicsymFrequent\  \dicIPA{{r}{j}{u}{\textlengthmark}{v}{a}} \dicPos{v}[6] \dicFlx{(rýf, rauf, rufum, ryfi, rofið)}[105] \dicFlx{acc} \dicSynonym{sundra} \dicDirectTranslationCS{prolomit, prorazit, probořit} \dicExampleIS{rjúfa einangrun e‑rs} \dicExampleCS{prolomit (čí) izolaci};  \dicPhraseIS{rjúfa þögnina} \dicDirectTranslationCS{prolomit ticho}
\dicEntry[rjúka] \dicTerm{rjúka} \dicsymFrequent\  \dicIPA{{r}{j}{u}{\textlengthmark}{\r{g}}{a}} \dicPos{v}[6] \dicFlx{(rýk, rauk, rukum, ryki, rokið)}[102] \textbf{1.} \dicPhraseIS{það rýkur úr e‑u} \dicFlx{impers} \dicSynonym{ósa} \dicDirectTranslationCS{z~(čeho) se kouří\,/\addthin čoudí\,/\addthin čadí} \dicExampleIS{Það rýkur upp úr strompinum.} \dicExampleCS{Z~komínu se kouří.}  \textbf{2.} \dicDirectTranslationCS{rozviřovat, vířit (prach ap.)}  \textbf{3.} \dicSynonym{þjóta} \dicDirectTranslationCS{pádit, (vy)letět, (vy)řítit se} \dicExampleIS{rjúka út úr húsinu} \dicExampleCS{vyřítit se z~domu};  \dicIdiom{rjúka}[á]{ \dicPhraseIS{það rýkur á e‑ð}} \dicFlx{impers} \dicDirectTranslationCS{(co) propuká (nečas ap.)};  \dicIdiom{rjúka}[í]{ \dicPhraseIS{rjúka í að (gera e‑ð)}} \dicLangCat{přen.} \dicDirectTranslationCS{vletět po hlavě do (dělání (čeho))};  \dicIdiom{rjúka}[upp]{ \dicPhraseIS{rjúka upp}} {\textbf{a.}} \dicDirectTranslationCS{vzkypět (hněvem), vyletět (z~kůže)};  {\textbf{b.}} \dicDirectTranslationCS{vyletět, dramaticky vzrůst};  \dicIdiom{rjúka}[til]{ \dicPhraseIS{rjúka til}} \dicDirectTranslationCS{hodit sebou, mrsknout sebou}
\dicEntry[rjúpa] \dicTerm{rjúp|a} \dicIPA{{r}{j}{u}{\textlengthmark}{\textsubring{b}}{a}} \dicPos{f}[1] \dicFlx{(‑u, ‑ur)}[13] \dicFieldCat{zool.} \dicDirectTranslationCS{bělokur, bělokur horský} \textit{(l.~{\textLA{Lagopus muta}})}  \dicsymPhoto\ ;  \dicIdiom{rjúpa}{ \dicPhraseIS{rembast eins og rjúpan við staurinn}} \dicFlx{refl} \dicLangCat{přen.} \dicDirectTranslationCS{snažit se ze všech sil} \dicIndirectTranslationCS{(většinou bez velkého úspěchu)}
\dicFigure{20809.jpg}{Rjúpa}{Rjúpa - Bollinger, Karen, Biolib, PD}
\dicEntry[roð] \dicTerm{roð} \dicIPA{{r}{\textopeno}{\textlengthmark}{\texttheta}} \dicPos{n}[2] \dicFlx{(‑s, ‑)}[5] \dicDirectTranslationCS{(rybí) kůže};  \dicIdiom{roð}{ \dicPhraseIS{e‑að er rýrt í roðinu}} \dicDirectTranslationCS{(co) je slabé\,/\addthin nevalné (návrh ap.)}; { \dicPhraseIS{hafa ekki roð við e‑m}} \dicLangCat{přen.} \dicDirectTranslationCS{nemít na (koho)};  \dicPhraseIS{vera eins og snúið roð í hundskjaft\,/\addthin hund} \dicLangCat{přen.} \dicDirectTranslationCS{vypadat, jako by (komu) uletěly včely};  \dicPhraseIS{vera tvöfaldur í roðinu} \dicDirectTranslationCS{být licoměrný\,/\addthin pokrytecký}
\dicEntry[roði] \dicTerm{roð|i} \dicIPA{{r}{\textopeno}{\textlengthmark}{ð}{\textsci}} \dicPos{m}[1] \dicFlx{(‑a, ‑ar)}[1] \dicDirectTranslationCS{zardění, začervenání} \dicExampleIS{roði í kinnum} \dicExampleCS{zardění tváří}
\dicEntry[roðið] \dicTerm{roðið} \dicIPA{{r}{\textopeno}{\textlengthmark}{ð}{\textsci}{\texttheta}} \dicPos{v} \dicFlx{supin} \dicLink{rjóða}
\dicEntry[roðna] \dicTerm{roðn|a} \dicIPA{{r}{\textopeno}{ð}{n}{a}} \dicPos{v}[1] \dicFlx{(‑aði)}[44] \dicDirectTranslationCS{(za)červenat se, (za)rdít se} \dicExampleIS{roðna út undir eyru} \dicExampleCS{začervenat se až po uši}
\dicEntry[rof] \dicTerm{rof} \dicIPA{{r}{\textopeno}{\textlengthmark}{f}} \dicPos{n}[2] \dicFlx{(‑s, ‑)}[5] \textbf{1.} \dicSynonym*{rauf} \dicDirectTranslationCS{trhlina, prasklina}  \textbf{2.} \dicDirectTranslationCS{porušení (smlouvy ap.)} \dicExampleIS{rof á samningi} \dicExampleCS{porušení smlouvy}  \textbf{3.} \dicDirectTranslationCS{přerušení} \dicExampleIS{rof á straumi} \dicExampleCS{přerušení proudu}  \textbf{4.} \dicFieldCat{geol.} \dicSynonym{veðrun} \dicDirectTranslationCS{zvětrávání, eroze}
\dicEntry[rofi] \dicTerm{rof|i} \dicIPA{{r}{\textopeno}{\textlengthmark}{v}{\textsci}} \dicPos{m}[1] \dicFlx{(‑a, ‑ar)}[1] \dicFieldCat{elek.} \dicDirectTranslationCS{vypínač, spínač, přepínač}
\dicEntry[rofið] \dicTerm{rofið} \dicIPA{{r}{\textopeno}{\textlengthmark}{v}{\textsci}{\texttheta}} \dicPos{v} \dicFlx{supin} \dicLink{rjúfa}
\dicEntry[rofinn] \dicTerm{rofinn} \dicIPA{{r}{\textopeno}{\textlengthmark}{v}{\textsci}{\textsubring{n}}} \dicPos{adj}[6]\dicFlx{}[-6] \textbf{1.} \dicDirectTranslationCS{prolomený, probořený}  \textbf{2.} \dicDirectTranslationCS{přerušený} \dicExampleIS{Sambandið er rofið.} \dicExampleCS{Spojení je přerušené.}  \textbf{3.} \dicDirectTranslationCS{porušený (smlouva ap.)}
\dicEntry[rofna] \dicTerm{rofn|a} \dicIPA{{r}{\textopeno}{\textsubring{b}}{n}{a}} \dicPos{v}[1] \dicFlx{(‑aði)}[34] \textbf{1.} \dicSynonym{slitna} \dicDirectTranslationCS{přerušit se, být přerušen} \dicExampleIS{Vegasamband rofnaði í flóðunum.} \dicExampleCS{Silniční spojení bylo přerušeno během povodní.}  \textbf{2.} \dicDirectTranslationCS{zvětrat, erodovat}
\dicEntry[rogastans] \dicTerm{roga··stans} \dicIPA{{r}{\textopeno}{\textlengthmark}{\textbabygamma}{a}{s}{\textsubring{d}}{a}{n}{s}} \dicPos{n}[2] \dicFlx{(‑)}[24] \dicPhraseIS{e‑n rekur í rogastans} \dicFlx{impers} \dicDirectTranslationCS{(komu) to vyráží dech, (kdo) je ohromen}
\dicEntry[rogginn] \dicTerm{rogginn} \dicIPA{{r}{\textopeno}{\r{\textObardotlessj}}{\textlengthmark}{\textsci}{\textsubring{n}}} \dicPos{adj}[6]\dicFlx{}[-2] \dicSynonym{hreykinn} \dicDirectTranslationCS{pyšný, domýšlivý}
\dicEntry[rok] \dicTerm{rok} \dicIPA{{r}{\textopeno}{\textlengthmark}{\r{g}}} \dicPos{n}[2] \dicFlx{(‑s, ‑)}[5] \dicDirectTranslationCS{silná vichřice} \dicIndirectTranslationCS{(Beaufort 10) (24,5--28,4~m/s)} \dicExampleIS{Það er öskrandi rok.} \dicExampleCS{Je strašná vichřice.}
\dicEntry[rokið] \dicTerm{rokið} \dicIPA{{r}{\textopeno}{\textlengthmark}{\r{\textObardotlessj}}{\textsci}{\texttheta}} \dicPos{v} \dicFlx{supin} \dicLink{rjúka}
\dicEntry[rokk] \dicTerm{rokk} \dicIPA{{r}{\textopeno}{h}{\r{g}}} \dicPos{n}[2] \dicFlx{(‑s)}[2] \dicFieldCat{hud.} \dicDirectTranslationCS{rock}
\dicEntry[rokkari] \dicTerm{rokk··ar|i} \dicIPA{{r}{\textopeno}{h}{\r{g}}{a}{r}{\textsci}} \dicPos{m}[1] \dicFlx{(‑a, ‑ar)}[13] \dicDirectTranslationCS{rocker(ka)}
\dicEntry[rokktónlist] \dicTerm{rokk··tón·list} \dicIPA{{r}{\textopeno}{h}{\r{g}}{t\smash{\textsuperscript{h}}}{ou}{n}{l}{\textsci}{s}{\textsubring{d}}} \dicPos{f}[7] \dicFlx{(‑ar)}[3] \dicFieldCat{hud.} \dicDirectTranslationCS{rock, rocková hudba}
\dicEntry[rokkur] \dicTerm{rokk|ur} \dicIPA{{r}{\textopeno}{h}{\r{g}}{\textscy}{\textsubring{r}}} \dicPos{m}[6] \dicFlx{(‑s, ‑ar)}[22] \dicDirectTranslationCS{kolovrat, kolovrátek}
\dicEntry[roksala] \dicTerm{rok··|sala} \dicIPA{{r}{\textopeno}{\textlengthmark}{\r{g}}{s}{a}{l}{a}} \dicPos{f}[1] \dicFlx{(‑sölu, ‑sölur)}[20] \dicDirectTranslationCS{totální výprodej}
\dicEntry[rola] \dicTerm{rol|a} \dicIPA{{r}{\textopeno}{\textlengthmark}{l}{a}} \dicPos{f}[1] \dicFlx{(‑u, ‑ur)}[19] \dicSynonym*{daufingi} \dicDirectTranslationCS{měkota, padavka, chcípáček}
\dicEntry[rolla] \dicTerm{roll|a} \dicIPA{{r}{\textopeno}{\textsubring{d}}{l}{a}} \dicPos{f}[1] \dicFlx{(‑u, ‑ur)}[19] \textbf{1.} \dicSynonym{ær\smash{\textsuperscript{1}}} \dicDirectTranslationCS{ovce}  \textbf{2.} \dicSynonym{romsa\smash{\textsuperscript{1}}} \dicDirectTranslationCS{obšírný text}  \textbf{3.} \dicSynonym*{bókarvöndull} \dicDirectTranslationCS{svitek, rulička}
\dicEntry[romm] \dicTerm{romm} \dicIPA{{r}{\textopeno}{m}{\textlengthmark}} \dicPos{n}[2] \dicFlx{(‑s)}[2] \dicDirectTranslationCS{rum}
\dicEntry[romsa] \dicTerm{roms|a\smash{\textsuperscript{1}}} \dicIPA{{r}{\textopeno}{m}{s}{a}} \dicPos{f}[1] \dicFlx{(‑u, ‑ur)}[7] \dicSynonym{langloka} \dicDirectTranslationCS{litanie, obšírný výklad}
\dicEntry[romsa] \dicTerm{roms|a\smash{\textsuperscript{2}}} \dicIPA{{r}{\textopeno}{m}{s}{a}} \dicPos{v}[1] \dicFlx{(‑aði)}[1] \dicDirectTranslationCS{řečnit, plkat};  \dicIdiom{romsa}[upp]{ \dicPhraseIS{romsa upp e‑u}} \dicDirectTranslationCS{odříkávat (co), vyjmenovávat (co)};  \dicPhraseIS{romsa upp úr sér e‑u} \dicDirectTranslationCS{vychrlit (co) ze sebe}
\dicEntry[ropa] \dicTerm{rop|a} \dicIPA{{r}{\textopeno}{\textlengthmark}{\textsubring{b}}{a}} \dicPos{v}[1] \dicFlx{(‑aði)}[1] \textbf{1.} \dicDirectTranslationCS{(od)říhnout si, říhat, krknout, krkat}  \textbf{2.} \dicDirectTranslationCS{krákat} \dicIndirectTranslationCS{(jako bělokur)}
\dicEntry[ropi] \dicTerm{rop|i} \dicIPA{{r}{\textopeno}{\textlengthmark}{\textsubring{b}}{\textsci}} \dicPos{m}[1] \dicFlx{(‑a, ‑ar)}[1] \dicDirectTranslationCS{říhnutí, krknutí}
\dicEntry[rosalega] \dicTerm{rosa··lega} \dicIPA{{r}{\textopeno}{\textlengthmark}{s}{a}{l}{\textepsilon}{\textbabygamma}{a}} \dicPos{adv} \dicSynonym{ofsalega} \dicDirectTranslationCS{moc, fantasticky} \dicExampleIS{rosalega gott} \dicExampleCS{moc dobré}
\dicEntry[rosalegur] \dicTerm{rosa··legur} \dicIPA{{r}{\textopeno}{\textlengthmark}{s}{a}{l}{\textepsilon}{\textbabygamma}{\textscy}{\textsubring{r}}} \dicPos{adj}[1]\dicFlx{}[-8] \dicSynonym{geysilegur} \dicDirectTranslationCS{obrovský, kolosální, mega-} \dicExampleIS{rosalegt atvinnuleysi} \dicExampleCS{obrovská nezaměstnanost}
\dicEntry[roskinn] \dicTerm{roskinn} \dicIPA{{r}{\textopeno}{s}{\r{\textObardotlessj}}{\textsci}{\textsubring{n}}} \dicPos{adj}[6]\dicFlx{}[-2] \dicDirectTranslationCS{starší, postarší} \dicIndirectTranslationCS{(o~člověku)}
\dicEntry[rosti] \dicTerm{rost|i} \dicIPA{{r}{\textopeno}{s}{\textsubring{d}}{\textsci}} \dicPos{m}[1] \dicFlx{(‑a)}[3] \dicSynonym{hroki} \dicDirectTranslationCS{nadutost, arogance} \dicExampleIS{vera með rosta} \dicExampleCS{být arogantní};  \dicPhraseIS{lækka\,/\addthin lægja rostann í e‑m} \dicLangCat{přen.} \dicDirectTranslationCS{srazit (komu) hřebínek, naučit (koho) slušnému chování}
\dicEntry[rostungur] \dicTerm{rost··ung|ur} \dicIPA{{r}{\textopeno}{s}{\textsubring{d}}{u}{\ng}{\r{g}}{\textscy}{\textsubring{r}}} \dicPos{m}[6] \dicFlx{(‑s, ‑ar)}[8] \dicFieldCat{zool.} \dicDirectTranslationCS{mrož} \textit{(l.~{\textLA{Odobenus rosmarus}})}  \dicsymPhoto\ 
\dicFigure{21298.jpg}{Rostungur}{Rostungur - U.S. Fish and Wildlife Service, Biolib, PD}
\dicEntry[rot] \dicTerm{rot} \dicIPA{{r}{\textopeno}{\textlengthmark}{\textsubring{d}}} \dicPos{n}[2] \dicFlx{(‑s, ‑)}[5] \dicSynonym{meðvitundarleysi} \dicDirectTranslationCS{omráčení, bezvědomí} \dicIndirectTranslationCS{(po úderu do hlavy)};  \dicPhraseIS{liggja í roti} \dicDirectTranslationCS{ležet v~bezvědomí};  \dicPhraseIS{rakna úr roti\,/\addthin rotinu} \dicSynonym*{fá meðvitund} \dicDirectTranslationCS{přijít ke smyslům, probrat se z~bezvědomí};  \dicPhraseIS{slá e‑n í rot} \dicSynonym{rota} \dicDirectTranslationCS{omráčit (koho)}
\dicEntry[rota] \dicTerm{rot|a} \dicIPA{{r}{\textopeno}{\textlengthmark}{\textsubring{d}}{a}} \dicPos{v}[1] \dicFlx{(‑aði)}[1] \dicFlx{acc} \dicDirectTranslationCS{omráčit (úderem do hlavy)};  \dicIdiom{rotast}{ \dicPhraseIS{rotast}} \dicFlx{refl} \dicDirectTranslationCS{omráčit se, být omráčen}
\dicEntry[rothögg] \dicTerm{rot··högg} \dicIPA{{r}{\textopeno}{\textlengthmark}{\textsubring{d}}{h}{\oe}{\r{g}}} \dicPos{n}[2] \dicFlx{(‑s, ‑)}[5] \textbf{1.} \dicSynonym{högg} \dicDirectTranslationCS{knokaut}  \textbf{2.} \dicLangCat{přen.} \dicSynonym{skaðvaldur} \dicDirectTranslationCS{knokaut, smrtelná rána} \dicExampleIS{rothögg fyrir félagið} \dicExampleCS{smrtelná rána pro společnost}
\dicEntry[rotinn] \dicTerm{rotinn} \dicIPA{{r}{\textopeno}{\textlengthmark}{\textsubring{d}}{\textsci}{\textsubring{n}}} \dicPos{adj}[6]\dicFlx{}[-2] \textbf{1.} \dicSynonym{úldinn} \dicDirectTranslationCS{shnilý, prohnilý, zetlelý}  \textbf{2.} \dicSynonym{spilltur} \dicDirectTranslationCS{zkažený, prohnilý (společnost ap.)}
\dicEntry[rotna] \dicTerm{rotn|a} \dicIPA{{r}{\textopeno}{h}{\textsubring{d}}{n}{a}} \dicPos{v}[1] \dicFlx{(‑aði)}[1] \dicSynonym*{brotna niður} \dicDirectTranslationCS{hnít, tlít, rozkládat se}
\dicEntry[rotnun] \dicTerm{rotn|un} \dicIPA{{r}{\textopeno}{h}{\textsubring{d}}{n}{\textscy}{\textsubring{n}}} \dicPos{f}[7] \dicFlx{(‑unar)}[9] \dicDirectTranslationCS{hnití, uhnívání, tlení, trouchnivění}
\dicEntry[rotta] \dicTerm{rott|a} \dicIPA{{r}{\textopeno}{h}{\textsubring{d}}{a}} \dicPos{f}[1] \dicFlx{(‑u, ‑ur)}[19] \textbf{1.} \dicFieldCat{zool.} \dicDirectTranslationCS{krysa} \textit{(l.~{\textLA{Rattus}})}  \textbf{2.} \dicSynonym*{ódrenglyndur maður} \dicDirectTranslationCS{krysa} \dicIndirectTranslationCS{(bezcharakterní člověk)}
\dicEntry[rotvarnarefni] \dicTerm{rot·varnar··efni} \dicIPA{{r}{\textopeno}{\textlengthmark}{\textsubring{d}}{v}{a}{r}{\textsubring{d}}{n}{a}{r}{\textepsilon}{\textsubring{b}}{n}{\textsci}} \dicPos{n}[2] \dicFlx{(‑s, ‑)}[14] \dicDirectTranslationCS{konzervant}
\dicEntry[ró] \dicTerm{ró\smash{\textsuperscript{1}}} \dicsymFrequent\  \dicIPA{{r}{ou}{\textlengthmark}} \dicPos{f}[9] \dicFlx{(‑ar)}[3] \dicSynonym{kyrrð} \dicDirectTranslationCS{klid, pokoj, mír};  \dicIdiom{ró}{ \dicPhraseIS{halda ró sinni}} \dicDirectTranslationCS{zůstat v~klidu};  \dicPhraseIS{í ró og næði} \dicFlx{adv} \dicDirectTranslationCS{v~klidu a~míru};  \dicPhraseIS{raska ró e‑rs} \dicDirectTranslationCS{narušit (čí) klid};  \dicPhraseIS{taka e‑u með ró} \dicDirectTranslationCS{brát (co) s~klidem};  \dicPhraseIS{vera ekki í rónni} \dicDirectTranslationCS{nebýt v~klidu, být rozrušený}
\dicEntry[ró] \dicTerm{ró\smash{\textsuperscript{2}}} \dicIPA{{r}{ou}{\textlengthmark}} \dicPos{f}[9] \dicFlx{(‑ar, rær)}[1] \dicDirectTranslationCS{matice, matička} \dicExampleIS{ró og bolti} \dicExampleCS{matička a~šroubek}
\dicEntry[róa] \dicTerm{ró|a\smash{\textsuperscript{1}}} \dicsymFrequent\  \dicIPA{{r}{ou}{\textlengthmark}{a}} \dicPos{v}[1] \dicFlx{(‑aði)}[1] \dicFlx{acc} \dicSynonym*{gera rólegan} \dicDirectTranslationCS{utišit, upokojit, ukonejšit, zklidnit} \dicExampleIS{Ég reyndi að róa hana.} \dicExampleCS{Zkoušel jsem ji utišit.};  \dicIdiom{róa}[niður]{ \dicPhraseIS{róa e‑n niður}} \dicDirectTranslationCS{uklidnit (koho), zklidnit (koho)} \dicExampleIS{róa sig niður} \dicExampleCS{uklidnit se};  \dicIdiom{róast}{ \dicPhraseIS{róast}} \dicFlx{refl} \dicSynonym*{sefast} \dicDirectTranslationCS{utišit se, upokojit se, uklidnit se}
\dicEntry[róa] \dicTerm{róa\smash{\textsuperscript{2}}} \dicsymFrequent\  \dicIPA{{r}{ou}{\textlengthmark}{a}} \dicPos{v}[5] \dicFlx{(ræ, reri\,/\addthin réri, rerum\,/\addthin rérum, reri\,/\addthin réri, róið)}[2] \dicFlx{dat} \dicDirectTranslationCS{veslovat, pádlovat} \dicExampleIS{róa bátnum yfir fjörðinn} \dicExampleCS{veslovat s~člunem přes fjord};  \dicIdiom{róa}[fram]{ \dicPhraseIS{róa fram í gráðið}} \dicDirectTranslationCS{sedět neklidně na židli a~houpat se};  \dicIdiom{róa}[til]{ \dicPhraseIS{róa til fiskjar}} \dicDirectTranslationCS{jet na rybolov};  \dicIdiom{róa}{ \dicPhraseIS{láta e‑n róa}} \dicLangCat{přen.} \dicDirectTranslationCS{zbavit se (koho), nechat (koho) plavat}; { \dicPhraseIS{það er ekki á vísan að róa}} \dicLangCat{přen.} \dicDirectTranslationCS{není to jisté}
\dicEntry[róandi] \dicTerm{ró··andi} \dicIPA{{r}{ou}{\textlengthmark}{a}{n}{\textsubring{d}}{\textsci}} \dicPos{adj}[13] \dicFlx{indecl}[1] \dicDirectTranslationCS{konejšivý, utišující, uklidňující} \dicExampleIS{róandi áhrif} \dicExampleCS{uklidňující vliv};  \dicPhraseIS{róandi lyf} \dicDirectTranslationCS{utišující lék, sedativum}
\dicEntry[róbot] \dicTerm{róbot} \dicIPA{{r}{ou}{\textlengthmark}{\textsubring{b}}{\textopeno}{\textsubring{d}}} \dicPos{m}[4] \dicFlx{(‑s, ‑ar)}[10] \dicSynonym{vélmenni} \dicDirectTranslationCS{robot}
\dicEntry[ródín] \dicTerm{ródín} \dicIPA{{r}{ou}{\textlengthmark}{\textsubring{d}}{i}{\textsubring{n}}} \dicPos{n}[2] \dicFlx{(‑s)}[2] \dicFieldCat{chem.} \dicDirectTranslationCS{rhodium} \textit{(l.~{\textLA{Rh, Rhodium}})}
\dicEntry[róði] \dicTerm{róð|i} \dicIPA{{r}{ou}{\textlengthmark}{ð}{\textsci}} \dicPos{m}[1] \dicFlx{(‑a, ‑ar)}[1] \dicPhraseIS{kasta\,/\addthin varpa e‑u fyrir róða} \dicLangCat{přen.} \dicDirectTranslationCS{pověsit (co) na hřebík}
\dicEntry[róðukross] \dicTerm{róðu··kross} \dicIPA{{r}{ou}{\textlengthmark}{ð}{\textscy}{k\smash{\textsuperscript{h}}}{r}{\textopeno}{s}} \dicPos{m}[4] \dicFlx{(‑, ‑ar)}[28] \dicDirectTranslationCS{krucifix}
\dicEntry[róður] \dicTerm{róð|ur} \dicIPA{{r}{ou}{\textlengthmark}{ð}{\textscy}{\textsubring{r}}} \dicPos{m}[5] \dicFlx{(‑rar\,/\addthin ‑urs, ‑rar)}[9] \dicDirectTranslationCS{veslování, pádlování};  \dicIdiom{róður}{ \dicPhraseIS{þetta er þungur róður}} \dicLangCat{přen.} \dicDirectTranslationCS{je to dřina};  \dicProverb\  \dicPhraseIS{Þungur er þegjandi róður.} \dicLangCat{přís.} \dicDirectTranslationCS{Těžké je tiché veslování.} \dicIndirectTranslationCS{(ve významu -- bez dobré nálady je každá práce ještě více vyčerpávající)}
\dicEntry[rófa] \dicTerm{róf|a} \dicIPA{{r}{ou}{\textlengthmark}{a}} \dicPos{f}[1] \dicFlx{(‑u, ‑ur)}[19] \textbf{1.} \dicSynonym{skott} \dicDirectTranslationCS{ocas, ohon, oháňka} \dicIndirectTranslationCS{(kočkovitých šelem, opice a~medvěda)}  \textbf{2.} \dicFieldCat{bot.} \dicSynonym{gulrófa} \dicDirectTranslationCS{tuřín, brukev řepka tuřín} \textit{(l.~{\textLA{Brassica napus var. napobrassica}})}
\dicEntry[rófubein] \dicTerm{rófu··bein} \dicIPA{{r}{ou}{\textlengthmark}{\textscy}{\textsubring{b}}{ei}{\textsubring{n}}} \dicPos{n}[2] \dicFlx{(‑s, ‑)}[5] \dicFieldCat{anat.} \dicDirectTranslationCS{kostrč}
\dicEntry[rógbar] \dicTerm{róg··bar} \dicIPA{{r}{ou}{\textlengthmark}{\textsubring{b}}{a}{\textsubring{r}}} \dicPos{v} \dicFlx{ind pf sg 1 pers} \dicLink{rógbera}
\dicEntry[rógbárum] \dicTerm{róg··bárum} \dicIPA{{r}{ou}{\textlengthmark}{\textsubring{b}}{au}{r}{\textscy}{\textsubring{m}}} \dicPos{v} \dicFlx{ind pf pl 1 pers} \dicLink{rógbera}
\dicEntry[rógber] \dicTerm{róg··ber} \dicIPA{{r}{ou}{\textlengthmark}{\textsubring{b}}{\textepsilon}{\textsubring{r}}} \dicPos{v} \dicFlx{ind praes sg 1 pers} \dicLink{rógbera}
\dicEntry[rógbera] \dicTerm{róg··|bera} \dicIPA{{r}{ou}{\textlengthmark}{\textsubring{b}}{\textepsilon}{r}{a}} \dicPos{v}[6] \dicFlx{(‑ber, ‑barði, ‑börðum, ‑bæri, ‑barið)}[5] \dicFlx{acc} \dicSynonym{rægja} \dicDirectTranslationCS{pomluvit, pomlouvat, očernit, hanobit, (po)hanět}
\dicEntry[rógberi] \dicTerm{róg··ber|i} \dicIPA{{r}{ou}{\textlengthmark}{\textsubring{b}}{\textepsilon}{r}{\textsci}} \dicPos{m}[1] \dicFlx{(‑a, ‑ar)}[1] \dicDirectTranslationCS{pomlouvač(ka), hanobitel(ka)}
\dicEntry[rógborið] \dicTerm{róg··borið} \dicIPA{{r}{ou}{\textlengthmark}{\textsubring{b}}{\textopeno}{r}{\textsci}{\texttheta}} \dicPos{v} \dicFlx{supin} \dicLink{rógbera}
\dicEntry[rógburður] \dicTerm{róg··burð|ur} \dicIPA{{r}{ou}{\textlengthmark}{\textsubring{b}}{\textscy}{r}{ð}{\textscy}{\textsubring{r}}} \dicPos{m}[10] \dicFlx{(‑ar)}[7] \dicDirectTranslationCS{pomluva, pohana, očernění}
\dicEntry[rógbæri] \dicTerm{róg··bæri} \dicIPA{{r}{ou}{\textlengthmark}{\textsubring{b}}{a}{i}{r}{\textsci}} \dicPos{v} \dicFlx{con pf sg 1 pers} \dicLink{rógbera}
\dicEntry[rógur] \dicTerm{róg|ur} \dicIPA{{r}{ou}{\textlengthmark}{\textscy}{\textsubring{r}}} \dicPos{m}[6] \dicFlx{(‑s)}[7] \dicDirectTranslationCS{pomluva, pohana, očernění}
\dicEntry[róið] \dicTerm{róið} \dicIPA{{r}{ou}{\textlengthmark}{\textsci}{\texttheta}} \dicPos{v} \dicFlx{supin} \dicLink{róa\smash{\textsuperscript{2}}}
\dicEntry[ról] \dicTerm{ról} \dicIPA{{r}{ou}{\textlengthmark}{\textsubring{l}}} \dicPos{n}[2] \dicFlx{(‑s, ‑)}[5] \dicSynonym{rölt} \dicDirectTranslationCS{procházka, vycházka, špacír};  \dicPhraseIS{vera á róli} \dicDirectTranslationCS{být na procházce, procházet se};  \dicPhraseIS{vera á góðu róli með e‑ð} \dicLangCat{přen.} \dicDirectTranslationCS{být na dobré cestě s~(čím)}
\dicEntry[róla] \dicTerm{ról|a\smash{\textsuperscript{1}}} \dicIPA{{r}{ou}{\textlengthmark}{l}{a}} \dicPos{f}[1] \dicFlx{(‑u, ‑ur)}[19] \dicDirectTranslationCS{houpačka}
\dicEntry[róla] \dicTerm{ról|a\smash{\textsuperscript{2}}} \dicIPA{{r}{ou}{\textlengthmark}{l}{a}} \dicPos{v}[1] \dicFlx{(‑aði)}[1] \dicFlx{dat} \dicPhraseIS{róla sér} {\textbf{a.}} \dicDirectTranslationCS{procházet se, špacírovat};  {\textbf{b.}} \dicDirectTranslationCS{houpat se} \dicIndirectTranslationCS{(na houpačce)}
\dicEntry[rólega] \dicTerm{ró··leg|a} \dicsymFrequent\  \dicIPA{{r}{ou}{\textlengthmark}{l}{\textepsilon}{\textbabygamma}{a}} \dicPos{adv} \dicFlx{(comp ‑ar, sup ‑ast)} \dicSynonym{stillilega} \dicDirectTranslationCS{klidně, pomalu, zvolna};  \dicIdiom{rólega}{ \dicPhraseIS{anda rólega}} \dicDirectTranslationCS{dýchat zhluboka};  \dicPhraseIS{hægt og rólega} \dicFlx{adv} \dicDirectTranslationCS{pomalu a~klidně};  \dicPhraseIS{taka það\,/\addthin því rólega} \dicDirectTranslationCS{být v~klidu, odpočívat}
\dicEntry[rólegheit] \dicTerm{ró·leg··heit} \dicIPA{{r}{ou}{\textlengthmark}{l}{\textepsilon}{\textbabygamma}{h}{ei}{\textsubring{d}}} \dicPos{n}[2] \dicFlx{pl}[1] \dicSynonym{kyrrð} \dicDirectTranslationCS{(po)klid};  \dicPhraseIS{í rólegheitum} \dicFlx{adv} \dicDirectTranslationCS{v~klidu}
\dicEntry[rólegur] \dicTerm{ró··legur} \dicsymFrequent\  \dicIPA{{r}{ou}{\textlengthmark}{l}{\textepsilon}{\textbabygamma}{\textscy}{\textsubring{r}}} \dicPos{adj}[1]\dicFlx{}[-8] \textbf{1.} \dicSynonym{stilltur} \dicDirectTranslationCS{klidný, pohodový};  \dicPhraseIS{bíddu rólegur} \dicFlx{imper} \dicDirectTranslationCS{počkej chvilku};  \dicPhraseIS{vertu rólegur} \dicFlx{imper} \dicDirectTranslationCS{uklidni se, buď v~klidu}  \textbf{2.} \dicSynonym*{viðburðalítill} \dicDirectTranslationCS{pokojný, klidný} \dicExampleIS{rólegt hverfi} \dicExampleCS{klidná čtvrť}
\dicEntry[rólfær] \dicTerm{ról··fær} \dicIPA{{r}{ou}{\textlengthmark}{l}{f}{a}{i}{\textsubring{r}}} \dicPos{adj}[5] \dicFlx{(f ‑)}[8] \dicDirectTranslationCS{schopný pohybu (pacient po nemoci ap.)}
\dicEntry[rólyndur] \dicTerm{ró··lyndur} \dicIPA{{r}{ou}{\textlengthmark}{l}{\textsci}{n}{\textsubring{d}}{\textscy}{\textsubring{r}}} \dicPos{adj}[2]\dicFlx{}[-14] \dicSynonym{stilltur} \dicDirectTranslationCS{klidný, (jsoucí) klidné mysli, (duševně) vyrovnaný}
\dicEntry[Róm] \dicTerm{Róm}\dicTerm{, Rómaborg} \dicIPA{{r}\-{ou}\-{\textlengthmark}\-{\textsubring{m}}\-} \dicPos{f}[4] \dicFlx{(‑ar)}[4] \dicFieldCat{geog.} \dicDirectTranslationCS{Řím} \dicIndirectTranslationCS{(hlavní město Itálie)}
\dicEntry[róma] \dicTerm{róm|a} \dicIPA{{r}{ou}{\textlengthmark}{m}{a}} \dicPos{v}[1] \dicFlx{(‑aði)}[1] \dicFlx{acc} \dicSynonym{lofa} \dicDirectTranslationCS{(vy)chválit, vychvalovat} \dicExampleIS{róma þjónustuna} \dicExampleCS{vychvalovat servis}
\dicEntry[Rómaborg] \dicTerm{Róma··borg} \dicIPA{{r}{ou}{\textlengthmark}{m}{a}{\textsubring{b}}{\textopeno}{r}{\r{g}}} \dicPos{f}[7] \dicFlx{(‑ar)}[4] \dicLink{Róm}
\dicEntry[rómaður] \dicTerm{róm··|aður} \dicIPA{{r}{ou}{\textlengthmark}{m}{a}{ð}{\textscy}{\textsubring{r}}} \dicPos{adj}[3] \dicFlx{(f ‑uð)}[3] \dicSynonym{nafnfrægur} \dicDirectTranslationCS{vyhlášený, věhlasný, legendární} \dicExampleIS{rómaður fyrir e‑ð} \dicExampleCS{vyhlášený (čím)}
\dicEntry[róman] \dicTerm{róman} \dicIPA{{r}{ou}{\textlengthmark}{m}{a}{\textsubring{n}}} \dicPos{m}[4] \dicFlx{(‑s, ‑ar)}[11] \dicLangCat{hovor.} \dicSynonym{skáldsaga} \dicDirectTranslationCS{román}
\dicEntry[rómanskur] \textls[15]{\dicTerm{rómanskur} \dicIPA{{r}{ou}{\textlengthmark}{m}{a}{n}{s}{\r{g}}{\textscy}{\textsubring{r}}} \dicPos{adj}[1] \dicFlx{(f rómönsk)}[3] \dicDirectTranslationCS{románský};  \dicPhraseIS{Rómanska Ameríka, Rómanska-Ameríka} \dicFieldCat{geog.} \dicDirectTranslationCS{La\-tin\-ská Amerika};  \dicPhraseIS{rómanskur stíll} \dicDirectTranslationCS{románský styl};  \dicPhraseIS{ró\-mönsk mál} \dicDirectTranslationCS{románské jazyky}}
\dicEntry[rómantík] \dicTerm{rómantík} \dicIPA{{r}{ou}{\textlengthmark}{m}{a}{\textsubring{n}}{\textsubring{d}}{i}{\r{g}}} \dicPos{f}[10] \dicFlx{(‑ur)}[2] \textbf{1.} \dicDirectTranslationCS{romantismus}  \textbf{2.} \dicDirectTranslationCS{romantika, romantičnost}
\dicEntry[rómantískur] \dicTerm{rómantískur} \dicIPA{{r}{ou}{\textlengthmark}{m}{a}{\textsubring{n}}{\textsubring{d}}{i}{s}{\r{g}}{\textscy}{\textsubring{r}}} \dicPos{adj}[1]\dicFlx{}[-1] \dicDirectTranslationCS{romantický};  \dicPhraseIS{rómantíska stefnan} \dicDirectTranslationCS{romantické hnutí}
\dicEntry[rómur] \dicTerm{róm|ur} \dicsymFrequent\  \dicIPA{{r}{ou}{\textlengthmark}{m}{\textscy}{\textsubring{r}}} \dicPos{m}[6] \dicFlx{(‑s, ‑ar)}[24] \textbf{1.} \dicSynonym{raust} \dicDirectTranslationCS{hlas};  \dicPhraseIS{e‑m liggur hátt\,/\addthin lágt rómur} \dicFlx{impers} \dicDirectTranslationCS{(kdo) má silný\,/\addthin slabý hlas}  \textbf{2.} \dicSynonym{viðurkenning} \dicDirectTranslationCS{ohlas, přijetí};  \dicPhraseIS{gera góðan róm að e‑u} \dicDirectTranslationCS{uvítat (co), (po)chválit (co)};  \dicPhraseIS{einum rómi} \dicFlx{adv} \dicDirectTranslationCS{jednohlasně, jedním hlasem}
\begin{xtolerant}{}{1pt}
\dicEntry[rómversk-kaþólskur] \dicTerm{róm·versk-kaþólskur}\addthinS\dicIPA{{r}{ou}{m}{v}{\textepsilon}{\textsubring{r}}{s}{\r{g}}{k\smash{\textsuperscript{h}}}{a}{\texttheta}{ou}{l}{s}{\r{g}}{\textscy}{\textsubring{r}}}\addthinS\dicPos{adj}[1]\dicFlx{}[-1] \dicFieldCat{náb.} \dicDirectTranslationCS{římskokatolický}
\end{xtolerant}
\dicEntry[rómverskur] \dicTerm{róm··verskur} \dicIPA{{r}{ou}{m}{v}{\textepsilon}{\textsubring{r}}{s}{\r{g}}{\textscy}{\textsubring{r}}} \dicPos{adj}[1]\dicFlx{}[-1] \dicDirectTranslationCS{římský};  \dicPhraseIS{rómversk tala} \dicDirectTranslationCS{římská číslice}
\dicEntry[róni] \dicTerm{rón|i} \dicIPA{{r}{ou}{\textlengthmark}{n}{\textsci}} \dicPos{m}[1] \dicFlx{(‑a, ‑ar)}[1] \dicLangCat{hovor.} \dicDirectTranslationCS{pobuda, ožrala, ožralka}
\dicEntry[rór] \dicTerm{rór} \dicIPA{{r}{ou}{\textlengthmark}{\textsubring{r}}} \dicPos{adj}[4]\dicFlx{}[-4] \dicDirectTranslationCS{(po)klidný, pokojný, tichý};  \dicPhraseIS{e‑m er ekki rótt} \dicFlx{impers} \dicDirectTranslationCS{(kdo) je neklidný}
\dicEntry[rós] \dicTerm{rós} \dicsymFrequent\  \dicIPA{{r}{ou}{\textlengthmark}{s}} \dicPos{f}[7] \dicFlx{(‑ar, ‑ir)}[1] \dicFieldCat{bot.} \dicDirectTranslationCS{růže} \textit{(l.~{\textLA{Rosa}})}  \dicsymPhoto\  \dicExampleIS{rauðar rósir} \dicExampleCS{červené růže};  \dicIdiom{rós}{ \dicPhraseIS{fá rós í hnappagatið}} \dicLangCat{přen.} \dicDirectTranslationCS{získat uznání};  \dicPhraseIS{gera rósir} \dicLangCat{přen.} \dicDirectTranslationCS{dělat blbosti};  \dicPhraseIS{tala undir rós} \dicLangCat{přen.} \dicDirectTranslationCS{mluvit v~náznacích, mluvit v~hádankách}; { \dicPhraseIS{þetta er enginn dans á rósum}} \dicLangCat{přen.} \dicDirectTranslationCS{to není (žádná) procházka růžovým sadem};  \dicProverb\  \dicPhraseIS{Engin rós án þyrna.} \dicLangCat{přís.} \dicDirectTranslationCS{Není růže bez trní.}
\dicFigure{6907.jpg}{Rós}{Rós - Deml Miroslav, Biolib, Copyright/CC-BY}
\dicEntry[rósakál] \dicTerm{rósa··kál} \dicIPA{{r}{ou}{\textlengthmark}{s}{a}{k\smash{\textsuperscript{h}}}{au}{\textsubring{l}}} \dicPos{n}[2] \dicFlx{(‑s)}[2] \dicFieldCat{bot.} \dicDirectTranslationCS{růžičková kapusta} \textit{(l.~{\textLA{Brassica oleracea var. gemmifera}})}
\dicEntry[rósavín] \dicTerm{rósa··vín} \dicIPA{{r}{ou}{\textlengthmark}{s}{a}{v}{i}{\textsubring{n}}} \dicPos{n}[2] \dicFlx{(‑s, ‑)}[5] \dicDirectTranslationCS{růžové víno}
\dicEntry[rósemd] \dicTerm{ró··semd} \dicIPA{{r}{ou}{\textlengthmark}{s}{\textepsilon}{m}{\textsubring{d}}} \dicPos{f}[7] \dicFlx{(‑ar)}[3] \dicDirectTranslationCS{klid, rozvaha, rozvážnost}
\dicEntry[rósemi] \dicTerm{ró··sem|i} \dicIPA{{r}{ou}{\textlengthmark}{s}{\textepsilon}{m}{\textsci}} \dicPos{f}[2] \dicFlx{(‑i)}[2] \dicSynonym*{rólyndi} \dicDirectTranslationCS{klid, (duševní) vyrovnanost}
\dicEntry[rósmarín] \dicTerm{rós··marín} \dicIPA{{r}{ou}{s}{m}{a}{r}{i}{\textsubring{n}}} \dicPos{n}[2] \dicFlx{(‑s)}[2] \dicFieldCat{bot.} \dicDirectTranslationCS{rozmarýn(a), rozmarýn lékařský} \textit{(l.~{\textLA{Rosmarinus officinalis}})}  \dicsymPhoto\ 
\dicFigure{ds_image_rosmarin_0_2.jpg}{Rósmarín}{Rósmarín - Fir0002, GFDL}
\dicEntry[rósóttur] \dicTerm{rós··óttur} \dicIPA{{r}{ou}{\textlengthmark}{s}{ou}{h}{\textsubring{d}}{\textscy}{\textsubring{r}}} \dicPos{adj}[1]\dicFlx{}[-10] \dicDirectTranslationCS{květovaný, kytičkovaný}
\dicEntry[rósrauður] \dicTerm{rós··rauður} \dicIPA{{r}{ou}{\textlengthmark}{s}{r}{\oe i}{ð}{\textscy}{\textsubring{r}}} \dicPos{adj}[2]\dicFlx{}[-6] \dicSynonym*{bleikrauður} \dicDirectTranslationCS{růžový, narůžovělý}
\dicEntry[róstugur] \dicTerm{róst··ugur} \dicIPA{{r}{ou}{s}{\textsubring{d}}{\textscy}{\textbabygamma}{\textscy}{\textsubring{r}}} \dicPos{adj}[1]\dicFlx{}[-8] \dicDirectTranslationCS{bouřlivý, neklidný, nepokojný}
\dicEntry[róstusamur] \dicTerm{róstu··|samur} \dicIPA{{r}{ou}{s}{\textsubring{d}}{\textscy}{s}{a}{m}{\textscy}{\textsubring{r}}} \dicPos{adj}[1] \dicFlx{(f ‑söm)}[2] \dicSynonym{róstugur} \dicDirectTranslationCS{neklidný, bouřlivý, pohnutý}
\dicEntry[rót] \dicTerm{rót\smash{\textsuperscript{1}}} \dicsymFrequent\  \dicIPA{{r}{ou}{\textlengthmark}{\textsubring{d}}} \dicPos{f}[8] \dicFlx{(‑ar, rætur)}[5] \textbf{1.} \dicSynonym*{plönturót} \dicDirectTranslationCS{kořen, kořínek} \dicIndirectTranslationCS{(podzemní část rostliny)};  \dicPhraseIS{e‑að festir rætur} \dicDirectTranslationCS{(co) zakořeňuje (strom ap.)}  \textbf{2.} \dicSynonym{upptök} \dicDirectTranslationCS{počátek, kořen, původ, příčina};  \dicPhraseIS{frá rótum} \dicFlx{adv} \dicDirectTranslationCS{od základu, od počátku};  \dicPhraseIS{eiga rætur sínar að rekja til e‑s} \dicDirectTranslationCS{mít kořeny v~(čem), odvozovat svoje kořeny od (čeho)}  \textbf{3.} \dicFieldCat{jaz.} \dicDirectTranslationCS{kořen}  \textbf{4.} \dicFieldCat{mat.} \dicSynonym{kvaðratrót} \dicDirectTranslationCS{odmocnina}  \textbf{5.} \dicFieldCat{poč.} \dicDirectTranslationCS{kořenový adresář};  \dicIdiom{rót}{ \dicPhraseIS{festa ekki rætur e‑s staðar}} \dicLangCat{přen.} \dicDirectTranslationCS{nezakořenit (kde), nezapustit (kde) kořeny};  \dicPhraseIS{e‑að stendur djúpum rótum} \dicLangCat{přen.} \dicDirectTranslationCS{(co) má hluboké kořeny}
\dicEntry[rót] \dicTerm{rót\smash{\textsuperscript{2}}} \dicIPA{{r}{ou}{\textlengthmark}{\textsubring{d}}} \dicPos{n}[2] \dicFlx{(‑s)}[2] \dicDirectTranslationCS{(silný) pohyb, hnutí};  \dicPhraseIS{e‑að kemur róti á e‑ð} \dicLangCat{přen.} \dicDirectTranslationCS{(co) uvádí do pohybu (co)}
\dicEntry[róta] \dicTerm{rót|a} \dicIPA{{r}{ou}{\textlengthmark}{\textsubring{d}}{a}} \dicPos{v}[1] \dicFlx{(‑aði)}[1] \dicSynonym{hreyfa} \dicDirectTranslationCS{hrabat se, přehrabovat se} \dicIndirectTranslationCS{(rozrývat rukama)} \dicExampleIS{róta í e‑u} \dicExampleCS{přehrabovat se v~(čem)};  \dicIdiom{róta}[upp]{ \dicPhraseIS{róta upp e‑u}} \dicDirectTranslationCS{obracet\,/\addthin převracet (co) (zemi ap.)};  \dicIdiom{rótast}{ \dicPhraseIS{rótast}} \dicFlx{refl} {\textbf{a.}} \dicDirectTranslationCS{zakořenit, zapustit kořeny (květina ap.)};  {\textbf{b.}} \dicDirectTranslationCS{zneklidnět, být neklidný}
\dicEntry[rótari] \dicTerm{rót··ar|i} \dicIPA{{r}{ou}{\textlengthmark}{\textsubring{d}}{a}{r}{\textsci}} \dicPos{m}[1] \dicFlx{(‑a, ‑ar)}[13] \dicDirectTranslationCS{bedňák} \dicIndirectTranslationCS{(pomocník hudební kapely)}
\dicEntry[rótarlegur] \dicTerm{rótar··legur} \dicIPA{{r}{ou}{\textlengthmark}{\textsubring{d}}{a}{r}{l}{\textepsilon}{\textbabygamma}{\textscy}{\textsubring{r}}} \dicPos{adj}[1]\dicFlx{}[-8] \dicSynonym{meinlegur} \dicDirectTranslationCS{zlomyslný, zlovolný}
\dicEntry[róteind] \dicTerm{rót··eind} \dicIPA{{r}{ou}{\textlengthmark}{\textsubring{d}}{ei}{n}{\textsubring{d}}} \dicPos{f}[7] \dicFlx{(‑ar, ‑ir)}[1] \dicFieldCat{fyz.} \dicDirectTranslationCS{proton}
\dicEntry[rótgróinn] \dicTerm{rót··gróinn} \dicIPA{{r}{ou}{\textlengthmark}{\textsubring{d}}{\r{g}}{r}{ou}{\textsci}{\textsubring{n}}} \dicPos{adj}[6]\dicFlx{}[-2] \dicDirectTranslationCS{zakořeněný, vžitý, tradiční} \dicExampleIS{rótgróin venja} \dicExampleCS{zakořeněný zvyk}
\dicEntry[rótlaus] \dicTerm{rót··laus} \dicIPA{{r}{ou}{\textlengthmark}{\textsubring{d}}{l}{\oe i}{s}} \dicPos{adj}[5]\dicFlx{}[-1] \dicDirectTranslationCS{(jsoucí) bez kořenů, vykořeněný, bezprizorní}
\dicEntry[róttæklingur] \dicTerm{rót··tæk·ling|ur} \dicIPA{{r}{ou}{\textlengthmark}{t\smash{\textsuperscript{h}}}{a}{i}{\r{g}}{l}{i}{\ng}{\r{g}}{\textscy}{\textsubring{r}}} \dicPos{m}[6] \dicFlx{(‑s, ‑ar)}[8] \dicDirectTranslationCS{radikál(ka), extremista, extremistka}
\dicEntry[róttækni] \dicTerm{rót··tækn|i} \dicIPA{{r}{ou}{\textlengthmark}{t\smash{\textsuperscript{h}}}{a}{i}{h}{\r{g}}{n}{\textsci}} \dicPos{f}[3] \dicFlx{(‑i)}[3] \dicDirectTranslationCS{radikalismus, extremismus}
\dicEntry[róttækur] \dicTerm{rót··tækur} \dicIPA{{r}{ou}{\textlengthmark}{t\smash{\textsuperscript{h}}}{a}{i}{\r{g}}{\textscy}{\textsubring{r}}} \dicPos{adj}[1]\dicFlx{}[-1] \dicDirectTranslationCS{radikální, extremistický}
\dicEntry[rubba] \dicTerm{rubb|a} \dicIPA{{r}{\textscy}{\textsubring{b}}{\textlengthmark}{a}} \dicPos{v}[1] \dicFlx{(‑aði)}[44] \dicFlx{dat} \dicPhraseIS{rubba e‑u af} \dicDirectTranslationCS{odbýt (co), odfláknout (co)}
\dicEntry[ruddalegur] \dicTerm{rudda··legur} \dicIPA{{r}{\textscy}{\textsubring{d}}{\textlengthmark}{a}{l}{\textepsilon}{\textbabygamma}{\textscy}{\textsubring{r}}} \dicPos{adj}[1]\dicFlx{}[-8] \dicDirectTranslationCS{hrubý, sprostý}
\dicEntry[ruddaskapur] \dicTerm{rudda··skap|ur} \dicIPA{{r}{\textscy}{\textsubring{d}}{\textlengthmark}{a}{s}{\r{g}}{a}{\textsubring{b}}{\textscy}{\textsubring{r}}} \dicPos{m}[10] \dicFlx{(‑ar)}[15] \dicDirectTranslationCS{hrubost, hrubiánství, neotesanost}
\dicEntry[ruddi] \dicTerm{rudd|i\smash{\textsuperscript{1}}} \dicIPA{{r}{\textscy}{\textsubring{d}}{\textlengthmark}{\textsci}} \dicPos{m}[1] \dicFlx{(‑a, ‑ar)}[1] \dicSynonym{dóni} \dicDirectTranslationCS{hulvát, hrubián, hrubec}
\dicEntry[ruddi] \dicTerm{ruddi\smash{\textsuperscript{2}}} \dicIPA{{r}{\textscy}{\textsubring{d}}{\textlengthmark}{\textsci}} \dicPos{v} \dicFlx{ind pf sg 1 pers} \dicLink{ryðja}
\dicEntry[ruddum] \dicTerm{ruddum} \dicIPA{{r}{\textscy}{\textsubring{d}}{\textlengthmark}{\textscy}{\textsubring{m}}} \dicPos{v} \dicFlx{ind pf pl 1 pers} \dicLink{ryðja}
\dicEntry[ruðningur] \dicTerm{ruð··ning|ur} \dicIPA{{r}{\textscy}{ð}{n}{i}{\ng}{\r{g}}{\textscy}{\textsubring{r}}} \dicPos{m}[6] \dicFlx{(‑s, ‑ar)}[8] \textbf{1.} \dicSynonym*{það að ryðja} \dicDirectTranslationCS{klučení, mýcení (lesa ap.)}  \textbf{2.} \dicFieldCat{sport.} \dicDirectTranslationCS{americký fotbal}  \textbf{3.} \dicFieldCat{sport.} \dicSynonym{rúgbý} \dicDirectTranslationCS{ragby}  \textbf{4.} \dicFieldCat{sport.} \dicDirectTranslationCS{prorážení (v~házené ap.)} \dicIndirectTranslationCS{(nepovolená hra)}
\dicEntry[ruðum] \dicTerm{ruðum} \dicIPA{{r}{\textscy}{\textlengthmark}{ð}{\textscy}{\textsubring{m}}} \dicPos{v} \dicFlx{ind pf pl 1 pers} \dicLink{rjóða}
\dicEntry[rufum] \dicTerm{rufum} \dicIPA{{r}{\textscy}{\textlengthmark}{v}{\textscy}{\textsubring{m}}} \dicPos{v} \dicFlx{ind pf pl 1 pers} \dicLink{rjúfa}
\dicEntry[rugg] \dicTerm{rugg} \dicIPA{{r}{\textscy}{\r{g}}{\textlengthmark}} \dicPos{n}[2] \dicFlx{(‑s)}[2] \textbf{1.} \dicSynonym*{það að rugga} \dicDirectTranslationCS{kolébání, houpání (kolébky ap.)}  \textbf{2.} \dicSynonym*{veltingur} \dicDirectTranslationCS{kolébání, houpání (lodi ap.)}
\dicEntry[rugga] \dicTerm{rugg|a} \dicIPA{{r}{\textscy}{\r{g}}{\textlengthmark}{a}} \dicPos{v}[1] \dicFlx{(‑aði)}[1] \dicFlx{dat} \textbf{1.} \dicFlx{dat} \dicSynonym{vagga\smash{\textsuperscript{2}}} \dicDirectTranslationCS{houpat, kolébat} \dicExampleIS{rugga barni} \dicExampleCS{kolébat dítě};  \dicPhraseIS{rugga sér} \dicDirectTranslationCS{houpat se} \dicExampleIS{rugga sér í stólnum} \dicExampleCS{houpat se na židli}  \textbf{2.} \dicDirectTranslationCS{houpat se, kolébat se (o~lodi ap.)}
\dicEntry[rugguhestur] \dicTerm{ruggu··hest|ur} \dicIPA{{r}{\textscy}{\r{g}}{\textlengthmark}{\textscy}{h}{\textepsilon}{s}{\textsubring{d}}{\textscy}{\textsubring{r}}} \dicPos{m}[6] \dicFlx{(‑s, ‑ar)}[4] \dicDirectTranslationCS{houpací kůň}
\dicEntry[ruggustóll] \dicTerm{ruggu··stól|l} \dicIPA{{r}{\textscy}{\r{g}}{\textlengthmark}{\textscy}{s}{\textsubring{d}}{ou}{\textsubring{d}}{\textsubring{l}}} \dicPos{m}[6] \dicFlx{(‑s, ‑ar)}[48] \dicDirectTranslationCS{houpací křeslo\,/\addthin židle}
\dicEntry[rugl] \dicTerm{rugl} \dicIPA{{r}{\textscy}{\r{g}}{\textsubring{l}}} \dicPos{n}[2] \dicFlx{(‑s)}[2] \textbf{1.} \dicSynonym{þvættingur} \dicDirectTranslationCS{nesmysl, blbost}  \textbf{2.} \dicSynonym{ringulreið} \dicDirectTranslationCS{zmatek, chaos}
\dicEntry[rugla] \dicTerm{rugl|a} \dicIPA{{r}{\textscy}{\r{g}}{l}{a}} \dicPos{v}[1] \dicFlx{(‑aði)}[1] \dicFlx{acc\,/\addthin dat} \textbf{1.} \dicFlx{dat} \dicDirectTranslationCS{splést si, poplést si, zaměnit};  \dicPhraseIS{rugla e‑u saman við e‑ð} \dicDirectTranslationCS{poplést si (co) s~(čím)};  \dicPhraseIS{rugla e‑u\,/\addthin e‑m saman} \dicDirectTranslationCS{splést si (co\,/\addthin koho)}  \textbf{2.} \dicFlx{dat} \dicSynonym{hræra\smash{\textsuperscript{2}}} \dicDirectTranslationCS{pomíchat, zpřeházet (boty ap.)}  \textbf{3.} \dicFlx{acc} \dicSynonym{trufla} \dicDirectTranslationCS{zmást, poplést, splést};  \dicPhraseIS{rugla e‑n með e‑u} \dicDirectTranslationCS{poplést (koho čím)};  \dicPhraseIS{rugla e‑n í ríminu} \dicDirectTranslationCS{poplést (koho), uvést (koho) v~omyl}  \textbf{4.} \dicDirectTranslationCS{plácat nesmysly, plést páté přes deváté};  \dicIdiom{ruglast}{ \dicPhraseIS{ruglast}} \dicFlx{refl} {\textbf{a.}} \dicSynonym*{verða ringlaður} \dicDirectTranslationCS{poplést se, splést se};  {\textbf{b.}} \dicSynonym*{komast í ólag} \dicDirectTranslationCS{pomíchat se, zpřeházet se};  \dicIdiom{ruglast}[á]{ \dicPhraseIS{ruglast á e‑u}} \dicFlx{refl} \dicDirectTranslationCS{splést si (co)};  \dicIdiom{ruglast}[í]{ \dicPhraseIS{ruglast í ríminu}} \dicFlx{refl} \dicDirectTranslationCS{poplést se, splést se}
\dicEntry[ruglaður] \dicTerm{rugl··|aður} \dicIPA{{r}{\textscy}{\r{g}}{l}{a}{ð}{\textscy}{\textsubring{r}}} \dicPos{adj}[3] \dicFlx{(f ‑uð)}[3] \dicSynonym{ringlaður} \dicDirectTranslationCS{zmatený, popletený}
\dicEntry[ruglingur] \dicTerm{rugl··ing|ur} \dicIPA{{r}{\textscy}{\r{g}}{l}{i}{\ng}{\r{g}}{\textscy}{\textsubring{r}}} \dicPos{m}[6] \dicFlx{(‑s)}[9] \dicSynonym{óreiða} \dicDirectTranslationCS{zmatek, nepořádek, chaos} \dicExampleIS{vera á ruglingi} \dicExampleCS{být v~nepořádku}
\dicEntry[rugludallur] \dicTerm{ruglu··dall|ur} \dicIPA{{r}{\textscy}{\r{g}}{l}{\textscy}{\textsubring{d}}{a}{\textsubring{d}}{l}{\textscy}{\textsubring{r}}} \dicPos{m}[6] \dicFlx{(‑s, ‑ar)}[5] \dicDirectTranslationCS{pomatenec, ztřeštěnec, ztřeštěnkyně, ztřeštěnka}
\dicEntry[rukka] \dicTerm{rukk|a} \dicIPA{{r}{\textscy}{h}{\r{g}}{a}} \dicPos{v}[1] \dicFlx{(‑aði)}[1] \dicFlx{acc} \dicSynonym{innheimta\smash{\textsuperscript{2}}} \dicDirectTranslationCS{inkasovat, kasírovat, vybírat peníze};  \dicPhraseIS{rukka fyrir e‑ð} \dicDirectTranslationCS{vybírat peníze za (co) (časopis ap.)};  \dicPhraseIS{rukka e‑n} \dicDirectTranslationCS{inkasovat od (koho)};  \dicPhraseIS{rukka e‑n um e‑ð} \dicDirectTranslationCS{inkasovat od (koho co)}
\dicEntry[rukkari] \dicTerm{rukk··ar|i} \dicIPA{{r}{\textscy}{h}{\r{g}}{a}{r}{\textsci}} \dicPos{m}[1] \dicFlx{(‑a, ‑ar)}[13] \dicLangCat{hovor.} \dicDirectTranslationCS{výběrčí (peněz), inkasista, inkasistka}
\dicEntry[rukkun] \dicTerm{rukk|un} \dicIPA{{r}{\textscy}{h}{\r{g}}{\textscy}{\textsubring{n}}} \dicPos{f}[7] \dicFlx{(‑unar, ‑anir)}[8] \dicDirectTranslationCS{inkasování, kasírování, vybírání peněz}
\dicEntry[rukum] \dicTerm{rukum} \dicIPA{{r}{\textscy}{\textlengthmark}{\r{g}}{\textscy}{\textsubring{m}}} \dicPos{v} \dicFlx{ind pf pl 1 pers} \dicLink{rjúka}
\dicEntry[rulla] \dicTerm{rull|a} \dicIPA{{r}{\textscy}{\textsubring{d}}{l}{a}} \dicPos{f}[1] \dicFlx{(‑u, ‑ur)}[19] \textbf{1.} \dicSynonym*{taurulla} \dicDirectTranslationCS{mandl}  \textbf{2.} \dicSynonym{hlutverk} \dicDirectTranslationCS{(herecká) role}
\dicEntry[rumdi] \dicTerm{rumdi} \dicIPA{{r}{\textscy}{m}{\textsubring{d}}{\textsci}} \dicPos{v} \dicFlx{ind pf sg 1 pers} \dicLink{rymja}
\dicEntry[rumdum] \dicTerm{rumdum} \dicIPA{{r}{\textscy}{m}{\textsubring{d}}{\textscy}{\textsubring{m}}} \dicPos{v} \dicFlx{ind pf pl 1 pers} \dicLink{rymja}
\dicEntry[rumið] \dicTerm{rumið} \dicIPA{{r}{\textscy}{\textlengthmark}{m}{\textsci}{\texttheta}} \dicPos{v} \dicFlx{supin} \dicLink{rymja}
\dicEntry[rumska] \dicTerm{rumsk|a} \dicIPA{{r}{\textscy}{m}{s}{\r{g}}{a}} \dicPos{v}[1] \dicFlx{(‑aði)}[1] \dicDirectTranslationCS{probudit se, probouzet se, probrat se, probírat se} \dicExampleIS{Hann er farinn að rumska.} \dicExampleCS{Začíná se probouzet.}
\dicEntry[runa] \textls[15]{\dicTerm{run|a} \dicIPA{{r}{\textscy}{\textlengthmark}{n}{a}} \dicPos{f}[1] \dicFlx{(‑u, ‑ur)}[7] \textbf{1.} \dicSynonym{röð} \dicDirectTranslationCS{řada, posloupnost} \dicExampleIS{löng runa af nöfnum} \dicExampleCS{dlouhá řada jmen}  \textbf{2.} \dicFieldCat{mat.} \mbox{\dicDirectTranslationCS{posloupnost}}}
\dicEntry[runka] \dicTerm{runk|a} \dicIPA{{r}{u}{\r{\ng}}{\r{g}}{a}} \dicPos{v}[1] \dicFlx{(‑aði)}[1] \dicPhraseIS{runka sér, runka sig} \dicLangCat{hrub.} \dicDirectTranslationCS{masturbovat, onanovat}
\dicEntry[runni] \dicTerm{runn|i} \dicIPA{{r}{\textscy}{n}{\textlengthmark}{\textsci}} \dicPos{m}[1] \dicFlx{(‑a, ‑ar)}[1] \dicDirectTranslationCS{keř, křoví}
\dicEntry[runnið] \dicTerm{runnið} \dicIPA{{r}{\textscy}{n}{\textlengthmark}{\textsci}{\texttheta}} \dicPos{v} \dicFlx{supin} \dicLink{renna\smash{\textsuperscript{3}}}
\dicEntry[runnum] \dicTerm{runnum} \dicIPA{{r}{\textscy}{n}{\textlengthmark}{\textscy}{\textsubring{m}}} \dicPos{v} \dicFlx{ind pf pl 1 pers} \dicLink{renna\smash{\textsuperscript{3}}}
\dicEntry[runurím] \dicTerm{runu··rím} \dicIPA{{r}{\textscy}{\textlengthmark}{n}{\textscy}{r}{i}{\textsubring{m}}} \dicPos{n}[2] \dicFlx{(‑s, ‑)}[5] \dicFieldCat{lit.} \dicDirectTranslationCS{sdružený rým}
\dicEntry[rusl] \dicTerm{rusl} \dicsymFrequent\  \dicIPA{{r}{\textscy}{s}{\textsubring{d}}{\textsubring{l}}} \dicPos{n}[2] \dicFlx{(‑s)}[2] \textbf{1.} \dicSynonym{úrgangur} \dicDirectTranslationCS{odpadky, smetí} \dicExampleIS{fleygja rusli} \dicExampleCS{vyhodit odpadky};  \dicPhraseIS{fara út með ruslið} \dicDirectTranslationCS{jít vynést odpadky, jít s~odpadky};  \dicPhraseIS{henda e‑u í ruslið} \dicDirectTranslationCS{vyhodit (co) (do smetí)};  \dicPhraseIS{setja e‑ð í ruslið} \dicDirectTranslationCS{vyhodit (co), dát (co) do koše}  \textbf{2.} \dicSynonym{drasl} \dicDirectTranslationCS{smetí, šunt, brak} \dicIndirectTranslationCS{(něco bezcenného)};  \dicIdiom{rusl}{ \dicPhraseIS{vera í rusli}} \dicLangCat{hovor.} \dicDirectTranslationCS{být mimo, být zničený}
\dicEntry[rusla] \dicTerm{rusl|a} \dicIPA{{r}{\textscy}{s}{\textsubring{d}}{l}{a}} \dicPos{v}[1] \dicFlx{(‑aði)}[1] \dicIdiom{rusla}[af]{ \dicPhraseIS{rusla e‑u af}} \dicDirectTranslationCS{odbýt (co), odfláknout (co)};  \dicIdiom{rusla}[til\,/\addthin út]{ \dicPhraseIS{rusla til, rusla út}} \dicDirectTranslationCS{udělat nepořádek, zaneřádit} \dicExampleIS{rusla til í herberginu sínu} \dicExampleCS{zaneřádit si pokoj}
\dicEntry[ruslafata] \dicTerm{rusla··|fata} \dicIPA{{r}{\textscy}{s}{\textsubring{d}}{l}{a}{f}{a}{\textsubring{d}}{a}} \dicPos{f}[1] \dicFlx{(‑fötu, ‑fötur)}[20] \dicDirectTranslationCS{(odpadkový) koš}
\dicEntry[ruslagámur] \dicTerm{rusla··gám|ur} \dicIPA{{r}{\textscy}{s}{\textsubring{d}}{l}{a}{\r{g}}{au}{m}{\textscy}{\textsubring{r}}} \dicPos{m}[6] \dicFlx{(‑s, ‑ar)}[10] \dicDirectTranslationCS{odpadní kontejner}
\dicEntry[ruslakarfa] \dicTerm{rusla··|karfa} \dicIPA{{r}{\textscy}{s}{\textsubring{d}}{l}{a}{k\smash{\textsuperscript{h}}}{a}{r}{v}{a}} \dicPos{f}[1] \dicFlx{(‑körfu, ‑körfur)}[20] \dicDirectTranslationCS{(odpadkový) koš (na papír ap.)}
\dicEntry[ruslakompa] \dicTerm{rusla··komp|a} \dicIPA{{r}{\textscy}{s}{\textsubring{d}}{l}{a}{k\smash{\textsuperscript{h}}}{\textopeno}{\textsubring{m}}{\textsubring{b}}{a}} \dicPos{f}[1] \dicFlx{(‑u, ‑ur)}[19] \dicDirectTranslationCS{komora, sklípek}
\dicEntry[ruslapoki] \dicTerm{rusla··pok|i} \dicIPA{{r}{\textscy}{s}{\textsubring{d}}{l}{a}{p\smash{\textsuperscript{h}}}{\textopeno}{\r{\textObardotlessj}}{\textsci}} \dicPos{m}[1] \dicFlx{(‑a, ‑ar)}[1] \dicDirectTranslationCS{pytel na odpadky}
\dicEntry[ruslatunna] \dicTerm{rusla··tunn|a} \dicIPA{{r}{\textscy}{s}{\textsubring{d}}{l}{a}{t\smash{\textsuperscript{h}}}{\textscy}{n}{a}} \dicPos{f}[1] \dicFlx{(‑u, ‑ur)}[7] \dicSynonym{sorptunna} \dicDirectTranslationCS{popelnice}
\dicEntry[ruslpóstur] \dicTerm{rusl··póst|ur} \dicIPA{{r}{\textscy}{s}{\textsubring{d}}{\textsubring{l}}{p\smash{\textsuperscript{h}}}{ou}{s}{\textsubring{d}}{\textscy}{\textsubring{r}}} \dicPos{m}[6] \dicFlx{(‑s)}[7] \dicFieldCat{poč.} \dicDirectTranslationCS{spam, nevyžádaná pošta}
\dicEntry[rutl] \dicTerm{rutl} \dicIPA{{r}{\textscy}{h}{\textsubring{d}}{\textsubring{l}}} \dicPos{n}[2] \dicFlx{(‑s)}[2] \dicSynonym{ruglingur} \dicDirectTranslationCS{pomatení, vyšinutí}
\dicEntry[rutt] \dicTerm{rutt} \dicIPA{{r}{\textscy}{h}{\textsubring{d}}} \dicPos{v} \dicFlx{supin} \dicLink{ryðja}
\dicEntry[rú] \dicTerm{rú} \dicIPA{{r}{u}{\textlengthmark}} \dicPos{n}[2] \dicFlx{(‑s)}[2] \dicPhraseIS{það er allt á rúi og stúi} \dicLangCat{přen.} \dicDirectTranslationCS{všechno je vzhůru nohama}
\dicEntry[Rúanda] \dicTerm{Rúanda} \dicIPA{{r}{u}{\textlengthmark}{a}{n}{\textsubring{d}}{a}} \dicPos{n}[4] \dicFlx{indecl}[2] \dicFieldCat{geog.} \dicDirectTranslationCS{Rwanda}
\dicEntry[Rúandamaður] \textls[15]{\dicTerm{Rúanda··|maður} \dicIPA{{r}{u}{\textlengthmark}{a}{n}{\textsubring{d}}{a}{m}{a}{ð}{\textscy}{\textsubring{r}}} \dicPos{m}[13] \dicFlx{(‑manns, ‑menn)}[2]} \dicDirectTranslationCS{Rwanďan(ka)}
\dicEntry[rúandíska] \dicTerm{rúandísk|a} \dicIPA{{r}{u}{\textlengthmark}{a}{n}{\textsubring{d}}{i}{s}{\r{g}}{a}} \dicPos{f}[1] \dicFlx{(‑u)}[5] \dicDirectTranslationCS{rwandština}
\dicEntry[rúandskur] \dicTerm{rúandskur}\dicTerm{, rúanskur} \dicIPA{{r}\-{u}\-{\textlengthmark}\-{a}\-{n}\-{\textsubring{d}}\-{s}\-{\r{g}}\-{\textscy}\-{\textsubring{r}}\-} \dicPos{adj}[1] \dicFlx{(f rúöndsk)}[3] \dicDirectTranslationCS{rwandský}
\dicEntry[rúanskur] \dicTerm{rúanskur} \dicIPA{{r}{u}{\textlengthmark}{a}{n}{s}{\r{g}}{\textscy}{\textsubring{r}}} \dicPos{adj}[1] \dicFlx{(f rúönsk)}[3] \dicLink{rúandskur}
\dicEntry[rúbidín] \dicTerm{rúbidín} \dicIPA{{r}{u}{\textlengthmark}{\textsubring{b}}{\textsci}{\textsubring{d}}{i}{\textsubring{n}}} \dicPos{n}[2] \dicFlx{(‑s)}[2] \dicFieldCat{chem.} \dicDirectTranslationCS{rubidium} \textit{(l.~{\textLA{Rb, Rubidium}})}
\dicEntry[rúbín] \dicTerm{rúbín} \dicIPA{{r}{u}{\textlengthmark}{\textsubring{b}}{i}{\textsubring{n}}} \dicPos{m}[4] \dicFlx{(‑s, ‑ar)}[10] \dicFieldCat{geol.} \dicDirectTranslationCS{rubín}
\dicEntry[rúð] \dicTerm{rúð} \dicIPA{{r}{u}{\textlengthmark}{\texttheta}} \dicPos{v} \dicFlx{supin} \dicLink{rýja\smash{\textsuperscript{2}}}
\dicEntry[rúða] \dicTerm{rúð|a\smash{\textsuperscript{1}}} \dicsymFrequent\  \dicIPA{{r}{u}{\textlengthmark}{ð}{a}} \dicPos{f}[1] \dicFlx{(‑u, ‑ur)}[19] \dicDirectTranslationCS{okenní tabulka} \dicExampleIS{brjóta rúðu í bifreiðinni} \dicExampleCS{rozbít okenní tabulku v~autobuse}
\dicEntry[Rúða] \dicTerm{Rúð|a\smash{\textsuperscript{2}}}\dicTerm{, Rúðuborg} \dicIPA{{r}\-{u}\-{\textlengthmark}\-{ð}\-{a}\-} \dicPos{f}[1] \dicFlx{(‑u)}[6] \dicFieldCat{geog.} \dicDirectTranslationCS{Rouen} \dicIndirectTranslationCS{(město ve Francii)}
\dicEntry[rúði] \dicTerm{rúði} \dicIPA{{r}{u}{\textlengthmark}{ð}{\textsci}} \dicPos{v} \dicFlx{ind\,/\addthin con pf sg 1 pers} \dicLink{rýja\smash{\textsuperscript{2}}}
\dicEntry[Rúðuborg] \dicTerm{Rúðu··borg} \dicIPA{{r}{u}{\textlengthmark}{ð}{\textscy}{\textsubring{b}}{\textopeno}{r}{\r{g}}} \dicPos{f}[7] \dicFlx{(‑ar)}[4] \dicLink{Rúða\smash{\textsuperscript{2}}}
\dicEntry[rúðugler] \dicTerm{rúðu··gler} \dicIPA{{r}{u}{\textlengthmark}{ð}{\textscy}{\r{g}}{l}{\textepsilon}{\textsubring{r}}} \dicPos{n}[2] \dicFlx{(‑s, ‑)}[6] \dicDirectTranslationCS{okenní sklo}
\dicEntry[rúðum] \dicTerm{rúðum} \dicIPA{{r}{u}{\textlengthmark}{ð}{\textscy}{\textsubring{m}}} \dicPos{v} \dicFlx{ind pf pl 1 pers} \dicLink{rýja\smash{\textsuperscript{2}}}
\dicEntry[rúðuþurrka] \dicTerm{rúðu··þurrk|a} \dicIPA{{r}{u}{\textlengthmark}{ð}{\textscy}{\texttheta}{\textscy}{\textsubring{r}}{\r{g}}{a}} \dicPos{f}[1] \dicFlx{(‑u, ‑ur)}[13] \dicDirectTranslationCS{stěrač, stírač} \dicIndirectTranslationCS{(zvláště v~pl)}
\dicEntry[rúgbrauð] \dicTerm{rúg··brauð} \dicIPA{{r}{u}{\textlengthmark}{\textsubring{b}}{r}{\oe i}{\texttheta}} \dicPos{n}[2] \dicFlx{(‑s, ‑)}[5] \dicDirectTranslationCS{žitný\,/\addthin černý chléb}
\dicEntry[rúgbý] \dicTerm{rúgbý} \dicIPA{{r}{u}{\textsubring{b}}{i}} \dicPos{n}[2] \dicFlx{(‑s, ‑)}[13] \dicFieldCat{sport.} \dicDirectTranslationCS{ragby}
\dicEntry[rúgmjöl] \dicTerm{rúg··mjöl} \dicIPA{{r}{u}{\textlengthmark}{m}{j}{\oe}{\textsubring{l}}} \dicPos{n}[2] \dicFlx{(‑s)}[2] \dicDirectTranslationCS{žitná mouka}
\dicEntry[rúgur] \dicTerm{rúg|ur} \dicIPA{{r}{u}{\textlengthmark}{\textscy}{\textsubring{r}}} \dicPos{m}[6] \dicFlx{(‑s)}[26] \dicFieldCat{bot.} \dicDirectTranslationCS{žito} \textit{(l.~{\textLA{Secale cereale}})}  \dicsymPhoto\ 
\dicFigure{ds_image_rugur_0_2.jpg}{Rúgur}{Rúgur - LSDSL, GFDL}
\dicEntry[rúið] \dicTerm{rúið} \dicIPA{{r}{u}{\textlengthmark}{\textsci}{\texttheta}} \dicPos{v} \dicFlx{supin} \dicLink{rýja\smash{\textsuperscript{2}}}
\dicEntry[rúkragi] \dicTerm{rú··krag|i} \dicIPA{{r}{u}{\textlengthmark}{k\smash{\textsuperscript{h}}}{r}{a}{i}{j}{\textsci}} \dicPos{m}[1] \dicFlx{(‑a, ‑ar)}[8] \dicFieldCat{zool.} \dicDirectTranslationCS{jespák bojovný} \textit{(l.~{\textLA{Philomachus pugnax}})}  \dicsymPhoto\ 
\dicFigure{ds_image_rukragi_0_1.jpg}{Rúkragi}{Rúkragi - Arjan Haverkamp, CC BY 2.0}
\dicEntry[rúlla] \dicTerm{rúll|a\smash{\textsuperscript{1}}} \dicIPA{{r}{u}{l}{\textlengthmark}{a}} \dicPos{f}[1] \dicFlx{(‑u, ‑ur)}[19] \textbf{1.} \dicDirectTranslationCS{role, rolička, rulička (papíru ap.)}  \textbf{2.} \dicSynonym*{hárrúlla} \dicDirectTranslationCS{natáčka}  \textbf{3.} \dicSynonym*{málningarrúlla} \dicDirectTranslationCS{(malířský) váleček}
\dicEntry[rúlla] \dicTerm{rúll|a\smash{\textsuperscript{2}}} \dicIPA{{r}{u}{l}{\textlengthmark}{a}} \dicPos{v}[1] \dicFlx{(‑aði)}[1] \dicFlx{dat} \textbf{1.} \dicSynonym{vinda\smash{\textsuperscript{2}}} \dicDirectTranslationCS{stočit, srolovat, svinout} \dicExampleIS{rúlla pappírnum upp} \dicExampleCS{svinout papír}  \textbf{2.} \dicSynonym{velta\smash{\textsuperscript{3}}} \dicDirectTranslationCS{kutálet se, koulet se, valit se} \dicExampleIS{Boltinn rúllaði niður brekkuna.} \dicExampleCS{Míč se kutálel dolů z~kopce.};  \dicIdiom{rúlla}[upp]{ \dicPhraseIS{rúlla upp e‑m}} \dicLangCat{slang.} \dicDirectTranslationCS{převálcovat (koho), roznést (koho) na kopytech}
\dicEntry[rúllugardína] \dicTerm{rúllu··gardín|a} \dicIPA{{r}{u}{l}{\textlengthmark}{\textscy}{\r{g}}{a}{r}{\textsubring{d}}{i}{n}{a}} \dicPos{f}[1] \dicFlx{(‑u, ‑ur)}[7] \dicDirectTranslationCS{roleta}
\dicEntry[rúllukragapeysa] \dicTerm{rúllu·kraga··peys|a} \dicIPA{{r}{u}{l}{\textlengthmark}{\textscy}{k\smash{\textsuperscript{h}}}{r}{a}{\textbabygamma}{a}{p\smash{\textsuperscript{h}}}{ei}{s}{a}} \dicPos{f}[1] \dicFlx{(‑u, ‑ur)}[7] \dicDirectTranslationCS{rolák}
\dicEntry[rúllustigi] \dicTerm{rúllu··stig|i} \dicIPA{{r}{u}{l}{\textlengthmark}{\textscy}{s}{\textsubring{d}}{i}{j}{\textsci}} \dicPos{m}[1] \dicFlx{(‑a, ‑ar)}[1] \dicDirectTranslationCS{jezdicí schody}
\dicEntry[rúlluterta] \dicTerm{rúllu··tert|a} \dicIPA{{r}{u}{l}{\textlengthmark}{\textscy}{t\smash{\textsuperscript{h}}}{\textepsilon}{\textsubring{r}}{\textsubring{d}}{a}} \dicPos{f}[1] \dicFlx{(‑u, ‑ur)}[19] \dicFieldCat{kulin.} \dicDirectTranslationCS{roláda}
\dicEntry[rúm] \dicTerm{rúm} \dicsymFrequent\  \dicIPA{{r}{u}{\textlengthmark}{\textsubring{m}}} \dicPos{n}[2] \dicFlx{(‑s, ‑)}[5] \textbf{1.} \dicSynonym{rekkja} \dicDirectTranslationCS{postel, lůžko};  \dicPhraseIS{búa um rúm} \dicDirectTranslationCS{ustlat, stlát};  \dicPhraseIS{fara í rúmið} \dicDirectTranslationCS{jít si lehnout, jít do postele};  \dicPhraseIS{liggja í rúminu} {\textbf{a.}} \dicDirectTranslationCS{ležet v~posteli};  {\textbf{b.}} \dicDirectTranslationCS{být nemocný (a~ležet v~posteli)}  \textbf{2.} \dicSynonym{rými} \dicDirectTranslationCS{prostor, místo (pro změnu ap.)}  \textbf{3.} \dicSynonym{vídd} \dicDirectTranslationCS{prostor} \dicIndirectTranslationCS{(rozměrově nevymezené prostředí)};  \dicIdiom{rúm}{ \dicPhraseIS{láta sér e‑ð í léttu rúmi liggja}} \dicDirectTranslationCS{nezáležet na (čem), nepřikládat (čemu) váhu}; { \dicPhraseIS{ryðja sér til rúms}} \dicSynonym*{breiðast út} \dicDirectTranslationCS{rozšířit se, rozšiřovat se}
\dicEntry[rúma] \dicTerm{rúm|a} \dicIPA{{r}{u}{\textlengthmark}{m}{a}} \dicPos{v}[1] \dicFlx{(‑aði)}[1] \dicFlx{acc} \dicDirectTranslationCS{pojmout, obsáhnout, mít místo\,/\addthin kapacitu} \dicExampleIS{Taskan rúmaði bara fötin.} \dicExampleCS{Taška pojala jen oblečení.}
\dicEntry[rúmba] \dicTerm{rúmb|a} \dicIPA{{r}{u}{m}{\textsubring{b}}{a}} \dicPos{f}[1] \dicFlx{(‑u, ‑ur)}[13] \dicFieldCat{hud.} \dicDirectTranslationCS{rumba}
\dicEntry[Rúmeni] \dicTerm{Rúmen|i} \dicIPA{{r}{u}{\textlengthmark}{m}{\textepsilon}{n}{\textsci}} \dicPos{m}[1] \dicFlx{(‑a, ‑ar)}[1] \dicDirectTranslationCS{Rumun(ka)}
\dicEntry[Rúmenía] \dicTerm{Rúmení|a} \dicIPA{{r}{u}{\textlengthmark}{m}{\textepsilon}{n}{i}{j}{a}} \dicPos{f}[1] \dicFlx{(‑u)}[6] \dicFieldCat{geog.} \dicDirectTranslationCS{Rumunsko}
\dicEntry[rúmenska] \dicTerm{rúmensk|a} \dicIPA{{r}{u}{\textlengthmark}{m}{\textepsilon}{n}{s}{\r{g}}{a}} \dicPos{f}[1] \dicFlx{(‑u)}[5] \dicDirectTranslationCS{rumunština}
\dicEntry[rúmenskur] \dicTerm{rúmenskur} \dicIPA{{r}{u}{\textlengthmark}{m}{\textepsilon}{n}{s}{\r{g}}{\textscy}{\textsubring{r}}} \dicPos{adj}[1]\dicFlx{}[-6] \dicDirectTranslationCS{rumunský}
\dicEntry[rúmfastur] \dicTerm{rúm··|fastur} \dicIPA{{r}{u}{m}{f}{a}{s}{\textsubring{d}}{\textscy}{\textsubring{r}}} \dicPos{adj}[1] \dicFlx{(f ‑föst)}[12] \dicDirectTranslationCS{upoutaný na lůžko}
\dicEntry[rúmfræði] \dicTerm{rúm··fræð|i} \dicIPA{{r}{u}{m}{f}{r}{a}{i}{ð}{\textsci}} \dicPos{f}[3] \dicFlx{(‑i)}[3] \dicFieldCat{mat.} \dicDirectTranslationCS{geometrie}
\dicEntry[rúmföt] \dicTerm{rúm··föt} \dicIPA{{r}{u}{m}{f}{\oe}{\textsubring{d}}} \dicPos{n}[2] \dicFlx{pl}[9] \dicSynonym{sængurfatnaður} \dicDirectTranslationCS{ložní prád\-lo, lůžkoviny}
\dicEntry[rúmgóður] \dicTerm{rúm··|góður} \dicIPA{{r}{u}{m}{\r{g}}{ou}{ð}{\textscy}{\textsubring{r}}} \dicPos{adj}[11] \dicFlx{(comp ‑betri, sup ‑bestur)}[1] \dicDirectTranslationCS{prostorný, prostranný} \dicExampleIS{rúmgóður salur} \dicExampleCS{prostorný sál}
\dicEntry[rúml.] \dicTerm{rúml.} \dicPos{zkr} \dicPhraseIS{rúmlega} \dicFlx{adv} \dicDirectTranslationCS{přes, více než}
\dicEntry[rúmlega] \dicTerm{rúm··leg|a\smash{\textsuperscript{1}}} \dicIPA{{r}{u}{m}{l}{\textepsilon}{\textbabygamma}{a}} \dicPos{f}[1] \dicFlx{(‑u, ‑ur)}[19] \dicDirectTranslationCS{pobyt na lůžku}
\dicEntry[rúmlega] \dicTerm{rúm··lega\smash{\textsuperscript{2}}} \dicsymFrequent\  \dicIPA{{r}{u}{m}{l}{\textepsilon}{\textbabygamma}{a}} \dicPos{adv} \dicSynonym*{rúmt} \dicDirectTranslationCS{přes, více než} \dicExampleIS{Hann er rúmlega fimmtugur.} \dicExampleCS{Je mu přes padesát let.}
\dicEntry[rúmlest] \dicTerm{rúm··lest} \dicIPA{{r}{u}{m}{l}{\textepsilon}{s}{\textsubring{d}}} \dicPos{f}[7] \dicFlx{(‑ar, ‑ir)}[1] \dicFieldCat{nám.} \dicDirectTranslationCS{rejstříková tuna}
\dicEntry[rúmliggjandi] \dicTerm{rúm··liggj·andi} \dicIPA{{r}{u}{m}{l}{\textsci}{\r{\textObardotlessj}}{a}{n}{\textsubring{d}}{\textsci}} \dicPos{adj}[13] \dicFlx{indecl}[1] \dicSynonym{rúmfastur} \dicDirectTranslationCS{upoutaný na lůžko (nemocí ap.)}
\dicEntry[rúmmál] \dicTerm{rúm··mál} \dicIPA{{r}{u}{m}{\textlengthmark}{au}{\textsubring{l}}} \dicPos{n}[2] \dicFlx{(‑s)}[2] \dicDirectTranslationCS{objem, kubatura}
\dicEntry[rúmmetri] \dicTerm{rúm··metr|i} \dicIPA{{r}{u}{m}{\textlengthmark}{\textepsilon}{\textsubring{d}}{r}{\textsci}} \dicPos{m}[1] \dicFlx{(‑a, ‑ar)}[1] \dicDirectTranslationCS{krychlový metr}
\dicEntry[rúmmynd] \dicTerm{rúm··mynd} \dicIPA{{r}{u}{m}{\textlengthmark}{\textsci}{n}{\textsubring{d}}} \dicPos{f}[7] \dicFlx{(‑ar, ‑ir)}[1] \dicFieldCat{mat.} \dicDirectTranslationCS{(prostorové\,/\addthin trojrozměrné) těleso}
\dicEntry[rúmstokkur] \dicTerm{rúm··stokk|ur} \dicIPA{{r}{u}{m}{s}{\textsubring{d}}{\textopeno}{h}{\r{g}}{\textscy}{\textsubring{r}}} \dicPos{m}[6] \dicFlx{(‑s, ‑ar)}[22] \dicDirectTranslationCS{strana postele}
\dicEntry[rúmur] \dicTerm{rúmur} \dicsymFrequent\  \dicIPA{{r}{u}{\textlengthmark}{m}{\textscy}{\textsubring{r}}} \dicPos{adj}[10] \dicFlx{(comp rýmri, sup rýmstur)}[13] \textbf{1.} \dicSynonym{rúmgóður} \dicDirectTranslationCS{prostorný, rozměrný};  \dicPhraseIS{það er rúmt um e‑n e‑s staðar} \dicDirectTranslationCS{(kdo) má (kde) dostatek místa (v~bytě ap.)}  \textbf{2.} \dicSynonym{ríflegur} \dicDirectTranslationCS{(jsoucí) více než\,/\addthin jak} \dicIndirectTranslationCS{(o~čase, velikosti ap.)} \dicExampleIS{rúm vika} \dicExampleCS{více jak týden}
\dicEntry[rún] \dicTerm{rún} \dicIPA{{r}{u}{\textlengthmark}{\textsubring{n}}} \dicPos{f}[7] \dicFlx{(‑ar, ‑ir)}[1] \dicSynonym*{rúnastafur} \dicDirectTranslationCS{runa}
\dicEntry[rúning] \dicTerm{rún··ing} \dicIPA{{r}{u}{\textlengthmark}{n}{i}{\ng}{\r{g}}} \dicPos{f}[4] \dicFlx{(‑ar)}[7] \dicLink{rúningur}
\dicEntry[rúningur] \dicTerm{rún··ing|ur}\dicTerm{, rúning} \dicIPA{{r}\-{u}\-{\textlengthmark}\-{n}\-{i}\-{\ng}\-{\r{g}}\-{\textscy}\-{\textsubring{r}}\-} \dicPos{m}[6] \dicFlx{(‑s)}[9] \dicDirectTranslationCS{stříhání ovcí}
\dicEntry[rúnnstykki] \dicTerm{rúnn··stykki} \dicIPA{{r}{u}{\textlengthmark}{\textsubring{n}}{s}{\textsubring{d}}{\textsci}{h}{\r{\textObardotlessj}}{\textsci}} \dicPos{n}[2] \dicFlx{(‑s, ‑)}[16] \dicDirectTranslationCS{houska, žemle} \dicExampleIS{rúnnstykki með osti} \dicExampleCS{houska se sýrem}
\dicEntry[rúsína] \dicTerm{rúsín|a} \dicIPA{{r}{u}{\textlengthmark}{s}{i}{n}{a}} \dicPos{f}[1] \dicFlx{(‑u, ‑ur)}[7] \dicSynonym*{þurrkað vínber} \dicDirectTranslationCS{(h)rozinka};  \dicPhraseIS{rúsínan í pylsuendanum} \dicLangCat{přen.} \dicDirectTranslationCS{třešnička na dortu}
\dicEntry[rúskinn] \dicTerm{rú··skinn} \dicIPA{{r}{u}{s}{\r{\textObardotlessj}}{\textsci}{\textsubring{n}}} \dicPos{n}[2] \dicFlx{(‑s, ‑)}[5] \dicDirectTranslationCS{semiš}
\dicEntry[Rússi] \dicTerm{Rúss|i} \dicIPA{{r}{u}{s}{\textlengthmark}{\textsci}} \dicPos{m}[1] \dicFlx{(‑a, ‑ar)}[1] \dicDirectTranslationCS{Rus(ka)}
\dicEntry[rússíbani] \dicTerm{rússí··ban|i} \dicIPA{{r}{u}{s}{\textlengthmark}{i}{\textsubring{b}}{a}{n}{\textsci}} \dicPos{m}[1] \dicFlx{(‑a, ‑ar)}[8] \dicDirectTranslationCS{horská dráha}
\dicEntry[Rússland] \dicTerm{Rúss··land} \dicIPA{{r}{u}{s}{l}{a}{n}{\textsubring{d}}} \dicPos{n}[2] \dicFlx{(‑s)}[4] \dicFieldCat{geog.} \dicDirectTranslationCS{Rusko}
\dicEntry[rússneska] \dicTerm{rúss··nesk|a} \dicIPA{{r}{u}{s}{\textsubring{d}}{n}{\textepsilon}{s}{\r{g}}{a}} \dicPos{f}[1] \dicFlx{(‑u)}[5] \dicDirectTranslationCS{ruština}
\dicEntry[rússneskur] \dicTerm{rúss··neskur} \dicIPA{{r}{u}{s}{\textsubring{d}}{n}{\textepsilon}{s}{\r{g}}{\textscy}{\textsubring{r}}} \dicPos{adj}[1]\dicFlx{}[-1] \dicDirectTranslationCS{ruský}
\dicEntry[rúst] \dicTerm{rúst} \dicIPA{{r}{u}{s}{\textsubring{d}}} \dicPos{f}[7] \dicFlx{(‑ar, ‑ir)}[1] \dicDirectTranslationCS{zřícenina, ruina, troska} \dicExampleIS{rústir af gömlum bæ} \dicExampleCS{ruiny staré farmy}
\dicEntry[rústa] \dicTerm{rúst|a} \dicIPA{{r}{u}{s}{\textsubring{d}}{a}} \dicPos{v}[1] \dicFlx{(‑aði)}[1] \dicFlx{dat\,/\addthin acc} \dicDirectTranslationCS{zdevastovat, zpustošit, zničit} \dicExampleIS{Hann rústaði allt sem í kringum hann var.} \dicExampleCS{Zdevastoval vše kolem sebe.}
\dicEntry[rúta] \dicTerm{rút|a} \dicIPA{{r}{u}{\textlengthmark}{\textsubring{d}}{a}} \dicPos{f}[1] \dicFlx{(‑u, ‑ur)}[7] \dicDirectTranslationCS{(dálkový) autobus, autokar}
\dicEntry[rútíneraður] \dicTerm{rútíner··|aður} \dicIPA{{r}{u}{\textlengthmark}{\textsubring{d}}{i}{n}{\textepsilon}{r}{a}{ð}{\textscy}{\textsubring{r}}} \dicPos{adj}[3] \dicFlx{(f ‑uð)}[3] \dicDirectTranslationCS{procvičený}
\dicEntry[rútubíll] \dicTerm{rútu··bíl|l} \dicIPA{{r}{u}{\textlengthmark}{\textsubring{d}}{\textscy}{\textsubring{b}}{i}{\textsubring{d}}{\textsubring{l}}} \dicPos{m}[6] \dicFlx{(‑s, ‑ar)}[48] \dicDirectTranslationCS{(dálkový) autobus}
\dicEntry[rútubílstjóri] \dicTerm{rútu·bíl··stjór|i} \dicIPA{{r}{u}{\textlengthmark}{\textsubring{d}}{\textscy}{\textsubring{b}}{i}{l}{s}{\textsubring{d}}{j}{ou}{r}{\textsci}} \dicPos{m}[1] \dicFlx{(‑a, ‑ar)}[1] \dicDirectTranslationCS{řidič(ka) (dálkového) autobusu}
\dicEntry[rútuferð] \dicTerm{rútu··ferð} \dicIPA{{r}{u}{\textlengthmark}{\textsubring{d}}{\textscy}{f}{\textepsilon}{r}{\texttheta}} \dicPos{f}[7] \dicFlx{(‑ar, ‑ir)}[1] \dicDirectTranslationCS{jízda\,/\addthin cesta autobusem, autobusový výlet}
\dicEntry[RÚV] \dicTerm{RÚV} \dicPos{zkr} \dicPhraseIS{Ríkisútvarpið} \dicDirectTranslationCS{Islandský státní rozhlas}
\dicEntry[rúþen] \dicTerm{rúþen} \dicIPA{{r}{u}{\textlengthmark}{\texttheta}{\textepsilon}{\textsubring{n}}} \dicPos{n}[2] \dicFlx{(‑s)}[33] \dicFieldCat{chem.} \dicDirectTranslationCS{ruthenium} \textit{(l.~{\textLA{Ru, Ruthenium}})}
\dicEntry[Rvík] \dicTerm{Rvík} \dicPos{zkr.} \dicPhraseIS{Reykjavík} \dicDirectTranslationCS{Reykjavík}
\dicEntry[Rvk] \dicTerm{Rvk} \dicPos{zkr.} \dicPhraseIS{Reykjavík} \dicDirectTranslationCS{Reykjavík}
\dicEntry[ryddi] \dicTerm{ryddi} \dicIPA{{r}{\textsci}{\textsubring{d}}{\textlengthmark}{\textsci}} \dicPos{v} \dicFlx{con pf sg 1 pers} \dicLink{ryðja}
\dicEntry[ryð] \dicTerm{ryð\smash{\textsuperscript{1}}} \dicIPA{{r}{\textsci}{\textlengthmark}{\texttheta}} \dicPos{n}[2] \dicFlx{(‑s)}[2] \dicDirectTranslationCS{rez}
\dicEntry[ryð] \dicTerm{ryð\smash{\textsuperscript{2}}} \dicIPA{{r}{\textsci}{\textlengthmark}{\texttheta}} \dicPos{v} \dicFlx{ind praes sg 1 pers} \dicLink{ryðja}
\dicEntry[ryðfrír] \dicTerm{ryð··frír} \dicIPA{{r}{\textsci}{ð}{f}{r}{i}{\textsubring{r}}} \dicPos{adj}[4]\dicFlx{}[-1] \dicDirectTranslationCS{nerezový} \dicExampleIS{ryðfrítt stál} \dicExampleCS{nerezová ocel}
\dicEntry[ryðga] \dicTerm{ryðg|a} \dicIPA{{r}{\textsci}{ð}{\r{g}}{a}} \dicPos{v}[1] \dicFlx{(‑aði)}[1] \textbf{1.} \dicSynonym*{fá á sig ryð} \dicDirectTranslationCS{(z)rezavět, (z)rezivět}  \textbf{2.} \dicDirectTranslationCS{zapomenout (znalosti ap.)} \dicExampleIS{ryðga í ensku} \dicExampleCS{zapomenout angličtinu}
\dicEntry[ryðgaður] \dicTerm{ryðg··|aður} \dicIPA{{r}{\textsci}{ð}{\r{g}}{a}{ð}{\textscy}{\textsubring{r}}} \dicPos{adj}[3] \dicFlx{(f ‑uð)}[3] \textbf{1.} \dicDirectTranslationCS{rezavý, zrezavělý} \dicExampleIS{ryðgaður nagli} \dicExampleCS{rezavý hřebík}  \textbf{2.} \dicDirectTranslationCS{zapomínající (určité znalosti ap.)} \dicExampleIS{ryðgaður í enskunni} \dicExampleCS{(jsoucí) s~chatrnými znalostmi angličtiny}
\dicEntry[ryði] \dicTerm{ryði} \dicIPA{{r}{\textsci}{\textlengthmark}{ð}{\textsci}} \dicPos{v} \dicFlx{con pf sg 1 pers} \dicLink{rjóða}
\dicEntry[ryðja] \dicTerm{ryðja} \dicsymFrequent\  \dicIPA{{r}{\textsci}{ð}{j}{a}} \dicPos{v}[4] \dicFlx{(ryð, ruddi, ruddum, ryddi, rutt)}[42] \dicFlx{acc\,/\addthin dat} \textbf{1.} \dicFlx{acc} \dicDirectTranslationCS{vyklidit, vyklízet, uvolnit, uvolňovat} \dicExampleIS{ryðja salinn} \dicExampleCS{uvolnit sál}  \textbf{2.} \dicPhraseIS{ryðja e‑u burt} \dicFlx{dat} \dicSynonym{fjarlægja} \dicDirectTranslationCS{odstranit\,/\addthin odstraňovat (co)} \dicExampleIS{ryðja steini burt} \dicExampleCS{odstranit kámen};  \dicIdiom{ryðja}[sér]{ \dicPhraseIS{ryðja sér braut\,/\addthin leið (gegnum e‑ð)}} \dicDirectTranslationCS{razit si cestu (skrz (co)) (prales ap.)};  \dicIdiom{ryðja}[til]{ \dicPhraseIS{ryðja sér til rúms}} \dicLangCat{přen.} \dicDirectTranslationCS{prorazit si cestu, rozšířit se (idea ap.)};  \dicIdiom{ryðja (út) úr\,/\addthin upp úr}{ \dicPhraseIS{ryðja (út) úr sér e‑u, ryðja upp úr sér e‑u}} \dicDirectTranslationCS{vychrlit (co) ze sebe};  \dicIdiom{ryðjast}{ \dicPhraseIS{ryðjast}} \dicFlx{refl} \dicDirectTranslationCS{(pro)tlačit se, prodírat se} \dicExampleIS{ryðjast í gegnum mannþröngina} \dicExampleCS{tlačit se davem}
\dicEntry[ryfi] \dicTerm{ryfi} \dicIPA{{r}{\textsci}{\textlengthmark}{v}{\textsci}} \dicPos{v} \dicFlx{con pf sg 1 pers} \dicLink{rjúfa}
\dicEntry[ryk] \dicTerm{ryk} \dicsymFrequent\  \dicIPA{{r}{\textsci}{\textlengthmark}{\r{g}}} \dicPos{n}[2] \dicFlx{(‑s)}[2] \dicDirectTranslationCS{prach} \dicExampleIS{Það sest ryk á glerið.} \dicExampleCS{Na skle se usazuje prach.};  \dicIdiom{ryk}{ \dicPhraseIS{dusta rykið af e‑u}} \dicLangCat{přen.} \dicDirectTranslationCS{oprášit (co) (znalosti ap.)}
\dicEntry[rykfallinn] \dicTerm{ryk··fallinn} \dicIPA{{r}{\textsci}{\textlengthmark}{\r{g}}{f}{a}{\textsubring{d}}{l}{\textsci}{\textsubring{n}}} \dicPos{adj}[6]\dicFlx{}[-4] \dicDirectTranslationCS{zaprášený}
\dicEntry[ryki] \dicTerm{ryki} \dicIPA{{r}{\textsci}{\textlengthmark}{\r{\textObardotlessj}}{\textsci}} \dicPos{v} \dicFlx{con pf sg 1 pers} \dicLink{rjúka}
\dicEntry[rykkja] \dicTerm{rykk|ja} \dicIPA{{r}{\textsci}{h}{\r{\textObardotlessj}}{a}} \dicPos{v}[2] \dicFlx{(‑ti, ‑t)}[25] \dicFlx{dat\,/\addthin acc} \textbf{1.} \dicFlx{dat} \dicSynonym{kippa\smash{\textsuperscript{2}}} \dicDirectTranslationCS{škubnout, trhnout} \dicExampleIS{rykkja í tauminn} \dicExampleCS{škubnout otěžemi}  \textbf{2.} \dicFlx{acc} \dicDirectTranslationCS{nařasit, nabrat (sukni ap.)} \dicExampleIS{rykkja pils} \dicExampleCS{nařasit sukni};  \dicPhraseIS{rykkjast} \dicFlx{refl} \dicDirectTranslationCS{kodrcat se, nadskakovat, být házený ze strany na stranu}
\begin{xtolerant}{}{1pt}
\dicEntry[rykkjóttur] \dicTerm{rykkj··óttur}\addthinS\dicIPA{{r}{\textsci}{h}{\r{\textObardotlessj}}{ou}{h}{\textsubring{d}}{\textscy}{\textsubring{r}}}\addthinS\dicPos{adj}[1]\dicFlx{}[-10]\addthinS\dicDirectTranslationCS{trhavý, přerývaný}\addthinS\dicExampleIS{rykkjóttur andardráttur} \dicExampleCS{přerývaný dech}
\end{xtolerant}
\dicEntry[rykkur] \dicTerm{rykk|ur} \dicIPA{{r}{\textsci}{h}{\r{g}}{\textscy}{\textsubring{r}}} \dicPos{m}[9] \dicFlx{(‑jar\,/\addthin ‑s, ‑ir)}[26] \dicSynonym{kippur} \dicDirectTranslationCS{trhnutí, škubnutí};  \dicPhraseIS{í einum rykk} \dicFlx{adv} \dicDirectTranslationCS{na jeden zátah} \dicExampleIS{fara alla leiðina í einum rykk} \dicExampleCS{ujet celou cestu na jeden zátah}
\dicEntry[ryksuga] \dicTerm{ryk··sug|a\smash{\textsuperscript{1}}} \dicIPA{{r}{\textsci}{\textlengthmark}{\r{g}}{s}{\textscy}{\textbabygamma}{a}} \dicPos{f}[1] \dicFlx{(‑u, ‑ur)}[13] \dicDirectTranslationCS{lux, vysavač}
\dicEntry[ryksuga] \dicTerm{ryk··sug|a\smash{\textsuperscript{2}}} \dicIPA{{r}{\textsci}{\textlengthmark}{\r{g}}{s}{\textscy}{\textbabygamma}{a}} \dicPos{v}[1] \dicFlx{(‑aði)}[1] \dicFlx{acc} \dicDirectTranslationCS{(vy)luxovat, vysát, vysávat} \dicExampleIS{ryksuga teppið} \dicExampleCS{luxovat koberec}
\dicEntry[rykugur] \dicTerm{ryk··ugur} \dicIPA{{r}{\textsci}{\textlengthmark}{\r{g}}{\textscy}{\textbabygamma}{\textscy}{\textsubring{r}}} \dicPos{adj}[1]\dicFlx{}[-8] \dicDirectTranslationCS{zaprášený, prašný}
\dicEntry[rym] \dicTerm{rym} \dicIPA{{r}{\textsci}{\textlengthmark}{\textsubring{m}}} \dicPos{v} \dicFlx{ind praes sg 1 pers} \dicLink{rymja}
\dicEntry[rymdi] \dicTerm{rymdi} \dicIPA{{r}{\textsci}{m}{\textsubring{d}}{\textsci}} \dicPos{v} \dicFlx{con pf sg 1 pers} \dicLink{rymja}
\dicEntry[rymja] \dicTerm{rymja} \dicIPA{{r}{\textsci}{m}{j}{a}} \dicPos{v}[4] \dicFlx{(rym, rumdi, rumdum, rymdi, rumið)}[35] \textbf{1.} \dicDirectTranslationCS{(za)mručet, (za)bručet (člověk ap.)}  \textbf{2.} \dicDirectTranslationCS{(za)hýkat (osel ap.)}  \textbf{3.} \dicDirectTranslationCS{(za)mručet, (za)bručet (medvěd ap.)}
\dicEntry[rynni] \dicTerm{rynni} \dicIPA{{r}{\textsci}{n}{\textlengthmark}{\textsci}} \dicPos{v} \dicFlx{con pf sg 1 pers} \dicLink{renna\smash{\textsuperscript{3}}}
\dicEntry[rysjóttur] \dicTerm{rysj··óttur} \dicIPA{{r}{\textsci}{s}{j}{ou}{h}{\textsubring{d}}{\textscy}{\textsubring{r}}} \dicPos{adj}[1]\dicFlx{}[-10] \dicSynonym{óstöðugur} \dicDirectTranslationCS{nejistý, problematický (doba ap.)} \dicExampleIS{rysjótt tíð} \dicExampleCS{nejistá doba}
\dicEntry[rytjulegur] \dicTerm{rytju··legur} \dicIPA{{r}{\textsci}{\textlengthmark}{\textsubring{d}}{j}{\textscy}{l}{\textepsilon}{\textbabygamma}{\textscy}{\textsubring{r}}} \dicPos{adj}[1]\dicFlx{}[-8] \dicSynonym{ræfilslegur} \dicDirectTranslationCS{mizerný, ubohý, nevalný} \dicExampleIS{rytjuleg blóm} \dicExampleCS{ubohé květy}
\dicEntry[rytmi] \dicTerm{rytm|i} \dicIPA{{r}{\textsci}{h}{\textsubring{d}}{m}{\textsci}} \dicPos{m}[1] \dicFlx{(‑a)}[3] \dicLink{ritmi}
\dicEntry[rý] \dicTerm{rý} \dicIPA{{r}{i}{\textlengthmark}} \dicPos{v} \dicFlx{ind praes sg 1 pers} \dicLink{rýja\smash{\textsuperscript{2}}}
\dicEntry[rýð] \dicTerm{rýð} \dicIPA{{r}{i}{\textlengthmark}{\texttheta}} \dicPos{v} \dicFlx{ind praes sg 1 pers} \dicLink{rjóða}
\dicEntry[rýf] \dicTerm{rýf} \dicIPA{{r}{i}{\textlengthmark}{f}} \dicPos{v} \dicFlx{ind praes sg 1 pers} \dicLink{rjúfa}
\dicEntry[rýi] \dicTerm{rýi} \dicIPA{{r}{i}{j}{\textsci}} \dicPos{v} \dicFlx{ind praes sg 1 pers} \dicLink{rýja}
\dicEntry[rýja] \dicTerm{rýj|a\smash{\textsuperscript{1}}} \dicIPA{{r}{i}{j}{\textlengthmark}{a}} \dicPos{f}[1] \dicFlx{(‑u, ‑ur)}[7] \textbf{1.} \dicSynonym{tuska} \dicDirectTranslationCS{hadr, hadřík}  \textbf{2.} \dicLangCat{zast.} \dicDirectTranslationCS{chudáček, ubožátko} \dicIndirectTranslationCS{(v~milém oslovení)} \dicExampleIS{rýjan mín} \dicExampleCS{chudáčku můj}
\dicEntry[rýja] \dicTerm{rýja\smash{\textsuperscript{2}}} \dicIPA{{r}{i}{j}{\textlengthmark}{a}} \dicPos{v}[4] \dicFlx{(rý\,/\addthin rýi, rúði, rúðum, rúði, rúið\,/\addthin rúð)}[55] \dicFlx{acc} \dicDirectTranslationCS{(o)stříhat (ovci)};  \dicIdiom{rýja}[inn]{ \dicPhraseIS{rýja e‑n inn að skinninu\,/\addthin skyrtunni}} \dicLangCat{přen.} \dicDirectTranslationCS{oškubat (koho), sedřít (koho) z~kůže}
\dicEntry[rýk] \dicTerm{rýk} \dicIPA{{r}{i}{\textlengthmark}{\r{g}}} \dicPos{v} \dicFlx{ind praes sg 1 pers} \dicLink{rjúka}
\dicEntry[rýma] \dicTerm{rým|a} \dicIPA{{r}{i}{\textlengthmark}{m}{a}} \dicPos{v}[2] \dicFlx{(‑di, ‑t)}[141] \dicFlx{acc} \textbf{1.} \dicDirectTranslationCS{vyklidit, vyklízet, uvolnit, uvolňovat} \dicExampleIS{rýma herbergið} \dicExampleCS{uvolnit pokoj}  \textbf{2.} \dicDirectTranslationCS{vyklidit, vyklízet, evakuovat} \dicExampleIS{rýma flugvöllinn} \dicExampleCS{vyklidit letiště};  \dicIdiom{rýma}[til]{ \dicPhraseIS{rýma til}} \dicDirectTranslationCS{udělat místo}
\dicEntry[rými] \dicTerm{rými} \dicIPA{{r}{i}{\textlengthmark}{m}{\textsci}} \dicPos{n}[2] \dicFlx{(‑s, ‑)}[14] \dicSynonym{rúm} \dicDirectTranslationCS{prostor, místo} \dicExampleIS{nýta rýmið vel} \dicExampleCS{dobře využít prostor}
\dicEntry[rýmingarsala] \dicTerm{rýmingar··|sala} \dicIPA{{r}{i}{\textlengthmark}{m}{i}{\ng}{\r{g}}{a}{\textsubring{r}}{s}{a}{l}{a}} \dicPos{f}[1] \dicFlx{(‑sölu, ‑sölur)}[20] \dicDirectTranslationCS{(úplný) výprodej\,/\addthin doprodej}
\dicEntry[rýmka] \dicTerm{rýmk|a} \dicIPA{{r}{i}{\textsubring{m}}{\r{g}}{a}} \dicPos{v}[1] \dicFlx{(‑aði)}[1] \dicFlx{acc} \dicSynonym*{auka rúm} \dicDirectTranslationCS{rozšířit, rozšiřovat, uvolnit, uvolňovat, liberalizovat} \dicExampleIS{rýmka kosningaréttinn} \dicExampleCS{liberalizovat volební právo};  \dicIdiom{rýmka}[til]{ \dicPhraseIS{rýmka til e‑s staðar}} \dicDirectTranslationCS{udělat (kde) místo};  \dicIdiom{rýmka}[um]{ \dicPhraseIS{rýmka um e‑ð}} \dicDirectTranslationCS{zvětšit místo pro (co) (nohy v~letadle ap.)};  \dicIdiom{rýmkast}{ \dicPhraseIS{rýmkast}} \dicFlx{refl} \dicDirectTranslationCS{rozšířit se, zvětšit se, expandovat}; { \dicPhraseIS{það rýmkast um e‑ð}} \dicFlx{refl impers} \dicDirectTranslationCS{je větší prostor pro (co)}
\dicEntry[rýmstur] \dicTerm{rýmstur} \dicIPA{{r}{i}{m}{s}{\textsubring{d}}{\textscy}{\textsubring{r}}} \dicPos{adj} \dicFlx{m sg nom sup} \dicLink{rúmur}
\dicEntry[rýna] \dicTerm{rýn|a} \dicsymFrequent\  \dicIPA{{r}{i}{\textlengthmark}{n}{a}} \dicPos{v}[2] \dicFlx{(‑di, ‑t)}[141] \dicSynonym{skoða} \dicDirectTranslationCS{(pro)zkoumat, (pečlivě) prohlédnout\,/\addthin prohlížet} \dicExampleIS{rýna í e‑ð} \dicExampleCS{prozkoumat (co)}
\dicEntry[rýr] \dicTerm{rýr} \dicIPA{{r}{i}{\textlengthmark}{\textsubring{r}}} \dicPos{adj}[5] \dicFlx{(f ‑)}[8] \textbf{1.} \dicSynonym{lélegur} \dicDirectTranslationCS{nevydatný, slabý, skrovný (úroda ap.)} \dicExampleIS{rýr heyskapur} \dicExampleCS{nevydatná senoseč}  \textbf{2.} \dicSynonym{magur} \dicDirectTranslationCS{vychrtlý, vyzáblý} \dicExampleIS{rýr kind} \dicExampleCS{vychrtlá ovce}
\dicEntry[rýra] \dicTerm{rýr|a} \dicIPA{{r}{i}{\textlengthmark}{r}{a}} \dicPos{v}[2] \dicFlx{(‑ði, ‑t)}[99] \dicFlx{acc} \dicSynonym*{draga úr} \dicDirectTranslationCS{oslabit, oslabovat, podkopat, podkopávat (důvěru ap.)} \dicExampleIS{rýra trúverðugleika} \dicExampleCS{oslabit důvěryhodnost}
\dicEntry[rýrna] \dicTerm{rýrn|a} \dicIPA{{r}{i}{r}{\textsubring{d}}{n}{a}} \dicPos{v}[1] \dicFlx{(‑aði)}[34] \dicSynonym{minnka} \dicDirectTranslationCS{(z)menšit se, ubýt, ubývat} \dicExampleIS{Vöðvarnir hafa rýrnað við fótbrotið.} \dicExampleCS{Kvůli zlomenině nohy ubylo svalů.}
\dicEntry[rýrnun] \dicTerm{rýrn|un} \dicIPA{{r}{i}{r}{\textsubring{d}}{n}{\textscy}{\textsubring{n}}} \dicPos{f}[7] \dicFlx{(‑unar)}[9] \textbf{1.} \dicSynonym{hnignun} \dicDirectTranslationCS{úbytek, snížení, zmenšení} \dicExampleIS{rýrnun jökla} \dicExampleCS{úbytek ledovců}  \textbf{2.} \dicFieldCat{med.} \dicDirectTranslationCS{úbytek, atrofie} \dicExampleIS{vöðvarýrnun} \dicExampleCS{svalová atrofie}
\dicEntry[rýta] \dicTerm{rýt|a} \dicIPA{{r}{i}{\textlengthmark}{\textsubring{d}}{a}} \dicPos{v}[2] \dicFlx{(‑ti, ‑t)}[62] \dicDirectTranslationCS{(za)chrochtat, (za)kvičet} \dicExampleIS{Gölturinn rýtir.} \dicExampleCS{Kanec chrochtá.}
\dicEntry[ræ] \dicTerm{ræ} \dicIPA{{r}{a}{i}{\textlengthmark}} \dicPos{v} \dicFlx{ind praes sg 1 pers} \dicLink{róa\smash{\textsuperscript{2}}}
\dicEntry[ræð] \dicTerm{ræð} \dicIPA{{r}{a}{i}{\textlengthmark}{\texttheta}} \dicPos{v} \dicFlx{ind praes sg 1 pers} \dicLink{ráða}
\dicEntry[ræða] \dicTerm{ræð|a\smash{\textsuperscript{1}}} \dicsymFrequent\  \dicIPA{{r}{a}{i}{\textlengthmark}{ð}{a}} \dicPos{f}[1] \dicFlx{(‑u, ‑ur)}[19] \dicSynonym{tal} \dicDirectTranslationCS{řeč, projev, proslov} \dicExampleIS{halda ræðu} \dicExampleCS{mít projev};  \dicPhraseIS{bein ræða} \dicFieldCat{jaz.} \dicDirectTranslationCS{přímá řeč};  \dicPhraseIS{óbein ræða} \dicFieldCat{jaz.} \dicDirectTranslationCS{nepřímá řeč}
\dicEntry[ræða] \dicTerm{ræ|ða\smash{\textsuperscript{2}}} \dicsymFrequent\  \dicIPA{{r}{a}{i}{\textlengthmark}{ð}{a}} \dicPos{v}[2] \dicFlx{(‑ddi, ‑tt)}[167] \dicFlx{acc} \dicDirectTranslationCS{(po)hovořit, (pro)diskutovat, (pro)mluvit} \dicExampleIS{ræða málið við e‑n} \dicExampleCS{prodiskutovat s~(kým) záležitost};  \dicPhraseIS{ræða við e‑n} \dicDirectTranslationCS{(po)hovořit s~(kým)};  \dicPhraseIS{ræða um e‑ð} \dicDirectTranslationCS{diskutovat o~(čem), hovořit o~(čem)};  \dicIdiom{ræða}[um]{ \dicPhraseIS{það er ekki um annað að ræða en að (gera e‑ð)}} \dicDirectTranslationCS{nezbývá nic jiného než ((co) udělat)}; { \dicPhraseIS{það er um (tvennt) að ræða}} \dicDirectTranslationCS{jsou (dvě) možnosti};  \dicIdiom{ræðast}[við]{ \dicPhraseIS{ræðast við}} \dicFlx{refl} \dicDirectTranslationCS{hovořit, diskutovat, mluvit} \dicExampleIS{Þeir ræddust við undir fjögur augu.} \dicExampleCS{Diskutovali mezi čtyřma očima.}
\dicEntry[ræðinn] \dicTerm{ræðinn} \dicIPA{{r}{a}{i}{\textlengthmark}{ð}{\textsci}{\textsubring{n}}} \dicPos{adj}[6]\dicFlx{}[-2] \dicDirectTranslationCS{hovorný, povídavý, upovídaný}
\dicEntry[ræðismaður] \dicTerm{ræðis··|maður} \dicIPA{{r}{a}{i}{\textlengthmark}{ð}{\textsci}{s}{m}{a}{ð}{\textscy}{\textsubring{r}}} \dicPos{m}[13] \dicFlx{(‑manns, ‑menn)}[2] \dicDirectTranslationCS{konzul(ka)}
\dicEntry[ræðisskrifstofa] \dicTerm{ræðis··skrif·stof|a} \dicIPA{{r}{a}{i}{\textlengthmark}{ð}{\textsci}{s}{\r{g}}{r}{\textsci}{f}{s}{\textsubring{d}}{\textopeno}{v}{a}} \dicPos{f}[1] \dicFlx{(‑u, ‑ur)}[7] \dicDirectTranslationCS{konzulát}
\dicEntry[ræðumaður] \dicTerm{ræðu··|maður} \dicIPA{{r}{a}{i}{\textlengthmark}{ð}{\textscy}{m}{a}{ð}{\textscy}{\textsubring{r}}} \dicPos{m}[13] \dicFlx{(‑manns, ‑menn)}[2] \dicDirectTranslationCS{řečník, řečnice, mluvčí}
\dicEntry[ræðustóll] \dicTerm{ræðu··stól|l} \dicIPA{{r}{a}{i}{\textlengthmark}{ð}{\textscy}{s}{\textsubring{d}}{ou}{\textsubring{d}}{\textsubring{l}}} \dicPos{m}[6] \dicFlx{(‑s, ‑ar)}[48] \dicDirectTranslationCS{řečnický stolek}
\dicEntry[ræfill] \dicTerm{ræf|ill} \dicIPA{{r}{a}{i}{\textlengthmark}{v}{\textsci}{\textsubring{d}}{\textsubring{l}}} \dicPos{m}[6] \dicFlx{(‑ils, ‑lar)}[35] \textbf{1.} \dicSynonym*{óreiðumaður} \dicDirectTranslationCS{mizera, lotr(yně)}  \textbf{2.} \dicSynonym{vesalingur} \dicDirectTranslationCS{nešťastník, nešťastnice, ubožák, ubožačka} \dicIndirectTranslationCS{(nejčastěji s~určitým členem)}  \textbf{3.} \dicSynonym*{lélegur hlutur} \dicDirectTranslationCS{haraburdí, harampádí, krám}
\dicEntry[ræfilsháttur] \dicTerm{ræfils··hátt|ur} \dicIPA{{r}{a}{i}{\textlengthmark}{v}{\textsci}{l}{s}{h}{au}{h}{\textsubring{d}}{\textscy}{\textsubring{r}}} \dicPos{m}[12] \dicFlx{(‑ar)}[8] \dicSynonym{vesaldómur} \dicDirectTranslationCS{ubohost, bídnost}
\dicEntry[ræfilslegur] \dicTerm{ræfils··legur} \dicIPA{{r}{a}{i}{\textlengthmark}{v}{\textsci}{l}{s}{l}{\textepsilon}{\textbabygamma}{\textscy}{\textsubring{r}}} \dicPos{adj}[1]\dicFlx{}[-8] \textbf{1.} \dicSynonym{tuskulegur} \dicDirectTranslationCS{potrhaný, rozedraný}  \textbf{2.} \dicSynonym{veikindalegur} \dicDirectTranslationCS{neduživý, churavý}
\dicEntry[rægja] \dicTerm{ræg|ja} \dicIPA{{r}{a}{i}{j}{\textlengthmark}{a}} \dicPos{v}[2] \dicFlx{(‑ði, ‑t)}[89] \dicFlx{acc} \dicSynonym{niðra} \dicDirectTranslationCS{očernit, očerňovat, zostudit, ostouzet} \dicExampleIS{rægja e‑n} \dicExampleCS{ostouzet (koho)}
\dicEntry[ræki] \dicTerm{ræki} \dicIPA{{r}{a}{i}{\textlengthmark}{\r{\textObardotlessj}}{\textsci}} \dicPos{v} \dicFlx{con pf sg 1 pers} \dicLink{reka\smash{\textsuperscript{2}}}
\dicEntry[rækilega] \dicTerm{ræki··lega} \dicsymFrequent\  \dicIPA{{r}{a}{i}{\textlengthmark}{\r{\textObardotlessj}}{\textsci}{l}{\textepsilon}{\textbabygamma}{a}} \dicPos{adv} \dicSynonym{vandlega} \dicDirectTranslationCS{důkladně, pořádně} \dicExampleIS{skoða e‑ð rækilega} \dicExampleCS{prozkoumat (co) důkladně}
\dicEntry[rækilegur] \dicTerm{ræki··legur} \dicIPA{{r}{a}{i}{\textlengthmark}{\r{\textObardotlessj}}{\textsci}{l}{\textepsilon}{\textbabygamma}{\textscy}{\textsubring{r}}} \dicPos{adj}[1]\dicFlx{}[-8] \dicSynonym{vandlegur} \dicDirectTranslationCS{důkladný, zevrubný} \dicExampleIS{rækilegar leiðbeiningar} \dicExampleCS{důkladné instrukce}
\dicEntry[rækja] \dicTerm{rækj|a\smash{\textsuperscript{1}}} \dicIPA{{r}{a}{i}{\textlengthmark}{\r{\textObardotlessj}}{a}} \dicPos{f}[1] \dicFlx{(‑u, ‑ur)}[25] \dicFieldCat{zool.} \dicDirectTranslationCS{kreveta, krevetka severní} \textit{(l.~{\textLA{Pandalus borealis}})}  \dicsymPhoto\ 
\dicFigure{ds_image_raekja_1_1.jpg}{Rækja}{Rækja - U.S. National Oceanic and Atmospheric Administration, PD}
\dicEntry[rækja] \dicTerm{ræk|ja\smash{\textsuperscript{2}}} \dicIPA{{r}{a}{i}{\textlengthmark}{\r{\textObardotlessj}}{a}} \dicPos{v}[2] \dicFlx{(‑ti, ‑t)}[33] \dicFlx{acc} \dicSynonym{stunda} \dicDirectTranslationCS{dbát, hledět si} \dicExampleIS{rækja e‑ð} \dicExampleCS{dbát na (co)};  \dicPhraseIS{rækja skyldur sínar} \dicDirectTranslationCS{hledět si svých povinností}
\dicEntry[rækt] \dicTerm{rækt} \dicIPA{{r}{a}{i}{x}{\textsubring{d}}} \dicPos{f}[4] \dicFlx{(‑ar)}[3] \textbf{1.} \dicSynonym{ræktun} \dicDirectTranslationCS{kultivace, pěstování, obdělávání, chov} \dicExampleIS{garðrækt} \dicExampleCS{obdělávání zahrady}  \textbf{2.} \dicSynonym{umhyggja} \dicDirectTranslationCS{péče, starost};  \dicPhraseIS{leggja rækt við e‑ð} \dicDirectTranslationCS{dávat si záležet na (čem), dbát na (co)}  \textbf{3.} \dicPhraseIS{ræktin} \dicFlx{def} \dicDirectTranslationCS{posilovna, tělocvična} \dicExampleIS{fara í ræktina} \dicExampleCS{jít do posilovny}
\dicEntry[rækta] \dicTerm{rækt|a} \dicIPA{{r}{a}{i}{x}{\textsubring{d}}{a}} \dicPos{v}[1] \dicFlx{(‑aði)}[1] \dicFlx{acc} \textbf{1.} \dicSynonym*{yrkja} \dicDirectTranslationCS{obdělávat (půdu ap.)} \dicExampleIS{rækta landið} \dicExampleCS{obdělávat půdu}  \textbf{2.} \dicSynonym*{framleiða með ræktun} \dicDirectTranslationCS{pěstovat (obilí ap.)} \dicExampleIS{rækta kartöflur} \dicExampleCS{pěstovat brambory}  \textbf{3.} \dicDirectTranslationCS{chovat, pěstovat (zvířata ap.)}  \textbf{4.} \dicSynonym*{hirða um} \dicDirectTranslationCS{dbát, pečovat, starat se}
\begin{xtolerant}{}{1pt}
\dicEntry[ræktanlegur] \dicTerm{ræktan··legur}\addthinS\dicIPA{{r}{a}{i}{x}{\textsubring{d}}{a}{n}{l}{\textepsilon}{\textbabygamma}{\textscy}{\textsubring{r}}}\addthinS\dicPos{adj}[1]\dicFlx{}[-8]\addthinS\dicDirectTranslationCS{obdělávatelný, upravovatelný}
\end{xtolerant}
\dicEntry[ræktarlaus] \dicTerm{ræktar··laus} \dicIPA{{r}{a}{i}{x}{\textsubring{d}}{a}{r}{l}{\oe i}{s}} \dicPos{adj}[5]\dicFlx{}[-1] \textbf{1.} \dicSynonym*{umhyggjulaus} \dicDirectTranslationCS{netečný, lhostejný}  \textbf{2.} \dicSynonym*{óræktaður} \dicDirectTranslationCS{neobdělávatelný}
\dicEntry[ræktarsemi] \dicTerm{ræktar··sem|i} \dicIPA{{r}{a}{i}{x}{\textsubring{d}}{a}{\textsubring{r}}{s}{\textepsilon}{m}{\textsci}} \dicPos{f}[2] \dicFlx{(‑i)}[2] \dicDirectTranslationCS{náklonnost, přízeň} \dicExampleIS{ræktarsemi við landið} \dicExampleCS{přízeň k~zemi}
\dicEntry[ræktun] \dicTerm{rækt|un} \dicIPA{{r}{a}{i}{x}{\textsubring{d}}{\textscy}{\textsubring{n}}} \dicPos{f}[7] \dicFlx{(‑unar)}[9] \textbf{1.} \dicDirectTranslationCS{obdělávání (půdy ap.)}  \textbf{2.} \dicDirectTranslationCS{pěstování, kultivace (obilí ap.)} \dicExampleIS{ræktun kartaflna} \dicExampleCS{pěstování brambor};  \dicPhraseIS{lífræn ræktun} \dicDirectTranslationCS{organické pěstování}  \textbf{3.} \dicDirectTranslationCS{pěstování, chov (dobytka ap.)}  \textbf{4.} \dicFieldCat{med.} \dicDirectTranslationCS{kultivace}
\dicEntry[ræma] \dicTerm{ræm|a} \dicIPA{{r}{a}{i}{\textlengthmark}{m}{a}} \dicPos{f}[1] \dicFlx{(‑u, ‑ur)}[19] \textbf{1.} \dicSynonym{lengja\smash{\textsuperscript{1}}} \dicDirectTranslationCS{pruh, pás(ka)}  \textbf{2.} \dicSynonym*{landræma} \dicDirectTranslationCS{pruh\,/\addthin pás země}
\dicEntry[ræna] \dicTerm{ræn|a\smash{\textsuperscript{1}}} \dicIPA{{r}{a}{i}{\textlengthmark}{n}{a}} \dicPos{f}[1] \dicFlx{(‑u)}[5] \dicSynonym{meðvitund} \dicDirectTranslationCS{vědomí} \dicIndirectTranslationCS{(stav člověka, v~kterém si uvědomuje svou existenci)} \dicExampleIS{vera með rænu} \dicExampleCS{být při vědomí};  \dicPhraseIS{hafa ekki rænu á að (gera e‑ð)} \dicDirectTranslationCS{nemít sílu ((co) udělat), nenasbírat dostatek sil na (udělání (čeho))}
\dicEntry[ræna] \dicTerm{ræn|a\smash{\textsuperscript{2}}} \dicsymFrequent\  \dicIPA{{r}{a}{i}{\textlengthmark}{n}{a}} \dicPos{v}[2] \dicFlx{(‑di, ‑t)}[155] \dicFlx{dat\,/\addthin acc} \textbf{1.} \dicFlx{dat} \dicDirectTranslationCS{(u)krást, odcizit, odcizovat, (u)loupit} \dicExampleIS{ræna málverki} \dicExampleCS{ukrást obraz}  \textbf{2.} \dicFlx{acc} \dicDirectTranslationCS{(vy)loupit, (vy)krást} \dicExampleIS{ræna banka} \dicExampleCS{vyloupit banku}
\dicEntry[ræningi] \dicTerm{ræn··ing|i} \dicIPA{{r}{a}{i}{\textlengthmark}{n}{i}{\textltailn}{\r{\textObardotlessj}}{\textsci}} \dicPos{m}[1] \dicFlx{(‑ja, ‑jar)}[14] \dicDirectTranslationCS{lupič(ka), zloděj(ka)}
\dicEntry[rænulaus] \dicTerm{rænu··laus} \dicIPA{{r}{a}{i}{\textlengthmark}{n}{\textscy}{l}{\oe i}{s}} \dicPos{adj}[5]\dicFlx{}[-1] \dicSynonym{meðvitundarlaus} \dicDirectTranslationCS{bezvědomý, (jsoucí) v~bezvědomí}
\dicEntry[rær] \dicTerm{rær} \dicIPA{{r}{a}{i}{\textlengthmark}{\textsubring{r}}} \dicPos{f} \dicFlx{pl nom} \dicLink{ró\smash{\textsuperscript{1}}}
\dicEntry[ræsa] \dicTerm{ræs|a} \dicsymFrequent\  \dicIPA{{r}{a}{i}{\textlengthmark}{s}{a}} \dicPos{v}[2] \dicFlx{(‑ti, ‑t)}[63] \dicFlx{acc} \dicSynonym*{setja í gang} \dicDirectTranslationCS{spustit, spouštět, (na)startovat, (z)aktivovat} \dicExampleIS{ræsa vélina} \dicExampleCS{spustit stroj};  \dicIdiom{ræsa}[fram]{ \dicPhraseIS{ræsa fram e‑ð}} \dicDirectTranslationCS{odvodnit (co), vysušit (co)} \dicExampleIS{ræsa fram mýrina} \dicExampleCS{odvodnit mokřiny};  \dicIdiom{ræsa}[út]{ \dicPhraseIS{ræsa út e‑ð\,/\addthin e‑n}} \dicDirectTranslationCS{přivolat\,/\addthin zavolat\,/\addthin zalarmovat (co\,/\addthin koho) (požárníky ap.)}
\dicEntry[ræsi] \dicTerm{ræsi} \dicIPA{{r}{a}{i}{\textlengthmark}{s}{\textsci}} \dicPos{n}[2] \dicFlx{(‑s, ‑)}[14] \dicDirectTranslationCS{strouha, stoka, (odpadní) kanál}
\dicEntry[ræsir] \dicTerm{ræs|ir} \dicIPA{{r}{a}{i}{\textlengthmark}{s}{\textsci}{\textsubring{r}}} \dicPos{m}[7] \dicFlx{(‑is, ‑ar)}[1] \dicFieldCat{techn.} \dicDirectTranslationCS{startér}
\dicEntry[ræskja] \dicTerm{ræsk|ja} \dicsymFrequent\  \dicIPA{{r}{a}{i}{s}{\r{\textObardotlessj}}{a}} \dicPos{v}[2] \dicFlx{(‑ti, ‑t)}[34] \dicPhraseIS{ræskja sig} \dicDirectTranslationCS{odkašlat si} \dicExampleIS{Hann ræskti sig áður en hann hélt áfram.} \dicExampleCS{Dříve než pokračoval, tak si odkašlal.}
\dicEntry[ræsta] \dicTerm{ræst|a} \dicIPA{{r}{a}{i}{s}{\textsubring{d}}{a}} \dicPos{v}[2] \dicFlx{(‑i, ‑)}[17] \dicFlx{acc} \dicDirectTranslationCS{uklidit, uklízet, poklízet} \dicExampleIS{ræsta skrifstofur} \dicExampleCS{uklízet kanceláře}
\dicEntry[ræstiduft] \dicTerm{ræsti··duft} \dicIPA{{r}{a}{i}{s}{\textsubring{d}}{\textsci}{\textsubring{d}}{\textscy}{f}{\textsubring{d}}} \dicPos{n}[2] \dicFlx{(‑s)}[2] \dicDirectTranslationCS{čisticí prostředek}
\dicEntry[ræsting] \dicTerm{ræst··ing} \dicIPA{{r}{a}{i}{s}{\textsubring{d}}{i}{\ng}{\r{g}}} \dicPos{f}[4] \dicFlx{(‑ar, ‑ar)}[5] \dicDirectTranslationCS{úklid, uklízení (domů ap.)} \dicExampleIS{vinna við ræstingar} \dicExampleCS{živit se uklízením}
\dicEntry[ræstitæknir] \dicTerm{ræsti··tækn|ir} \dicIPA{{r}{a}{i}{s}{\textsubring{d}}{\textsci}{t\smash{\textsuperscript{h}}}{a}{i}{h}{\r{g}}{n}{\textsci}{\textsubring{r}}} \dicPos{m}[7] \dicFlx{(‑is, ‑ar)}[1] \dicDirectTranslationCS{uklízeč(ka)}
\dicEntry[rætast] \dicTerm{rætast} \dicsymFrequent\  \dicIPA{{r}{a}{i}{\textlengthmark}{\textsubring{d}}{a}{s}{\textsubring{d}}} \dicPos{v}[2] \dicFlx{(rættist, ræst)}[190] \dicFlx{refl} \dicSynonym*{koma fram} \dicDirectTranslationCS{vyplnit se, splnit se} \dicExampleIS{Spáin hefur ræst.} \dicExampleCS{Věštba se splnila.};  \dicIdiom{rætast}[úr]{ \dicPhraseIS{það rætist úr e‑m}} \dicFlx{impers} \dicDirectTranslationCS{(kdo) předčí očekávání};  \dicPhraseIS{það rætist úr e‑u} \dicFlx{impers} \dicDirectTranslationCS{(čemu) se daří lépe, než se očekávalo}
\dicEntry[rætinn] \dicTerm{rætinn} \dicIPA{{r}{a}{i}{\textlengthmark}{\textsubring{d}}{\textsci}{\textsubring{n}}} \dicPos{adj}[6]\dicFlx{}[-2] \dicDirectTranslationCS{podlý, zlomyslný}
\dicEntry[rætni] \dicTerm{rætn|i} \dicIPA{{r}{a}{i}{h}{\textsubring{d}}{n}{\textsci}} \dicPos{f}[3] \dicFlx{(‑i)}[3] \dicDirectTranslationCS{podlost, zlomyslnost}
\dicEntry[rætur] \dicTerm{rætur} \dicIPA{{r}{a}{i}{\textlengthmark}{\textsubring{d}}{\textscy}{\textsubring{r}}} \dicPos{f} \dicFlx{pl nom} \dicLink{rót\smash{\textsuperscript{1}}}
\dicEntry[rödd] \dicTerm{rödd} \dicsymFrequent\  \dicIPA{{r}{\oe}{\textsubring{d}}{\textlengthmark}} \dicPos{f}[7] \dicFlx{(raddar, raddir)}[17] \dicDirectTranslationCS{hlas} \dicExampleIS{fjarlægar raddir} \dicExampleCS{vzdálené hlasy};  \dicPhraseIS{brýna röddina} \dicDirectTranslationCS{zvýšit hlas}
\dicEntry[rödduð] \dicTerm{rödduð} \dicIPA{{r}{\oe}{\textsubring{d}}{\textlengthmark}{\textscy}{\texttheta}} \dicPos{adj} \dicFlx{f sg nom pos} \dicLink{raddaður}
\dicEntry[röddun] \dicTerm{rödd|un} \dicIPA{{r}{\oe}{\textsubring{d}}{\textlengthmark}{\textscy}{\textsubring{n}}} \dicPos{f}[7] \dicFlx{(‑unar)}[12] \dicFieldCat{jaz.} \dicDirectTranslationCS{znělost}
\dicEntry[röð] \dicTerm{röð} \dicsymFrequent\  \dicIPA{{r}{\oe}{\textlengthmark}{\texttheta}} \dicPos{f}[7] \dicFlx{(raðar, raðir)}[16] \textbf{1.} \dicDirectTranslationCS{řada, linie} \dicExampleIS{í þriðju röð} \dicExampleCS{ve třetí řadě};  \dicPhraseIS{ganga á röðina} \dicDirectTranslationCS{jít po řadě};  \dicPhraseIS{röðin kemur að e‑m} \dicDirectTranslationCS{řada přichází na (koho)};  \dicPhraseIS{(standa) í röð} \dicDirectTranslationCS{(stát) v~řadě, (stát) ve frontě}  \textbf{2.} \dicFieldCat{mat.} \dicDirectTranslationCS{řada};  \dicIdiom{röð}{ \dicPhraseIS{í fremstu röð}} \dicFlx{adv} \dicDirectTranslationCS{v~první řadě, především}; { \dicPhraseIS{í röð og reglu}} \dicFlx{adv} \dicDirectTranslationCS{v~nejlepším pořádku}; { \dicPhraseIS{vera einstakur í sinni röð}} \dicLangCat{přen.} \dicDirectTranslationCS{být výjimečný, být jediný svého druhu}
\dicEntry[röðun] \dicTerm{röð|un} \dicIPA{{r}{\oe}{\textlengthmark}{ð}{\textscy}{\textsubring{n}}} \dicPos{f}[7] \dicFlx{(‑unar)}[12] \dicDirectTranslationCS{řazení, seřazení, třídění, klasifikace} \dicExampleIS{röðun í hópa} \dicExampleCS{seřazení do skupin}
\dicEntry[röf] \dicTerm{röf} \dicIPA{{r}{\oe}{\textlengthmark}{f}} \dicPos{n} \dicFlx{pl nom} \dicLink{raf\smash{\textsuperscript{1}}}
\dicEntry[röfl] \dicTerm{röfl} \dicIPA{{r}{\oe}{\textsubring{b}}{\textsubring{l}}} \dicPos{n}[2] \dicFlx{(‑s)}[2] \dicLink{rövl}
\dicEntry[röfla] \dicTerm{röfl|a} \dicIPA{{r}{\oe}{\textsubring{b}}{l}{a}} \dicPos{v}[1] \dicFlx{(‑aði)}[34] \dicLink{rövla}
\dicEntry[rög] \dicTerm{rög} \dicIPA{{r}{\oe}{\textlengthmark}{x}} \dicPos{adj} \dicFlx{f sg nom pos} \dicLink{ragur}
\dicEntry[rögg] \dicTerm{rögg} \dicIPA{{r}{\oe}{\r{g}}{\textlengthmark}} \dicPos{f}[6] \dicFlx{(‑var, ‑var)}[1] \dicPhraseIS{taka á sig rögg} \dicLangCat{přen.} \dicDirectTranslationCS{vzchopit se, dát se do kupy}
\dicEntry[röggsamur] \dicTerm{rögg··|samur} \dicIPA{{r}{\oe}{\r{g}}{s}{a}{m}{\textscy}{\textsubring{r}}} \dicPos{adj}[1] \dicFlx{(f ‑söm)}[2] \dicDirectTranslationCS{rázný, energický, průbojný} \dicExampleIS{vera röggsamur stjórnandi} \dicExampleCS{být rázným velitelem}
\dicEntry[rök] \dicTerm{rök\smash{\textsuperscript{1}}} \dicsymFrequent\  \dicIPA{{r}{\oe}{\textlengthmark}{\r{g}}} \dicPos{n}[2] \dicFlx{pl}[9] \dicDirectTranslationCS{důvod, důkaz, argument};  \dicPhraseIS{e‑að á ekki við rök að styðjast} \dicDirectTranslationCS{(co) se neopírá o~argumenty, (co) není zdůvodněné};  \dicPhraseIS{e‑að er á rökum reist} \dicDirectTranslationCS{(co) je odůvodněné};  \dicPhraseIS{færa rök fyrir e‑u, leiða rök að e‑u} \dicDirectTranslationCS{zdůvodnit (co), odůvodnit (co)}
\dicEntry[rök] \dicTerm{rök\smash{\textsuperscript{2}}} \dicIPA{{r}{\oe}{\textlengthmark}{\r{g}}} \dicPos{adj} \dicFlx{f sg nom pos} \dicLink{rakur}
\dicEntry[rökfastur] \dicTerm{rök··|fastur} \dicIPA{{r}{\oe}{\textlengthmark}{\r{g}}{f}{a}{s}{\textsubring{d}}{\textscy}{\textsubring{r}}} \dicPos{adj}[1] \dicFlx{(f ‑föst)}[11] \dicDirectTranslationCS{logický, racionální}
\dicEntry[rökfesta] \dicTerm{rök··fest|a} \dicIPA{{r}{\oe}{\textlengthmark}{\r{g}}{f}{\textepsilon}{s}{\textsubring{d}}{a}} \dicPos{f}[1] \dicFlx{(‑u)}[5] \dicDirectTranslationCS{logičnost}
\dicEntry[rökfræði] \dicTerm{rök··fræð|i} \dicIPA{{r}{\oe}{\textlengthmark}{\r{g}}{f}{r}{a}{i}{ð}{\textsci}} \dicPos{f}[3] \dicFlx{(‑i)}[3] \dicFieldCat{filos.} \dicDirectTranslationCS{logika} \dicExampleIS{stærðfræðileg rökfræði} \dicExampleCS{matematická logika}
\dicEntry[rökhyggja] \dicTerm{rök··hyggj|a} \dicIPA{{r}{\oe}{\textlengthmark}{k\smash{\textsuperscript{h}}}{\textsci}{\r{\textObardotlessj}}{a}} \dicPos{f}[1] \dicFlx{(‑u)}[5] \dicFieldCat{filos.} \dicDirectTranslationCS{racionalismus}
\dicEntry[rökkur] \dicTerm{rökkur} \dicsymFrequent\  \dicIPA{{r}{\oe}{h}{\r{g}}{\textscy}{\textsubring{r}}} \dicPos{n}[2] \dicFlx{(‑s)}[28] \dicSynonym{húm} \dicDirectTranslationCS{soumrak, stmívání, přítmí} \dicExampleIS{í rökkri stofunnar} \dicExampleCS{v~přítmí obývacího pokoje}
\dicEntry[rökkva] \dicTerm{rökkv|a} \dicIPA{{r}{\oe}{h}{\r{g}}{v}{a}} \dicPos{v}[1] \dicFlx{(‑aði)}[78] \dicFlx{impers} \dicPhraseIS{það rökkvar} \dicSynonym*{það dimmir} \dicDirectTranslationCS{stmívá se, smráká se}
\dicEntry[röklegur] \dicTerm{rök··legur} \dicIPA{{r}{\oe}{\textlengthmark}{\r{g}}{l}{\textepsilon}{\textbabygamma}{\textscy}{\textsubring{r}}} \dicPos{adj}[1]\dicFlx{}[-8] \dicDirectTranslationCS{logický} \dicExampleIS{rökleg skýring} \dicExampleCS{logické vysvětlení}
\dicEntry[rökréttur] \dicTerm{rök··réttur} \dicIPA{{r}{\oe}{\textlengthmark}{\r{g}}{r}{j}{\textepsilon}{h}{\textsubring{d}}{\textscy}{\textsubring{r}}} \dicPos{adj}[1]\dicFlx{}[-10] \dicDirectTranslationCS{logický} \dicExampleIS{rökrétt afleiðing} \dicExampleCS{logický důsledek}
\dicEntry[rökræða] \dicTerm{rök··ræ|ða} \dicIPA{{r}{\oe}{\textlengthmark}{\r{g}}{r}{a}{i}{ð}{a}} \dicPos{v}[2] \dicFlx{(‑ddi, ‑tt)}[167] \dicFlx{acc} \dicDirectTranslationCS{(pro)diskutovat, (pro)debatovat, rokovat} \dicExampleIS{rökræða mál} \dicExampleCS{prodiskutovat záležitost};  \dicPhraseIS{rökræða um e‑ð} \dicDirectTranslationCS{debatovat o~(čem)}
\dicEntry[röksemd] \dicTerm{rök··semd} \dicIPA{{r}{\oe}{\textlengthmark}{\r{g}}{s}{\textepsilon}{m}{\textsubring{d}}} \dicPos{f}[7] \dicFlx{(‑ar, ‑ir)}[1] \dicDirectTranslationCS{odůvodnění, důvod, argument} \dicExampleIS{röksemd fyrir þessari skoðun} \dicExampleCS{důvod pro tento názor}
\dicEntry[röksemdafærsla] \dicTerm{rök·semda··færsl|a} \dicIPA{{r}{\oe}{\textlengthmark}{\r{g}}{s}{\textepsilon}{m}{\textsubring{d}}{a}{f}{a}{i}{\textsubring{r}}{s}{\textsubring{d}}{l}{a}} \dicPos{f}[1] \dicFlx{(‑u)}[5] \dicDirectTranslationCS{odůvodnění, argumentace}
\dicEntry[rökstuddi] \dicTerm{rök··studdi} \dicIPA{{r}{\oe}{\textlengthmark}{\r{g}}{s}{\textsubring{d}}{\textscy}{\textsubring{d}}{\textsci}} \dicPos{v} \dicFlx{ind pf sg 1 pers} \dicLink{rökstyðja}
\dicEntry[rökstuddum] \dicTerm{rök··studdum} \dicIPA{{r}{\oe}{\textlengthmark}{\r{g}}{s}{\textsubring{d}}{\textscy}{\textsubring{d}}{\textscy}{\textsubring{m}}} \dicPos{v} \dicFlx{ind pf pl 1 pers} \dicLink{rökstyðja}
\dicEntry[rökstuddur] \dicTerm{rök··studdur} \dicIPA{{r}{\oe}{\textlengthmark}{\r{g}}{s}{\textsubring{d}}{\textscy}{\textsubring{d}}{\textscy}{\textsubring{r}}} \dicPos{adj}[2]\dicFlx{}[-18] \dicDirectTranslationCS{důvodný, odůvodněný, opodstatněný} \dicExampleIS{rökstuddur grunnur} \dicExampleCS{důvodné podezření}
\dicEntry[rökstuðningur] \dicTerm{rök··stuð·ning|ur} \dicIPA{{r}{\oe}{\textlengthmark}{\r{g}}{s}{\textsubring{d}}{\textscy}{ð}{n}{i}{\ng}{\r{g}}{\textscy}{\textsubring{r}}} \dicPos{m}[6] \dicFlx{(‑s)}[9] \dicDirectTranslationCS{důvody, argumenty}
\dicEntry[rökstutt] \dicTerm{rök··stutt} \dicIPA{{r}{\oe}{\textlengthmark}{\r{g}}{s}{\textsubring{d}}{\textscy}{h}{\textsubring{d}}} \dicPos{v} \dicFlx{supin} \dicLink{rökstyðja}
\dicEntry[rökstyddi] \dicTerm{rök··styddi} \dicIPA{{r}{\oe}{\textlengthmark}{\r{g}}{s}{\textsubring{d}}{\textsci}{\textsubring{d}}{\textsci}} \dicPos{v} \dicFlx{con pf sg 1 pers} \dicLink{rökstyðja}
\dicEntry[rökstyð] \dicTerm{rök··styð} \dicIPA{{r}{\oe}{\textlengthmark}{\r{g}}{s}{\textsubring{d}}{\textsci}{\texttheta}} \dicPos{v} \dicFlx{ind praes sg 1 pers} \dicLink{rökstyðja}
\dicEntry[rökstyðja] \dicTerm{rök··|styðja} \dicIPA{{r}{\oe}{\textlengthmark}{\r{g}}{s}{\textsubring{d}}{\textsci}{ð}{j}{a}} \dicPos{v}[4] \dicFlx{(‑styð, ‑studdi, ‑studdum, ‑styddi, ‑stutt)}[41] \dicFlx{acc} \dicDirectTranslationCS{zdůvodnit, zdůvodňovat, odůvodnit, odůvodňovat, opodstatnit, opodstatňovat} \dicExampleIS{rökstyðja mál sitt} \dicExampleCS{zdůvodnit své stanovisko}
\dicEntry[röktum] \dicTerm{röktum} \dicIPA{{r}{\oe}{x}{\textsubring{d}}{\textscy}{\textsubring{m}}} \dicPos{v} \dicFlx{ind pf pl 1 pers} \dicLink{rekja}
\dicEntry[röll] \dicTerm{röll} \dicIPA{{r}{\oe}{\textsubring{d}}{\textsubring{l}}} \dicPos{n} \dicFlx{pl nom} \dicLink{rall}
\dicEntry[rölt] \dicTerm{rölt} \dicIPA{{r}{\oe}{\textsubring{l}}{\textsubring{d}}} \dicPos{n}[2] \dicFlx{(‑s)}[2] \dicSynonym{labb} \dicDirectTranslationCS{procházka, špacír} \dicExampleIS{vera á rölti um bæinn} \dicExampleCS{procházet se po městě}
\dicEntry[rölta] \dicTerm{rölt|a} \dicsymFrequent\  \dicIPA{{r}{\oe}{\textsubring{l}}{\textsubring{d}}{a}} \dicPos{v}[2] \dicFlx{(‑i, ‑)}[17] \dicDirectTranslationCS{projít se, procházet se} \dicExampleIS{rölta um bæinn} \dicExampleCS{projít se po městě}
\dicEntry[römm] \dicTerm{römm} \dicIPA{{r}{\oe}{m}{\textlengthmark}} \dicPos{adj} \dicFlx{f sg nom pos} \dicLink{rammur}
\dicEntry[rönd] \dicTerm{rönd} \dicsymFrequent\  \dicIPA{{r}{\oe}{n}{\textsubring{d}}} \dicPos{f}[8] \dicFlx{(randar, rendur\,/\addthin randir)}[9] \textbf{1.} \dicSynonym{brún} \dicDirectTranslationCS{hrana, okraj}  \textbf{2.} \dicSynonym{rák} \dicDirectTranslationCS{pruh, proužek, obruba} \dicExampleIS{Rendurnar á peysunni eru rauðar.} \dicExampleCS{Proužky na svetru jsou červené.};  \dicIdiom{rönd}{ \dicPhraseIS{elta e‑n á röndum}} \dicLangCat{přen.} \dicDirectTranslationCS{šlapat (komu) na paty}
\dicEntry[röndóttur] \dicTerm{rönd··óttur} \dicIPA{{r}{\oe}{n}{\textsubring{d}}{ou}{h}{\textsubring{d}}{\textscy}{\textsubring{r}}} \dicPos{adj}[1]\dicFlx{}[-10] \dicDirectTranslationCS{pruhovaný, proužkovaný}
\dicEntry[röng] \dicTerm{röng} \dicIPA{{r}{\oe i}{\ng}{\r{g}}} \dicPos{adj} \dicFlx{f sg nom pos} \dicLink{rangur}
\dicEntry[röntgen] \dicTerm{röntgen} \dicIPA{{r}{\oe}{\textsubring{n}}{\textsubring{d}}{\r{\textObardotlessj}}{\textepsilon}{\textsubring{n}}} \dicPos{n}[2] \dicFlx{(‑s)}[33] \dicFieldCat{fyz.} \dicDirectTranslationCS{rentgen} \dicIndirectTranslationCS{(jednotka ozáření)}
\dicEntry[röntgenfræði] \dicTerm{röntgen··fræð|i} \dicIPA{{r}{\oe}{\textsubring{n}}{\textsubring{d}}{\r{\textObardotlessj}}{\textepsilon}{n}{f}{r}{a}{i}{ð}{\textsci}} \dicPos{f}[3] \dicFlx{(‑i)}[3] \dicDirectTranslationCS{radiologie}
\dicEntry[röntgengeisli] \dicTerm{röntgen··geisl|i} \dicIPA{{r}{\oe}{\textsubring{n}}{\textsubring{d}}{\r{\textObardotlessj}}{\textepsilon}{n}{\r{\textObardotlessj}}{ei}{s}{\textsubring{d}}{l}{\textsci}} \dicPos{m}[1] \dicFlx{(‑a, ‑ar)}[1] \dicDirectTranslationCS{rentgenový paprsek, rentgenové záření} \dicIndirectTranslationCS{(zvláště v~pl)}
\dicEntry[röntgenín] \dicTerm{röntgenín} \dicIPA{{r}{\oe}{\textsubring{n}}{\textsubring{d}}{\r{\textObardotlessj}}{\textepsilon}{n}{i}{\textsubring{n}}} \dicPos{n}[2] \dicFlx{(‑s)}[2] \dicFieldCat{chem.} \dicDirectTranslationCS{roentgenium} \textit{(l.~{\textLA{Rg, Roentgenium}})}
\dicEntry[röntgenmynd] \dicTerm{röntgen··mynd} \dicIPA{{r}{\oe}{\textsubring{n}}{\textsubring{d}}{\r{\textObardotlessj}}{\textepsilon}{n}{m}{\textsci}{n}{\textsubring{d}}} \dicPos{f}[7] \dicFlx{(‑ar, ‑ir)}[1] \dicDirectTranslationCS{rentgenový snímek}
\dicEntry[röntgenskoðun] \dicTerm{röntgen··skoð|un} \dicIPA{{r}{\oe}{\textsubring{n}}{\textsubring{d}}{\r{\textObardotlessj}}{\textepsilon}{n}{s}{\r{g}}{\textopeno}{ð}{\textscy}{\textsubring{n}}} \dicPos{f}[7] \dicFlx{(‑unar, ‑anir)}[8] \dicDirectTranslationCS{rentgenové vyšetření}
\dicEntry[röntgentæki] \dicTerm{röntgen··tæki} \dicIPA{{r}{\oe}{\textsubring{n}}{\textsubring{d}}{\r{\textObardotlessj}}{\textepsilon}{n}{t\smash{\textsuperscript{h}}}{a}{i}{\r{\textObardotlessj}}{\textsci}} \dicPos{n}[2] \dicFlx{(‑s, ‑)}[16] \dicDirectTranslationCS{rentgen} \dicIndirectTranslationCS{(přístroj na prosvěcování hmoty)}
\dicEntry[rör] \dicTerm{rör} \dicIPA{{r}{\oe}{\textlengthmark}{\textsubring{r}}} \dicPos{n}[2] \dicFlx{(‑s, ‑)}[5] \dicDirectTranslationCS{(instalatérská) trubka, roura, potrubí}
\dicEntry[röskleiki] \dicTerm{rösk··leik|i} \dicIPA{{r}{\oe}{s}{\r{g}}{l}{ei}{\r{\textObardotlessj}}{\textsci}} \dicPos{m}[1] \dicFlx{(‑a)}[3] \dicDirectTranslationCS{síla, energie, agilnost}
\dicEntry[röskur] \dicTerm{röskur} \dicIPA{{r}{\oe}{s}{\r{g}}{\textscy}{\textsubring{r}}} \dicPos{adj}[1]\dicFlx{}[-1] \textbf{1.} \dicSynonym{vaskur\smash{\textsuperscript{1}}} \dicDirectTranslationCS{aktivní, činorodý, čilý}  \textbf{2.} \dicSynonym{rúmur} \dicDirectTranslationCS{(jsoucí) více než, dobrý (časově)} \dicExampleIS{röskur tveggja tíma gangur} \dicExampleCS{více než dvouhodinový pochod}
\dicEntry[rövl] \dicTerm{rövl}\dicTerm{, röfl} \dicIPA{{r}\-{\oe}\-{v}\-{\textsubring{l}}\-} \dicPos{n}[2] \dicFlx{(‑s)}[2] \dicSynonym{þvaður} \dicDirectTranslationCS{nesmysl, kec, žvást} \dicExampleIS{Vertu ekki með neitt rövl.} \dicExampleCS{Nech těch nesmyslů.}
\dicEntry[rövla] \dicTerm{rövl|a}\dicTerm{, röfla} \dicIPA{{r}\-{\oe}\-{v}\-{l}\-{a}\-} \dicPos{v}[1] \dicFlx{(‑aði)}[34] \dicSynonym{mögla} \dicDirectTranslationCS{stěžovat si, reptat} \dicExampleIS{rövla um e‑ð} \dicExampleCS{stěžovat si na (co)}

\else
\makeatletter
\@for\next:=\alphabetF\do{\input{letters/\next.tex}}
\makeatother 
\fi
\restoregeometry
\pagestyle{plain}

\fi

% ================================================================ PHONETICS =

\ifinputphon

\cleardoublepage

\chapter{Seznam islandských fonémů}              \label{sec:phon_phonems}

\begin{table}[h]
\begin{tabular}{lllll} \toprule

\textbf{foném} & \textbf{vysvětlivka} & \textbf{příklad} & \textbf{fonetický přepis} & \textbf{překlad} \\
\midrule
{\textipa{[{a}]}} & otevřená přední nezaokrouhlená samohl. & rata & {\textipa{[{r}{a}{\textlengthmark}{\textsubring{d}}{a}]}} & trefit \\ 
{\textipa{[{\textepsilon}]}} & polootevřená přední nezaokrouhlená samohl. & meta & {\textipa{[{m}{\textepsilon}{\textlengthmark}{\textsubring{d}}{a}]}} & ocenit \\ 
{\textipa{[{e}]}} & polozavřená přední nezaokrouhlená samohl. & heiður & {\textipa{[{h}{ei}{\textlengthmark}{ð}{\textscy}{\textsubring{r}}]}} & čest \\
{\textipa{[{\textsci}]}} & téměř zavřená téměř přední nezaokrouhlená samohl. & vinna & {\textipa{[{v}{\textsci}{n}{\textlengthmark}{a}]}} & pracovat \\ 
{\textipa{[{i}]}} & zavřená přední nezaokrouhlená samohl. & fínn & {\textipa{[{f}{i}{\textsubring{d}}{\textsubring{n}}]}} & jemný \\ 
{\textipa{[{\textopeno}]}} & polootevřená zadní zaokrouhlená samohl. & lofa & {\textipa{[{l}{\textopeno}{\textlengthmark}{v}{a}]}} & slibovat \\ 
{\textipa{[{o}]}} & polozavřená zadní zaokrouhlená samohl. & bóndi & {\textipa{[{\textsubring{b}}{ou}{n}{\textsubring{d}}{\textsci}]}} & farmář \\ 
{\textipa{[{\textscy}]}} & téměř zavřená téměř přední zaokrouhlená samohl. & hundur & {\textipa{[{h}{\textscy}{n}{\textsubring{d}}{\textscy}{\textsubring{r}}]}} & pes \\ 
{\textipa{[{u}]}} & zavřená zadní zaokrouhlená samohl. & súkkulaði & {\textipa{[{s}{u}{h}{\r{g}}{\textscy}{l}{a}{ð}{\textsci}]}} & čokoláda \\ 
{\textipa{[{\oe}]}} & téměř otevřená přední nezaokrouhlená samohl. & lönd & {\textipa{[{l}{\oe}{n}{\textsubring{d}}]}} & země \\ 
{\textipa{[{m}]}} & bilabiální nazála & mýri & {\textipa{[{m}{i}{\textlengthmark}{r}{\textsci}]}} & mokřina \\ 
{\textipa{[{\textsubring{m}}]}} & neznělá bilabiální nazála & þyrmt & {\textipa{[{\texttheta}{\textsci}{\textsubring{r}}{\textsubring{m}}{\textsubring{d}}]}} & {\textit{$\shortrightarrow$ þyrma}} \\  
{\textipa{[{n}]}} & alveolární nazála & nemandi & {\textipa{[{n}{\textepsilon}{\textlengthmark}{m}{a}{n}{\textsubring{d}}{\textsci}]}} & žák \\ 
{\textipa{[{\textsubring{n}}]}} & neznělá alveolární nazála & mennt & {\textipa{[{m}{\textepsilon}{\textsubring{n}}{\textsubring{d}}]}} & vzdělání \\ 
{\textipa{[{\textltailn}]}} & palatální nazála & angi & {\textipa{[{au}{\textltailn}{\r{\textObardotlessj}}{\textsci}]}} & drobek \\ 
{\textipa{[{\r{\textltailn}}]}} & neznělá palatální nazála & banki & {\textipa{[{\textsubring{b}}{au}{\r{\textltailn}}{\r{\textObardotlessj}}{\textsci}]}} & banka \\
{\textipa{[{\ng}]}} & velární nazála & skyggndur & {\textipa{[{s}{\r{\textObardotlessj}}{\textsci}{\ng}{\textsubring{d}}{\textscy}{\textsubring{r}}]}} & {\textit{$\shortrightarrow$ skyggna}} \\ 
{\textipa{[{\r{\ng}}]}} & neznělá velární nazála & skyggnt & {\textipa{[{s}{\r{\textObardotlessj}}{\textsci}{\r{\ng}}{\textsubring{d}}]}} & {\textit{$\shortrightarrow$ skyggna}} \\ 
{\textipa{[{\textsubring{b}}]}} & neznělá bilabiální ploziva & bara & {\textipa{[{\textsubring{b}}{a}{\textlengthmark}{r}{a}]}} & jenom \\ 
{\textipa{[{p\textsuperscript{h}}]}} & neznělá bilabiální ploziva s přídechem & prestur & {\textipa{[{p\textsuperscript{h}}{r}{\textepsilon}{s}{\textsubring{d}}{\textscy}{\textsubring{r}}]}} & kněz \\ 
{\textipa{[{\textsubring{d}}]}} & neznělá alveolární ploziva & dæmi & {\textipa{[{\textsubring{d}}{a}{i}{\textlengthmark}{m}{\textsci}]}} & příklad \\ 
{\textipa{[{t\textsuperscript{h}}]}} & neznělá alveolární ploziva s přídechem & tómur & {\textipa{[{t\textsuperscript{h}}{ou}{\textlengthmark}{m}{\textscy}{\textsubring{r}}]}} & prázdný \\ 
{\textipa{[{\r{\textObardotlessj}}]}} & neznělá palatální ploziva & gefa & {\textipa{[{\r{\textObardotlessj}}{\textepsilon}{\textlengthmark}{v}{a}]}} & dát \\ 
{\textipa{[{c\textsuperscript{h}}]}} & neznělá palatální ploziva s přídechem & kenna & {\textipa{[{c\textsuperscript{h}}{\textepsilon}{n}{\textlengthmark}{a}]}} & učit \\
{\textipa{[{\r{g}}]}} & neznělá velární ploziva & gata & {\textipa{[{\r{g}}{a}{\textlengthmark}{\textsubring{d}}{a}]}} & ulice \\ 
{\textipa{[{k\textsuperscript{h}}]}} & neznělá velární ploziva s přídechem & koma & {\textipa{[{k\textsuperscript{h}}{\textopeno}{\textlengthmark}{m}{a}]}} & přijít \\
{\textipa{[{v}]}} & znělá labiodentální frikativa & hafa & {\textipa{[{h}{a}{\textlengthmark}{v}{a}]}} & mít \\ 
{\textipa{[{f}]}} & neznělá labiodentální frikativa & finna & {\textipa{[{f}{\textsci}{n}{\textlengthmark}{a}]}} & najít \\ 
{\textipa{[{ð}]}} & znělá dentální frikativa & borða & {\textipa{[{\textsubring{b}}{\textopeno}{r}{ð}{a}]}} & jíst \\ 
{\textipa{[{\texttheta}]}} & neznělá dentální frikativa & þola & {\textipa{[{\texttheta}{\textopeno}{\textlengthmark}{l}{a}]}} & snést \\
{\textipa{[{s}]}} & neznělá alveolární frikativa & systir & {\textipa{[{s}{\textsci}{s}{\textsubring{d}}{\textsci}{\textsubring{r}}]}} & sestra \\ 
{\textipa{[{\c{c}}]}} & neznělá palatální frikativa & hjóla & {\textipa{[{\c{c}}{ou}{\textlengthmark}{l}{a}]}} & jet na kole \\ 
{\textipa{[{\textbabygamma}]}} & znělá velární frikativa & sagði & {\textipa{[{s}{a}{\textbabygamma}{ð}{\textsci}]}} & řekl \\ 
{\textipa{[{x}]}} & neznělá velární frikativa & lag & {\textipa{[{l}{a}{\textlengthmark}{x}]}} & píseň \\
{\textipa{[{h}]}} & neznělá glotální frikativa & heima & {\textipa{[{h}{ei}{\textlengthmark}{m}{a}]}} & doma \\ 
{\textipa{[{l}]}} & alveolární laterální aproximanta & loforð & {\textipa{[{l}{\textopeno}{\textlengthmark}{v}{\textopeno}{r}{\texttheta}]}} & slib \\ 
{\textipa{[{\textsubring{l}}]}} & neznělá alveolární laterální aproximanta & hlutur & {\textipa{[{\textsubring{l}}{\textscy}{\textlengthmark}{\textsubring{d}}{\textscy}{\textsubring{r}}]}} & věc \\  
{\textipa{[{j}]}} & palatální aproximanta & jól & {\textipa{[{j}{ou}{\textlengthmark}{\textsubring{l}}]}} & Vánoce \\ 
{\textipa{[{r}]}} & alveolární vibranta & reynsla & {\textipa{[{r}{ei}{n}{s}{\textsubring{d}}{l}{a}]}} & zkušenost \\ 
{\textipa{[{\textsubring{r}}]}} & neznělá alveolární vibranta & hraun & {\textipa{[{\textsubring{r}}{\oe i}{\textlengthmark}{\textsubring{n}}]}} & láva \\ 
\bottomrule
\end{tabular}
\end{table}

\clearpage

%\section{Výslovnost islandských samohlásek}      \label{sec:phon_vowels}
%\LTXtable{\columnwidth}{phon/vowels}
%\clearpage

%\section{Výslovnost islandských souhlásek}       \label{sec:phon_consonants}
%\LTXtable{\columnwidth}{phon/consonants}
%\clearpage
\fi

% =============================================================== MORPHOLOGY =

\ifinputmorpho

\cleardoublepage

\chapter{Morfologický klíč}                      \label{sec:morpho}

\section{Podstatné jméno, rod mužský}            \label{sec:morpho_m}
{\small\LTXfw{morpho/m}}
\clearpage

\section{Podstatné jméno, rod ženský}            \label{sec:morpho_f}
{\small\LTXfw{morpho/f}}
\clearpage

\section{Podstatné jméno, rod střední}           \label{sec:morpho_n}
{\small\LTXfw{morpho/n}}
\clearpage

\section{Přídavné jméno}                         \label{sec:morpho_adj}
{\small\LTXfw{morpho/adj}}
\clearpage

\section{Zájmeno}                                \label{sec:morpho_pron}
{\small\LTXfw{morpho/pron}}
 \begin{center}
\begin{minipage}[t]{.45\textwidth}
\begin{tabular}{l>{\footnotesize\itshape}lll}
 {\textbf{\textit{třída}}} & {\textit{pád}}   & \textit{sg} & \textit{pl}  \\ 
 \hline
\multirow{3}{*}{{{\textbf{pron} \Large{\textbf{17}}}}}  &  nom & ég & við   \\*
 & acc &  mig  & okkur  \\*
 & dat & mér & okkur   \\*
 & gen & mín  & okkar  \\
\hline

\multirow{3}{*}{{{\textbf{pron} \Large{\textbf{17}}}}}  &  nom & hann & þeir   \\*
 & acc &  hann  & þá  \\*
 & dat & honum & þeim   \\*
 & gen & hans  & þeirra  \\
\hline

\multirow{3}{*}{{{\textbf{pron} \Large{\textbf{17}}}}}  &  nom & hún & þær   \\*
 & acc &  hana  & þær  \\*
 & dat & henni & þeim   \\*
 & gen & hennar  & þeirra  \\
\hline

\multirow{3}{*}{{{\textbf{pron} \Large{\textbf{17}}}}}  &  nom &  & vér   \\*
 & acc &    & oss  \\*
 & dat &  & oss   \\*
 & gen &   & vor  \\
\hline

\multirow{3}{*}{{{\textbf{pron} \Large{\textbf{17}}}}}  &  nom &  & þér   \\*
 & acc &    & yður  \\*
 & dat &  & yður   \\*
 & gen &   & yðar  \\
\hline

\multirow{3}{*}{{{\textbf{pron} \Large{\textbf{17}}}}}  &  nom & þú & þið   \\*
 & acc &  þig  & ykkur  \\*
 & dat & þér & ykkur   \\*
 & gen & þín  & ykkar  \\
\hline

\multirow{3}{*}{{{\textbf{pron} \Large{\textbf{17}}}}}  &  nom & það & þau   \\*
 & acc &  það  & þau  \\*
 & dat & því & þeim   \\*
 & gen & þess  & þeirra  \\
\hline
\end{tabular}
\end{minipage}\hfil
\begin{minipage}[t]{.45\textwidth}
\begin{tabular}{l>{\footnotesize\itshape}lll}
\hline
\multirow{3}{*}{{{\textbf{pron} \Large{\textbf{18}}}}}  &  nom &  &    \\*
 & acc &  sig  &   \\*
 & dat & sér &    \\*
 & gen & sín  &   \\
 \hline
\end{tabular}
\end{minipage}
\end{center}

\clearpage

\section{Číslovka}                               \label{sec:morpho_num}
 {\small\LTXfw{morpho/num}}
\clearpage

\begin{landscape}
\section{Sloveso}                                \label{sec:morpho_v}
{\small\LTXfw{morpho/v}}
\end{landscape}

\clearpage

\fi

% ============================================================== BACK MATTER =

\ifmakebackmatter

% ------------------------------------------------ Photographs authors index -

\cleardoublepage
%\pagestyle{empty}
%\pagestyle{indexstyle}
%\phantomsection
%\addcontentsline{toc}{chapter}
%  {Seznam autorů fotografií a ilustrací}        
\label{sec:photo}
\printindex[figures]

% ------------------------------------------------------------- Bibliography -
\cleardoublepage
\chapter{Bibliografie}
%\twocolumn
%\phantomsection
%\addcontentsline{toc}{chapter}{Bibliografie}
\begin{multicols}{2}
\nocite{*}\printbibliography[heading=none]
\end{multicols}

% ------------------------------------------------------------------ Apendix -
\cleardoublepage
\onecolumn
\ifPDF
\else
\chapter{Apendix: Islandská gramatika}
\clearpage
\covergeometry  
\makeatletter
\includepdf[
  pages={1,2,3,4,5,6,7,8,9,10,11,12,13,14,15,16,17,18,19,20,21,22,23,24,25,26,27,28,29,30,31,32,33,34,35,36,37,38,39,40,41,42,43,44,45,46,47,48,49,50,51,52,53,54,55,56,57,58,59,60,61,62,63,64,65,66,67,68,69,70,71,72,73},
  turn=false,
  scale=0.79,
  keepaspectratio,
  offset={\dimexpr(\paperwidth-\csname Gm@layoutwidth\endcsname)/2-\csname Gm@layouthoffset\endcsname\relax}
  {\dimexpr(\paperheight-\csname  Gm@layoutheight\endcsname)/2-\csname Gm@layoutvoffset\endcsname\relax}
   ]{gram.pdf}
\makeatother
\restoregeometry  
\fi

% ------------------------------------------------------------------ Barcode -
\pagenumbering{gobble}
\ifPDF
\else
\cleardoublepage\null\clearpage % clear to verso
\vspace*\fill
\EANisbn[SC3]
\fi
\fi

% --------------------------------------------------------------- Back cover -

\ifmakecovers
  \cleardoublepage\null\clearpage % inner side of the cover
\covergeometry  
  \makecoverwith\backcoverimages
\restoregeometry  
\fi

\end{document}

